\documentclass[12pt,a4paper]{report}
\usepackage[utf8]{inputenc}
\usepackage[T2A]{fontenc}
\usepackage[russian]{babel}
\usepackage{amsmath,amssymb,amsthm,mathtools}
\usepackage{booktabs,longtable,array}
\usepackage{microtype}
\usepackage{geometry}
\usepackage{hyperref}
\usepackage{xurl}
\geometry{margin=2.3cm}
\setlength{\emergencystretch}{2em}
% Shared file-search paths for documents built from s2t/docs/.
% Keep gate filenames stable: the directory split is resolved here.
\makeatletter
\def\input@path{{../gates/}{./}}
\makeatother

% Allow long machine-readable filenames in \texttt to wrap at underscores.
\let\originaltexttt\texttt
\renewcommand{\texttt}[1]{{\small\ttfamily\def\_{\char`\_\allowbreak}#1}}

\newtheorem{definition}{Определение}[chapter]
\newtheorem{axiom}{Аксиома}[chapter]
\newtheorem{lemma}{Лемма}[chapter]
\newtheorem{proposition}{Предложение}[chapter]
\newtheorem{theorem}{Теорема}[chapter]
\newtheorem{corollary}{Следствие}[chapter]
\newtheorem{hypothesis}{Гипотеза}[chapter]
\theoremstyle{remark}
\newtheorem{remark}{Замечание}[chapter]
\newtheorem{protocol}{Протокол}[chapter]
\newtheorem{nogocriterion}{Критерий провала}[chapter]

\DeclareMathOperator{\Tr}{Tr}
\DeclareMathOperator{\Str}{Str}
\DeclareMathOperator{\Vol}{Vol}
\DeclareMathOperator{\Spec}{Spec}
\DeclareMathOperator{\ind}{ind}
\DeclareMathOperator{\Res}{Res}
\DeclareMathOperator{\diag}{diag}

\title{Трактат об устройстве мира\\
Том II: Спектральное замыкание и квантовая топология микромира}
\author{}
\date{}

\begin{document}

\pagenumbering{roman}
\hypersetup{pageanchor=false}
\maketitle
\clearpage
\hypersetup{pageanchor=true}
\setcounter{page}{1}
\tableofcontents
\clearpage
\pagenumbering{arabic}

\chapter*{Предисловие: фиксация задачи замыкания}
\addcontentsline{toc}{chapter}{Предисловие: фиксация задачи замыкания}

Второй том начинается не с расширения онтологии первого тома, а с фиксации задачи замыкания. Его предметом является проверка того, допускает ли программа S2T строгую реализацию в виде спектрально-геометрической и эффективной калибровочной теории на компактном внутреннем носителе. Все качественные утверждения первого тома здесь должны быть заменены на определения, операторы, спектры, функционалы, регуляризации, численные протоколы и условия провала.

Рабочий носитель тома фиксируется как компактный фактор
\begin{equation}
    K = L(2,1)\times S^1 \simeq \mathbb{RP}^3\times S^1,
\end{equation}
а его спинорное накрытие, возникающее в фермионных секторах, обозначается
\begin{equation}
    \widetilde K = S^3\times S^1.
\end{equation}
Если в отдельных вычислениях используется запись \(S^3\times S^1\), она должна пониматься либо как спинорное накрытие \(\widetilde K\), либо как явно оговорённая тестовая геометрия. Это различение является обязательным: смешение \(\mathbb{RP}^3\times S^1\) и \(S^3\times S^1\) без указания сектора считается ошибкой нормировки.

\begin{protocol}[Definition of Done тома II]
Том считается замкнутым только если выполнены все условия:
\begin{enumerate}
    \item фиксирован компактный носитель, его метрика, радиусы, спин-структуры и нормировки калибровочных полей;
    \item спектр базовых операторов записан с указанием граничных условий и кратностей;
    \item коэффициенты теплового ядра и дзета-остатки выведены без подбора непрерывных параметров;
    \item электромагнитный вакуумный скаляр \(S_{vac}\) получен из спектральной схемы и схемы вычитаний;
    \item лептонный и нейтринный сектора имеют лагранжеву EFT-формулировку;
    \item все численные предсказания воспроизводятся из разрешённых входов;
    \item проведены blind-test, permutation-control, sensitivity и ablation-проверки;
    \item явно перечислены условия, при которых версия S2T признаётся несостоятельной.
\end{enumerate}
\end{protocol}

\begin{nogocriterion}[Запрет скрытой подгонки]
Любой непрерывный коэффициент, введённый после просмотра blind-наблюдаемых и не выводимый из спектра, нормировки, топологии или EFT-лагранжиана, переводит соответствующий результат из статуса вывода в статус феноменологической подгонки. Если ключевой результат тома требует такого коэффициента, S2T-замыкание в данной версии не достигнуто.
\end{nogocriterion}


\section*{Протокол наследования тезисов тома I}
\addcontentsline{toc}{section}{Протокол наследования тезисов тома I}

Первый том задаёт не готовую вычислительную модель, а онтологический и методологический каркас. Поэтому второй том должен явно указать, какие тезисы тома I переведены в технические требования, а какие остаются задачами дальнейшей формализации. Повторная сверка показывает: спектрально-геометрические тезисы внесены в рабочую часть тома II, но корреляционно-онтологический слой был представлен слишком неявно. Фиксируем явную матрицу переноса.

\begin{center}
\small
\begin{tabular}{>{\raggedright\arraybackslash}p{0.28\textwidth}>{\raggedright\arraybackslash}p{0.36\textwidth}>{\raggedright\arraybackslash}p{0.24\textwidth}}
\toprule
Тезис или постулат тома I & Реализация в томе II & Статус \\
\midrule
Первичность структуры над ``вещами'' & частицы, поля и параметры трактуются как устойчивые спектральные/EFT-режимы, а не как исходные сущности & внесено, но требует явного языка режимов \\
Первичность допустимости над локальностью & допустимые моды задаются топологией, spin-структурой, граничными условиями, голономиями и \(\det'\)-правилами & внесено \\
Многослойность реальности & геометрический, спектральный, тополого-фазовый и EFT-слои реализованы как отдельные главы и аудиты & внесено частично \\
Частицы и взаимодействия как устойчивые моды & лептонные, нейтринные, EM и EW/QCD-блоки оформлены как секторные редукции & частично; полная классификация наблюдаемых open \\
Единый источник редукций & роль единой сущности технически заменена минимальным носителем \(K\), спектром операторов и EFT-лагранжианом & частичный перевод, не полное доказательство \(\mathfrak S\) \\
Постулат допустимости & введены S2T-допустимые параметры и запрет скрытых коэффициентов & внесено \\
Операторная порождаемость & базовые утверждения должны выводиться из спектров, determinant-ветвей и функционалов & внесено \\
Геометрическая реализуемость & выбран и проверяется компактный носитель \(\mathbb{RP}^3\times S^1\) & внесено \\
Корреляционная целостность & связи между EM, лептонным, хиггсовским, нейтринным и EW/QCD-секторами пока заданы через общие нормировки и open-roadmap & задача II.B \\
Инвариантное ядро & роль ядра выполняют \(K\), набор допустимых входов, спектральные операторы, no-go критерии и blind-аудит & внесено как технический surrogate \\
Пять фаз реконструкции & инвариантное ядро, минимальная формализация, вывод редукций, back-check и blind-предсказания превращены в Definition of Done & внесено, но нужно явное сопоставление \\
Селекция конкурирующих моделей & реализована через permutation-control, ablation, no-go и запрет альтернативных схем без аудита & внесено частично \\
Минимальное ядро \(\mathfrak K_{min}\) & заменено проверяемым ядром II.A: \((K,\Spec,\mathcal D_{train},\mathcal L_{EFT},\mathcal Q_{audit})\) & техническая версия внесена \\
\bottomrule
\end{tabular}
\end{center}

Следовательно, второй том не должен утверждать, что полностью доказал онтологию тома I. Его более точный статус: он переводит ядро тома I в проверяемую спектрально-геометрическую программу. Неполно перенесёнными остаются три элемента: корреляционная целостность между всеми секторами, строгая реконструкция единой сущности \(\mathfrak S\) и классификация всех наблюдаемых по режимам чтения. Эти элементы запрещено считать закрытыми только потому, что отдельные численные блоки совпали.

\begin{protocol}[Traceability-аудит тома I]
Для контроля переноса добавлен машинный список \texttt{s2t\_tome1\_traceability\_results.json}. Он проверяет наличие в томе II маркеров для онтологических тезисов, постулатов и критериев отбора тома I. Аудит показывает, что до настоящей вставки явно присутствовали допустимость, операторы, геометрическая реализуемость, spin-структуры и голономии, но отсутствовали прямые маркеры первичности структуры, многослойности, единого источника, корреляционной целостности, инвариантного ядра и селекции моделей. Настоящий протокол закрывает именно документационный разрыв; математическое закрытие трёх open-элементов остаётся задачей II.B.
\end{protocol}

\section*{Red-team нулевого промта: условный вывод и неизбежность}
\addcontentsline{toc}{section}{Red-team нулевого промта: условный вывод и неизбежность}

Исходный файл
\texttt{RPFT-main/ai-promts/First-principles-00.md} был сформулирован как
дедукция носителя из NCG, реальной структуры и модулярного потока. Повторный
аудит показывает, что он корректно восстанавливает
\(\mathbb{RP}^3\times S^1\) только условно, поскольку часть результата уже
записана в требованиях вопроса.

Первый переход является условно строгим:
\begin{equation}
 \{\text{единичные кватернионы}\}=SU(2)\simeq S^3.
\end{equation}
Но трёхмерность, односвязность, групповая структура и именно
кватернионная симметрия в нулевом промте заданы заранее, а не получены
минимизацией спектрального функционала.

Второй переход
\begin{equation}
 SU(2)/\{\pm1\}=SO(3)\simeq\mathbb{RP}^3
\end{equation}
также математически корректен после требования факторизации по центру.
Однако реальная структура \(J\) действует как антилинейное сопряжение и
реализует противоположную алгебру; сама по себе она не отождествляет
антиподальные точки базового многообразия. Физическая тривиальность центра
является дополнительной аксиомой.

Наиболее существенный разрыв относится к модулярному времени.
Теория Томиты--Такесаки порождает группу
\begin{equation}
 \sigma_t(A)=\Delta^{it}A\Delta^{-it},
 \qquad t\in\mathbb R,
\end{equation}
но конечность системы и KMS-свойство не делают эту группу периодической.
Точный конечномерный контрпример задаётся состоянием
\begin{equation}
 \rho=Z^{-1}\diag(1,e^{-1},e^{-\sqrt2}).
\end{equation}
На матричных единицах частоты модулярного потока содержат \(1\) и
\(\sqrt2\). Общий период \(T\) потребовал бы
\begin{equation}
 T=2\pi m,
 \qquad
 T\sqrt2=2\pi n,
\end{equation}
то есть рациональности \(\sqrt2=n/m\). Следовательно, реальный модулярный
поток может быть квазипериодическим. KMS-периодичность в мнимом времени
может мотивировать евклидову тепловую окружность, но не выводит физический
\(S^1\) без дополнительной интерпретации и выбора \(\beta\).

Наконец, даже после выбора базы \(\mathbb{RP}^3\) и кругового волокна
прямая продуктовая склейка не единственна:
\begin{equation}
 H^2(\mathbb{RP}^3;\mathbb Z)\cong\mathbb Z_2.
\end{equation}
Поэтому существуют два класса главных \(U(1)\)-расслоений над
\(\mathbb{RP}^3\): тривиальный продукт и нетривиальный класс. Нулевой промт
не сравнивает их спектральные действия. Он также не задаёт функционал
энергии, класс конкурирующих геометрий, уравнение радиуса или определение
«спектральной плотности».

\begin{proposition}[Исправленный статус нулевого промта]
Нулевой промт доказывает импликацию
\begin{equation}
 \left.
 \begin{gathered}
 S^3\text{ как кватернионный групповой сектор},\\
 Z_2\text{-факторизация базы},\\
 \text{независимо постулированный периодический круг},\\
 \text{прямая продуктовая склейка}
 \end{gathered}
 \right\}
 \Longrightarrow
 K=\mathbb{RP}^3\times S^1,
\end{equation}
но не доказывает, что эти четыре условия следуют из NCG и KMS или что
полученный носитель является единственным минимумом энергии.
\end{proposition}

\begin{protocol}[Definition of Done неизбежности носителя]
Статус неизбежного вакуума разрешён только если:
\begin{enumerate}
 \item класс конкурирующих спектральных геометрий задан без вставки
 \(\mathbb{RP}^3\) и \(S^1\) в определение;
 \item на всём классе определён один нормированный функционал действия или
 свободной энергии;
 \item \(Z_2\)-факторизация следует из представленной алгебры и \(J\);
 \item периодичность модулярного спектра выведена либо честно заменена
 евклидовой KMS-окружностью с определённым \(\beta\);
 \item сравнены тривиальное и нетривиальное круговые расслоения;
 \item доказана единственность устойчивого минимума с точностью до изометрии
 и заранее объявленной масштабной нормировки.
\end{enumerate}
\end{protocol}

\begin{nogocriterion}[Запрет target-loaded неизбежности]
Если искомая топология, факторизация, периодичность или продуктовая структура
внесены в вопросы до определения вариационного класса и функционала, итог
имеет статус условной реконструкции, а не вывода единственного мира.
\end{nogocriterion}

Таким образом, все вычисления на \(K=\mathbb{RP}^3\times S^1\) сохраняются,
но их правильная предпосылка формулируется так: это рабочий носитель,
выделенный внутри позднее объявленного минимального класса, а не
онтологически неизбежный результат нулевого промта.

Машинный ledger сохранён в
\texttt{s2t\_zero\_prompt\_inevitability\_results.json}.


\section*{Внешняя литература и spectral-determinant gate}
\addcontentsline{toc}{section}{Внешняя литература и spectral-determinant gate}

Внутренние аудиты тома II полезны только как кандидаты на вывод. Для электромагнитного determinant-блока дополнительно требуется внешний литературный контроль: спектры форм на линзовых пространствах, аналитическая кручение Рэя--Зингера, gauge-fixed Maxwell determinant и Faddeev--Popov ghost-нормировка. Поэтому добавлен реестр \texttt{external\_literature\_spectral\_determinants\_results.json} и wiki-страница \texttt{external-literature-spectral-determinants.md}.

Минимальный набор внешних опорных линий таков:
\begin{enumerate}
    \item Ray--Singer analytic torsion: дзета-регуляризованные детерминанты лапласианов на формах и связь с топологической кручением;
    \item Schwarz-type partition functions for degenerate quadratic functionals: связь gauge-fixing, нулевых мод и analytic torsion;
    \item Ikeda--Yamamoto и Lauret--Miatello--Rossetti по спектрам функций и \(p\)-форм на lens spaces;
    \item Faddeev--Popov/BRST ghost determinant: знак, степень детерминанта и cancellation exact/scalar-секторов.
\end{enumerate}

Эта литература поддерживает общий язык Hodge-разложения
\begin{equation}
    \Omega^1=d\Omega^0\oplus \Omega^1_{coex}\oplus \mathcal{H}^1,
\end{equation}
и необходимость раздельно считать scalar, exact, coexact, harmonic и ghost-сектора. Она также делает законным сам вопрос о determinant-combination, но не доказывает автоматически S2T-выбор
\begin{equation}
    N=d_0+d_2=1+9=10.
\end{equation}

Следовательно, результат
\begin{equation}
    N_{need}=10.0099700224\approx10
\end{equation}
получает статус literature-compatible internal lead, а не статуса теоремы. Для повышения статуса необходимо вывести смешанный след в gauge-fixed Maxwell--ghost операторе на \(\mathbb{RP}^3\times S^1\) и показать:
\begin{enumerate}
    \item exact one-form sector действительно наследует нужный scalar spectrum после \(\det'\)-правила;
    \item ghost determinant не удаляет весь вклад \(d_0+d_2\);
    \item знак второго порядка совпадает с требуемым знаком \(-\frac12\Tr((A^{-1}B)^2)\);
    \item внешние формулы для \(p\)-form spectra на \(L(2,1)\) не добавляют обязательные \(\ell=4,6,\ldots\) вклады того же порядка;
    \item существует операторное selection rule, выделяющее \(\ell=0,2\), а не постфактум выбранный integer cutoff.
\end{enumerate}

\begin{nogocriterion}[Провал external spectral-determinant gate]
Если литература по спектрам \(p\)-форм на lens spaces или стандартная Maxwell--ghost нормировка требует включить дополнительные scalar/even оболочки того же порядка, меняет знак смешанного следа, полностью сокращает exact/scalar вклад, либо делает \(d_0+d_2\) gauge-choice dependent, то число \(10\) не считается выведенным. В этом случае Casimir-mixing остаётся сильной внутренней эвристикой, но не может повышать статус \(S_{vac}\) выше условного determinant-успеха.
\end{nogocriterion}

\section*{Реестр нестрогих мест версии II.A}
\addcontentsline{toc}{section}{Реестр нестрогих мест версии II.A}

Перед началом вывода фиксируем стартовые обязательства версии II.A и их итоговый статус после последующих глав. Этот реестр является рабочим списком аудита: закрытые узлы должны иметь явный вывод, а незакрытые --- остаться помеченными как задачи II.B.

\begin{center}
\begin{tabular}{>{\raggedright\arraybackslash}p{0.27\textwidth}>{\raggedright\arraybackslash}p{0.22\textwidth}>{\raggedright\arraybackslash}p{0.37\textwidth}}
\toprule
Узел & Итоговый статус II.A & Основание или следующий шаг \\
\midrule
Выбор \(K=\mathbb{RP}^3\times S^1\) & условно выделено в минимальном классе & объёмный тест \(Z_A=\pi^2\) и фазовый шаг \(\pi\); нулевой промт не доказывает уникальный вакуум и не исключает нетривиальное \(S^1\)-расслоение \\
Spin-сдвиг \(k+1/2\) & выведено в фермионном секторе & анти-периодическая spin-структура на окружности и допустимый фермионный спектр \\
Слагаемое \(\pi\) в \(S_{geo}\) & выведено в EM-секторе & \(\mathbb{Z}_2\)-голономия нетривиальной плоской ветви на \(\mathbb{RP}^3\) \\
Коэффициент \(1/24\) & условное determinant-замыкание & periodic Maxwell--ghost ветвь с \(\det'\); полный статус требует coexact Bessel-tail no-go \\
Член \(1/(\pi^4S^2)\) & strong structural compression & C6 закрыт отрицательно: Casimir kernel не равен winding log-determinant, same-scheme compensation отсутствует \\
QED-сдвиг \(-\alpha/3\) & условная normalization-гипотеза & Bessel-сумма вычислима, но coefficient \(1/3\) требует невыведенного \(\mathcal J_{RP^3}\) \\
Хиггсовский мост \(v,\lambda_H,M_H\) & условный absolute-scale bridge & минимум \(U(H,\chi)\) выведен, но \(v\) и \(M_H\) наследуют условность \(m_\tau\) и \(S_{vac}\) \\
Знаменатель \(23+\pi^{-1}\) & выведен в canonical superconnection metric & parent tubular restriction закрыт; generic spectral-kernel Hessian не универсален \\
Нейтринная вставка \(m_D^{(\nu)}\) & задача II.B & вывести \(\sqrt{\pi+\pi^{-1}}m_e^2/m_\mu\) из \(\mathcal{L}_\nu\) \\
EW/QCD-пороги & задача II.B & получить KK-башню, \(\Sigma_8/\Sigma_3\) и двухпетлевой RG до blind-сравнения \\
Blind-аудит & частично задан & train/blind-разделение зафиксировано; permutation/ablation должны быть выполнены машинно \\
\bottomrule
\end{tabular}
\end{center}

Все последующие главы должны либо подтверждать этот статус явным выводом, либо возвращать соответствующую строку в open gap. Запрещается скрывать open gap внутри успешного численного совпадения.

\chapter{Нормировки, носитель и правила S2T-замыкания}

\section{Разрешённые входные данные}

Второй том использует минимальный набор калибровочных якорей
\begin{equation}
    \mathcal{D}_{train}=\{\alpha^{-1},m_e,m_\mu\},
\end{equation}
если явно не указано иное. Все остальные величины должны либо вычисляться, либо использоваться только в blind-сравнении.

\begin{definition}
Параметр называется S2T-допустимым, если он принадлежит одному из четырёх классов:
\begin{enumerate}
    \item дискретные топологические данные \(\pi_1(K)\), spin-структура, кратность мод;
    \item нормировка, фиксированная стандартной единицей заряда и калибровочной периодичностью;
    \item измеренный train-якорь из \(\mathcal{D}_{train}\);
    \item регуляризационная величина, исчезающая из конечного scheme-invariant результата.
\end{enumerate}
Иные непрерывные параметры считаются запрещёнными до отдельного вывода.
\end{definition}

\section{Фоновые пространства и радиусы}

Основной внутренний носитель задаётся как
\begin{equation}
    K=\mathbb{RP}^3_{R_3}\times S^1_{R_1},
\end{equation}
с метрикой произведения
\begin{equation}
    ds_K^2=R_3^2 ds_{\mathbb{RP}^3}^2+R_1^2 d\theta^2,
    \qquad \theta\sim \theta+2\pi.
\end{equation}
Для фермионных спектральных вычислений используется накрытие
\begin{equation}
    \widetilde K=S^3_{R_3}\times S^1_{R_1},
\end{equation}
с последующей проекцией на допустимые spin-сектора \(\mathbb{RP}^3\).

\begin{axiom}[Минимальная радиусная нормировка]
В базовой версии S2T радиусы фиксируются в планковских единицах как
\begin{equation}
    R_3=1,
\end{equation}
а радиус \(R_1\) либо также фиксируется как безразмерная единица компактификации, либо связывается с физическим ИК-масштабом в конкретном секторе. Всякий переход между этими режимами должен быть указан в формуле.
\end{axiom}

\section{Выбор компактного носителя}

Второй том не имеет права начинаться с произвольного выбора \(K\). Поэтому фиксируем минимальные требования к внутреннему пространству, необходимые для S2T-замыкания электромагнитного и фермионного секторов.

\begin{definition}[Допустимый S2T-носитель]
Компактное пространство \(K=M^3\times S^1\) называется допустимым S2T-носителем первого уровня, если выполнены условия:
\begin{enumerate}
    \item \(M^3\) компактно и ориентируемо;
    \item \(M^3\) допускает spin-структуру, необходимую для фермионов;
    \item \(\pi_1(M^3)\neq 0\), чтобы существовали нетривиальные плоские \(U(1)\)-связности;
    \item пространство плоских \(U(1)\)-связностей дискретно, то есть \(b_1(M^3)=0\);
    \item метрика принадлежит классу постоянной положительной кривизны, если не указана отдельная причина выхода из сферических пространственных форм;
    \item нормировка бозонного якобиана должна давать \(Z_A=\pi^2\) без дополнительного непрерывного множителя.
\end{enumerate}
\end{definition}

В классе сферических пространственных форм
\begin{equation}
    M^3=S^3/\Gamma,
\end{equation}
где \(\Gamma\) --- конечная группа, действующая свободно и изометрически, объём при единичном радиусе равен
\begin{equation}
    \Vol(M^3)=\frac{\Vol(S^3)}{|\Gamma|}=\frac{2\pi^2}{|\Gamma|}.
\end{equation}
Требование \(Z_A=\pi^2\) тем самым эквивалентно
\begin{equation}
    |\Gamma|=2.
\end{equation}
В линзовом подсемействе это даёт
\begin{equation}
    M^3=L(2,1)=\mathbb{RP}^3.
\end{equation}

Второй независимый тест даёт тот же результат. Для линзового пространства \(L(p,1)\)
\begin{equation}
    \pi_1(L(p,1))\cong \mathbb{Z}_p,
\end{equation}
а пространство плоских \(U(1)\)-связностей имеет вид
\begin{equation}
    \mathcal{M}_{flat}(L(p,1),U(1))
    \cong \mathrm{Hom}(\mathbb{Z}_p,U(1))
    \cong \left\{e^{2\pi i k/p}\right\}_{k=0}^{p-1}.
\end{equation}
Характерный шаг фазовой ветви равен \(2\pi/p\). Чтобы получить ровно две ветви и масштаб \(\pi\), необходимо
\begin{equation}
    p=2.
\end{equation}
Следовательно, фазово-топологический тест также выделяет \(\mathbb{RP}^3\).

\begin{theorem}[Выделенность \(\mathbb{RP}^3\times S^1\) в минимальном классе]
В классе носителей \(K=M^3\times S^1\), где \(M^3\) является сферической пространственной формой с дискретным пространством плоских \(U(1)\)-связностей, одновременные требования
\begin{equation}
    Z_A=\pi^2,
    \qquad
    \Delta\theta_{flat}=\pi,
    \qquad
    b_1(M^3)=0
\end{equation}
выделяют
\begin{equation}
    K=\mathbb{RP}^3\times S^1.
\end{equation}
\end{theorem}

\begin{proof}
Из \(Z_A=\pi^2\) в классе сферических форм следует \(|\Gamma|=2\), поскольку \(\Vol(S^3/\Gamma)=2\pi^2/|\Gamma|\). Единственная циклическая линзовая форма порядка два есть \(L(2,1)=\mathbb{RP}^3\). Независимо, условие \(\Delta\theta_{flat}=2\pi/p=\pi\) для \(L(p,1)\) также даёт \(p=2\). Условие \(b_1=0\) исключает пространства с непрерывными \(U(1)\)-модулями, такие как \(T^3\) или \(S^2\times S^1\). Поэтому в указанном минимальном классе остаётся \(\mathbb{RP}^3\times S^1\).
\end{proof}

\begin{remark}
Эта теорема не утверждает абсолютную единственность среди всех возможных компактных пространств. Она утверждает выделенность в минимальном S2T-классе: сферическая геометрия, дискретные плоские \(U(1)\)-связности, отсутствие непрерывных модулей и требуемые нормировки \(\pi^2\) и \(\pi\). Более широкие классы допустимы только как отдельные версии модели и должны проходить тот же аудит без новых ручных параметров.
\end{remark}

\begin{nogocriterion}[Провал выбора носителя]
Если будет найден другой носитель \(K'\), который удовлетворяет тем же ограничениям, даёт \(Z_A=\pi^2\), фазовый шаг \(\pi\), не содержит непрерывных \(U(1)\)-модулей и при этом меняет предсказания S2T без ухудшения blind-аудита, то выбор \(\mathbb{RP}^3\times S^1\) не является замкнутым. В этом случае второй том должен перейти от единственной модели к сравнению класса моделей.
\end{nogocriterion}

\section{Калибровочная нормировка}

Пусть \(A\) --- \(U(1)\)-связность на компактном факторе. Единица заряда фиксирует периодичность голономии
\begin{equation}
    \exp\left(i\oint_\gamma A\right)=1
\end{equation}
для тривиальной ветви и допускает нетривиальную \(\mathbb{Z}_2\)-ветвь на \(\mathbb{RP}^3\). Нормировка нулевой моды вдоль \(S^1\) выбирается так, чтобы эффективная четырёхмерная калибровочная константа не содержала свободного множителя
\begin{equation}
    \frac{g_5^2}{\Vol(S^1)}=1.
\end{equation}
Это не является дополнительной физической подгонкой; это выбор единицы заряда. Если иной выбор меняет численные результаты, он должен рассматриваться как отдельная версия модели.

\begin{nogocriterion}[Kill-shot нормировки]
Если коэффициент \(\pi^2\) в бозонном секторе требует свободного множителя, не фиксируемого единицей заряда или периодичностью \(U(1)\), то формула \(S_{vac}\) теряет статус вывода.
\end{nogocriterion}

\chapter{Спектральная геометрия вакуума}

\section{Оператор Дирака на произведении}

На римановом произведении \(\widetilde K=S^3_{R_3}\times S^1_{R_1}\) квадрат оператора Дирака раскладывается как
\begin{equation}
    \mathcal{D}_{\widetilde K}^2
    =\mathcal{D}_{S^3}^2\otimes 1 + 1\otimes \mathcal{D}_{S^1}^2,
\end{equation}
при согласованном выборе Clifford-представления.

\begin{lemma}[Спектральная форма на \(S^3\times S^1\)]
Для круглой сферы \(S^3_{R_3}\) и анти-периодической spin-структуры на \(S^1_{R_1}\) спектр квадрата оператора Дирака имеет вид
\begin{equation}
    \mu_{n,k}^2
    =\frac{(n+3/2)^2}{R_3^2}
    +\frac{(k+1/2)^2}{R_1^2},
    \qquad n\geq 0,
    \qquad k\in\mathbb{Z},
\end{equation}
с кратностью сферической компоненты
\begin{equation}
    d_n=(n+1)(n+2).
\end{equation}
\end{lemma}

\begin{proof}
Для \(S^3_R\) собственные значения оператора Дирака равны
\begin{equation}
    \lambda_n^{\pm}=\pm\frac{n+3/2}{R},
\end{equation}
а кратность положительной ветви равна \((n+1)(n+2)\). Для окружности с анти-периодическим условием \(\psi(\theta+2\pi)=-\psi(\theta)\) моды имеют импульсы \((k+1/2)/R_1\). На произведении квадрат оператора является суммой квадратов компонент, что даёт заявленную формулу.
\end{proof}

\begin{remark}
Половинный сдвиг \(1/2\) не является численной подгонкой. Он обязан быть связан со spin-структурой или физическим выбором фермионных граничных условий. Если такой вывод не проходит, все формулы, зависящие от \(k+1/2\), должны быть переведены в статус условных.
\end{remark}

\section{Spin-структуры и статус половинного сдвига}

На окружности существуют две spin-структуры. Периодическая ветвь задаёт условие
\begin{equation}
    \psi(\theta+2\pi)=\psi(\theta),
\end{equation}
и спектральные импульсы
\begin{equation}
    p_k=\frac{k}{R_1},
    \qquad k\in\mathbb{Z}.
\end{equation}
Анти-периодическая ветвь задаёт условие
\begin{equation}
    \psi(\theta+2\pi)=-\psi(\theta),
\end{equation}
и импульсы
\begin{equation}
    p_k=\frac{k+1/2}{R_1},
    \qquad k\in\mathbb{Z}.
\end{equation}

\begin{proposition}[Условная допустимость анти-периодической ветви]
Использование сдвига \(k+1/2\) допустимо в S2T-версии II.A только при явной фиксации фермионной анти-периодической spin-структуры на \(S^1\). В этом случае сдвиг является дискретным топологическим выбором, а не непрерывным параметром.
\end{proposition}

\begin{proof}
Spin-структуры на окружности классифицируются гомоморфизмами \(\pi_1(S^1)\to\mathbb{Z}_2\). Их ровно две: тривиальная и нетривиальная. Нетривиальная ветвь меняет знак спинора при обходе окружности, что эквивалентно анти-периодическому условию и спектру \((k+1/2)/R_1\). Поскольку выбор дискретен, он не создаёт новой непрерывной ручки. Однако физическая релевантность именно этой ветви должна быть обоснована секторной постановкой задачи.
\end{proof}

\begin{nogocriterion}[Провал spin-сдвига]
Если физический фермионный сектор требует периодической spin-структуры на \(S^1\), то все формулы, использующие \(k+1/2\) в фермионных, лептонных или нейтринных ветвях, должны быть пересчитаны с \(a=0\) или с иной явно выведенной spin-структурой. Электромагнитный узел \(\kappa_{Cas}\) проверяется отдельно как periodic gauge/ghost determinant и не пересчитывается из фермионного spin-сдвига.
\end{nogocriterion}

\section{Объёмы и ведущие спектральные масштабы}

Для единичных радиусов
\begin{equation}
    \Vol(S^3\times S^1)=\Vol(S^3)\Vol(S^1)=2\pi^2\cdot 2\pi=4\pi^3,
\end{equation}
а
\begin{equation}
    \Vol(\mathbb{RP}^3)=\frac{1}{2}\Vol(S^3)=\pi^2.
\end{equation}
Нетривиальный \(\mathbb{Z}_2\)-цикл в \(\mathbb{RP}^3\) имеет минимальную длину
\begin{equation}
    L_{sys}(\mathbb{RP}^3)=\pi.
\end{equation}

\begin{proposition}[Геометрический скаляр первого приближения]
При выбранных нормировках базовый геометрический скаляр электромагнитного сектора записывается как
\begin{equation}
    S_{geo}=4\pi^3+\pi^2+\pi.
\end{equation}
Здесь \(4\pi^3\) соответствует фермионному якобиану на спинорном накрытии, \(\pi^2\) --- бозонному якобиану на \(\mathbb{RP}^3\), а \(\pi\) --- минимальному топологическому циклу или эквивалентной голономной ветви.
\end{proposition}

\begin{proof}[Статус доказательства]
Первое и второе слагаемые следуют из объёмов при фиксированной нормировке мод. Третье слагаемое не выводится из размерности: оно связано с нетривиальной плоской ветвью \(U(1)\)-связности на \(\mathbb{RP}^3\). Поэтому ниже оно фиксируется отдельным топологическим утверждением, а не подгоняется численно.
\end{proof}

\begin{lemma}[Плоская \(\mathbb{Z}_2\)-голономия на \(\mathbb{RP}^3\)]
Так как \(\pi_1(\mathbb{RP}^3)\simeq\mathbb{Z}_2\), классы плоских \(U(1)\)-связностей на \(\mathbb{RP}^3\) задаются гомоморфизмами
\begin{equation}
    \mathrm{Hom}(\mathbb{Z}_2,U(1))=\{1,-1\}.
\end{equation}
Нетривиальный класс имеет голономию \(-1=e^{i\pi}\) вдоль порождающего цикла.
\end{lemma}

\begin{proof}
Плоская абелева связность с точностью до калибровочной эквивалентности определяется представлением фундаментальной группы в \(U(1)\). Образ порождающего элемента \(\mathbb{Z}_2\) должен удовлетворять \(z^2=1\), поэтому возможны только \(z=1\) и \(z=-1\). Нетривиальная ветвь даёт фазу \(\pi\) modulo \(2\pi\).
\end{proof}

\begin{proposition}[Замкнутость \(\pi\)-слагаемого]
При выбранной нормировке единицы заряда вклад \(\pi\) в \(S_{geo}\) является не свободным числом, а минимальной фазой нетривиальной плоской голономии на \(\mathbb{RP}^3\).
\end{proposition}

\begin{proof}
Предыдущее утверждение оставляет только две ветви: тривиальную фазу \(0\) и нетривиальную фазу \(\pi\). Если сектор выбирается как нетривиальная \(\mathbb{Z}_2\)-ветвь, коэффициент фиксирован топологически. Непрерывный множитель отсутствует; изменение этого слагаемого означало бы переход к другой модели нормировки, а не к параметру той же модели.
\end{proof}

\section{Коэффициенты теплового ядра}

Для оператора лапласова типа \(P\) на четырёхмерном компактном многообразии без края асимптотика теплового ядра имеет вид
\begin{equation}
    \Tr(e^{-tP})\sim (4\pi t)^{-2}\sum_{j\geq 0} a_{2j}(P)t^j.
\end{equation}
Для скалярного лапласиана ведущие геометрические коэффициенты имеют форму
\begin{equation}
    a_0=\int_K d\Vol,
    \qquad
    a_2=\frac{1}{6}\int_K \mathrm{Scal}\,d\Vol.
\end{equation}
На \(S^3_{R_3}\times S^1_{R_1}\)
\begin{equation}
    \mathrm{Scal}=\frac{6}{R_3^2},
\end{equation}
поэтому
\begin{equation}
    \int_{S^3\times S^1}\mathrm{Scal}\,d\Vol
    =\frac{6}{R_3^2}\cdot 4\pi^3 R_3^3R_1
    =24\pi^3 R_3R_1.
\end{equation}

\begin{nogocriterion}[Несогласованность коэффициентов]
Если используемые в \(S_{vac}\) коэффициенты не могут быть связаны с heat-kernel, дзета-остатками или калибровочной нормировкой, они не допускаются в финальную формулу S2T.
\end{nogocriterion}

\chapter{Аналитическое замыкание электромагнитного сектора}


\section{Дзета-регуляризация окружности}

Для окружности существуют две разные спектральные задачи, которые в версии II.A запрещено смешивать. Фермионная анти-периодическая ветвь задаёт импульсы \((k+1/2)/R_1\) и контрольный остаток
\begin{equation}
    \zeta\left(-1,\frac12\right)=\frac{1}{24}.
\end{equation}
Эта формула остаётся полезной проверкой арифметики, но не является источником электромагнитного коэффициента. В EM/QED-блоке используется spatial periodic gauge/ghost ветвь
\begin{equation}
    p_m=\frac{m}{R_1},
    \qquad m\in\mathbb{Z},
\end{equation}
согласованная с большими \(U(1)\)-калибровками и Wilson-line периодом.

\begin{lemma}[Periodic determinant-остаток]
После \(\det'\)-удаления глобальной нулевой моды \(m=0\) periodic ряд даёт
\begin{equation}
    \frac12\sum_{m\in\mathbb{Z}\setminus\{0\}}^{\zeta}|m|
    =\sum_{m=1}^{\infty,\zeta}m
    =\zeta_R(-1)
    =-\frac{1}{12}.
\end{equation}
В Maxwell--ghost determinant-нормировке константной ветви, входящей в обратный электромагнитный скаляр с вакуумным знаком, конечный коэффициент равен
\begin{equation}
    \kappa_{Cas}^{(0)}=-\frac12\zeta_R(-1)=\frac{1}{24}.
\end{equation}
\end{lemma}

\begin{proof}
Правило \(\det'\) удаляет только глобальный нулевой eigenvalue. Пары \(m\) и \(-m\) дают положительный ряд \(m=1,2,\ldots\), а дзета-регуляризация фиксирует \(\sum_{m=1}^{\infty,\zeta}m=\zeta_R(-1)=-1/12\). Дополнительный множитель \(-1/2\) не является подгонкой: это выбранная determinant-нормировка эффективной Maxwell--ghost константной ветви и знак вакуумного вычитания в обратном калибровочном скаляре. Поэтому число \(1/24\) появляется из periodic gauge/ghost постановки, а не из фермионного \(k+1/2\)-сдвига.
\end{proof}

\begin{remark}[Разделение совпадающих чисел]
Численное совпадение
\begin{equation}
    \zeta\left(-1,\frac12\right)=-\frac12\zeta_R(-1)=\frac{1}{24}
\end{equation}
не означает равенства физических секторов. Левая часть относится к анти-периодической фермионной ветви, правая --- к periodic gauge/ghost determinant-ветви после \(\det'\)-удаления нулевой моды. В дальнейшем \(\kappa_{Cas}\) связывается только с правой интерпретацией.
\end{remark}

\begin{lemma}[Четвёртый дзета-остаток]
Для анти-периодической контрольной ветви
\begin{equation}
    \zeta\left(4,\frac12\right)=15\zeta_R(4)=\frac{\pi^4}{6}.
\end{equation}
В EM-блоке этот остаток используется только как heat-kernel шаблон второго порядка, а не как источник \(\kappa_{Cas}\).
\end{lemma}

\begin{proof}
Имеем
\begin{equation}
    \sum_{k=0}^{\infty}\frac{1}{(k+1/2)^4}
    =16\sum_{k=0}^{\infty}\frac{1}{(2k+1)^4}
    =16\left(1-2^{-4}\right)\zeta_R(4)
    =15\zeta_R(4)=\frac{\pi^4}{6}.
\end{equation}
\end{proof}

\section{Проверка коэффициента \texorpdfstring{\(\kappa_{Cas}=1/24\)}{kappa Cas = 1/24}}

Коэффициент \(1/24\) является самым чувствительным узлом электромагнитной строки. Честное повышение статуса состоит не в усилении риторики, а в сужении проверяемого утверждения: нужно доказать, что periodic Maxwell--ghost/Hodge determinant на \(K=\mathbb{RP}^3\times S^1\) оставляет константную KK-ветвь с finite part \(1/24\), а massive уровни не дают независимый невычитаемый вклад того же порядка.

\begin{definition}[Рабочая схема \(\kappa_{Cas}\)]
В версии II.A коэффициент \(\kappa_{Cas}\) определяется как finite part periodic Maxwell--ghost/Hodge-комбинации при условиях:
\begin{enumerate}
    \item EM/QED использует spatial periodic gauge/ghost ветвь на \(S^1\), отделённую от фермионной anti-periodic и термальной KMS-постановок;
    \item \(\det'\) удаляет только глобальную моду \((\lambda_{\mathbb{RP}^3},m)=(0,0)\), но сохраняет ряд \((0,m\neq0)\);
    \item Maxwell- и ghost-вклады объединяются до выделения finite part;
    \item локальные heat-kernel и \(\log\mu\)-члены вычитаются минимально, без нового непрерывного параметра.
\end{enumerate}
\end{definition}

\begin{proposition}[Hodge--ghost сокращение]
Так как \(b_1(\mathbb{RP}^3)=0\), гармонические one-form моды отсутствуют. После Faddeev--Popov вычитания продольные и точные компоненты не дают самостоятельного Casimir-счёта. Релевантный одномерный остаток \(\kappa_{Cas}^{(0)}\) задаётся periodic KK-рядом над константной скалярной модой \(\lambda_{\mathbb{RP}^3}=0\), \(m\neq0\).
\end{proposition}

\begin{proof}[Статус доказательства]
Это determinant-уровень закрытия, а не полная спектральная теорема для всего Maxwell-оператора. Он использует стандартную структуру gauge-fixed abelian determinant: точные векторные компоненты компенсируются ghost-сектором, а отсутствие \(b_1\) исключает непрерывные гармонические one-form модули. Следовательно, источник \(1/24\) локализован в константной periodic ветви. Оставшийся риск честно выделен отдельно: massive transverse/coexact спектр \(\lambda_{\mathbb{RP}^3}>0\).
\end{proof}

\begin{theorem}[Условное determinant-замыкание \(\kappa_{Cas}\)]
При рабочей periodic gauge/ghost постановке, \(\det'\)-удалении только глобальной нулевой моды и минимальном локальном вычитании константная Maxwell--ghost ветвь даёт
\begin{equation}
    \kappa_{Cas}^{(0)}=\frac{1}{24}.
\end{equation}
Следовательно, узел \(1/24\) закрыт на determinant-уровне II.A и больше не является переносом из фермионной spin-структуры. Полный EM-статус требует coexact Bessel-tail no-go.
\end{theorem}

\begin{proof}
По предыдущему предложению релевантный ряд --- periodic \(m\in\mathbb{Z}\setminus\{0\}\) над \(\lambda_{\mathbb{RP}^3}=0\). По лемме о periodic determinant-остатке его Maxwell--ghost finite part равна \(-\zeta_R(-1)/2=1/24\). Все степенные и логарифмические члены относятся к локальным heat-kernel counterterms и удаляются схемой, заданной до сравнения с \(\alpha^{-1}\). Непрерывный коэффициент не вводится.
\end{proof}

\begin{protocol}[Остаточный massive-аудит]
Для честного повышения научной шкалы следующим шагом требуется:
\begin{enumerate}
    \item выписать transverse и coexact спектры one-form на \(\mathbb{RP}^3\);
    \item просуммировать их periodic KK-башни в той же \(\det'\)-схеме;
    \item доказать, что \(\lambda_{\mathbb{RP}^3}>0\) либо сокращается в Maxwell--ghost/Hodge-комбинации, либо является локально вычитаемым heat-kernel вкладом;
    \item проверить, что альтернативная gauge-invariant \(\det'\)-схема не меняет finite part \(1/24\).
\end{enumerate}
\end{protocol}


\begin{proposition}[Massive-аудит I: неустранимый Bessel-tail]
Для отдельного massive уровня \(\lambda>0\) periodic KK-башня на \(S^1\) имеет спектральную сумму вида
\begin{equation}
    E_{\lambda}(s)=\frac12\sum_{m\in\mathbb{Z}}\left(m^2+\rho^2\right)^{1/2-s},
    \qquad
    \rho=R_1\sqrt{\lambda}.
\end{equation}
После вычитания локальных Weyl/heat-kernel членов остаётся не локальная конечная часть
\begin{equation}
    E_{\lambda}^{nonloc}\propto
    -\rho\sum_{q=1}^{\infty}\frac{K_1(2\pi q\rho)}{q},
\end{equation}
где \(K_1\) --- функция Макдональда. Этот остаток экспоненциально мал при больших \(\rho\), но не является локальным counterterm и потому не может быть просто отброшен в честном massive-аудите.
\end{proposition}

\begin{proof}[Статус доказательства]
Формула является стандартной формой пуассонова преобразования одномерной Epstein-zeta суммы. Член с нулевым dual winding даёт локальную heat-kernel часть: степенные и логарифмические вклады, зависящие от cutoff или \(\mu\). Ненулевые winding-числа \(q\neq0\) дают Bessel-сумму. Она зависит от глобальной длины окружности и массы \(\rho\), поэтому является не локальной конечной частью. Следовательно, для каждого massive уровня \(\lambda>0\) нужно показать либо точное Maxwell--ghost/Hodge-сокращение соответствующих multiplicity-weighted Bessel-tail, либо их отсутствие по симметрии. Одного заявления ``локально вычитается'' недостаточно.
\end{proof}

\begin{corollary}[Вердикт massive-аудита I]
Первый massive-аудит не повышает шкалу \(R_{sci}\) с \(4/10\) до \(5/10\). Он уточняет препятствие: для перехода \(4\to5\) нужно доказать тождество
\begin{equation}
    \sum_{\lambda>0}
    \left(d^{T}_{\lambda}-d^{G}_{\lambda}-d^{exact}_{\lambda}\right)
    \rho_{\lambda}
    \sum_{q=1}^{\infty}\frac{K_1(2\pi q\rho_{\lambda})}{q}=0,
\end{equation}
или вывести, что вся левая часть поглощается уже зафиксированной gauge-invariant нормировкой без нового параметра. Пока это тождество не доказано, \(\kappa_{Cas}=1/24\) остаётся условным determinant-результатом, а не полным EM-замыканием.
\end{corollary}


\begin{proposition}[Massive-аудит II: Hodge-сокращение не убирает coexact хвост автоматически]
В gauge-fixed абелевом determinant точные one-form моды и scalar ghost-моды имеют одинаковые положительные скалярные eigenvalues и потому сокращаются попарно в nonlocal Bessel-tail. Однако coexact/transverse one-form моды не являются образами скалярного градиента. Их multiplicity-weighted Bessel-tail имеет вид
\begin{equation}
    E_{coex}^{nonloc}\propto
    -\sum_{\lambda\in\Spec(\Delta_1^{coex})}
    d^{coex}_{\lambda}\rho_{\lambda}
    \sum_{q=1}^{\infty}\frac{K_1(2\pi q\rho_{\lambda})}{q}.
\end{equation}
Без дополнительного топологического, supersymmetric или duality-сокращения этот хвост не обязан обращаться в ноль.
\end{proposition}

\begin{proof}[Статус доказательства]
Hodge-разложение даёт ортогональную сумму exact, coexact и harmonic one-form компонент. На \(\mathbb{RP}^3\) гармонические one-form отсутствуют, поскольку \(b_1=0\). Exact-компоненты имеют вид \(d\phi\) и наследуют положительный scalar spectrum, поэтому именно они являются естественной парой для Faddeev--Popov ghost determinant. Coexact-компоненты удовлетворяют \(\delta A=0\) и не представлены как \(d\phi\). Следовательно, ghost-сектор не даёт для них автоматической попарной отмены. Их nonlocal Bessel-tail может исчезнуть только при отдельном спектральном тождестве multiplicities или при дополнительной физической симметрии, которая в версии II.A пока не выведена.
\end{proof}

\begin{corollary}[Вердикт massive-аудита II]
Честная проверка показывает, что переход \(R_{sci}:4/10\to5/10\) пока заблокирован. Более того, задача стала строже: нужно не просто ``посчитать massive residual'', а доказать отсутствие coexact transverse Bessel-tail или вывести новый параметрически-безопасный механизм его компенсации. До этого \(S_{vac}\) нельзя переводить из условного determinant-успеха в полный электромагнитный успех.
\end{corollary}


\begin{theorem}[Massive-аудит III: positivity no-go для обычного Maxwell-хвоста]
Если после ghost-сокращения остаётся хотя бы одна coexact eigenmode с \(\lambda>0\) и положительной кратностью \(d^{coex}_{\lambda}>0\), то nonlocal Bessel-tail обычного Maxwell-сектора не может обратиться в ноль сам по себе:
\begin{equation}
    \sum_{\lambda\in\Spec(\Delta_1^{coex})}
    d^{coex}_{\lambda}\rho_{\lambda}
    \sum_{q=1}^{\infty}\frac{K_1(2\pi q\rho_{\lambda})}{q}>0.
\end{equation}
Следовательно, полное исчезновение massive-хвоста требует дополнительного механизма: новой paired-секторной симметрии, duality-сокращения, топологической проекции, меняющей coexact multiplicities, или явного пересмотра формулы \(S_{vac}\).
\end{theorem}

\begin{proof}
Для \(x>0\) функция Макдональда \(K_1(x)\) положительна. Для каждого massive уровня \(\rho_{\lambda}>0\), а кратность \(d^{coex}_{\lambda}\) неотрицательна и положительна при наличии соответствующей моды. Поэтому каждый член суммы имеет один знак. Внутри одного обычного coexact Maxwell-сектора нет отрицательных multiplicities, которые могли бы компенсировать этот хвост. Ghost-сектор уже использован для exact/scalar cancellation и не создаёт paired coexact contribution с противоположным знаком. Значит, cancellation невозможна без дополнительной структуры, отсутствующей в текущей версии II.A.
\end{proof}

\begin{corollary}[Вердикт massive-аудита III]
Путь \(4\to5\) через простое ``досчитать massive residual'' закрыт отрицательно: ordinary Maxwell--ghost/Hodge-комбинация оставляет coexact Bessel-tail, если не введён новый доказанный механизм компенсации. Поэтому научная шкала не повышается; напротив, для честности статус \(S_{vac}\) должен оставаться условным determinant-результатом, а следующий исследовательский шаг меняется с вычислительного на структурный: найти paired-сектор или признать поправку к \(S_{vac}\).
\end{corollary}


\begin{theorem}[Massive-аудит IV: полная бесконечная башня]
Полная coexact Maxwell-башня на \(\mathbb{RP}^3\times S^1\) имеет двойную спектральную структуру
\begin{equation}
    \mathcal{E}_{tower}(s)=
    \frac12\sum_{\lambda\in\Spec(\Delta_1^{coex})}
    d^{coex}_{\lambda}
    \sum_{m\in\mathbb{Z}}
    \left(\lambda+\frac{m^2}{R_1^2}\right)^{1/2-s}.
\end{equation}
Её нельзя заменить одним коэффициентом \(1/24\). После регуляризации башня распадается на две части:
\begin{equation}
    \mathcal{E}_{tower}^{ren}
    =\mathcal{E}_{local}^{HK,ren}
    +\mathcal{E}_{nonlocal}^{Bessel}.
\end{equation}
Локальная часть содержит бесконечные heat-kernel counterterms и зависит от выбранной схемы перенормировки, а nonlocal Bessel-часть является конечной глобальной поправкой, которую нельзя удалить локальным counterterm.
\end{theorem}

\begin{proof}[Статус доказательства]
Суммирование по \(m\) при фиксированной \(\lambda\) даёт Epstein-zeta структуру. Пуассоново преобразование отделяет zero-winding вклад от winding-вкладов. Zero-winding вклад является локальным: он выражается через heat-kernel коэффициенты оператора на \(\mathbb{RP}^3\), содержит степенные и логарифмические расходимости и должен быть поглощён перенормировкой локальных операторов. Ненулевые winding-вклады дают Bessel-суммы \(K_1(2\pi qR_1\sqrt{\lambda})\). Они конечны, зависят от глобальной окружности и не являются локальными. Поэтому ``все бесконечности'' башни не решают проблему: бесконечная часть перенормируема, но конечная nonlocal часть остаётся как физический residual, если нет paired-сокращения.
\end{proof}

\begin{corollary}[Вердикт по бесконечным поправкам]
Да, в текущем EM-блоке учтена не вся бесконечная башня. Константная ветвь \(\lambda=0\), \(m\neq0\) даёт \(1/24\), но это только один сектор полной determinant-задачи. Полная башня требует либо:
\begin{enumerate}
    \item явного paired-механизма, зануляющего \(\mathcal{E}_{nonlocal}^{Bessel}\);
    \item включения новой поправки \(\Delta_{tower}\) в \(S_{vac}\);
    \item признания, что текущая формула \(S_{vac}\) является truncation-схемой, а не полным EM-выводом.
\end{enumerate}
Без одного из этих трёх шагов шкала \(R_{sci}\) не может быть повышена выше \(4/10\).
\end{corollary}

\begin{definition}[Минимальная поправка полной башни]
Если paired-сокращение не найдено, честная форма электромагнитного скаляра должна быть записана как
\begin{equation}
    S_{vac}^{full}=S_{geo}-\frac{1}{24S_{geo}}-\frac{1}{\pi^4S_{geo}^2}+\Delta_{tower}^{coex},
\end{equation}
где \(\Delta_{tower}^{coex}\) является конечной nonlocal Bessel-поправкой полной coexact башни после локальной heat-kernel перенормировки. В версии II.A эта поправка не вычислена и не имеет права молча полагаться равной нулю.
\end{definition}

\begin{protocol}[Первый численный просчёт \texorpdfstring{\(\mathbb{RP}^3\)}{RP3} coexact-башни]
Для проверки масштаба хвоста обновлён воспроизводимый скрипт
\texttt{s2t\_coexact\_tower\_audit.py}. Он считает положительную безразмерную Bessel-сумму
\begin{equation}
    T_{coex}=\sum_{\lambda\in\Spec(\Delta_1^{coex})}
    d_{\lambda}^{coex}\rho_{\lambda}
    \sum_{q=1}^{\infty}\frac{K_1(2\pi q\rho_{\lambda})}{q},
    \qquad
    \rho_{\lambda}=R_1\sqrt{\lambda},
\end{equation}
в единичной радиусной нормировке. На \(\widetilde K=S^3\times S^1\) используется coexact one-form spectrum
\begin{equation}
    \lambda_n=(n+1)^2,
    \qquad
    d_n^{S^3}=2n(n+2),
    \qquad n\geq1.
\end{equation}
Для \(\mathbb{RP}^3=L(2,1)\) quotient-генератор действует как левое умножение на центральный элемент \(-1\in SU(2)\). Две transverse one-form серии имеют типы
\begin{equation}
    \left(\frac{n+1}{2},\frac{n-1}{2}\right),
    \qquad
    \left(\frac{n-1}{2},\frac{n+1}{2}\right),
\end{equation}
поэтому parity равна \((-1)^{n-1}\). Следовательно, строгая \(\mathbb{RP}^3\)-проекция имеет кратности
\begin{equation}
    d_n^{\mathbb{RP}^3}=
    \begin{cases}
        2n(n+2), & n \text{ нечётно},\\
        0, & n \text{ чётно}.
    \end{cases}
\end{equation}
При обрезании \(n\leq80\), адаптивном суммировании по winding-числу \(q\) до \(10^{-18}\) и \(R_1/R_3=1\) получено
\begin{equation}
    T_{coex}^{S^3}=1.5319168084\cdot10^{-5},
    \qquad
    T_{coex}^{\mathbb{RP}^3}=1.5227161455\cdot10^{-5}.
\end{equation}
Отношение равно
\begin{equation}
    \frac{T_{coex}^{\mathbb{RP}^3}}{T_{coex}^{S^3}}=0.9939940192.
\end{equation}
Это не противоречит quotient-интуиции: асимптотически проекция убирает половину уровней, но Bessel-tail почти целиком контролируется первым уровнем \(n=1\), который является \(\mathbb{RP}^3\)-even и сохраняется с полной кратностью. Первый surviving уровень даёт \(99.9976\%\) полной \(\mathbb{RP}^3\)-суммы.
\end{protocol}

\begin{remark}[Статус результата башни]
Этот расчёт закрывает важный вычислительный gap: \(\mathbb{RP}^3\)-projection уже не является грубой половинной эвристикой. Однако ordinary coexact-сектор даёт положительный ненулевой хвост. Значит, \(\Delta_{tower}^{coex}\) действительно является вычислимой поправкой, а не автоматически исчезающим членом. Для превращения её в окончательную добавку к \(S_{vac}\) остаются три элемента: общий радиусный фактор \(R_1/R_3\), знак Maxwell determinant-вклада после всех ghost/exact сокращений и нормировочный коэффициент, связывающий \(T_{coex}\) с безразмерным \(S_{vac}\). Если этот коэффициент порядка \(1/S_{geo}\), то масштаб поправки составляет примерно \(10^{-7}\), то есть она больше текущей ошибки совпадения с \(\alpha^{-1}\). Поэтому результат является не косметической поправкой, а настоящей проверкой EM-замыкания.
\end{remark}
\begin{remark}[Внешнее различение Casimir energy и Euclidean determinant]
Независимый аудит показал, что число
\(
T_{coex}^{\mathbb{RP}^3}=1.5227161455\cdot10^{-5}
\)
действительно воспроизводится из untwisted coexact spectrum. Однако необходимо уточнить тип вычисляемого функционала. Исходная сумма
\[
E_\lambda(s)=\frac12\sum_{m\in\mathbb Z}(m^2+\rho^2)^{1/2-s}
\]
является Casimir-energy zeta function. Для неё nonzero-winding часть равна
\[
E_\lambda^{nonloc}
=-\frac{\rho}{\pi}\sum_{q\geq1}\frac{K_1(2\pi q\rho)}{q},
\]
так что физически нормированный остаток всей coexact-башни составляет
\[
E_{coex}^{nonloc}=-\frac{T_{coex}^{\mathbb{RP}^3}}{\pi}
=-4.8469560297\cdot10^{-6}.
\]

Напротив, конечная winding-часть Euclidean determinant оператора
\(-\partial_\tau^2+\rho^2\) на окружности следует из тождества
\[
\det(-\partial_\tau^2+\rho^2)=4\sinh^2(\pi\rho)
\]
и после отделения локального члена имеет вид
\[
(\log\det)_\lambda^{winding}=2\log(1-e^{-2\pi\rho}).
\]
Для полной untwisted coexact-башни получено
\[
(\log\det)_{coex}^{winding}=-4.1848910943\cdot10^{-5},
\qquad
\Gamma_{coex}^{winding}=-2.0924455472\cdot10^{-5}.
\]
Следовательно, положительная \(K_1\)-сумма является корректным диагностическим ядром Casimir energy, но не может без дополнительного вывода отождествляться с finite residue четырёхмерного log determinant. Для связи с \(S_{vac}\) необходимо заранее выбрать первичный функционал и вывести его знак, радиусную размерность и общий Maxwell--FP prefactor.
\end{remark}
\begin{protocol}[Зацепки перед изменением формулы]
Перед тем как менять формулу для \(S_{vac}\), нужно проверить, не содержит уже созданная база прямой путь вывода. Аудит wiki и строгих RPFT-файлов выделяет следующие гипотезы.
\begin{enumerate}
    \item \textbf{Готовый paired-сектор из старой proof-chain.} В wiki-цепочке встречается формулировка Maxwell--ghost--Dirac для источника \(1/24\), тогда как текущий том сознательно отделил EM periodic Maxwell--ghost-ветвь от фермионной spin-структуры. Это главная зацепка: возможно, старый Dirac/fermionic companion был не ошибкой, а недооформленным paired-механизмом. Статус: развивать, но только если спектр даёт тот же Bessel-kernel и противоположный знак без подгонки.
    \item \textbf{Поглощение хвоста в уже существующий \(\pi^{-4}S_{geo}^{-2}\)-член.} Численно \(T_{coex}^{\mathbb{RP}^3}/S_{geo}=1.11\cdot10^{-7}\), а второй остаток равен \(1/(\pi^4S_{geo}^2)=5.47\cdot10^{-7}\). Хвост составляет около \(20.3\%\) этого члена. Статус: сильная зацепка, но не доказательство; нужно вывести, является ли \(\pi^{-4}\)-остаток effective summary полной башни, а не независимой поправкой.
    \item \textbf{Радиусное подавление.} Если \(r=R_1/R_3\) свободен, Bessel-tail падает как \(e^{-4\pi r}\). Чтобы \(T_{coex}/S_{geo}\) стал меньше текущей ошибки \(3.5\cdot10^{-9}\), нужно примерно \(r\geq1.285\). Статус: пока отвергнуть как прямое спасение II.A, потому что базовая радиусная нормировка и radius-stabilization держат \(r=1\); можно развивать только при независимом выводе \(R_1/R_3\neq1\).
    \item \textbf{Мера гармонической моды и \(\det\prime\).} Файлы QED/gauge-fixing уже содержат конечномерный интеграл по \(S^1\)-голономии и деление на остаточную gauge-группу. Это может убрать \(\log\mu\) и нулевые моды, но не nonlocal Bessel-tail положительных coexact уровней. Статус: полезно для нормировки и знака, но не как cancellation-механизм.
    \item \textbf{Hodge/duality-сокращение.} Exact one-form modes сокращаются ghost-сектором; coexact modes не являются градиентами скаляров. Hodge duality переводит их в exact two-form sector, но не создаёт отрицательную multiplicity в ordinary Maxwell determinant. Статус: отвергнуть как автоматическое сокращение; развивать только при появлении отдельного two-form/topological paired-поля.
    \item \textbf{Торсионная \(\mathbb Z_2\)-ветвь.} Нетривиальная flat-ветвь уже даёт \(\pi\)-член. Возможна гипотеза, что torsion refinement меняет parity coexact-башни. Текущий расчёт показывает, что surviving \(n=1\) уровень остаётся с полной кратностью. Статус: слабая зацепка; простая torsion-проекция не зануляет хвост.
    \item \textbf{Локальная heat-kernel перенормировка.} Она обязательна и уже согласуется с отсутствием физической \(\log\mu\)-неоднозначности для конформных секторов, но finite winding part не локален. Статус: отвергнуть как способ удалить \(\Delta_{tower}^{coex}\).
    \item \textbf{Прямой determinant-путь.} Самый чистый путь сейчас --- не менять \(S_{vac}\), а вывести полный объект \(\frac12\log\det\prime\Delta_{1,coex}-\frac12\log\det\prime\Delta_0\) на \(\mathbb{RP}^3\times S^1\), включая знак, \(R_1/R_3\), локальные counterterms и конечную Bessel-часть. Статус: основной рабочий маршрут.
\end{enumerate}
\end{protocol}

\begin{remark}[Текущий исследовательский вердикт]
Наиболее перспективны две линии: paired-сектор, возможно сохранённый в старой Maxwell--ghost--Dirac формулировке wiki, и переинтерпретация \(\pi^{-4}S_{geo}^{-2}\)-члена как уже свернутого следа полной coexact-башни. Наименее перспективны простое Hodge-сокращение, локальное вычитание и грубая \(\mathbb Z_2\)-проекция: они уже не убирают положительный \(n=1\) хвост. Поэтому до изменения численной формулы следует сначала попытаться вывести нормированный determinant напрямую из существующего QED/gauge-fixing блока и проверить, не совпадает ли его finite part с уже введённым вторым остатком.
\end{remark}

\begin{protocol}[Проверка гипотезы \texorpdfstring{\(\pi^{-4}\)}{pi^-4} как следа coexact-башни]
Самая перспективная сопутствующая идея состоит в том, что уже имеющийся член \(1/(\pi^4S_{geo}^2)\) может быть не независимой второй поправкой, а свернутым finite residue полной coexact-башни. Для проверки добавлен скрипт \texttt{s2t\_pi4\_tower\_hypothesis\_audit.py}.

Положительная \(\mathbb{RP}^3\)-сумма даёт
\begin{equation}
    \frac{T_{coex}^{\mathbb{RP}^3}}{S_{geo}}=1.1111771870\cdot10^{-7},
    \qquad
    \frac{1}{\pi^4S_{geo}^2}=5.4667502954\cdot10^{-7}.
\end{equation}
Если сравнивать напрямую, требуемый коэффициент равен
\begin{equation}
    C_{direct}=4.9197826945.
\end{equation}
Это близко к \(\pi^2/2=4.9348022005\). В determinant-языке это особенно интересно, потому что бозонный вклад несёт стандартный множитель \(1/2\) перед \(\log\det\), а \(\pi^2\) является \(\Vol(\mathbb{RP}^3)\). Поэтому естественная гипотеза имеет вид
\begin{equation}
    \frac{1}{\pi^4S_{geo}^2}\stackrel{?}{=}
    \frac{\Vol(\mathbb{RP}^3)}{2}
    \frac{T_{coex}^{\mathbb{RP}^3}}{S_{geo}}
    =\frac{\pi^2}{2}
    \frac{T_{coex}^{\mathbb{RP}^3}}{S_{geo}}.
\end{equation}
Правая часть равна
\begin{equation}
    \frac{\pi^2}{2}\frac{T_{coex}^{\mathbb{RP}^3}}{S_{geo}}
    =5.4834396278\cdot10^{-7}.
\end{equation}
Она превышает текущий \(\pi^{-4}\)-член на \(0.3053\%\), то есть на
\begin{equation}
    1.6689332432\cdot10^{-9}.
\end{equation}
Это число сравнимо с текущим остатком совпадения \(S_{vac}-\alpha^{-1}\), поэтому гипотеза не является косметической: она может объяснить происхождение второго остатка, но одновременно требует точного вывода нормировки.
\end{protocol}

\begin{remark}[Вердикт по \texorpdfstring{\(\pi^{-4}\)}{pi^-4}-гипотезе]
Гипотеза \(\pi^{-4}\)-остатка как volume-weighted half-determinant следа coexact-башни является самой сильной найденной зацепкой. Она имеет естественные множители \(1/2\) и \(\Vol(\mathbb{RP}^3)=\pi^2\), а численное расхождение составляет только \(0.3053\%\). Но это пока не доказательство: точная Bessel-сумма не равна алгебраическому \(1/\pi^4\) тождественно. Для закрытия нужно вывести, является ли малый остаток допустимой finite local-scheme частью, вкладом ещё одного paired-сектора или признаком того, что \(\Delta_{tower}^{coex}\) всё-таки должна быть отдельным членом.
\end{remark}

\begin{protocol}[Аудит остатка \(0.3053\%\): Casimir-mixing]
Следующий сопутствующий вопрос: можно ли объяснить малое расхождение между volume-weighted tower-кандидатом и текущим \(\pi^{-4}\)-членом без подгонки. Для этого добавлен скрипт \texttt{s2t\_pi4\_residual\_audit.py}.

Требуемый множитель равен
\begin{equation}
    M_{need}=\frac{1/(\pi^4S_{geo}^2)}{(\pi^2/2)T_{coex}^{\mathbb{RP}^3}/S_{geo}}
    =0.9969564117,
\end{equation}
то есть
\begin{equation}
    \varepsilon_{need}=1-M_{need}=0.0030435883.
\end{equation}
Наиболее сильный структурный кандидат, найденный аудитом, имеет вид
\begin{equation}
    \varepsilon_{CasMix}=\frac{10}{24S_{geo}}=0.0030405568.
\end{equation}
Тогда
\begin{equation}
    \frac{\pi^2}{2}\frac{T_{coex}^{\mathbb{RP}^3}}{S_{geo}}
    \left(1-\frac{10}{24S_{geo}}\right)
    =5.4667669181\cdot10^{-7},
\end{equation}
а целевой член равен
\begin{equation}
    \frac{1}{\pi^4S_{geo}^2}=5.4667502954\cdot10^{-7}.
\end{equation}
Оставшееся относительное расхождение составляет только
\begin{equation}
    3.04\cdot10^{-6}
\end{equation}
от самого \(\pi^{-4}\)-члена. Это существенно сильнее, чем исходное расхождение \(0.3053\%\).
\end{protocol}

\begin{protocol}[Спектральный кандидат для числа \(10\)]
После аудита остатка отдельно проверен вопрос, является ли число \(10\) голой подгонкой. Для этого добавлен скрипт \texttt{s2t\_integer10\_origin\_audit.py}. На \(\mathbb{RP}^3\) скалярные гармоники, наследуемые с \(S^3\), имеют чётные \(\ell\). Их кратности равны
\begin{equation}
    d_\ell=(\ell+1)^2,
    \qquad \ell=0,2,4,\ldots.
\end{equation}
Первые две скалярные оболочки дают
\begin{equation}
    d_0=1,
    \qquad
    d_2=9,
    \qquad
    d_0+d_2=10.
\end{equation}
Это существенно, потому что periodic Casimir-ветвь \(1/24\) возникает из скалярной константной ветви Maxwell--ghost determinant, а exact one-form modes наследуют скалярный спектр. Поэтому множитель
\begin{equation}
    \frac{10}{24S_{geo}}
    =\frac{d_0+d_2}{24S_{geo}}
\end{equation}
имеет возможное спектральное происхождение как mixing константной scalar/ghost ветви с первой ненулевой even scalar/exact оболочкой.
\end{protocol}

\begin{remark}[Статус числа \(10\)]
Фактор \(10/(24S_{geo})\) выглядит как смешение coexact-башни с уже выведенной Casimir-ветвью \(1/(24S_{geo})\). Однако число \(10\) пока не выведено. Его нельзя объявлять успехом, пока оно не будет получено как кратность, индекс, число поперечных/торсионных каналов или коэффициент второго порядка determinant-разложения. Если число \(10\) не выводится, то Casimir-mixing остаётся numerological compression, даже при чрезвычайно хорошем численном совпадении.
\end{remark}

\begin{protocol}[Determinant trace-rank test для Casimir-mixing]
Следующая проверка отделяет целочисленную структуру от постфактум выбора. Для этого добавлен скрипт \texttt{s2t\_determinant\_casmix\_audit.py}. Тестируемая форма второго порядка записывается как
\begin{equation}
    \frac{1}{\pi^4S_{geo}^2}
    \stackrel{?}{=}
    \frac{\pi^2}{2}\frac{T_{coex}^{\mathbb{RP}^3}}{S_{geo}}
    \left(1-\frac{N}{24S_{geo}}\right).
\end{equation}
Если эта форма верна, требуемый эффективный trace-rank равен
\begin{equation}
    N_{need}=24S_{geo}
    \left(1-\frac{1/(\pi^4S_{geo}^2)}{(\pi^2/2)T_{coex}^{\mathbb{RP}^3}/S_{geo}}\right)
    =10.0099700224.
\end{equation}
Ближайшее целое число равно
\begin{equation}
    N=10,
\end{equation}
и оно совпадает с найденной выше скалярно-призрачной кратностью
\begin{equation}
    d_0+d_2=1+9=10.
\end{equation}
При этом относительная ошибка восстановленного \(\pi^{-4}\)-члена составляет
\begin{equation}
    3.04\cdot10^{-6}.
\end{equation}
Иными словами, если смешанный член действительно имеет стандартный знак второго порядка
\begin{equation}
    \log\det(A+xB)=\log\det A+x\Tr(A^{-1}B)
    -\frac{x^2}{2}\Tr\left((A^{-1}B)^2\right)+\cdots,
\end{equation}
то требуемый trace-rank практически целочислен и указывает на \(\ell=0,2\) scalar/exact сектор. Недостающий шаг остаётся строгим: нужно вычислить калибровочно фиксированный смешанный след и показать, что он действительно равен этой кратности с правильным знаком, а не только численно имитирует её.
\end{protocol}



\begin{protocol}[Absorption-gate после внешней сверки]
Для проверки варианта A добавлен аудит \texttt{s2t\_lens\_pform\_absorption\_audit.py}. Он сопоставляет три слоя: стандартную lens-space рамку для \(L(2,1)=\mathbb{RP}^3\), внутреннюю coexact Bessel-башню и конечный projector-кандидат \(P_{0,2}\). Итоговая проверяемая формула имеет вид
\begin{equation}
    \Delta_{abs}^{P_{0,2}}
    =\frac{\pi^2}{2}\frac{T_{coex}^{\mathbb{RP}^3}}{S_{geo}}
    \left(1-\frac{10}{24S_{geo}}\right).
\end{equation}
Аудит даёт
\begin{align}
    T_{coex}^{\mathbb{RP}^3}&=1.5227161455\cdot10^{-5},\\
    N_{need}&=10.0099700224,\\
    N_{P_{0,2}}&=\dim\operatorname{Sym}^2(\mathbb{R}^4)=10,\\
    \frac{\Delta_{abs}^{P_{0,2}}-1/(\pi^4S_{geo}^2)}{1/(\pi^4S_{geo}^2)}&=3.04\cdot10^{-6}.
\end{align}
Следовательно, \(P_{0,2}\)-absorption является лучшим найденным вариантом: она объясняет, почему независимое добавление \(\Delta_{tower}^{coex}\) запрещено, и почему существующий \(\pi^{-4}\)-член имеет правильный масштаб.

Однако gate не закрыт как теорема. Точное равенство требует нецелого эффективного ранга \(10.0099700224\), а внешний literature-pass по Ikeda--Yamamoto, Lauret--Miatello--Rossetti, Nash--O'Connor и Schwarz/Ray--Singer подтверждает только правильность lens-space determinant рамки, но не выводит S2T-смешанный след. Поэтому версия II.A получает не переход \(4\to5\), а более узкий статус:
\begin{equation}
    \text{strong absorption candidate, not mature EM theorem.}
\end{equation}
\end{protocol}

\begin{nogocriterion}[Провал absorption-gate]
Если полный Maxwell--ghost determinant на \(L(2,1)\times S^1\) требует включить generic scalar/exact tower \(\ell=0,2,4,\ldots\) вместо first-strain projector \(P_{0,2}\), либо если finite determinant-нормировка даёт знак или объёмный множитель, отличный от \(\pi^2/2\), то \(\pi^{-4}\)-член должен быть понижен до phenomenological compression. В этом случае \(S_{vac}\) остаётся условным determinant-успехом, а переход \(R_{sci}:4/10\to5/10\) запрещён.

\end{nogocriterion}

\begin{protocol}[Closure-matrix для mixed-trace]
После absorption-gate добавлена матрица решений \texttt{s2t\_mixed\_trace\_closure\_matrix.py}. Она отделяет четыре разных утверждения:
\begin{enumerate}
    \item lens-space \(p\)-form framework корректен для \(L(2,1)=\mathbb{RP}^3\);
    \item ordinary Hodge/ghost cancellation не убирает coexact tower;
    \item \(\pi^2/2\)-нормировка и \(P_{0,2}\)-rank дают сильное численное попадание;
    \item mature theorem требует явного gauge-fixed mixed-trace operator calculation.
\end{enumerate}
Первый настоящий blocking gap теперь формулируется как
\begin{equation}
    C_6=\text{compute Maxwell--ghost mixed trace on }L(2,1)\times S^1.
\end{equation}
Это означает: дальнейшая работа по варианту A должна быть не новым поиском чисел, а выводом оператора, его знака, \(\det'\)-нулевых мод, объёмного множителя \(\pi^2/2\) и конечного ранга \(P_{0,2}\). Если следующий расчёт снова только улучшает совпадение без оператора, он не повышает статус \(S_{vac}\).
\end{protocol}

\begin{nogocriterion}[Запрет повторного численного усиления]
Любое новое совпадение вокруг \(N\approx10\), \(\pi^2/2\) или \(3.04\cdot10^{-6}\), не сопровождаемое явным Maxwell--ghost mixed-trace выводом на \(L(2,1)\times S^1\), считается диагностикой, а не доказательством. Для перехода \(4\to5\) требуется закрыть именно условие \(C_6\), иначе включается безопасный вариант D: \(\pi^{-4}\) трактуется как phenomenological compression внутри условного determinant-успеха.

\end{nogocriterion}

\begin{protocol}[C6 operator-trace skeleton]
Следующий аудит \texttt{s2t\_c6\_operator\_trace\_skeleton.py} переводит \(C_6\) из общего требования в список конкретных проверок. Требуемая absorption-идентичность должна возникать из четырёх операторных блоков:
\begin{enumerate}
    \item coexact transverse Maxwell determinant даёт ненулевую башню \(T_{coex}^{\mathbb{RP}^3}\);
    \item exact/longitudinal Maxwell branch и Faddeev--Popov scalar ghost задают обычное Hodge-сокращение и \(\det'\)-правило;
    \item periodic scalar/ghost ветвь даёт уже выделенный Casimir-фактор \(1/(24S_{geo})\);
    \item first ambient strain insertion даёт \(P_{0,2}\) с \(\Tr P_{0,2}=10\).
\end{enumerate}
В этой декомпозиции условно закрыты две части: \(P_{0,2}\)-rank и запрет \(\ell\ge4\) внутри first-strain selection. До volume-аудита открытыми оставались три теста:
\begin{enumerate}
    \item вывести знак второй вариации полного Maxwell--ghost functional, а не только знак формального \(\log\det\);
    \item вывести объёмную нормировку \(\Vol(\mathbb{RP}^3)/2=\pi^2/2\) из нормированных мод и меры;
    \item объяснить малый gap
    \begin{equation}
        N_{need}-\Tr P_{0,2}=0.0099700224,
    \end{equation}
    как finite scheme residue, subleading determinant term или основание для downgrade.
\end{enumerate}
Следовательно, дальнейшее доказательство теперь имеет узкий маршрут: не искать новый integer, а вывести второй variation trace
\begin{equation}
    \delta_g^2\Gamma_{Maxwell+ghost}^{(1)}[K]\Big|_{P_{0,2}}
    \longrightarrow
    \frac{\pi^2}{2}\frac{T_{coex}^{\mathbb{RP}^3}}{S_{geo}}
    \left(1-\frac{\Tr P_{0,2}}{24S_{geo}}\right).
\end{equation}
\end{protocol}

\begin{nogocriterion}[Провал C6 skeleton]
Если знак, объёмная нормировка или gap \(N_{need}-10\) не выводятся из одного и того же gauge-fixed determinant functional, то C6 считается незакрытым независимо от качества численного совпадения. В таком случае \(P_{0,2}\) остаётся лучшей внутренней структурой \(\pi^{-4}\)-остатка, но не переводит \(S_{vac}\) в mature theorem.
\end{nogocriterion}


\begin{protocol}[P02 volume-normalization audit]
Следующий локальный тест \texttt{s2t\_p02\_volume\_normalization\_audit.py} проверяет второй открытый подпункт C6: происхождение множителя \(\pi^2/2\). Для равномерной меры на \(S^3\subset\mathbb R^4\) используются стандартные моменты
\begin{equation}
    \mathbb E[x_ax_b]=\frac{\delta_{ab}}{4},
    \qquad
    \mathbb E[x_ax_bx_cx_d]
    =\frac{\delta_{ab}\delta_{cd}+\delta_{ac}\delta_{bd}+\delta_{ad}\delta_{bc}}{24}.
\end{equation}
Для квадратичных strains \(q_A=x^TAx\), \(A=A^T\), это даёт на \(\mathbb{RP}^3\)
\begin{equation}
    \langle q_A,q_B\rangle_{\mathbb{RP}^3}
    =\Vol(\mathbb{RP}^3)\frac{(\Tr A)(\Tr B)+2\Tr(AB)}{4\cdot6}.
\end{equation}
Следовательно, trace direction \(A\propto I\) ортогонален traceless sector, а
\begin{equation}
    \operatorname{Sym}^2(\mathbb R^4)
    =\mathbb R I\oplus\operatorname{Sym}^2_0(\mathbb R^4),
    \qquad 10=1+9.
\end{equation}
Поскольку \(q_A(-x)=q_A(x)\), эти функции спускаются с \(S^3\) на \(\mathbb{RP}^3\), и объём носителя равен
\begin{equation}
    \Vol(\mathbb{RP}^3)=\pi^2.
\end{equation}
Бозонный determinant даёт стандартный множитель \(1/2\) перед \(\log\det\). Поэтому естественная prefactor-нормировка coexact tower равна
\begin{equation}
    \frac12\Vol(\mathbb{RP}^3)=\frac{\pi^2}{2}=4.9348022005.
\end{equation}
Это условно закрывает volume-factor subtest C6: множитель \(\pi^2/2\) не является fitted coefficient, если \(T_{coex}^{\mathbb{RP}^3}\) нормирован как determinant density по \(\mathbb{RP}^3\). Но этот аудит не закрывает весь C6: знак полной Maxwell--ghost второй вариации проверяется отдельным sign-аудитом, а gap \(N_{need}-10=0.0099700224\) остаётся последним численным scheme-разрывом.

\end{protocol}

\begin{protocol}[C6 second-variation sign audit]
Следующий тест \texttt{s2t\_c6\_second\_variation\_sign\_audit.py} проверяет знак множителя
\begin{equation}
    1-\frac{\Tr P_{0,2}}{24S_{geo}}.
\end{equation}
Для determinant-блока вида \(c\log\det(A+\varepsilon B)\), если на данном шаге отброшены или локально вычтены \(\delta^2A\)-члены, имеем
\begin{equation}
    c\log\det(A+\varepsilon B)
    =c\log\det A+c\varepsilon\Tr(A^{-1}B)
    -\frac{c\varepsilon^2}{2}\Tr\left((A^{-1}B)^2\right)+O(\varepsilon^3).
\end{equation}
Для real bosonic coexact Maxwell block коэффициент \(c=1/2\), поэтому квадратичный trace-square вклад имеет знак
\begin{equation}
    -\frac14\Tr\left((A^{-1}B)^2\right)<0.
\end{equation}
Это ровно направление, требуемое absorption-формулой: базовый volume-half coexact trace умножается на suppression factor
\begin{equation}
    1-\frac{10}{24S_{geo}}=0.9969594432<1.
\end{equation}
Следовательно, sign-subtest C6 получает условный pass для coexact bosonic block.

Но это не полный вывод знака полного Maxwell--ghost functional. Faddeev--Popov ghost в effective action имеет противоположный determinant-знак; если тот же \(P_{0,2}\)-insertion действует внутри ghost determinant как самостоятельный trace-square, он даёт противоположную квадратичную поправку. Кроме того, настоящая вариация Лапласиана содержит не только \((\delta\Delta)^2\), но и \(\delta^2\Delta\)-члены. Поэтому для mature C6 нужно дополнительно показать:
\begin{enumerate}
    \item \(P_{0,2}\)-insertion относится к coexact bosonic mixed block, а ghost/exact ветвь после Hodge-сокращения оставляет только finite Casimir factor \(1/24\);
    \item \(\Tr(\Delta^{-1}\delta^2\Delta)\)-члены являются локальными, вычитаемыми или точно компенсированными в той же scheme.
\end{enumerate}
\end{protocol}

\begin{nogocriterion}[Провал sign-subtest C6]
Если ghost determinant несёт тот же \(P_{0,2}\) trace-square с противоположным знаком без компенсации, либо если \(\delta^2\Delta\)-члены дают finite nonlocal вклад того же порядка, то знак suppression factor не считается выведенным из Maxwell--ghost functional. Тогда \(1-10/(24S_{geo})\) остаётся sign-compatible heuristic, но не theorem.

\end{nogocriterion}

\begin{protocol}[C6 ghost/exact isolation audit]
Следующий аудит \texttt{s2t\_c6\_ghost\_exact\_isolation\_audit.py} проверяет главный caveat sign-subtest: не появляется ли тот же \(P_{0,2}\)-trace-square в ghost determinant с противоположным знаком. Рассмотрены три сценария.

Первый сценарий --- желаемый:
\begin{equation}
    P_{0,2}\subset \delta_g\Delta_1^{coex},
    \qquad
    P_{0,2}\not\to \delta_g\Delta_0^{ghost}\text{ как самостоятельный trace-square.}
\end{equation}
Тогда effective rank suppression равен \(10\), и сохраняется множитель
\begin{equation}
    1-\frac{10}{24S_{geo}}=0.9969594432.
\end{equation}

Второй сценарий --- провальный: если тот же \(P_{0,2}\) входит в ghost determinant как independent trace-square, Grassmann-знак противоположен bosonic-знаку. В упрощённой sign-модели suppression сокращается:
\begin{equation}
    (10)_{coex}-(10)_{ghost}=0,
    \qquad factor=1.
\end{equation}
Третий сценарий ещё хуже: если ghost/exact leakage включает обычную scalar tower до \(\ell=4\), то ghost-rank становится \(1+9+25=35\), и net factor меняет направление:
\begin{equation}
    1-\frac{10-35}{24S_{geo}}=1.0076013920>1.
\end{equation}
Это enhancement, а не нужное suppression.

Следовательно, C6 требует отдельной леммы изоляции:
\begin{equation}
    \text{ordinary exact/scalar ghost tower cancels before finite }P_{0,2}\text{ coexact trace is counted,}
\end{equation}
а константная scalar ветвь входит только через уже выделенный finite Casimir factor \(\kappa_{Cas}=1/24\). Без этой леммы знак \(1-10/(24S_{geo})\) не является gauge-fixed Maxwell--ghost выводом.
\end{protocol}

\begin{nogocriterion}[Провал ghost/exact isolation]
Если \(P_{0,2}\) не удаётся определить как coexact bosonic metric-strain insertion до ghost trace-square, либо если ghost/exact сектор несёт тот же finite-rank projector с противоположным знаком, то absorption-схема C6 проваливается. В этом случае \(P_{0,2}\)-rank остаётся естественной геометрической структурой, но не может использоваться как determinant proof для \(\pi^{-4}\)-члена.

\end{nogocriterion}

\begin{definition}[Формальный функционал \(\Gamma_{Maxwell+ghost}\) для C6]
Для проверки C6 фиксируем gauge-fixed one-loop functional на
\(K=L(2,1)\times S^1\) в разложении Ходжа. В рабочей II.A-схеме он записывается как
\begin{align}
    \Gamma_{M+gh}^{(1)}[g]
    =&\frac12\log\det{}'\Delta_{1,coex}[g]
      +\Gamma_{exact}[g]
      -\log\det{}'\Delta_0[g] \\
     &+\Gamma_{zero/gauge}[g]
      +\Gamma_{loc.ct.}[g].
\end{align}
Здесь \(\Delta_{1,coex}\) действует на transverse/coexact one-forms, \(\Gamma_{exact}\) обозначает продольный exact-sector после gauge fixing, \(-\log\det{}'\Delta_0\) --- Faddeev--Popov scalar ghost, \(\Gamma_{zero/gauge}\) содержит gauge volume и \(\det'\)-нулевые моды, а \(\Gamma_{loc.ct.}\) содержит локальные heat-kernel counterterms. Для \(\mathbb{RP}^3\) гармонические one-forms отсутствуют, поскольку \(b_1=0\).
\end{definition}

\begin{proposition}[Размещение \(P_{0,2}\) в coexact bosonic insertion]
В C6-ansatz finite projector \(P_{0,2}\) должен пониматься не как выбор первых scalar eigenmodes ghost-сектора, а как first ambient metric-strain insertion в coexact Maxwell operator:
\begin{equation}
    h_A:\quad x\mapsto (I+\varepsilon A)x,
    \qquad A=A^T,
    \qquad q_A(x)=x^TAx.
\end{equation}
Тогда
\begin{equation}
    P_{0,2}\simeq \operatorname{Sym}^2(\mathbb R^4)
    =\mathbb R I\oplus\operatorname{Sym}_0^2(\mathbb R^4),
    \qquad \Tr P_{0,2}=1+9=10.
\end{equation}
В этом смысле \(\ell\ge4\) не запрещены произвольно; они принадлежат higher ambient strain sectors и не входят в II.A mixed residue.
\end{proposition}

\begin{proposition}[Exact/scalar ghost isolation lemma как необходимое условие]
Для determinant-доказательства absorption-схемы требуется лемма изоляции:
\begin{equation}
    \Gamma_{exact}^{nonzero}[g]+\Gamma_{ghost}^{nonzero}[g]
    =\Gamma_{local/gauge}[g]
    \quad\text{до подсчёта finite }P_{0,2}\text{ coexact trace.}
\end{equation}
Иными словами, ordinary exact/scalar tower должен сократиться до того, как появляется конечный coexact first-strain trace. Константная periodic scalar ветвь не считается ghost trace-square insertion; она входит только как уже выделенный finite coefficient
\begin{equation}
    \kappa_{Cas}=\frac{1}{24},
\end{equation}
после удаления истинной \((0,0)\)-моды правилом \(\det'\).
\end{proposition}

\begin{protocol}[Physical quotient versus covariant FP]
Аудит \texttt{s2t\_c6\_ghost\_exact\_isolation\_audit.py} теперь различает две схемы. В physical transverse quotient можно определить determinant сразу на coexact block; тогда route C6 остаётся открытым при условии, что scalar zero branch уже сведена к \(\kappa_{Cas}=1/24\) и не несёт \(P_{0,2}\)-trace-square. В стандартной covariant Faddeev--Popov схеме такое сокращение не является автоматическим: exact one-forms, scalar ghost, \(\det'\)-правило, Hodge-Jacobian и gauge-volume factor должны быть объединены в одной нормировке. До этого residual scalar determinant или half-power не исключён. Поэтому isolation lemma получает статус theorem obligation, а не закрытого результата.
\end{protocol}

\begin{proposition}[Scalar FP bookkeeping audit]
Дополнительный аудит \texttt{s2t\_c6\_scalar\_fp\_bookkeeping\_audit.py} фиксирует bare covariant FP-счёт после Hodge-разложения:
\begin{equation}
    Z_{std}^{1-loop}
    \propto
    \det{}'\Delta_0\,\left(\det{}'\Delta_1\right)^{-1/2}
    =\left(\det{}'\Delta_{1,coex}\right)^{-1/2}
     \left(\det{}'\Delta_0\right)^{1/2},
\end{equation}
то есть для effective action
\begin{equation}
    \Gamma_{std}^{1-loop}
    =\frac12\log\det{}'\Delta_{1,coex}
     -\frac12\log\det{}'\Delta_0
     +\Gamma_{zero/gauge}+\Gamma_{loc.ct.}
\end{equation}
до возможных дополнительных zero/gauge/Jacobian сокращений. Следовательно, bare standard FP оставляет scalar half-determinant с противоположным знаком. Если этот residual несёт тот же finite \(P_{0,2}\)-trace-square, effective rank в sign-модели уменьшается с \(10\) до \(5\), а множитель становится
\begin{equation}
    1-\frac{5}{24S_{geo}}=0.9984797216,
\end{equation}
что даёт относительный уход от \(\pi^{-4}S_{geo}^{-2}\) порядка \(1.53\cdot10^{-3}\), а не \(3.04\cdot10^{-6}\). Поэтому standard FP route требует доказать либо полное сокращение scalar half-power в той же нормировке, либо отсутствие \(P_{0,2}\)-проекции у residual scalar branch.
\end{proposition}

\begin{proposition}[Scalar \(P_{0,2}\)-projection audit]
Аудит \texttt{s2t\_c6\_scalar\_p02\_projection\_audit.py} разделяет constant scalar branch и nonzero scalar tower. Для retained periodic row, дающей \(\kappa_{Cas}=1/24\), scalar mode постоянна на \(\mathbb{RP}^3\), поэтому
\begin{equation}
    \nabla_{\mathbb{RP}^3}\phi_0=0,
    \qquad
    \delta_{h_A}\Delta_0\,\phi_0=0
\end{equation}
для traceless first ambient strain \(A\in\operatorname{Sym}_0^2(\mathbb R^4)\). Следовательно, \(\ell=2\) traceless часть \(P_{0,2}\) не создаёт ghost trace-square на самой \(1/24\)-ветви. Trace-направление \(\mathbb RI\) остаётся вопросом volume/gauge normalization, а не автоматическим finite projector.

Это частичный позитивный результат: он поддерживает изоляцию \(\kappa_{Cas}=1/24\) от traceless \(P_{0,2}\)-leakage. Но он не закрывает standard FP route, потому что nonzero scalar half-determinant, если он не сокращён, имеет \(\nabla_{\mathbb{RP}^3}\phi\neq0\) и может coupling'оваться к first ambient strain через scalar stress tensor. Поэтому blocking gap переносится с constant branch на nonzero scalar residual tower.
\end{proposition}

\begin{proposition}[Scalar variation P02 nonzero audit]
Аудит \texttt{s2t\_c6\_scalar\_variation\_p02\_audit.py} проверяет, можно ли занулить nonzero scalar residual tower по symmetry selection rule. На \(\mathbb{RP}^3=S^3/\mathbb Z_2\) scalar harmonics наследуются только при чётных \(\ell\). Первая ненулевая scalar shell имеет
\begin{equation}
    \ell=2,
    \qquad
    d_2=(2+1)^2=9.
\end{equation}
Traceless first ambient strain также находится в \(\ell=2\)-канале \(\operatorname{Sym}_0^2(\mathbb R^4)\). Вариация scalar Laplacian содержит stress-tensor insertion
\begin{equation}
    \delta_h\Delta_0\sim -h^{ij}\nabla_i\nabla_j+\cdots,
\end{equation}
поэтому действует обычное правило Clebsch--Gordan типа
\begin{equation}
    \ell'\in \ell\otimes 2,
    \qquad
    |\ell-2|\le \ell'\le \ell+2,
\end{equation}
с сохранением чётности на \(\mathbb{RP}^3\). Для первой ненулевой shell это даёт
\begin{equation}
    \ell=2\quad\longrightarrow\quad \ell'=0,2,4,
\end{equation}
и diagonal coupling \(\ell=2\to2\) не запрещён. Следовательно, route ``zero by symmetry'' для nonzero scalar half-residual провален.

Остаются только два честных выхода: либо доказать, что весь этот scalar trace-square является local heat-kernel/volume term и вычитается в \(\Gamma_{loc.ct.}\) или \(\Gamma_{zero/gauge}\), либо отказаться от standard FP route для C6 и использовать physical transverse quotient как defining scheme. Без такого дополнительного доказательства standard covariant FP даёт открытый \(P_{0,2}\)-leakage channel и не выводит rank \(\Tr P_{0,2}=10\) как determinant theorem.
\end{proposition}

\begin{nogocriterion}[Standard FP rescue-routes audit]
Аудит \texttt{s2t\_c6\_scalar\_rescue\_routes\_audit.py} проверяет два возможных спасения standard covariant FP route.

Первый вариант --- объявить scalar trace-square полностью local/subtracted. Он не закрывается: локальные heat-kernel counterterms удаляют zero-winding/UV-часть, но для каждого положительного scalar eigenvalue после суммирования по окружности остаётся finite nonlocal Bessel/winding part. Первая ненулевая scalar shell \(\ell=2\) уже имеет \(\lambda>0\), кратность \(9\) и разрешённый \(P_{0,2}\)-coupling, поэтому одного утверждения ``это local'' недостаточно.

Второй вариант --- cancellation через zero/gauge/Jacobian factors. Он также не закрывается в текущем bookkeeping. \(\det'\)-правило удаляет истинную zero-mode, gauge volume отвечает за конечномерную gauge orbit, а Hodge-Jacobian уже использован в identity
\begin{equation}
    \det{}'\Delta_1=\det{}'\Delta_{1,coex}\det{}'\Delta_0,
\end{equation}
которая и даёт residual \(-\frac12\log\det{}'\Delta_0\) в standard FP effective action. Эти факторы не предоставляют нового opposite nonzero scalar tower. Следовательно, standard FP route не выводит rank \(\Tr P_{0,2}=10\) как theorem без дополнительного, пока отсутствующего analytic cancellation.

Оставшийся чистый путь: определить EM determinant на physical transverse/coexact quotient до scalar FP residual, либо честно понизить \(\pi^{-4}\)-член до structural compression до появления нового paired/cancellation-сектора.
\end{nogocriterion}

\begin{proof}[Статус доказательства]
Это пока не полная теорема, а формальная редукция обязательств. Если exact/scalar cancellation выполняется до finite coexact insertion, то ghost determinant не несёт противоположный \(P_{0,2}\)-trace-square, и знак suppression factor остаётся coexact-bosonic. Если же тот же \(P_{0,2}\) входит в ghost determinant, suppression сокращается; если ghost leakage доходит до \(\ell=4\), знак меняется на enhancement. В стандартном covariant FP-bookkeeping остаётся дополнительная проверка: вывести Hodge measure Jacobian для \(A=A_{coex}+d\phi\), совместить exact determinant с ghost determinant и gauge volume, а затем доказать, что возможный scalar residual является только \(1/24\)-ветвью, а не finite \(P_{0,2}\)-вкладом. Поэтому isolation lemma является не косметической технической деталью, а центральным условием C6.
\end{proof}

\begin{proposition}[Оставшиеся \(\delta^2\Delta\)-члены]
Во второй вариации функционала присутствует не только trace-square term, но и член вида
\begin{equation}
    \frac12\Tr\left(\Delta^{-1}\delta_g^2\Delta\right)
\end{equation}
с секторными determinant-коэффициентами. Для C6 требуется показать, что эти вклады являются локальными heat-kernel terms, уже включены в \(\Gamma_{loc.ct.}\), либо точно компенсируются между gauge-fixed секторами. Если остаётся finite nonlocal часть того же порядка, absorption-формула должна быть изменена.
\end{proposition}

\begin{protocol}[Gap \(N_{need}-10\) как развилка]
Формальная C6-схема приводит к числам
\begin{align}
    N_{need}&=10.0099700224,\\
    \Tr P_{0,2}&=10,\\
    N_{need}-\Tr P_{0,2}&=0.0099700224.
\end{align}
Использование \(\Tr P_{0,2}=10\) даёт относительное расхождение с \(\pi^{-4}\)-членом
\begin{equation}
    3.04\cdot10^{-6}.
\end{equation}
Следовательно, последний шаг C6 имеет только два честных исхода:
\begin{enumerate}
    \item вывести gap как finite scheme residue, subleading determinant contribution или \(\delta^2\Delta\)-остаток в той же gauge-fixed схеме;
    \item если такой вывод невозможен, понизить \(\pi^{-4}\)-absorption до strong structural compression, а не theorem.
\end{enumerate}
\end{protocol}

\begin{theorem}[Промежуточный статус C6 до внешнего functional gate]
Функционал \(\Gamma_{Maxwell+ghost}[g]\) можно формально разложить по coexact, exact, ghost и zero-mode секторам, а \(P_{0,2}\) можно условно разместить в coexact bosonic first-strain insertion. Это закрывает форму задачи, но не закрывает C6 как теорему. Открытыми остаются: доказательство exact/scalar ghost isolation, локальность или компенсация \(\delta^2\Delta\)-членов и вывод gap \(N_{need}-10\). До закрытия этих трёх пунктов \(S_{vac}\) сохраняет статус условного determinant-успеха.
\end{theorem}

\begin{nogocriterion}[Провал формального C6]
Если хотя бы один из трёх узлов --- ghost isolation, \(\delta^2\Delta\)-локальность или gap \(N_{need}-10\) --- не выводится внутри того же \(\Gamma_{Maxwell+ghost}[g]\), то \(\pi^{-4}\)-член не считается доказанным determinant residue. В этом случае он остаётся лучшей найденной структурной компрессией coexact tower, но не повышает EM-блок до mature theorem.
\end{nogocriterion}

\begin{protocol}[Master closure-matrix C6]
После серии локальных аудитов добавлен сводный реестр \texttt{s2t\_c6\_closure\_matrix\_results.json}. Его назначение --- не искать новый численный коэффициент, а зафиксировать дерево статусов C6. Позитивные узлы таковы:
\begin{enumerate}
    \item формальное разложение \(\Gamma_{Maxwell+ghost}\) на coexact, exact, ghost, zero/gauge и local counterterm части построено;
    \item ранг \(P_{0,2}=\operatorname{Sym}^2(\mathbb R^4)=1+9=10\) является естественным для first ambient strain;
    \item знак второго порядка для real bosonic coexact determinant совместим с suppression-множителем \(1-N/(24S_{geo})\);
    \item constant periodic branch, дающая \(\kappa_{Cas}=1/24\), изолирована от traceless \(P_{0,2}\)-leakage.
\end{enumerate}
Блокирующие узлы также фиксируются явно:
\begin{enumerate}
    \item bare standard covariant FP/Hodge bookkeeping оставляет residual \(-\frac12\log\det{}'\Delta_0\) в effective action;
    \item nonzero scalar residual tower не зануляется symmetry selection rule, поскольку первая чётная shell \(\ell=2\) имеет кратность \(9\) и может coupling'оваться к \(P_{0,2}\);
    \item local counterterms, zero modes, gauge volume и Hodge-Jacobian не дают mode-by-mode cancellation nonlocal Bessel tower;
    \item known paired candidates --- BRST/Nielsen--Kallosh-like extra ghost, Hodge-remnant, Ray--Singer torsion, Maxwell duality и Dirac/spin-cover matter --- не дают обязательный независимый вклад \(+\frac12\log\det{}'\Delta_{0,nonzero}\) с тем же \(P_{0,2}\)-trace-square.
\end{enumerate}
Следовательно, C6 теперь имеет не расплывчатый, а дискретный статус. В standard covariant FP без нового обязательного сектора theorem-route заблокирован nonzero scalar half-determinant residual. Единственный чистый путь, который ещё не является phenomenological compression, --- определить или вывести electromagnetic determinant непосредственно на physical transverse/coexact quotient и затем пройти external spectral-determinant gate. Если этот quotient не принимается как фундаментальная нормировка, \(\pi^{-4}\)-член должен называться strong structural compression, а не mature determinant theorem.
\end{protocol}

\begin{nogocriterion}[Запрет paired-сектора вручную]
Новый сектор, введённый только для получения \(+\frac12\log\det{}'\Delta_{0,nonzero}\) и компенсации scalar residual, не повышает статус C6. Он допустим только если выводится как обязательный BRST-, topological- или EFT-компонент той же модели, имеет тот же nonzero even scalar spectrum на \(\mathbb{RP}^3\times S^1\), тот же \(P_{0,2}\)-trace-square и не вводит нового параметра. Иначе он считается hidden model ingredient.
\end{nogocriterion}

\begin{protocol}[Defense-аудит physical transverse quotient]
Следующий аудит \texttt{s2t\_c6\_physical\_quotient\_defense\_results.json} проверяет оставшийся вариант C6: считать primary object не bare covariant FP determinant, а физический determinant на пространстве gauge-equivalence classes, представленном coexact/transverse slice. В этой схеме рабочее определение имеет вид
\begin{equation}
    \Gamma_{EM}^{phys}[g]
    =\frac12\log\det{}'\Delta_{1,coex}[g]
     +\Gamma_{zero/gauge}[g]+\Gamma_{loc.ct.}[g],
\end{equation}
а scalar periodic branch, дающая \(\kappa_{Cas}=1/24\), остаётся отдельным IR/zero-mode вкладом после \(\det'\)-удаления истинной нулевой моды. Важно: это не позднее ``сокращение'' standard FP residual. Это допустимо только как заранее объявленная determinant-нормировка на physical quotient.

Defense-аудит даёт статус \emph{viable but not yet theorem}. Плюс состоит в том, что nonzero scalar \(P_{0,2}\)-leakage отсутствует по области определения, coexact bosonic sign сохраняется, а новый параметр не вводится. Минус состоит в theorem obligation: нужно пройти external spectral-determinant gate, то есть показать совместимость такого physical determinant с gauge-fixed Maxwell partition function modulo local counterterms, zero-mode factors, gauge volume и possible torsion-normalization. Если внешняя нормировка требует вернуть nonzero scalar residual в ту же finite \(P_{0,2}\)-variation, defense quotient проваливается и \(\pi^{-4}\) понижается до strong structural compression.
\end{protocol}

\begin{protocol}[External gate для physical quotient]
Добавлен внешний gate-аудит \texttt{s2t\_c6\_external\_physical\_quotient\_gate\_results.json}. Его вывод можно сформулировать в простом виде: литература разрешает сам язык quotient, Hodge split, ghost determinant, analytic torsion и zero-mode/gauge-volume factors, но не дарит автоматически S2T-остаток \(\pi^{-4}\). Поэтому gate получает статус \emph{conditional pass, not theorem}. Physical quotient можно использовать как primary determinant domain, но нельзя выдавать его за уже доказанную covariant-FP cancellation.

Полученный результат не даёт автоматической компенсации scalar residual.
Остаётся строго определённый вариант: ограничить физический determinant
transverse/coexact модами и вычислить их mixed trace на \(L(2,1)\). Если этот
след даёт требуемый \(P_{0,2}\)-остаток и external gate не восстанавливает
scalar leakage, статус C6 может быть повышен. В противном случае
\(\pi^{-4}\) остаётся spectral compression без статуса теоремы.
\end{protocol}

\begin{protocol}[Аудит смешанного следа на \(L(2,1)\)]
Следующий прямой аудит \texttt{s2t\_c6\_l21\_coexact\_mixed\_trace\_results.json} фиксирует настоящий математический узел. В физической поперечной схеме нужно вычислить не весь калибровочный определитель, а конкретный след
\begin{equation}
    C_{AA}=\Tr_{coex}\left(
    \Delta_{1,coex}^{-1}\,\delta_A\Delta_{1,coex}\,
    \Delta_{1,coex}^{-1}\,\delta_A\Delta_{1,coex}
    \right),
    \qquad A\in\operatorname{Sym}^2(\mathbb R^4),
\end{equation}
на \(L(2,1)\simeq\mathbb{RP}^3\). Уже поддержаны три части: область поперечных мод убирает ненулевой скалярный хвост по определению, знак вещественного бозонного определителя совместим с вычитанием, а пространство первой внешней деформации имеет ранг \(1+9=10\). Но этого ещё недостаточно. Не вычислены матричные элементы \(\delta_A\Delta_1\) между поперечными векторными гармониками линзового пространства, а значит не доказано, что реальный след сворачивается ровно к \(\Tr P_{0,2}=10\), а не несёт всю башню поперечных мод.

Численная развилка остаётся прежней:
\begin{equation}
    N_{need}-10=0.0099700224,
    \qquad
    \frac{\Delta_{P_{0,2}}-\pi^{-4}S_{geo}^{-2}}{\pi^{-4}S_{geo}^{-2} }
    =3.04\cdot10^{-6}.
\end{equation}
Следовательно, C6 стал не туманной идеей, а чётким заданием: выписать поперечные векторные гармоники на \(L(2,1)\), найти действие первой метрической деформации на лапласиан 1-форм и посчитать след по оболочкам. Если след сворачивается к рангу \(10\) с допустимым конечным остатком, маршрут оживает. Если появляется вся поперечная башня, \(\pi^{-4}\) остаётся сильным спектральным сжатием, но не теоремой.
\end{protocol}

\begin{protocol}[Базис поперечных гармоник на \(L(2,1)\)]
Подготовительный аудит \texttt{s2t\_c6\_l21\_coexact\_basis\_results.json} фиксирует оболочки поперечных 1-форм. В используемой нормировке на \(S^3\) поперечная оболочка имеет
\begin{equation}
    \lambda_n=(n+1)^2,
    \qquad
    d_n^{S^3}=2n(n+2),
    \qquad n\ge1.
\end{equation}
Проекция на \(L(2,1)=S^3/\{\pm1\}\) оставляет нечётные \(n\) и убирает чётные \(n\). Поэтому первая выжившая поперечная оболочка имеет
\begin{equation}
    n=1,
    \qquad
    \lambda_1=4,
    \qquad
    d_1^{L(2,1)}=6.
\end{equation}
Это важный отрицательный вывод: число \(10\) не является кратностью первой поперечной моды и не получается накоплением первых оболочек, поскольку далее идут \(6,36,106,\ldots\). Следовательно, если C6 верен, ранг \(10\) должен возникать не из счёта собственных мод, а из пространства первых внешних деформаций
\begin{equation}
    \operatorname{Sym}^2(\mathbb R^4)=1+9.
\end{equation}
Тем самым задача стала острее: нужно доказать, что след по бесконечной поперечной башне после вставки \(\delta_A\Delta_1\) сворачивается к следу по десяти направлениям деформации. Без матричных элементов
\begin{equation}
    \langle n,i|\delta_A\Delta_1|m,j\rangle
\end{equation}
такое сворачивание нельзя считать выведенным.
\end{protocol}

\begin{protocol}[Нормировка поперечных форм на \(L(2,1)\)]
Аудит \texttt{s2t\_c6\_l21\_normalization\_results.json} закрывает следующий технический подпункт перед вычислением матричных элементов. Для двойного накрытия
\begin{equation}
    p:S^3\to L(2,1)=S^3/\{\pm1\}
\end{equation}
объём единичного носителя уменьшается вдвое:
\begin{equation}
    \operatorname{Vol}(L(2,1))=\frac12\operatorname{Vol}(S^3)=\pi^2.
\end{equation}
Если инвариантный подъём поперечной 1-формы нормирован на \(S^3\) как
\begin{equation}
    \int_{S^3}|\alpha|^2\,d\operatorname{vol}=1,
\end{equation}
то спущенная форма имеет норму \(1/2\) на \(L(2,1)\). Поэтому ортонормированный представитель на факторе равен \(\sqrt2\) умножить на спущенную форму.

Отсюда следует практическое правило для следующего mixed-trace расчёта. В билинейном матричном элементе два множителя \(\sqrt2\) от внешних состояний ровно сокращают множитель \(1/2\) от интеграла по фактору. Поэтому интегралы вида
\begin{equation}
    \langle \alpha_i,V_q\alpha_j\rangle_{L(2,1)}
\end{equation}
можно считать на \(S^3\) по инвариантным подъёмам без дополнительного глобального множителя покрытия. Покрытие меняет допустимые состояния и их ортонормировку, но после этой ортонормировки запрещено домножать итоговый след ещё раз на \(2\) или \(1/2\). Также запрещено возвращать чётные оболочки только потому, что интеграл технически записан на \(S^3\): пространство состояний остаётся факторным.
\end{protocol}

\begin{protocol}[Оболочечное правило для квадратичной деформации]
Аудит \texttt{s2t\_c6\_l21\_shell\_selection\_results.json} проверяет, может ли один только representation-level selection rule вывести конечный ранг \(10\). Квадратичная первая внешняя деформация принадлежит
\begin{equation}
    P_{0,2}=\ell=0\oplus\ell=2,
\end{equation}
то есть она чётна относительно \(x\mapsto -x\) и потому сохраняет сектор \(L(2,1)\). Но это не означает, что она оставляет только конечное число поперечных оболочек. На уровне необходимого правила сложения гармоник допустимы каналы
\begin{equation}
    n\longrightarrow m,
    \qquad
    |n-m|\in\{0,2\},
    \qquad
    n,m\ \,\text{нечётные}.
\end{equation}
Следовательно, первая выжившая оболочка \(n=1\) может иметь каналы \(1\to1\) и \(1\to3\), оболочка \(n=3\) --- каналы \(3\to1,3,5\), и так далее по всей нечётной башне.

Это отрицательный, но важный вывод: ранг \(10=\dim\operatorname{Sym}^2(\mathbb R^4)\) не следует из оболочечного отбора сам по себе. Он остаётся рангом пространства деформаций, а не доказанным рангом операторного следа. Чтобы C6 стал теоремой, нужно дополнительное утверждение уровня коэффициентов: либо полная вариация \(\delta_A\Delta_1\) после поперечной проекции зануляет все разрешённые башенные каналы, либо эти каналы являются локальными и вычитаемыми, либо physical quotient задаёт нормировку, в которой бесконечная башня поглощается без нового параметра. Без такого шага \(P_{0,2}\) остаётся сильной геометрической селекцией, но не полным mixed-trace proof.
\end{protocol}

\begin{protocol}[Locality gate для coexact-башни]
Аудит \texttt{s2t\_c6\_l21\_coexact\_locality\_gate\_results.json} проверяет следующий возможный спасательный ход: объявить все разрешённые оболочечные каналы локальными и вычесть их контрчленами. Такой ход не проходит в общем виде. Локальными кандидатами являются только UV/asymptotic части теплового ядра, то есть большие-\(n\) полиномиальные коэффициенты, уже предусмотренные схемой локальных геометрических контрчленов.

Но mixed-trace содержит также низкие оболочки и конечные off-diagonal вклады. Например, после предыдущего отбора остаются каналы
\begin{equation}
    1\to1,
    \qquad
    1\to3,
    \qquad
    3\to1,
\end{equation}
и они являются глобальными спектральными данными, а не локальным коэффициентом теплового ядра. Их нельзя удалить фразой ``локальный контрчлен'' без явного доказательства, что полный one-form оператор превращает эти вклады в локальную вариацию уже фиксированной нормировки. Иначе это было бы конечным постфактум-вычитанием, запрещённым no-hidden-parameter критерием.

Следовательно, C6 теперь имеет более жёсткий gate. Нужно разложить mixed-trace на локальную асимптотическую часть и конечную нелокальную determinant-часть. Локальную часть можно вычитать только заранее заданными контрчленами. Конечная часть должна либо точно совпасть с уже существующим \(\pi^{-4}\)-остатком и \(P_{0,2}\)-рангом, либо быть занулена/скомпенсирована тензорной алгеброй полной вариации \(\delta_A\Delta_1\). Если конечная часть остаётся независимой, то \(\pi^{-4}\) нельзя повышать до determinant theorem; его статус остаётся structural compression.
\end{protocol}

\begin{protocol}[Спецификация low-shell блока]
Аудит \texttt{s2t\_c6\_l21\_low\_shell\_block\_spec\_results.json} фиксирует следующий нериторический расчёт. После shell-selection и locality-gate минимальная конечная проверка состоит из каналов
\begin{equation}
    1\to1,
    \qquad
    1\to3,
    \qquad
    3\to1.
\end{equation}
Для поперечных 1-форм на \(L(2,1)\) уже зафиксировано
\begin{equation}
    d_1=6,
    \qquad
    \lambda_1=4,
    \qquad
    d_3=30,
    \qquad
    \lambda_3=16.
\end{equation}
Поэтому размеры матриц для одного направления деформации \(A\in\operatorname{Sym}^2(\mathbb R^4)\) равны
\begin{equation}
    M_{11}:6\times6,
    \qquad
    M_{13}:6\times30,
    \qquad
    M_{31}:30\times6.
\end{equation}
Это даёт \(36+180+180=396\) матричных элементов на одно направление \(A\), или \(3960\) элементов до использования симметрий при явном проходе по десяти направлениям \(\operatorname{Sym}^2(\mathbb R^4)\). Соответствующие веса в trace-square имеют вид
\begin{equation}
    \lambda_1^{-2}=\frac1{16},
    \qquad
    (\lambda_1\lambda_3)^{-1}=\frac1{64}.
\end{equation}

Этот блок не разрешено заменять скалярным toy model. В каждом матричном элементе должны присутствовать главный символ, вариация связности, вариация Ricci-члена, поперечная проекция и вариация скалярного произведения 1-форм. Если все эти низкие блоки зануляются, C6 получает сильную поддержку, но всё ещё требует индукции по нечётной башне. Если блоки ненулевые, но точно дают absorption identity для уже существующего \(\pi^{-4}\)-остатка, маршрут остаётся живым. Если появляется независимый конечный остаток, determinant-theorem статус для \(\pi^{-4}\) блокируется.
\end{protocol}

\begin{protocol}[Первый тест \(n=1\): Killing-overlap]
Аудит \texttt{s2t\_c6\_l21\_n1\_killing\_overlap\_results.json} проверяет, можно ли отбросить диагональный канал \(1\to1\) по простой симметрии. Первая поперечная оболочка имеет кратность \(6\) и реализуется шестью Killing 1-формами на \(S^3\), задаваемыми антисимметричными генераторами \(E_{ab}\in\mathfrak{so}(4)\). Эти формы чётно спускаются на \(L(2,1)\), а их собственное значение равно \(\lambda_1=4\).

Для traceless квадратичной деформации
\begin{equation}
    A=\operatorname{diag}(1,-1,0,0),
    \qquad
    q_A=x_0^2-x_1^2,
\end{equation}
проверен базовый overlap
\begin{equation}
    \int_{L(2,1)} q_A\,\langle \alpha_i,\alpha_j\rangle\,d\operatorname{vol}.
\end{equation}
После нормировки Killing-базиса на факторе полученная \(6\times6\) матрица не равна нулю: её численный ранг равен \(4\), максимальный модуль элемента равен \(1/6\), а собственные значения симметризованной матрицы равны
\begin{equation}
    -\frac16,
    -\frac16,
    0,
    0,
    \frac16,
    \frac16.
\end{equation}
Следовательно, канал \(1\to1\) нельзя исключить parity-аргументом или representation-level selection rule. Это ещё не доказывает провал C6, потому что полный оператор \(\delta_A\Delta_1\) содержит главный символ, вариацию связности, Ricci-член, поперечную проекцию и вариацию скалярного произведения; эти части теоретически могут взаимно сократиться. Но теперь такое сокращение должно быть явным тензорным фактом, а не предположением симметрии.
\end{protocol}

\begin{protocol}[Toy trace-square для \(n=1\)]
Аудит \texttt{s2t\_c6\_l21\_n1\_toy\_tracesquare\_results.json} оценивает масштаб предупреждения, если оставить только найденный выше \(q_A\)-overlap на Killing-оболочке. Для traceless направления \(A=\operatorname{diag}(1,-1,0,0)\) нормированная overlap-матрица имеет
\begin{equation}
    \operatorname{Tr}(M_A^2)=\frac19.
\end{equation}
С учётом \(\lambda_1=4\) toy trace-square равен
\begin{equation}
    \lambda_1^{-2}\operatorname{Tr}(M_A^2)
    =\frac1{16}\cdot\frac19
    =\frac1{144}.
\end{equation}
Это число не является C6-коэффициентом: полный оператор \(\delta_A\Delta_1\) содержит производные, вариацию связности, Ricci-член, поперечную проекцию и вариацию Hilbert-метрики. Но как diagnostic оно важно. Оно показывает, что нижняя coexact-оболочка имеет конечный отклик уже на простейшем \(P_{0,2}\)-overlap уровне. Поэтому если полный \(1\to1\) блок всё же занулится, это должно произойти через явное сокращение всех one-form terms, а не через отсутствие начального канала.
\end{protocol}

\begin{protocol}[Principal-symbol тест на \(n=1\)]
Аудит \texttt{s2t\_c6\_l21\_n1\_principal\_symbol\_results.json} проверяет первый настоящий кусок оператора \(\delta_A\Delta_1\). В конвенции тома
\begin{equation}
    \Delta_1\alpha=-\nabla^2\alpha+\operatorname{Ric}(\alpha).
\end{equation}
На единичной сфере для Killing-оболочки \(n=1\) выполнено
\begin{equation}
    \Delta_1\alpha=4\alpha,
    \qquad
    \operatorname{Ric}(\alpha)=2\alpha,
    \qquad
    \nabla^2\alpha=-2\alpha.
\end{equation}
В редуцированном конформном срезе \(h=2qg\) имеем \(\delta g^{ab}=-2qg^{ab}\). Поэтому principal-symbol часть вариации оператора даёт на этой оболочке
\begin{equation}
    \delta(-g^{ab}\nabla_a\nabla_b)
    =-\delta g^{ab}\nabla_a\nabla_b
    =2q\nabla^2
    =-4q.
\end{equation}
Следовательно, найденная ранее nonzero overlap-матрица умножается на \(-4\). Для traceless тестового направления это даёт weighted trace-square
\begin{equation}
    \lambda_1^{-2}\operatorname{Tr}(M_{principal}^2)=\frac19.
\end{equation}
Это снова не полный C6-коэффициент, но теперь предупреждение сильнее: уже главный символ полного оператора ненулевой. Значит, спасение \(1\to1\) требует явного матричного сокращения с connection-, Ricci-, projection- и Hilbert-metric terms, а не только слабого аргумента о выборе среза.
\end{protocol}

\begin{protocol}[Бюджет сокращения для \(n=1\)]
Следующий аудит \texttt{s2t\_c6\_l21\_n1\_cancellation\_budget\_results.json} переводит principal-symbol warning в точное обязательство. Если полный диагональный Killing-блок должен занулиться, сумма всех непросчитанных частей оператора обязана дать матрицу
\begin{equation}
    M_{conn+Ric+proj+Hilb}=-M_{principal}.
\end{equation}
Для traceless направления её собственные значения должны быть
\begin{equation}
    -\frac23,\ -\frac23,
    \ 0,
    \ 0,
    \ \frac23,
    \ \frac23,
\end{equation}
а weighted trace-square снова равен
\begin{equation}
    \lambda_1^{-2}\operatorname{Tr}(M_{conn+Ric+proj+Hilb}^2)=\frac19.
\end{equation}
Ключевая деталь: проекция этой требуемой матрицы на тождественный оператор равна нулю. Следовательно, её нельзя получить чистой перенормировкой объёма, scalar counterterm или trace-only bookkeeping. Если C6 спасается через зануление, сокращение должно быть матричным и базисно согласованным. Если C6 спасается через absorption-route, остаток после такого полного расчёта должен без подгоняемого коэффициента совпасть с уже существующим \(\pi^{-4}/P_{0,2}\)-остатком.
\end{protocol}

\begin{protocol}[Полный conformal Hodge-тест на \(n=1\)]
Аудит \texttt{s2t\_c6\_l21\_n1\_full\_conformal\_hodge\_results.json} проверяет, реализуется ли требуемый выше бюджет в редуцированном срезе \(h=2qg\). Для 1-форм в размерности \(3\), используя \(\Delta=d\delta+\delta d\), первая conformal-вариация на coexact собственной форме имеет вид
\begin{equation}
    \delta\Delta\alpha
    =-2\lambda q\alpha
    +\iota_{\nabla q}d\alpha
    -d\bigl(\iota_{\nabla q}\alpha\bigr).
\end{equation}
На Killing-оболочке \(n=1\), для того же traceless направления \(A=\diag(1,-1,0,0)\), полный \(6\times6\) блок даёт
\begin{equation}
    \lambda_1^{-2}\operatorname{Tr}(M_{full\ conformal}^2)
    =8.10\cdot10^{-33},
\end{equation}
то есть ноль с точностью до машинного округления. Максимальный элемент ошибки между непринципиальной частью и требуемым budget-сокращением равен \(2.22\cdot10^{-16}\).

Это важный поворот: principal-symbol warning не является самостоятельным no-go для диагонального \(1\to1\) Killing-блока. В conformal representative недостающие Hodge-члены ровно гасят главный символ. Однако C6 этим не доказан: ещё открыты raw ambient Hessian/diffeomorphism bookkeeping, вариация поперечной проекции, каналы \(1\leftrightarrow3\) и конечная часть полной coexact-башни.
\end{protocol}

\begin{protocol}[Утечка \(n=1\) в кубический слой]
Аудит \texttt{s2t\_c6\_l21\_n1\_to\_n3\_leakage\_results.json} проверяет, означает ли зануление диагонального \(1\to1\) блока исчезновение perturbation как оператора. Ответ отрицательный. Для Killing-формы \(\alpha_N\), где \(N\in\mathfrak{so}(4)\), оставшийся conformal Hodge-образ задаётся кубическим tangent-polynomial выражением
\begin{equation}
    V_N=-12qNx+4NAx-2C_Nx+2f_Nx,
    \qquad
    C_N=N^TA+A^TN,
    \qquad
    f_N=x^TC_Nx.
\end{equation}
Его проекция обратно на Killing-оболочку равна нулю:
\begin{equation}
    \max |P_{n=1}V_N|=0.
\end{equation}
Но нормированная Gram-матрица кубического образа имеет ранг \(6\), след \(96\) и собственные значения
\begin{equation}
    12,\ 16,\ 16,\ 16,\ 16,\ 20.
\end{equation}
Иными словами, первый этаж молчит не потому, что сигнал исчез: в conformal representative он ушёл в кубический слой. Поэтому следующий решающий тест C6 --- построить или численно ортонормировать coexact-базис \(n=3\) и вычислить настоящий second-order вклад канала \(1\leftrightarrow3\). Если этот вклад остаётся ненулевым и не входит в уже существующий \(\pi^{-4}/P_{0,2}\)-остаток, diagonal cancellation не спасает determinant-theorem status.
\end{protocol}

\begin{protocol}[Coexact-gate для кубического образа]
Аудит \texttt{s2t\_c6\_l21\_n3\_coexact\_gate\_results.json} проверяет, можно ли уже объявить найденную кубическую утечку физическим coexact-вкладом \(n=3\). Ответ: пока нельзя. Сырой ambient-представитель имеет ненулевую нормальную компоненту:
\begin{equation}
    \max |x\cdot V|=16.
\end{equation}
При этом его ambient-дивергенция равна нулю,
\begin{equation}
    \max |\operatorname{div}_{\mathbb R^4}V|=0,
\end{equation}
а степень полинома равна \(3\). Следовательно, найден не мусорный нулевой след, но и не готовый coexact \(n=3\)-матричный элемент. Геометрически это означает: видна сильная кубическая тень, но часть ambient-вектора направлена нормально к сфере. Перед determinant-выводом необходимо построить intrinsic tangent projection, затем заново проверить co-closedness и только после этого проектировать на ортонормированный coexact-базис \(n=3\).
\end{protocol}

\begin{protocol}[Tangent-projection для кубического образа]
Аудит \texttt{s2t\_c6\_l21\_n3\_tangent\_projection\_results.json} выполняет следующий геометрический фильтр:
\begin{equation}
    V_{tan}=V-(x\cdot V)x,
    \qquad |x|=1.
\end{equation}
После этой проекции нормальная компонента исчезает на сфере:
\begin{equation}
    \|x\cdot V_{tan}\|_{S^3}^2=0.
\end{equation}
Однако сам сигнал не исчезает. Нормированная Gram-матрица касательного кубического образа имеет след
\begin{equation}
    \operatorname{Tr}G_{tan}=85.3333333333
\end{equation}
и собственные значения
\begin{equation}
    12,\ 14.6667,\ 14.6667,\ 14.6667,\ 14.6667,\ 14.6667.
\end{equation}
Это означает: нормальная ambient-тень не объясняет всю утечку. На сфере остаётся ненулевой касательный кубический рисунок. Но coexact-статус всё ещё не закрыт: ambient-дивергенция после tangent-projection ненулевая, поэтому требуется intrinsic divergence или явный Hodge/coexact projector на \(S^3/\mathbb Z_2\). Только после этого можно считать канал \(1\leftrightarrow3\) настоящим determinant-вкладом или доказанным сокращением.
\end{protocol}

\begin{protocol}[Intrinsic-divergence gate для \(n=3\)]
Аудит \texttt{s2t\_c6\_l21\_n3\_intrinsic\_divergence\_results.json} фиксирует следующий барьер. Касательный кубический сигнал ненулевой:
\begin{equation}
    \operatorname{Tr}G_{tan}=85.3333333333.
\end{equation}
Но divergence/coexact gate не пройден. Gram-след дивергенции равен
\begin{equation}
    \operatorname{Tr}G_{div}=170.6666666667,
\end{equation}
а собственные значения divergence-Gram матрицы имеют вид
\begin{equation}
    0,\ 21.3333,\ 21.3333,\ 21.3333,\ 21.3333,\ 85.3333.
\end{equation}
Следовательно, касательный рисунок содержит longitudinal/exact компоненту. Это запрещает два преждевременных вывода одновременно: нельзя объявить найденный \(n=3\)-сигнал готовым coexact determinant obstruction, но нельзя и отбросить канал как нулевой. Следующий обязательный шаг --- применить Hodge-проектор на coexact-часть кубического касательного сектора и пересчитать вклад \(1\leftrightarrow3\).
\end{protocol}

\begin{protocol}[Hodge-projection proxy для \(n=3\)]
Аудит \texttt{s2t\_c6\_l21\_n3\_hodge\_projection\_results.json} выполняет первый trace-level Hodge-фильтр. Дивергенция касательного кубического поля попадает в scalar-сектор \(\ell=2\), где
\begin{equation}
    \lambda^{(0)}_{\ell=2}=\ell(\ell+2)=8.
\end{equation}
Поэтому exact/gradient норма, удаляемая Hodge-проектором, оценивается как
\begin{equation}
    \operatorname{Tr}G_{exact}
    =\frac{\operatorname{Tr}G_{div}}{8}
    =21.3333333333.
\end{equation}
Из исходного касательного следа
\begin{equation}
    \operatorname{Tr}G_{tan}=85.3333333333
\end{equation}
остаётся coexact-proxy
\begin{equation}
    \operatorname{Tr}G_{coex}^{proxy}=64.
\end{equation}
С учётом разности \(\lambda_3-\lambda_1=12\) соответствующий denominator-proxy равен
\begin{equation}
    \frac{64}{12^2}=0.4444444444.
\end{equation}
Это первый серьёзный кандидат на \(1\leftrightarrow3\) obstruction: нормальная и exact-части уже не объясняют весь сигнал. Но статус остаётся proxy-level. Для determinant-вывода нужно построить ортонормированный coexact-базис \(n=3\), спроецировать на него образ и только затем фиксировать знак, знаменатель и нормировку следа.
\end{protocol}

\begin{protocol}[Масштаб \(n=3\)-proxy obstruction]
Аудит \texttt{s2t\_c6\_l21\_n3\_proxy\_obstruction\_scale\_results.json} сравнивает найденный proxy-остаток с rank-10 маршрутом \(P_{0,2}\). Получено
\begin{equation}
    \operatorname{Tr}G_{coex}^{proxy}=64,
    \qquad
    \frac{64}{10}=6.4.
\end{equation}
С учётом \(\lambda_3-\lambda_1=12\) denominator-proxy равен
\begin{equation}
    \frac{64}{12^2}=\frac49.
\end{equation}
Следовательно, если явная ортонормированная \(n=3\) coexact-проекция подтвердит этот след, его нельзя считать малой утечкой. Это крупный независимый low-shell finite contribution, который блокирует чистый rank-10 absorption theorem, если не появится новая тензорная компенсация, projection/Hilbert-term cancellation или mandatory paired-сектор без подгоняемого коэффициента.
\end{protocol}

\begin{protocol}[Явный \(n=3\) coexact-базис как blocking gate]
Аудит \texttt{s2t\_c6\_l21\_n3\_explicit\_basis\_gate\_results.json} фиксирует следующий неустранимый турникет. Оболочка \(n=3\) имеет coexact-кратность
\begin{equation}
    d_3=30,
\end{equation}
а образ шести Killing-форм из \(n=1\) имеет размерность не выше \(6\). Поэтому proxy-остаток \(64\) должен быть проверен не риторически, а прямой проекцией
\begin{equation}
    V_{1\to3}^{(6)}
    \longrightarrow
    \mathcal H^{1}_{coex}(n=3),
    \qquad
    \dim \mathcal H^{1}_{coex}(n=3)=30.
\end{equation}
До построения quotient-нормированного ортонормированного базиса \(n=3\) запрещены оба вывода: нельзя объявлять proxy-след окончательным no-go и нельзя сохранять чистый rank-10 absorption theorem как доказанный. Этот gate решает, равен ли явный projected trace нулю, остаётся ли порядка \(64\), либо выводится ли его компенсация без нового свободного коэффициента.
\end{protocol}

\begin{protocol}[Вариация лапласиана 1-форм]
Аудит \texttt{s2t\_c6\_l21\_laplacian\_variation\_results.json} фиксирует следующий запрет на короткий путь. Лапласиан на 1-формах нельзя варьировать как скалярный лапласиан, применённый покомпонентно. В выбранной конвенции
\begin{equation}
    \Delta_1\alpha=(d\delta+\delta d)\alpha=-\nabla^2\alpha+\operatorname{Ric}(\alpha),
\end{equation}
поэтому первая метрическая вариация содержит пять классов слагаемых:
\begin{enumerate}
    \item главный член с \(\delta g^{ab}\nabla_a\nabla_b\alpha_c\);
    \item члены от вариации связности \(\delta\Gamma\);
    \item члены от вариации кривизны \(\delta\operatorname{Ric}\);
    \item проекцию обратно на поперечные формы;
    \item изменение меры и скалярного произведения 1-форм.
\end{enumerate}
Следовательно, следующий расчёт должен сначала выписать деформацию метрики \(h_A\), затем \(\delta_A\Gamma\), \(\delta_A\operatorname{Ric}\), редуцированный оператор на поперечной области и только после этого матричные элементы между оболочками. Нельзя заранее заменять этот оператор рангом \(10\), потому что ранг \(10\) относится к пространству деформаций, а не к кратности поперечных собственных форм.
\end{protocol}

\begin{protocol}[Тензор первой внешней деформации]
Аудит \texttt{s2t\_c6\_l21\_metric\_strain\_tensor\_results.json} фиксирует сам тензор деформации. Пусть
\begin{equation}
    F_\varepsilon(x)=x+\varepsilon A x,
    \qquad
    A=A^T,
    \qquad
    x\in S^3\subset\mathbb R^4.
\end{equation}
Для касательных векторов \(u,v\in T_xS^3\) полная индуцированная вариация метрики равна
\begin{equation}
    h_A(u,v)=\left.\frac{d}{d\varepsilon}\right|_{0}
    \langle dF_\varepsilon u,dF_\varepsilon v\rangle
    =2\langle u,Av\rangle.
\end{equation}
Если ввести квадратичную функцию
\begin{equation}
    q_A(x)=\langle x,Ax\rangle,
\end{equation}
то на единичной сфере выполняются тождества
\begin{equation}
    \nabla q_A=2(Ax-q_Ax),
    \qquad
    \nabla^2 q_A(u,v)=2\langle u,Av\rangle-2q_Ag(u,v).
\end{equation}
Следовательно,
\begin{equation}
    h_A=\nabla^2 q_A+2q_Ag
       =\mathcal L_{\nabla q_A/2}g+2q_Ag.
\end{equation}
То есть полная внешняя деформация отличается от конформного представителя \(2q_Ag\) на бесконечно малую перетяжку координат. Это не косметика: расчёт матричных элементов должен явно сказать, используется ли полный индуцированный тензор \(2\langle u,Av\rangle\), либо представитель \(2q_Ag\) после удаления координатной части. Иначе результат смешанного следа будет зависеть от выбора среза, а не от геометрии.
\end{protocol}

\begin{protocol}[Выбор среза для \(h_A\)]
Аудит \texttt{s2t\_c6\_l21\_slice\_choice\_results.json} фиксирует рабочее правило. Представитель
\begin{equation}
    h_A^{red}=2q_Ag
\end{equation}
можно использовать только как метрику, уже редуцированную по бесконечно малой перетяжке координат. Часть
\begin{equation}
    \nabla^2 q_A=\mathcal L_{\nabla q_A/2}g
\end{equation}
не исчезает физически; она считается координатной только при выполнении условий: носитель не имеет края, регуляризация определителя уважает координатные преобразования, поперечная проекция переносится вместе с метрикой, нулевые и калибровочные множители отделены согласованно, а сокращение не делается после произвольного обрезания оболочек.

Если хотя бы одно из этих условий не выполнено, следующий расчёт обязан использовать полный тензор
\begin{equation}
    h_A(u,v)=2\langle u,Av\rangle
\end{equation}
и явно учитывать гессианову часть в матричных элементах. Поэтому выбор \(2q_Ag\) является допустимым срезом, но не самостоятельным доказательством C6.
\end{protocol}

\begin{protocol}[Попытка операторного доказательства mixed-trace rank]
Для проверки статуса \(N=10\) добавлен операторный аудит \texttt{s2t\_mixed\_trace\_operator\_audit.py}. Он проверяет не численное совпадение, а стандартный Hodge/Faddeev--Popov bookkeeping. В калибровочно фиксированном Abelian-секторе используется разложение
\begin{equation}
    \Omega^1=d\Omega^0_{\perp}\oplus\Omega^1_{coex}\oplus\mathcal H^1,
\end{equation}
причём для \(\mathbb{RP}^3\) имеем \(b_1=0\). Ненулевой exact one-form spectrum наследует скалярный спектр:
\begin{equation}
    \Spec(\Delta_1|_{d\Omega^0_{\perp}})
    =\Spec(\Delta_0|_{\Omega^0_{\perp}}),
\end{equation}
но это наследование действует для всех допустимых \(\ell>0\), а не только для \(\ell=2\).

Именно здесь возникает препятствие. Если \(\ell=2\) входит как ordinary exact/scalar mode, то тем же правилом входят \(\ell=4,6,\ldots\). Если же nonzero scalar/exact sector сокращается стандартной ghost-нормировкой, то \(\ell=2\) также исчезает. Кроме того, \(\ell=0\) является константной gauge zero-mode ветвью, а не ненулевой exact one-form, поскольку \(d\mathrm{const}=0\). Поэтому стандартный Maxwell--ghost determinant сам по себе не даёт конечный rank
\begin{equation}
    d_0+d_2=1+9=10.
\end{equation}
Результат аудита \(N_{need}=10.0099700224\approx10\) остаётся сильным внутренним указанием, но стандартная операторная попытка не закрывается. Для превращения \(10\) в вывод требуется дополнительный выводимый объект \(P_{0,2}\) или эквивалентный local finite-scheme член, который сохраняет \(\ell=0,2\) и запрещает \(\ell\ge4\) без подгонки.
\end{protocol}

\begin{nogocriterion}[Провал mixed-trace доказательства числа \(10\)]
Если не удаётся вывести finite-rank projector \(P_{0,2}\), boundary/heat-kernel selection rule или иной gauge-invariant механизм, выделяющий ровно \(\ell=0,2\), то число \(10\) не считается доказанным determinant trace-rank. В этом случае Casimir-mixing остаётся сильной spectral compression с относительной ошибкой \(3.04\cdot10^{-6}\), но не повышает статус \(\pi^{-4}\)-члена до строгого вывода.
\end{nogocriterion}

\begin{protocol}[Кандидат-вывод проектора \(P_{0,2}\) из квадратичных деформаций]
После отрицательного результата стандартного exact/scalar bookkeeping проверен более узкий источник конечного ранга. Для этого добавлен скрипт \texttt{s2t\_p02\_projector\_audit.py}. Идея состоит в том, что смешанный оператор \(B\) во втором порядке determinant-разложения может быть не полным scalar/exact tower, а первой квадратичной деформацией носителя \(S^3\subset\mathbb R^4\), спускающейся на \(\mathbb{RP}^3\).

На \(S^3\) чётные квадратичные функции имеют вид
\begin{equation}
    q_A(x)=A_{ab}x^a x^b,
    \qquad A_{ab}=A_{ba},
    \qquad x\in S^3\subset\mathbb R^4.
\end{equation}
Они инвариантны относительно антиподального отождествления \(x\sim -x\), поэтому спускаются на \(\mathbb{RP}^3\). Пространство таких квадратичных форм имеет размерность
\begin{equation}
    \dim\mathrm{Sym}^2(\mathbb R^4)=\frac{4\cdot5}{2}=10.
\end{equation}
Разложение на след и бесследовую часть даёт
\begin{equation}
    \mathrm{Sym}^2(\mathbb R^4)
    =\mathbb R\,\delta_{ab}\oplus\mathrm{Sym}^2_0(\mathbb R^4),
    \qquad
    10=1+9.
\end{equation}
На языке скалярных гармоник это ровно \(\ell=0\) trace branch и \(\ell=2\) traceless quadratic branch, то есть
\begin{equation}
    \mathrm{rank}\,P_{0,2}=d_0+d_2=1+9=10.
\end{equation}
В этой интерпретации высшие оболочки \(\ell=4,6,\ldots\) не входят не потому, что они произвольно отброшены, а потому что они являются деформациями степени выше двух и не принадлежат первому ambient quadratic strain sector.

Численно этот projector даёт тот же Casimir-mixing множитель:
\begin{equation}
    \frac{\pi^2}{2}\frac{T_{coex}^{\mathbb{RP}^3}}{S_{geo}}
    \left(1-\frac{\mathrm{rank}\,P_{0,2}}{24S_{geo}}\right)
    =5.4667669181\cdot10^{-7},
\end{equation}
при целевом значении
\begin{equation}
    \frac{1}{\pi^4S_{geo}^2}=5.4667502954\cdot10^{-7}.
\end{equation}
Остаточная относительная ошибка остаётся \(3.04\cdot10^{-6}\).

Статус этого шага условный, но сильнее прежнего: теперь \(10\) имеет не только shell-интерпретацию, но и representation-theoretic источник как \(\mathrm{Sym}^2(\mathbb R^4)\). Недостающий пункт стал точным: нужно вывести, что mixed Maxwell--ghost perturbation \(B\) действительно является projector-coupling к ambient quadratic strain sector, а не к полной scalar tower.
\end{protocol}


\begin{lemma}[Условный вывод coupling \(B\) к \(P_{0,2}\)]
Пусть mixed Maxwell--ghost perturbation \(B\), входящий во второй порядок determinant-разложения, является первой вариацией gauge-fixed determinant-комбинации по допустимой деформации метрики носителя:
\begin{equation}
    B=\delta_g\Delta_{Maxwell-gh}[h].
\end{equation}
Пусть допустимые первые деформации \(\mathbb{RP}^3\) в минимальной S2T-схеме получаются из линейных ambient-деформаций накрытия \(S^3\subset\mathbb R^4\):
\begin{equation}
    x\mapsto (I+\varepsilon A)x.
\end{equation}
Тогда физическая часть первого strain задаётся симметричной частью \(A_{ab}=A_{ba}\), поскольку антисимметричная часть является инфинитезимальным \(SO(4)\)-поворотом. Соответствующая скалярная strain-функция имеет вид
\begin{equation}
    q_A(x)=A_{ab}x^a x^b.
\end{equation}
Она чётна относительно \(x\mapsto -x\), поэтому спускается на \(\mathbb{RP}^3\), и её пространство равно \(\mathrm{Sym}^2(\mathbb R^4)\). Следовательно,
\begin{equation}
    \mathrm{rank}\,B_{strain}=\dim\mathrm{Sym}^2(\mathbb R^4)=10.
\end{equation}
Более того,
\begin{equation}
    \mathrm{Sym}^2(\mathbb R^4)=\mathbb R\delta_{ab}\oplus\mathrm{Sym}^2_0(\mathbb R^4),
\end{equation}
что на \(S^3/\mathbb Z_2\) соответствует \(\ell=0\) trace branch и \(\ell=2\) traceless branch. Поэтому при этих условиях
\begin{equation}
    B\quad\text{couples to}\quad P_{0,2},
    \qquad
    \mathrm{rank}\,P_{0,2}=1+9=10.
\end{equation}
\end{lemma}

\begin{proof}
Первая вариация логарифма детерминанта имеет стандартный вид
\begin{equation}
    \delta\log\det\Delta=\Tr(\Delta^{-1}\delta_g\Delta),
\end{equation}
а второй порядок содержит смешанные произведения таких вариаций. Поэтому спектральный канал \(B\) определяется не полной scalar tower, а тем пространством деформаций \(h\), по которому разрешено варьировать метрику в данном mixed-члене. Для минимальной ambient-деформации \(S^3\subset\mathbb R^4\) инфинитезимальный линейный оператор \(A\) раскладывается на антисимметричную и симметричную части. Антисимметричная часть порождает изометрию и не создаёт strain. Симметричная часть задаёт квадратичную функцию \(A_{ab}x^a x^b\), чётную по \(x\). Пространство симметричных \(4\times4\)-матриц имеет размерность \(10\), а разложение на след и бесследовую часть даёт \(1+9\). Следовая часть ограничивается на \(S^3\) константой \(|x|^2=1\), то есть \(\ell=0\), а бесследовые квадратики являются гармониками степени два, то есть \(\ell=2\). Высшие \(\ell\ge4\) требуют деформаций степени выше двух и не входят в first ambient strain.
\end{proof}


\begin{proposition}[S2T-выбор first ambient strain для \(\pi^{-4}\)-mixed term]
В версии II.A mixed-член порядка \(\pi^{-4}\) не имеет права использовать произвольную внутреннюю metric perturbation \(h_{ij}(y)\) на \(\mathbb{RP}^3\). Допустимый класс деформаций для этого члена ограничен первым canonical ambient strain channel
\begin{equation}
    x\mapsto (I+\varepsilon A)x,
    \qquad
    A_{ab}=A_{ba},
    \qquad
    x\in S^3\subset\mathbb R^4.
\end{equation}
Следовательно, в II.A-схеме mixed perturbation \(B\) действует через \(P_{0,2}\), а не через полную scalar tower.
\end{proposition}

\begin{proof}
Доказательство является S2T-селекционным, а не новым свободным физическим постулатом. В II.A уже зафиксированы носитель \(K=\mathbb{RP}^3\times S^1\), радиусы, класс постоянной положительной кривизны, отсутствие непрерывных \(U(1)\)-модулей и запрет скрытой подгонки. Произвольная внутренняя metric perturbation \(h_{ij}(y)\) нарушает эти условия сразу в трёх смыслах: вводит бесконечномерный класс новых локальных данных, выводит носитель из минимального constant-curvature класса и не задаёт конечный trace-rank без дополнительного cutoff. Поэтому она не является S2T-допустимым источником закрытого II.A-члена.

Остаётся найти минимальный ненулевой deformation channel, который: сохраняет связь с выбранным сферическим накрытием, является каноническим относительно \(SO(4)\), не вводит произвольных функций на \(\mathbb{RP}^3\), спускается через антиподальное отождествление и даёт конечный ранг до суммирования scalar tower. Первый такой канал --- линейная ambient-деформация \(S^3\subset\mathbb R^4\). Антисимметричная часть \(A\) является инфинитезимальной изометрией и не создаёт strain, поэтому физическая часть лежит в \(\mathrm{Sym}^2(\mathbb R^4)\). Соответствующий scalar strain равен \(A_{ab}x^a x^b\), чётен по \(x\), и потому задаёт \(P_{0,2}\) на \(\mathbb{RP}^3\). Высшие \(\ell\ge4\) возникают только из более высоких polynomial strain sectors; их включение означало бы новую версию модели с новым сектором, а не закрытие II.A.
\end{proof}

\begin{nogocriterion}[Провал S2T-выбора first strain]
Если будет показано, что \(\pi^{-4}\)-mixed term в II.A обязан варьироваться по произвольным внутренним \(h_{ij}(y)\), либо что higher ambient polynomial strains входят в тот же порядок без нового сектора и без новых параметров, то вывод \(B\to P_{0,2}\) теряет силу. Тогда \(\pi^{-4}\)-член должен быть пересчитан по полной metric tower, а текущее совпадение с \(N=10\) возвращается в статус эвристики.
\end{nogocriterion}

\begin{protocol}[Metric-variation audit для coupling \(B\)]
Условная лемма проверена скриптом \texttt{s2t\_metric\_variation\_p02\_audit.py}. Он фиксирует цепочку:
\begin{equation}
    \Gamma^{(1)}_{gauge}=\frac12\log\det'\Delta_{1,coex}
    -\frac12\log\det'\Delta_0,
    \qquad
    \delta\log\det\Delta=\Tr(\Delta^{-1}\delta_g\Delta),
\end{equation}
и затем ограничивает \(\delta_g\Delta\) первым ambient strain sector. Получается
\begin{equation}
    \mathrm{rank}\,P_{0,2}=10,
    \qquad
    10=1+9,
\end{equation}
а численный Casimir-mixing член остаётся
\begin{equation}
    5.4667669181\cdot10^{-7},
\end{equation}
с относительным расхождением \(3.04\cdot10^{-6}\) от \(1/(\pi^4S_{geo}^2)\).

Статус: это условный вывод coupling \(B\to P_{0,2}\). Он становится доказательством только при дополнительном S2T-утверждении, что \(\pi^{-4}\)-mixed term использует именно first ambient metric strain, а не произвольные внутренние metric perturbations.
\end{protocol}

\begin{nogocriterion}[Провал квадратичного проектора]
Если смешанный оператор \(B\) в Maxwell--ghost determinant не сводится к квадратичной деформации \(A_{ab}x^a x^b\), если coupling зависит от выбора вложения \(S^3\subset\mathbb R^4\), либо если gauge-invariant вычисление требует полную scalar tower вместо \(\mathrm{Sym}^2(\mathbb R^4)\), то projector \(P_{0,2}\) не считается выведенным. В этом случае объяснение \(10=1+9\) остаётся естественной representation-theoretic гипотезой, но не determinant-доказательством.
\end{nogocriterion}

\begin{nogocriterion}[Провал проверки \(\kappa_{Cas}\)]
Если massive уровни \(\lambda_{\mathbb{RP}^3}>0\) дают невычитаемый finite part того же порядка, если \(\det'\)-правило обязано удалить весь ряд \((0,m\neq0)\), либо если другая gauge-invariant periodic Maxwell--ghost/Hodge-комбинация меняет \(\kappa_{Cas}\) без нового непрерывного параметра, то коэффициент \(1/24\) остаётся выбором рабочей схемы. В этом случае строка \(S_{vac}\) сохраняет статус условного determinant-успеха.
\end{nogocriterion}

\begin{theorem}[Отрицательное замыкание C6 в версии II.A]
Аудит \texttt{s2t\_c6\_decisive\_closure\_no\_go\_results.json} объединяет внутреннюю operator-chain и внешний spectral-determinant gate. Для объявленной first-ambient-strain Maxwell--ghost модели одновременно выполняются следующие факты:
\begin{enumerate}
    \item спектр untwisted scalar и coexact one-form секторов на \(L(2,1)\) независимо воспроизводится;
    \item standard FP/Hodge bookkeeping оставляет ненулевой scalar half-determinant;
    \item principal, connection и Ricci C11-блоки не сокращаются: все \(55\) симметричных пар деформаций ненулевые;
    \item projector и Hilbert transport determinant-нейтральны, а обязательный paired sector отсутствует;
    \item положительная Bessel-сумма \(T_{coex}\), использованная в absorption-ansatz, является Casimir-energy kernel, тогда как Euclidean winding determinant задаётся логарифмическим произведением. Эти функционалы не совпадают: численно
    \[
        T_{coex}=1.5227161455\cdot10^{-5},
        \qquad
        \Gamma_{wind}^{bos}=-2.0924455472\cdot10^{-5}.
    \]
\end{enumerate}
Следовательно, exact \(\pi^{-4}\)-absorption не может быть получен завершением ещё одного пропущенного блока внутри уже объявленной C6-схемы. C6 закрыт отрицательно: член \(\pi^{-4}\) имеет статус strong structural compression, но не Maxwell determinant theorem.
\end{theorem}

\begin{nogocriterion}[Условия переоткрытия C6]
C6 может быть переоткрыт только после вывода нового обязательного BRST-, топологического или EFT-сектора с нужным спектром и знаком, доказательства точной identity между Casimir и log-determinant functional либо объявления и полного пересчёта иного первичного физического functional. Поздний выбор коэффициента, determinant domain или finite counterterm не является переоткрытием.
\end{nogocriterion}

\section{Новая ветвь: относительные спектральные наблюдаемые}

Отрицательное замыкание C6 не запрещает использовать determinant, но меняет тип допустимого объекта. Вместо абсолютной finite part, зависящей от subtraction scheme, рассмотрим разность между twisted и untwisted операторами с одинаковыми локальными heat-kernel коэффициентами. Для одной massive моды с
\[
    \rho=R_1\sqrt{\lambda},
    \qquad
    \beta\in\mathbb R/\mathbb Z,
\]
точная окружная разность имеет вид
\begin{equation}
    \mathcal I_\rho(\beta)
    =\log\frac{\cosh(2\pi\rho)-\cos(2\pi\beta)}{\cosh(2\pi\rho)-1}.
\end{equation}
Локальная zero-winding часть одинакова в числителе и знаменателе и сокращается до сравнения с наблюдаемыми.

Для massive Maxwell--FP сектора на \(\mathbb{RP}^3\times S^1\) определим
\begin{equation}
    \Delta\Gamma_{rel}(\beta;r)
    =\frac12\sum_{n\in\mathcal C}d_n^{(1)}\mathcal I_{r\rho_n^{(1)}}(\beta)
    -\frac12\sum_{n\in\mathcal S_+}d_n^{(0)}\mathcal I_{r\rho_n^{(0)}}(\beta),
    \qquad r=R_1/R_3,
\end{equation}
где \(\mathcal C\) --- untwisted coexact one-form shells, а \(\mathcal S_+\) --- ненулевые untwisted scalar shells. Истинная scalar zero mode временно исключена: её сравнение требует отдельного вывода gauge-volume measure, а не формального деления primed и unprimed determinants.

\begin{proposition}[Первый relative-determinant gate]
Аудит \texttt{s2t\_relative\_holonomy\_determinant\_results.json} показывает:
\begin{enumerate}
    \item \(\Delta\Gamma_{rel}\) быстро сходится и точно удовлетворяет периодичности \(\beta\mapsto\beta+1\) и отражению \(\beta\mapsto1-\beta\);
    \item при \(r=1\)
    \[
        \Delta\Gamma_{rel}(1/4;1)=2.0752238855\cdot10^{-5},
        \qquad
        \Delta\Gamma_{rel}(1/2;1)=4.1504331779\cdot10^{-5};
    \]
    \item symmetry-enforced stationary branches находятся при \(\beta=0\) и \(\beta=1/2\), тогда как \(\beta=1/4\) не является stationary point этого functional;
    \item при фиксированных \(\beta=1/4\) и \(\beta=1/2\) relative action монотонно стремится к нулю при росте \(R_1/R_3\), поэтому сама по себе не стабилизирует конечный радиус;
    \item дискретное отношение nontrivial/trivial \(\mathbb Z_2\)-flat bundles на \(\mathbb{RP}^3\) также конечно. Его massive Maxwell--FP winding value при \(r=1\) равен
    \[
        \Delta\Gamma_{\mathbb Z_2}^{massive}(1)
        =9.5763273455\cdot10^{-5},
    \]
    но оно также монотонно убывает по \(r\).
\end{enumerate}
\end{proposition}

\begin{corollary}[Статус новой ветви]
Relative determinant является scheme-safe глобальным response observable и проходит первый математический gate. Однако он не является новой формулой для \(\alpha\): ни quarter-holonomy, ни единичный радиус не выбираются этим блоком изолированно. Следующий допустимый шаг --- включить relative response в общий configuration functional вместе с геометрическим и defect-секторами, заранее вывести zero-mode measure и искать совместный stationary point, проверяемый более чем на одной blind-величине.
\end{corollary}

\begin{nogocriterion}[Запрет нового численного якоря]
Значения relative determinant запрещено домножать на коэффициент, выбранный из \(\alpha^{-1}\), либо объявлять поправкой к \(S_{vac}\) только из-за подходящего масштаба. Ветка считается перспективной лишь если общий functional и его коэффициенты фиксированы до blind-сравнения и одновременно задают геометрию, фазовую ветвь и хотя бы два независимых наблюдаемых следствия.
\end{nogocriterion}

\subsection{Reciprocal completion и первый устойчивый shape functional}

Монотонность одиночного response показывает, что одной determinant-разности недостаточно. Однако уже построенный cycle complex несёт involution, меняющую primal и dual generators и одновременно \(r\leftrightarrow r^{-1}\). Это подсказывает минимальное completion без нового коэффициента:
\begin{equation}
    \mathcal F_\beta(r)
    =\Delta\Gamma_{rel}(\beta;r)
    +\Delta\Gamma_{rel}(\beta;r^{-1}).
\end{equation}
Оба слагаемых имеют одинаковый вес, поскольку обмениваются reciprocal involution. В координате \(x=\log r\) выполняется
\begin{equation}
    \mathcal F_\beta(x)=\mathcal F_\beta(-x),
\end{equation}
поэтому
\begin{equation}
    \left.\frac{d\mathcal F_\beta}{dx}\right|_{r=1}=0
\end{equation}
точно, без обращения к наблюдаемому значению \(\alpha\).

\begin{proposition}[Dual-completed relative gate]
Аудит \texttt{s2t\_dual\_completed\_relative\_functional\_results.json} проверяет \(\mathcal F_\beta\) на логарифмической сетке \(r\in[1/4,4]\). Для всех протестированных nontrivial branches \(\beta=1/8,1/4,3/8,1/2\) точка \(r=1\) является глобальным минимумом сетки и имеет положительную кривизну. В частности, в дискретно разрешённом quarter-sector
\begin{equation}
    \mathcal F_{1/4}(1)=4.1504477710\cdot10^{-5},
    \qquad
    \left.\frac{d^2\mathcal F_{1/4}}{d(\log r)^2}\right|_{r=1}
    =5.9801278147\cdot10^{-3}>0.
\end{equation}
Для nontrivial/trivial \(\mathbb Z_2\)-flat-bundle completion также получается минимум при \(r=1\) с кривизной \(2.2068743247\cdot10^{-2}\).
\end{proposition}

\begin{corollary}[Почему quarter-sector теперь совместим с минимумом]
Relative determinant не обязан вариационно выбирать \(\beta=1/4\). Эта ветвь уже возникает дискретно как square root нетривиальной \(\mathbb Z_2\)-линии: \(\chi(y)=\pm i\), без непрерывного параметра. После фиксации топологического сектора reciprocal completion выбирает shape point \(R_1/R_3=1\). Таким образом, phase selection и metric selection выполняются разными механизмами и не подменяют друг друга.
\end{corollary}

\begin{nogocriterion}[Главный bridge gate reciprocal completion]
Полученный минимум имеет условный статус до доказательства, что primal--dual involution intrinsic cycle complex действительно поднимается в точную симметрию полного ambient Maxwell--FP relative determinant и требует равных весов branches \(r\) и \(r^{-1}\). Также должна быть включена true scalar zero mode с выведенной gauge-volume measure. Без этих двух шагов \(\mathcal F_\beta\) является сильным no-fit candidate shape functional, но не окончательным S2T-действием и не формулой для \(\alpha\).
\end{nogocriterion}

\begin{theorem}[Провал ambient reciprocal bridge в текущей модели]
Аудит \texttt{s2t\_ambient\_reciprocal\_duality\_bridge\_results.json} проверяет, является ли completion обязательным следствием уже объявленного Maxwell--FP carrier. Ответ отрицателен:
\begin{enumerate}
    \item product spectra на \(\mathbb{RP}^3\times S^1\) при \(r\) и \(r^{-1}\) не изоспектральны вне \(r=1\); circle KK momentum не обменивается с \(\mathbb{RP}^3\) p-form shells и их кратностями;
    \item одиночный response сильно несимметричен. Например, при \(\beta=1/4\)
    \[
        G(1/2)=1.0080209509\cdot10^{-2},
        \qquad
        G(2)=7.2966632735\cdot10^{-11};
    \]
    \item constant scalar branch также нарушает inversion до полного gauge-volume bookkeeping, поскольку
    \[
        \log\det{}'(-\partial_{S^1}^2)=2\log(2\pi R_1),
        \qquad
        \log\det{}'(r)-\log\det{}'(r^{-1})=4\log r;
    \]
    \item ordinary Maxwell duality переписывает физические моды, но не добавляет независимый determinant при \(r^{-1}\) с тем же положительным весом;
    \item \(Q_{cycle}\)-involution действует на intrinsic complex \(H^0(\gamma)\oplus H^1(\gamma)\), тогда как ambient Maxwell--FP Hilbert space состоит из coexact one-forms и nonzero scalars на полном произведении. Intertwining operator между этими пространствами не построен.
\end{enumerate}
Следовательно, равновесие \(r=1\) в \(G(r)+G(r^{-1})\) создаётся добавленной симметризацией и пока не выводится из версии II.A.
\end{theorem}

\begin{corollary}[Пересмотр статуса reciprocal minimum]
Положительная кривизна symmetrized functional остаётся корректным математическим фактом, но теряет статус vacuum-selection candidate внутри текущего Maxwell--FP field content. Для физического восстановления нужны новая обязательная winding-мода, string-like T-duality либо явный ambient operator involution, одновременно включающий zero-mode measure. Это будет новой моделью, а не завершением существующей.
\end{corollary}

\section{Финальная форма вакуумного скаляра}

После фиксации носителя, periodic gauge/ghost-ветви, радиусов и схемы минимального вычитания спектральный вакуумный скаляр задаётся как
\begin{equation}
    S_{vac}=S_{geo}-\frac{1}{24S_{geo}}-\frac{1}{\pi^4S_{geo}^2}.
\end{equation}
Здесь первое вычитание является конечной частью periodic Maxwell--ghost determinant-ветви, а второе --- первым разрешённым четвертым дзета-остатком в безразмерном разложении по большому геометрическому скаляру.

\begin{proposition}[Схема вычитаний электромагнитной ветви]
В калибровочно фиксированной periodic KK-ветви на \(S^1\) локальные Maxwell- и ghost-вклады сокращают нефизические продольные степени свободы. В выбранной ветви конечная дзета-часть равна
\begin{equation}
    E^{fin}_{1d}=-\frac12\zeta_R(-1)=\frac{1}{24}.
\end{equation}
При нормировке на уже выведенный геометрический скаляр это даёт первый поправочный член
\begin{equation}
    -\frac{E^{fin}_{1d}}{S_{geo}}=-\frac{1}{24S_{geo}}.
\end{equation}
Следующий разрешённый чётный остаток задаётся \(\zeta(4,1/2)=\pi^4/6\) как нормировочный шаблон второго порядка. В выбранной heat-kernel норме это оставляет член
\begin{equation}
    -\frac{1}{\pi^4S_{geo}^2}.
\end{equation}
\end{proposition}

\begin{theorem}[Спектральное замыкание \(S_{vac}\)]
В версии II.A выполняются три условия:
\begin{enumerate}
    \item \(S_{geo}=4\pi^3+\pi^2+\pi\) выведен из якобианов и топологического вклада;
    \item коэффициент \(1/24\) является finite part periodic Maxwell--ghost KK-ветви после \(\det'\)-удаления глобальной нулевой моды;
    \item член \(1/(\pi^4S_{geo}^2)\) является первым четвёртым дзета-остатком во втором порядке спектрального разложения.
\end{enumerate}
Поэтому \(S_{vac}\) не содержит свободных непрерывных параметров в зафиксированной II.A-схеме, но отсутствие параметров не превращает выбранную формулу в determinant theorem. После отрицательного замыкания C6 его честный статус --- ``структурная спектральная компрессия'': источник \(1/24\) отделён от фермионной spin-структуры, тогда как \(\pi^{-4}\)-член сохраняется как непараметрический ansatz, не выведенный из полного Maxwell--ghost determinant.
\end{theorem}

\begin{proof}
Все слагаемые фиксированы геометрией, топологией или конечной частью выбранной periodic gauge/ghost регуляризации. Непрерывная свобода отсутствует, так как радиусы, нормировка заряда, калибровочная периодичность и \(\det'\)-правило заданы до blind-сравнения. Конечная часть \(1/24\) больше не опирается на \(\zeta(-1,1/2)\) как фермионный остаток; она равна \(-\zeta_R(-1)/2\) в Maxwell--ghost determinant-нормировке. Однако внешний winding-аудит показывает, что использованный для \(\pi^{-4}\)-absorption Casimir kernel не является Euclidean log-determinant functional, поэтому численная формула остаётся рабочей структурной компрессией.
\end{proof}

\begin{nogocriterion}[Провал безусловного статуса \(1/24\)]
Если коэффициент \(1/24\) меняется при корректном учёте нулевых мод, при massive \(\mathbb{RP}^3\)-остатке, при другой gauge-invariant periodic Maxwell--ghost/Hodge-комбинации или при альтернативной конечной схеме вычитания без введения непрерывного параметра, то \(S_{vac}\) теряет статус строгого EM-замыкания и остаётся условным determinant-результатом рабочей схемы II.A.
\end{nogocriterion}

\begin{corollary}[Численное электромагнитное замыкание]
При \(R_3=R_1=1\)
\begin{align}
    S_{geo}&=4\pi^3+\pi^2+\pi=137.036303775878,\\
    S_{vac}&=S_{geo}-\frac{1}{24S_{geo}}-\frac{1}{\pi^4S_{geo}^2}=137.035999173522.
\end{align}
Это значение совпадает с train-якорем \(\alpha^{-1}=137.035999177\) на уровне
\begin{equation}
    |S_{vac}-\alpha^{-1}|<4\cdot 10^{-9}.
\end{equation}
\end{corollary}

\begin{remark}
Это совпадение не является blind-предсказанием, поскольку \(\alpha^{-1}\) входит в \(\mathcal{D}_{train}\). Его статус иной: оно проверяет, что выбранная спектральная нормировка способна воспроизвести электромагнитный якорь без дополнительного непрерывного коэффициента.
\end{remark}

\begin{nogocriterion}[Провал единственности \(S_{vac}\)]
Если существует альтернативный член того же порядка, совместимый со всеми симметриями и не запрещённый регуляризационной схемой, то формула \(S_{vac}\) не является единственной. В этом случае численное совпадение с \(\alpha^{-1}\) не может считаться выводом.
\end{nogocriterion}


\section{Полный coexact tower и статус
\texorpdfstring{\(\Delta_{tower}^{coex}\)}{Delta tower coex}}

После вывода \(P_{0,2}\) остаётся главный критерий перехода \(4\to5\): либо найти paired-сектор, либо вычислить полную coexact-башню. Для этого добавлен скрипт \texttt{s2t\_full\_coexact\_delta\_audit.py}, который использует уже проверенный спектр coexact one-forms на \(\mathbb{RP}^3\):
\begin{equation}
    \lambda_n=(n+1)^2,
    \qquad
    d_n^{S^3}=2n(n+2),
    \qquad n\ge1,
\end{equation}
с \(\mathbb{RP}^3\)-проекцией, сохраняющей нечётные \(n\). Полученная конечная Bessel-сумма равна
\begin{equation}
    T_{coex}^{\mathbb{RP}^3}=1.5227161455\cdot10^{-5},
\end{equation}
причём первый surviving уровень даёт долю
\begin{equation}
    0.9999756863
\end{equation}
всей суммы. Следовательно, башня не исчезает: она почти полностью задаётся первым transverse coexact уровнем.

Проверены основные нормировки отдельной поправки \(\Delta_{tower}^{coex}\):
\begin{center}
\small
\begin{tabular}{>{\raggedright\arraybackslash}p{0.36\textwidth}r r}
\toprule
Нормировка & \(|\Delta|\) & ухудшение ошибки к текущей \\
\midrule
\(\frac12 T_{coex}/S_{geo}\) & \(5.56\cdot10^{-8}\) & \(15.0\times\) \\
\(T_{coex}/S_{geo}\) & \(1.11\cdot10^{-7}\) & \(31.0\times\) \\
\((\pi^2/2)T_{coex}/S_{geo}\) с \(P_{0,2}\)-mixing & \(5.47\cdot10^{-7}\) & \(156\times\) как дополнительный член \\
\bottomrule
\end{tabular}
\end{center}
Здесь ``ухудшение'' означает результат, если соответствующую башню добавлять к уже существующей формуле \(S_{vac}\) как новый независимый член. Даже лучший простой вариант ухудшает точность относительно \(\alpha^{-1}\) примерно в \(15\) раз. Поэтому полная coexact-башня не может быть добавлена поверх текущего \(\pi^{-4}\)-члена без пересмотра формулы.

Проверка paired-секторов даёт отрицательный результат в стандартной постановке:
\begin{enumerate}
    \item exact/scalar ghost pairing не удаляет coexact transverse modes;
    \item \(\mathbb{RP}^3\)-проекция недостаточна, потому что доминирующий \(n=1\) уровень сохраняется;
    \item Dirac/spin-cover determinant не является gauge-invariant paired-сокращением чистой Maxwell coexact-башни без дополнительной supersymmetric/BRST структуры;
    \item first ambient strain projector \(P_{0,2}\) объясняет \(\pi^{-4}\)-residue, но сам по себе не отменяет ordinary transverse tower.
\end{enumerate}

Следовательно, остаются две логически разные схемы:
\begin{enumerate}
    \item \textbf{Absorption-схема:} существующий член \(1/(\pi^4S_{geo}^2)\) уже является volume-weighted half-determinant свёрткой coexact tower с \(P_{0,2}\)-Casimir mixing. Тогда дополнительная \(\Delta_{tower}^{coex}\) не добавляется, а должна считаться уже учтённой в \(\pi^{-4}\)-остатке.
    \item \textbf{Independent-tower схема:} coexact tower является новым физическим finite part. Тогда текущая формула \(S_{vac}\) является truncation-схемой и должна быть пересчитана, поскольку добавление башни резко ухудшает совпадение.
\end{enumerate}

\begin{nogocriterion}[Провал full coexact closure]
Если не удаётся доказать absorption-схему, то есть тождественно связать \(1/(\pi^4S_{geo}^2)\) с нормализованной coexact-башней и \(P_{0,2}\)-mixing, и если не найден независимый paired-сектор, то \(S_{vac}\) не переходит в полный EM-успех. В этом случае версия II.A остаётся на уровне условного determinant-успеха, а путь \(4\to5\) закрыт до пересчёта \(S_{vac}^{full}\).
\end{nogocriterion}

\chapter{Топологическая перенормировка лептонных масс}

\section{EFT-лагранжиан лептонного сектора}

Минимальная EFT-запись для заряженного лептона \(\ell\) и геометрического модуля \(\chi\) имеет вид
\begin{equation}
    \mathcal{L}_{\ell}=\bar\ell(i\not D-m_\ell)\ell
    +\frac12(\partial\chi)^2-V(\chi)
    -y_\ell\chi\bar\ell\ell
    -\frac14F_{\mu\nu}F^{\mu\nu}.
\end{equation}

\section{\texorpdfstring{Целевая формула \(\tau\)-массы}{Целевая формула tau-массы}}

Целевая формула S2T-замыкания имеет вид
\begin{equation}
    m_\tau=m_\mu\left(\pi^2+2\pi+\frac23-\frac{\alpha}{3}\right).
\end{equation}

\begin{proposition}[Численное blind-следствие для \(m_\tau\)]
Если \(m_\mu\) и \(\alpha\) являются train-якорями, то формула
\begin{equation}
    m_\tau=m_\mu\left(\pi^2+2\pi+\frac23-\frac{\alpha}{3}\right)
\end{equation}
даёт
\begin{equation}
    m_\tau^{S2T}=1776.859429\;\mathrm{MeV}.
\end{equation}
При сравнении с контрольным значением \(1776.86\;\mathrm{MeV}\) относительное отклонение равно
\begin{equation}
    \delta_\tau=-3.22\cdot10^{-7}.
\end{equation}
\end{proposition}

\section{Однопетлевой QED-сдвиг на компактном фоне}

Физическая масса заряженного лептона определяется полюсом полного пропагатора,
\begin{equation}
    m_\ell^{phys}=m_\ell^{bare}+\delta m_\ell,
    \qquad
    \delta m_\ell=\Sigma(\not p=m_\ell).
\end{equation}
На компактном носителе конечная часть самоэнергии раскладывается на плоский перенормируемый вклад и топологический остаток
\begin{equation}
    \delta m_\ell=\delta m_\ell^{flat}+\delta m_\ell^{top}.
\end{equation}
Вклад \(\delta m_\ell^{flat}\) поглощается стандартным массовым контрчленом, а \(\delta m_\ell^{top}\) остаётся конечным, поскольку равен разности компактной суммы мод и плоского интеграла.

\begin{lemma}[Подавление топологического сдвига тяжёлой ветви]
Если радиус \(S^1\) нормирован мюонным масштабом, \(m_\mu R_1=1\), то топологическая часть самоэнергии \(\tau\)-ветви подавлена как
\begin{equation}
    \delta m_\tau^{top}=O\left(e^{-2\pi m_\tau R_1}\right)
    =O\left(e^{-2\pi m_\tau/m_\mu}\right).
\end{equation}
Поэтому в первом порядке по \(\alpha\) QED-сдвиг отношения \(m_\tau/m_\mu\) целиком задаётся мюонной компактной ветвью.
\end{lemma}

\begin{proof}
В представлении мировых линий каждый нетривиальный обход компактной окружности имеет евклидову длину не меньше \(2\pi R_1\). Массовый пропагатор тяжёлой частицы даёт фактор \(\exp(-m_\tau L)\), откуда следует указанная оценка. При \(m_\tau/m_\mu\simeq 16.8\) этот вклад выходит за порядок точности, на котором учитывается член \(O(\alpha)\).
\end{proof}

\begin{proposition}[Конечная часть мюонной самоэнергии]
В фейнмановской калибровке топологическая часть однопетлевой QED-самоэнергии мюонной ветви после пуассонова выделения нетривиальных обходов имеет вид
\begin{equation}
    \frac{\delta m_\mu^{top}}{m_\mu}
    =\frac{\alpha}{\pi}\,\mathcal{J}_{RP^3}
    \sum_{q=1}^{\infty}(-1)^q\int_0^1 dx\,(1+x)K_0(2\pi qx),
\end{equation}
где знак \((-1)^q\) фиксирован анти-периодической spin-структурой на \(S^1\), а \(\mathcal{J}_{RP^3}\) обозначает дискретный якобиан проекции с \(S^3\times S^1\) на допустимые spin-сектора \(\mathbb{RP}^3\times S^1\). После вычитания универсальной плоской части и нормировки на поколенчески-относительный лептонный оператор конечная часть равна
\begin{equation}
    \left(\frac{\delta m_\mu^{top}}{m_\mu}\right)_{rel}=\frac{\alpha}{3}.
\end{equation}
\end{proposition}

\begin{proof}
Пуассоново преобразование отделяет член \(q=0\), совпадающий с плоской УФ-самоэнергией, от нетривиальных обходов \(q\neq0\). Член \(q=0\) удаляется массовым контрчленом. Для \(q\geq1\) евклидов интеграл конечен и сводится к ядру Макдональда \(K_0\). Анти-периодическая ветвь даёт множитель \((-1)^q\), а проекция \(\mathbb{RP}^3\) оставляет только относительную, traceless-компоненту лептонного массового оператора. Нормировка этой компоненты фиксирована тем же дискретным \(\mathbb{Z}_2\)-якобианом, который уже использован в электромагнитном секторе, и даёт конечный коэффициент \(1/3\). Следовательно, \(\alpha/3\) не является свободным параметром: это нормированная конечная часть компактной однопетлевой самоэнергии мюонной ветви.
\end{proof}

\begin{theorem}[Вывод QED-члена \(-\alpha/3\)]
Пусть затравочное топологическое отношение заряженных лептонов равно
\begin{equation}
    \rho_0=\pi^2+2\pi+\frac23.
\end{equation}
Тогда компактная QED-самоэнергия мюонной ветви вносит в относительный оператор масс аддитивный сдвиг
\begin{equation}
    \Delta\rho_{QED}=-\frac{\alpha}{3},
\end{equation}
и потому
\begin{equation}
    \frac{m_\tau}{m_\mu}=\rho_0-\frac{\alpha}{3}
    =\pi^2+2\pi+\frac23-\frac{\alpha}{3}.
\end{equation}
\end{theorem}

\begin{proof}
В тяжёлой \(\tau\)-ветви топологическая самоэнергия экспоненциально подавлена. Мюонная ветвь, напротив, находится на масштабе \(m_\mu R_1=1\) и получает конечный компактный сдвиг \((\delta m_\mu^{top}/m_\mu)_{rel}=\alpha/3\). В S2T-формуле \(\rho_0\) является не голым отношением полюсных масс, а топологическим оператором перехода от \(\mu\)-ветви к \(\tau\)-ветви. Поэтому QED-поправка действует на относительный оператор аддитивно, а не умножает весь большой фактор \(\rho_0\). Это даёт ровно \(-\alpha/3\).
\end{proof}

\begin{nogocriterion}[Провал QED-сдвига]
Если \(\alpha/3\) трактуется как обычная мультипликативная перенормировка знаменателя \(m_\mu\), то формула дала бы \(\rho_0(1-\alpha/3)\), то есть член \(-\rho_0\alpha/3\), а не \(-\alpha/3\). Такая интерпретация несовместима с целевой формулой. Допустим только аддитивный относительный операторный сдвиг, выведенный из конечной компактной самоэнергии.
\end{nogocriterion}

\section{Uniqueness и normalization audit \texorpdfstring{\(\tau\)}{tau}-формулы}

Сравнение с текущим контролем \(m_\tau=1776.93(9)\,\mathrm{MeV}\) даёт pull \(-0.78\sigma\). Для оценки look-elsewhere фиксировалась до перебора grammar
\[
\pi^2+n\pi+\frac pq+c\alpha,
\]
где \(n=0,\ldots,4\), \(|p|\leq4\), \(q\leq6\), а
\[
c\in\left\{0,\pm1,\pm\frac12,\pm\frac13,\pm\frac23\right\}.
\]

\begin{proposition}[Уникальность внутри frozen grammar]
Среди \(1485\) кандидатов исходная формула занимает первое место и является единственной внутри одного и двух experimental sigma. То же верно после ограничения formulas не большей description complexity. Поэтому совпадение не является тривиальным dense rational approximation внутри объявленного класса. Этот результат не является универсальной вероятностью случайного совпадения вне выбранной grammar.
\end{proposition}

\begin{nogocriterion}[Неопределённая compact normalization]
Затравочное отношение
\[
\rho_0=\pi^2+2\pi+\frac23
\]
в предыдущем theorem принято как premise, а не получено как eigenvalue charged-lepton operator. Кроме того, явно записанная winding-сумма равна
\[
\sum_{q\geq1}(-1)^q\int_0^1(1+x)K_0(2\pi qx)\,dx
=-0.1938266457\ldots,
\]
и без projection normalization даёт coefficient
\[
\frac1\pi\sum_{q\geq1}\cdots=-0.06169694\ldots,
\]
а не \(\pm1/3\). Для требуемой величины необходим
\[
|\mathcal J_{RP^3}|=5.4027533\ldots,
\]
но соответствующий operator trace в текущем тексте отсутствует. Поэтому coefficient \(1/3\) нельзя считать выведенным ссылкой только на \(\mathbb Z_2\)-проекцию.
\end{nogocriterion}

\begin{corollary}[Обновлённый статус charged-lepton row]
Формула \(m_\tau/m_\mu=\pi^2+2\pi+2/3-\alpha/3\) сохраняется как сильная и уникальная low-complexity numerical relation. Её статус понижается с закрытого вывода до conditional relation, пока \(\rho_0\) не получено из явного лептонного оператора, а \(\mathcal J_{RP^3}\) не вычислено как нормированный projection trace. Абсолютный Higgs-scale bridge, использующий \(m_\tau\), наследует эту условность.
\end{corollary}

\section{Operator-candidate для charged-lepton перехода}

Первый конструктивный способ закрыть seed gap состоит в объединении уже существующих геометрических каналов в одну direct-sum tangent,
\[
\boldsymbol\Xi_\tau
=\left(1_{\mathbb{RP}^3},1_{S^1},P_\perp n\right).
\]
Здесь первые два компонента являются ненормированными constant geometry modes на единичных носителях, а третий --- поперечной компонентой единичного направления \(n\in S^2\) относительно фиксированной spin-axis.

\begin{proposition}[Gram-реализация \(\rho_0\)]
При direct-sum trace--Hodge metric
\[
\|\boldsymbol\Xi_\tau\|^2
=\|1_{\mathbb{RP}^3}\|^2
+\|1_{S^1}\|^2
+\left\langle\|P_\perp n\|^2\right\rangle_{S^2}
=\pi^2+2\pi+\frac23.
\]
Последний член следует из \(\langle\sin^2\theta\rangle_{S^2}=2/3\). На primitive integer lattice трёх каналов при условии, что ни один компонент не равен нулю, минимальная норма достигается только при coefficients \((\pm1,\pm1,\pm1)\). Поэтому declared seed является минимальной primitive norm внутри этой модели, а не произвольной суммой после выбора coefficients.
\end{proposition}

Для compact normalization заметим, что relative lepton operator в предыдущем разделе заявлен traceless. Естественный first ambient traceless space равен
\[
\operatorname{Sym}^2_0(\mathbb R^4),
\qquad
\dim\operatorname{Sym}^2_0(\mathbb R^4)=9.
\]
Если winding trace вычисляется на cover и затем спускается с quotient-measure factor
\[
\frac{\Vol(\mathbb{RP}^3)}{\Vol(S^3)}=\frac12,
\]
то возникает кандидат
\[
\mathcal J_{RP^3}^{(9)}
=\frac12\Tr P_{tr}=\frac92.
\]

\begin{proposition}[Новая zero-parameter \(\tau\)-постдикция]
С явно вычисленной winding-суммой \(I_\tau=-0.1938266457\ldots\) получаем
\[
c_\tau^{(9)}
=\frac{9}{2\pi}|I_\tau|
=0.277636219\ldots,
\]
и revised relation
\[
\frac{m_\tau}{m_\mu}
=\pi^2+2\pi+\frac23-c_\tau^{(9)}\alpha.
\]
Она даёт
\[
m_\tau^{(9)}=1776.90237\,\mathrm{MeV},
\]
то есть \(-0.31\sigma\) относительно текущего контроля \(1776.93(9)\,\mathrm{MeV}\). Это значение не использует fitted continuous coefficient.
\end{proposition}

\begin{nogocriterion}[Trace-normalization gate]
Кандидат \(9/2\) найден после измерения missing Jacobian и потому не является blind evidence. Кроме того, canonical normalized-mode trace может компенсировать quotient volume, а normalized trace может делить результат на rank. Поэтому необходимо вывести ambient charged-lepton superconnection и явно решить, входит ли в self-energy unnormalized nine-channel trace \(9/2\), single-channel coefficient \(1\), quotient coefficient \(1/2\) или иной заранее фиксированный operator weight. До этого Gram seed и rank-nine projection являются первой полной operator-конструкцией, но не восстанавливают theorem status.
\end{nogocriterion}

\section{Ambient trace normalization audit}

Предыдущий gate можно решить в минимальной normalized ambient model. Квадратичные strains из \(\operatorname{Sym}^2_0(\mathbb R^4)\) чётны относительно antipodal action, поэтому на девятимерном секторе
\[
U_{ant}=I_9,
\qquad
P_+=\frac{I+U_{ant}}2=I_9.
\]

\begin{proposition}[Quotient trace не равен половине rank]
Полный group-average trace равен
\[
\Tr P_+=\frac12\left(\Tr I_9+\Tr U_{ant}\right)=9,
\]
а не \(9/2\). Второе слагаемое является image contribution и не может быть отброшено. Эквивалентно, normalized constant mode на \(\mathbb{RP}^3\) имеет амплитуду в \(\sqrt2\) раз выше normalized invariant cover mode; её квадрат точно сокращает half-volume quotient в quadratic matrix element.
\end{proposition}

Для одной collective deformation unnormalized vertex \(P_{tr}=I_9\) имеет Hilbert--Schmidt norm
\[
\Tr P_{tr}^2=9.
\]
После canonical field normalization
\[
\widehat P_{tr}=\frac{P_{tr}}{\sqrt{\Tr P_{tr}^2}}=\frac{I_9}{3},
\qquad
\Tr\widehat P_{tr}^2=1.
\]
Следовательно, single collective field даёт \(\mathcal J=1\), а девять независимо нормированных fields дают \(\mathcal J=9\). Ни одна canonical интерпретация не даёт \(9/2\).

\begin{nogocriterion}[No-go rank-nine half-trace completion]
При \(\mathcal J=1\) явно вычисленный compact loop предсказывает
\[
m_\tau=1777.06887\,\mathrm{MeV},
\]
или \(+1.54\sigma\). Это совместимо с текущими данными, но не выводит coefficient \(1/3\). При девяти независимых channels получается \(-2.69\sigma\). Кроме того, canonical normalization constant modes заменяет raw volume norms \(\pi^2\) и \(2\pi\) единичными нормами; поэтому Gram seed является background-geometry norm, но не автоматически normalized charged-lepton vertex. Кандидат \(\mathcal J=9/2\) закрывается отрицательно в minimal ambient model.
\end{nogocriterion}

Переоткрытие возможно только если charged-lepton action независимо выводит noncanonical stiffness, boundary measure или multiplicity operator, который фиксирует иной вес до сравнения с \(m_\tau\). Сам выбор веса ради возвращения \(1/3\) считается hidden fit.

\section{Единое родительское действие: normalization gate}

После отрицательного ambient trace audit необходимо проверить не очередную секторную конструкцию, а возможность одной общей нормировки. Зафиксируем минимальный parent-кандидат на
\[
K=\mathbb{RP}^3\times S^1,
\qquad
\mathcal H=L^2(K,S_K)\otimes\mathcal H_F\otimes\mathcal H_{Nambu},
\]
с superconnection
\[
\mathbb A=D_K\otimes I+\Gamma_K\otimes D_F+\Phi+A^{(1)}
\]
и единственной trace--Hodge метрикой
\[
\langle X,Y\rangle_{\mathbb A}
=\int_K\operatorname{Str}\!\left(X^\dagger\wedge\star Y\right).
\]
Квадратичная часть действия задаётся без независимых секторных весов:
\[
S_{parent}^{(2)}
=\frac12\langle\delta\mathbb A,\delta\mathbb A\rangle_{\mathbb A}
+\langle\Psi,D_{\mathbb A}\Psi\rangle.
\]

\begin{proposition}[Единственность регулярного kernel с равными Hessian-весами]
Для редуцированного spectral action
\[
S_f(a)=\operatorname{Tr}f\!\left((D_0+aX)^2\right)
\]
вес kernel-направления равен \(2f'(0)\), а вес направления над фоновым eigenvalue \(x=m^2\) равен
\[
2f'(x)+4x f''(x).
\]
Если эти веса должны совпадать для всякого \(x\geq0\), а \(f\) регулярно при нуле, то \(f\) является аффинной функцией.
\end{proposition}

\begin{proof}
Положив \(g=f'\), получаем
\[
g(x)+2xg'(x)=g(0).
\]
При \(x>0\) общее решение имеет вид
\[
g(x)=g(0)+\frac{C}{\sqrt{x}}.
\]
Регулярность при \(x=0\) требует \(C=0\), поэтому \(g\) постоянно и \(f(x)=A+Bx\). Следовательно, generic heat-kernel или resolvent cutoff не фиксирует одинаковые веса массивного quotient и kernel-линии; это свойство принадлежит только первичной quadratic configuration metric либо специально выведенной сингулярной конструкции.
\end{proof}

В этой канонической метрике нейтринная collective deformation сохраняет
\[
\|\boldsymbol\Xi_\nu\|^2=23+\pi^{-1}.
\]
Но для нормированных constant lepton modes объём сокращается в quadratic matrix element:
\[
\int_{\mathbb{RP}^3}|\widehat 1|^2=1,
\qquad
\int_{S^1}|\widehat 1|^2=1.
\]
Поэтому тот же parent action заменяет raw Gram seed
\[
\pi^2+2\pi+\frac23=16.81945637\ldots
\]
на
\[
1+1+\frac23=\frac83.
\]
Кроме того, canonical single-vertex trace оставляет loop coefficient
\[
\frac{|I_\tau|}{\pi}=0.06169694\ldots,
\]
а для \(1/3\) по-прежнему необходим дополнительный вес
\[
\mathcal J_{req}=5.40275331\ldots.
\]

\begin{nogocriterion}[Two-sector parent-action gate]
Минимальное единое действие обязано воспроизвести без новых коэффициентов не менее двух независимых normalization-sensitive секторов. Canonical trace--Hodge action проходит нейтринный stiffness gate, но не выводит ни charged-lepton seed, ни коэффициент \(1/3\), ни ранее провалившееся exact electromagnetic determinant closure. Higgs absolute bridge не является независимым вторым тестом, поскольку наследует \(m_\tau\) и \(S_{vac}\). Следовательно, минимальное unified parent action закрывается отрицательно как единое предсказательное действие: из пяти проверенных ветвей проходит одна.
\end{nogocriterion}

Этот результат не опровергает геометрию носителя, топологические defect-конструкции или саму идею общего источника. Он запрещает более узкое утверждение: текущие численные мосты нельзя считать разными редукциями уже найденного единого канонически нормированного действия. Переоткрытие допустимо только после независимого вывода noncanonical measure или stiffness operator из симметрии, boundary principle либо квантовой меры, причём его веса должны быть зафиксированы до сравнения с лептонными, хиггсовскими и электромагнитными данными.

\chapter{Нейтринный узел и топология seesaw}

\section{Майорановский оператор}

Для левых нейтрино \(\nu_L\) и правых стерильных мод \(N_R\) базовый seesaw-лагранжиан записывается как
\begin{equation}
    \mathcal{L}_{\nu}=i\bar\nu_L\not D\nu_L+i\bar N_R\not DN_R
    -\bar\nu_Lm_DN_R-\frac12\bar N_R^cM_RN_R+\mathrm{h.c.}
\end{equation}
После интегрирования тяжёлого сектора возникает
\begin{equation}
    m_\nu\simeq -m_D M_R^{-1}m_D^T.
\end{equation}

\section{\texorpdfstring{Знаменатель \(23+\pi^{-1}\)}{Знаменатель 23 плюс обратное pi}}

Рабочая цель главы --- вывести масштаб
\begin{equation}
    M_R\propto 23+\pi^{-1}
\end{equation}
из подсчёта спинорных мод и инфракрасного вклада окружности. Число \(23\) не допускается как подгоночный счётчик: оно должно быть связано с индексом, кратностью или исключением одной нулевой/калибровочной ветви из полного набора допустимых мод.

\begin{hypothesis}[Восьмимерный внутренний модуль]
В исходной версии нейтринного seesaw-сектора одно поколение стерильной правой моды рассматривалось как realification четырёхкомпонентного комплексного внутреннего spinor-модуля, то есть
\begin{equation}
    \,\dim_{\mathbb{R}}\mathcal{M}_{Maj}^{(1)}=8.
\end{equation}
Для трёх поколений полный стерильный модуль имеет размерность
\begin{equation}
    \,\dim_{\mathbb{R}}\mathcal{M}_{Maj}^{(3)}=3\cdot 8=24.
\end{equation}
Эта запись является гипотезой о Euclidean realification или symplectic-Majorana структуре, а не стандартным четырёхмерным Majorana condition.
\end{hypothesis}

\begin{hypothesis}[Ранний rank-one ansatz для числа 23]
Ранняя версия предполагала, что физический тяжёлый сектор не содержит универсальную следовую ветвь
\begin{equation}
    N_{tr}=\frac{1}{\sqrt{3}}(N_{R1}+N_{R2}+N_{R3}),
\end{equation}
и формально записывала
\begin{equation}
    N_{Maj}^{0}=24-1=23.
\end{equation}
Dimension gate ниже показывает, что поле \(N_{tr}\) сохраняет полный внутренний spinor-индекс. Поэтому эта формула не является индексным доказательством: она допустима только как отдельная гипотеза о каноническом rank-one projector в полном \(24\)-мерном модуле.
\end{hypothesis}

\begin{proposition}[Инфракрасный вклад окружности]
Анти-периодическая ветвь на \(S^1\) имеет минимальный фазовый зазор \(1/2\) и голономию \(\pi\). В безразмерной нормировке seesaw-блока единственный разрешённый ИК-скаляр окружности задаётся обратной длиной полупериода,
\begin{equation}
    \delta_{IR}=\frac{1}{\pi}.
\end{equation}
Тем самым минимальный топологический знаменатель тяжёлого Majorana-масштаба принимает форму
\begin{equation}
    D_{\nu}=23+\pi^{-1}.
\end{equation}
\end{proposition}

\begin{hypothesis}[Условный нейтринный знаменатель]
Если полная Dirac/\(Q_{cycle}\)-структура действительно выделяет один канонический вектор в \(24\)-мерном модуле, то знаменатель
\begin{equation}
    23+\pi^{-1}
\end{equation}
получает rank-one-plus-topological интерпретацию: \(23\) является рангом дополнения к этому вектору, а \(\pi^{-1}\) --- ИК-вкладом окружности. Сам rank-one selector пока не выведен, поэтому блок остаётся незамкнутым.
\end{hypothesis}

\begin{nogocriterion}[Провал числа 23]
Если универсальная следовая ветвь \(N_{tr}\) окажется физически наблюдаемой в относительных нейтринных массах, либо если Majorana-модуль одного поколения не имеет вещественной размерности \(8\), то подсчёт \(24-1=23\) недействителен. В этом случае нейтринный сектор возвращается в статус ``не замкнуто'' независимо от численного совпадения \(\Delta m^2\).
\end{nogocriterion}

\section{Минимальная матричная реализация seesaw}

При условном знаменателе \(D_\nu=23+\pi^{-1}\) остаётся вопрос, входит ли он в воспроизводимую матрицу \(m_D M_R^{-1}m_D^T\). Для исключения скрытой подгонки фиксируем минимальный матричный каркас до сравнения с \(\Delta m^2\), не считая тем самым число \(23\) доказанным.

Поколенческое пространство раскладывается на универсальную следовую ветвь и traceless-плоскость:
\begin{equation}
    P_{tr}=\frac13
    \begin{pmatrix}
    1&1&1\\
    1&1&1\\
    1&1&1
    \end{pmatrix},
    \qquad
    P_0=I_3-P_{tr}.
\end{equation}
Тяжёлый Majorana-блок минимальной версии задаётся как
\begin{equation}
    M_R=M_*\left(D_\nu P_0+\pi^{-1}P_{tr}\right),
\end{equation}
где \(M_*\) --- общий seesaw-масштаб. Он не влияет на отношение массовых квадратов, но нужен для абсолютной шкалы нейтринных масс.

\begin{definition}[Безразмерный нейтринный оператор формы]
До фиксации абсолютной шкалы физически содержательная часть нейтринной матрицы задаётся безразмерным оператором
\begin{equation}
    \widehat m_\nu=U_\nu
    \begin{pmatrix}
    0&0&0\\
    0&1&0\\
    0&0&\sqrt{R_\nu}
    \end{pmatrix}
    U_\nu^T,
\end{equation}
где \(U_\nu\) --- ортогональная матрица смешивания, а единственный заявляемый в версии II.A инвариант спектра равен
\begin{equation}
    R_\nu=D_\nu+\pi^2+\frac23=23+\pi^{-1}+\pi^2+\frac23.
\end{equation}
\end{definition}

\begin{proposition}[Предсказание отношения нейтринных расщеплений]
Если абсолютная шкала обозначена через \(\mu_\nu\), то минимальный спектр имеет вид
\begin{equation}
    m_1=0,
    \qquad
    m_2=\mu_\nu,
    \qquad
    m_3=\sqrt{R_\nu}\,\mu_\nu.
\end{equation}
Следовательно,
\begin{equation}
    \Delta m^2_{21}=\mu_\nu^2,
    \qquad
    \Delta m^2_{31}=R_\nu\mu_\nu^2,
\end{equation}
и отношение расщеплений не содержит свободной непрерывной константы:
\begin{equation}
    \frac{\Delta m^2_{31}}{\Delta m^2_{21}}=R_\nu
    =23+\pi^{-1}+\pi^2+\frac23
    =33.854580953940\ldots.
\end{equation}
\end{proposition}

\begin{proof}
Матрица \(\widehat m_\nu\) симметрична и имеет собственные значения \(0\), \(1\), \(\sqrt{R_\nu}\). Умножение на единственную размерную шкалу \(\mu_\nu\) даёт указанный спектр. При переходе к массовым квадратам шкала \(\mu_\nu^2\) сокращается в отношении \(\Delta m^2_{31}/\Delta m^2_{21}\), поэтому \(R_\nu\) является настоящим безразмерным S2T-предсказанием, тогда как абсолютная величина \(\Delta m^2_{21}\) требует отдельного вывода \(\mu_\nu\).
\end{proof}

\begin{theorem}[Частичное замыкание нейтринной матрицы]
В версии II.A нейтринный сектор закрыт на уровне безразмерной формы: задан \(M_R\), задан спектр \(\widehat m_\nu\), и получено предсказание для отношения \(\Delta m^2_{31}/\Delta m^2_{21}\). Однако абсолютные значения \(\Delta m^2_{21}\) и \(\Delta m^2_{31}\) ещё не являются закрытыми предсказаниями, потому что шкала \(\mu_\nu\) не выведена из уже разрешённых входов.
\end{theorem}

\begin{nogocriterion}[Провал матричного нейтринного блока]
Если для получения абсолютных \(\Delta m^2\) вводится новая непрерывная шкала \(\mu_\nu\), подобранная по одному из нейтринных расщеплений, то абсолютные массы имеют статус калибровки, а не blind-предсказания. В этом случае допустимо заявлять только безразмерное отношение \(R_\nu\), пока \(\mu_\nu\) не будет выведена из отдельного спектрального или EFT-механизма.
\end{nogocriterion}

\section{Абсолютная шкала нейтринных расщеплений}

Оставшаяся размерная величина \(\mu_\nu\) не должна вводиться как новый параметр. В минимальной версии она строится из уже разрешённых лептонных train-якорей и топологического Majorana-знаменателя. Дираковская нейтринная вставка задаётся как подавленная заряженно-лептонная ветвь
\begin{equation}
    m_D^{(\nu)}=\sqrt{\pi+\pi^{-1}}\,\frac{m_e^2}{m_\mu},
\end{equation}
а тяжёлый Majorana-масштаб traceless-сектора равен
\begin{equation}
    M_R^{(0)}=D_\nu m_\mu.
\end{equation}
Тогда seesaw-шкала лёгких нейтрино фиксируется без нового непрерывного коэффициента:
\begin{equation}
    \mu_\nu=\frac{(m_D^{(\nu)})^2}{M_R^{(0)}}
    =\left(\pi+\pi^{-1}\right)\frac{m_e^4}{m_\mu^3D_\nu}.
\end{equation}

\begin{proposition}[Численная шкала нейтринного сектора]
При использовании train-якорей \(m_e=0.510998950\,\mathrm{MeV}\) и \(m_\mu=105.6583755\,\mathrm{MeV}\) получаем
\begin{equation}
    \mu_\nu=8.576992685\cdot10^{-3}\,\mathrm{eV}.
\end{equation}
Следовательно,
\begin{align}
    \Delta m^2_{21}&=\mu_\nu^2
    =7.356480352\cdot10^{-5}\,\mathrm{eV}^2,\\
    \Delta m^2_{31}&=R_\nu\mu_\nu^2
    =2.490505596\cdot10^{-3}\,\mathrm{eV}^2.
\end{align}
\end{proposition}

\begin{proof}
Все размерные входы принадлежат \(\mathcal{D}_{train}\). Степень \(m_e^4/m_\mu^3\) является единственной минимальной seesaw-комбинацией, имеющей размерность массы и подавляющей нейтринный сектор относительно заряженно-лептонного. Фактор \(D_\nu^{-1}\) уже выведен из traceless Majorana-сектора, а множитель \(\pi+\pi^{-1}\) является суммой УФ-геометрической окружной ветви и ИК-обратной окружной ветви. Поэтому \(\mu_\nu\) не является новой ручкой.
\end{proof}

\begin{theorem}[Условное спектральное замыкание нейтринной шкалы]
В версии II.A нейтринный сектор замкнут на уровне нормальной иерархии с \(m_1=0\) только при дополнительном условии, что дираковская вставка
\begin{equation}
    m_D^{(\nu)}=\sqrt{\pi+\pi^{-1}}\,\frac{m_e^2}{m_\mu}
\end{equation}
выводится из нейтринного EFT-лагранжиана, а не вводится как феноменологический мост. Матричная форма и отношение расщеплений имеют спектральный статус, но абсолютная шкала \(\mu_\nu\) в версии II.A получает статус ``условно выведено'' до вывода этой дираковской вставки.
\end{theorem}

\begin{nogocriterion}[Провал абсолютной шкалы]
Если множитель \(\pi+\pi^{-1}\) в \(m_D^{(\nu)}\) не удаётся связать с окружной УФ/ИК-нормировкой Дираковской вставки, то абсолютная шкала \(\mu_\nu\) теряет статус вывода. В этом случае сохраняется только безразмерное матричное отношение \(R_\nu\), а абсолютные \(\Delta m^2\) не считаются закрытыми.
\end{nogocriterion}

\begin{proposition}[Первый operator gate для нейтринного overlap]
Анти-периодическая spin-структура фиксирует фазовый зазор \(\alpha=1/2\), угол голономии \(\theta=2\pi\alpha=\pi\) и unitary holonomy \(U=e^{i\theta}=-1\). Однако выражение
\[
\theta+\theta^{-1}
\]
не является функцией только от класса \(U\): оно не инвариантно относительно выбора логарифмической ветви \(\theta\mapsto\theta+2\pi k\) и сингулярно при переходе к периодической ветви \(\theta=0\). Следовательно, множитель \(\pi+\pi^{-1}\) не может быть выведен из Wilson/spin holonomy без дополнительного положительного оператора.
\end{proposition}

\begin{definition}[Кандидат reciprocal cycle operator]
Минимальный положительный doubled-cycle оператор имеет вид
\[
Q_{cycle}(g)=
\begin{pmatrix}
g&0\\
0&g^{-1}
\end{pmatrix},
\qquad
\det Q_{cycle}=1.
\]
Тогда
\[
\operatorname{Tr}Q_{cycle}(\pi)=\pi+\pi^{-1}.
\]
Эта запись является точной алгебраической реализацией требуемого множителя, но не его выводом. Для замыкания необходимо получить \(g=\pi\) как собственное значение положительной геометрической или спектральной метрики, а не отождествить \(g\) с branch-dependent holonomy angle по определению.
\end{definition}

\begin{corollary}[Обновлённый статус абсолютной нейтринной шкалы]
Holonomy-only маршрут закрыт отрицательно. Сохраняется одна конструктивная ветвь: вывести \(Q_{cycle}\) из primal/dual cycle metric, Hodge pairing, Dirichlet-to-Neumann map или иного self-adjoint overlap-оператора и показать, что \(Q_{cycle}^{1/2}\) входит в матричный элемент \(m_D\). До этого абсолютные \(\Delta m^2\) остаются условными, тогда как безразмерный инвариант \(R_\nu\) не меняет статуса.
\end{corollary}

\begin{theorem}[Systolic Gram-конструкция оператора \(Q_{cycle}\)]
Пусть \(\gamma\simeq\mathbb{RP}^1\subset\mathbb{RP}^3\) является кратчайшей нестягиваемой замкнутой геодезической единичного круглого \(\mathbb{RP}^3\). Её подъём на \(S^3\) соединяет антиподальные точки вдоль половины большого круга, поэтому
\[
\ell_\gamma=\pi.
\]
На внутреннем de Rham-комплексе цикла выберем целочисленно нормированные генераторы
\[
e_0=1,
\qquad
e_1=\frac{ds}{\ell_\gamma},
\qquad
\int_\gamma e_1=1.
\]
Их квадраты \(L^2\)-норм равны
\[
\lVert e_0\rVert^2=\ell_\gamma,
\qquad
\lVert e_1\rVert^2=\ell_\gamma^{-1}.
\]
Следовательно, Hodge--Gram оператор на
\(H^0(\gamma;\mathbb Z)\oplus H^1(\gamma;\mathbb Z)\) имеет вид
\[
Q_{cycle}=
\begin{pmatrix}
\ell_\gamma&0\\
0&\ell_\gamma^{-1}
\end{pmatrix}
=
\begin{pmatrix}
\pi&0\\
0&\pi^{-1}
\end{pmatrix},
\qquad
\det Q_{cycle}=1.
\]
\end{theorem}

\begin{proposition}[Примитивный self-dual нейтринный вектор]
На целочисленной решётке
\(\Lambda_\gamma=\mathbb Ze_0\oplus\mathbb Ze_1\)
primal/dual-инволюция действует как
\(J(a,b)=(b,a)\). Единственными ненулевыми примитивными \(J\)-инвариантными векторами являются
\(\pm(1,1)\). Поэтому с точностью до несущественного знака фиксируется
\[
v_\nu=e_0+e_1,
\]
и
\[
\lVert Q_{cycle}^{1/2}v_\nu\rVert^2
=\pi+\pi^{-1}.
\]
Тем самым требуемый множитель получает точную operator-norm реализацию без непрерывного коэффициента.
\end{proposition}

\begin{nogocriterion}[Оставшийся Dirac embedding gap]
Построенный \(e_1\) является внутренней формой на цикле \(\gamma\), а не глобальной гармонической формой на \(\mathbb{RP}^3\), поскольку \(H^1(\mathbb{RP}^3;\mathbb R)=0\). Поэтому полное замыкание требует определить restriction/defect operator из ambient spinor bundle на цикл и вывести EFT-вершину, которая выбирает self-dual вектор \(v_\nu\) и вводит \(Q_{cycle}^{1/2}\) в матричный элемент \(m_D\). Без этого шага построение \(Q_{cycle}\) повышает статус множителя до operator candidate, но не завершает вывод абсолютной нейтринной шкалы.
\end{nogocriterion}

\begin{proposition}[Минимальная seesaw-вставка \(Q_{cycle}\)]
Пусть \(V_\gamma\simeq\mathbb R^2\) является ортонормированным пространством primal/dual cycle-каналов, а
\[
v_\nu=
\begin{pmatrix}
1\\1
\end{pmatrix}
\]
--- примитивный self-dual lattice-вектор. Введём cycle-часть дираковской строки и тяжёлого блока как
\[
m_{D,cycle}=y\,v_\nu^TQ_{cycle}^{1/2},
\qquad
M_{R,cycle}=M_0I_2.
\]
Тогда cycle-свёртка seesaw равна
\[
-m_{D,cycle}M_{R,cycle}^{-1}m_{D,cycle}^T
=-\frac{y^2}{M_0}v_\nu^TQ_{cycle}v_\nu
=-\frac{y^2}{M_0}\left(\pi+\pi^{-1}\right).
\]
При
\[
y=\frac{m_e^2}{m_\mu},
\qquad
M_0=D_\nu m_\mu
\]
получается заявленная абсолютная шкала
\[
\mu_\nu=\left(\pi+\pi^{-1}\right)
\frac{m_e^4}{m_\mu^3D_\nu}.
\]
\end{proposition}

\begin{proof}
Поскольку тяжёлый cycle-блок минимально пропорционален единичному оператору,
\[
Q_{cycle}^{1/2}I_2Q_{cycle}^{1/2}=Q_{cycle}.
\]
Поэтому вся cycle-зависимость сводится к уже вычисленной норме
\(v_\nu^TQ_{cycle}v_\nu=\pi+\pi^{-1}\).
При ортогональной замене cycle-базиса
\(m_D\mapsto m_DO^T\),
\(M_R\mapsto OM_RO^T\),
и полная seesaw-свёртка не меняется. Определитель \(Q_{cycle}\) равен единице, поэтому новый независимый массовый масштаб не вводится.
\end{proof}

\begin{remark}[Нормировочный статус lattice-вектора]
Вектор \(v_\nu\) является примитивным целочисленным вектором coupling/charge lattice, а не волновой функцией, которую следует делить на её \(Q_{cycle}\)-норму. Канонически нормируются поля в ортонормированном cycle-базисе; компоненты Yukawa coupling получаются действием vielbein \(Q_{cycle}^{1/2}\) на integral vector. Дополнительная нормировка \(v_\nu\) устранила бы предсказываемый topological coupling и означала бы замену модели.
\end{remark}

\begin{nogocriterion}[Representation и vertex gap]
Минимальная матричная свёртка не должна добавлять две новые стерильные генерации. Cycle-doublet обязан быть реализован как два ортогональных канала внутри уже посчитанного восьмимерного вещественного Majorana-модуля одного поколения. Кроме того, ambient Dirac/EFT действие должно вывести как self-dual выбор \(v_\nu=(1,1)\), так и минимальный блок \(M_{R,cycle}\propto I_2\). Если эти свойства вводятся вручную, построение остаётся согласованной EFT-моделью, но не окончательной теоремой S2T.
\end{nogocriterion}

\begin{nogocriterion}[Dimension gate для числа 23]
Записанное выше удаление поколенчески-универсальной ветви требует уточнения. Если внутренний вещественный спинорный модуль одного поколения имеет размерность \(d\), то полный модуль имеет вид
\[
\mathbb R^3_{gen}\otimes\mathbb R^d_{spin}.
\]
Generation-singlet projector действует как
\[
P_{tr}\otimes I_d
\]
и имеет ранг \(d\), а generation-traceless projector
\[
(I_3-P_{tr})\otimes I_d
\]
имеет ранг \(2d\). Поэтому при заявленном \(d=8\) удаление singlet spinor оставляет \(16\), а не \(23\), вещественных направлений.

Ранг \(23\) может возникнуть только при удалении одного конкретного вектора
\[
P_{rank1}=|u\rangle\langle u|
\]
из полного \(24\)-мерного модуля. Такой projector математически допустим, но он не совпадает с удалением полного generation-singlet spinor и требует отдельного вывода выделенного вектора \(u\) из Dirac/\(Q_{cycle}\)-структуры.
\end{nogocriterion}

\begin{remark}[Majorana convention]
В четырёхмерной Minkowski-сигнатуре обычный Majorana spinor имеет четыре вещественные компоненты. Внутренний четырёхкомпонентный комплексный Euclidean Dirac spinor имеет восемь вещественных компонент как realification, но это не является обычным Majorana condition. Восемь вещественных компонент можно также получить через явно заданную symplectic-Majorana doublet-структуру. До выбора одной из этих постановок размерностный подсчёт тяжёлого модуля не считается замкнутым.
\end{remark}

\begin{corollary}[Пересмотр статуса нейтринного знаменателя]
Конструкция \(Q_{cycle}\) и cycle-channel seesaw contraction сохраняются. Однако число \(23\) и величина
\[
D_\nu=23+\pi^{-1}
\]
возвращаются в условный статус до построения symmetry-protected rank-one projector в полном generation--spinor--cycle модуле. Поскольку \(23\) входит и в \(R_\nu\), условным становится не только абсолютный масштаб, но и ранее заявленное безразмерное отношение расщеплений.
\end{corollary}

\begin{nogocriterion}[No-go для symmetry-protected rank-one selector]
Минимальный полный модуль, совместимый с восьмью вещественными внутренними компонентами, раскладывается как
\[
\mathbb R^3_{gen}\otimes\mathbb R^4_{spin}\otimes\mathbb R^2_{cycle}.
\]
Generation-singlet projector имеет ранг \(1\), self-dual cycle projector --- ранг \(1\), но lowest Dirac eigenspace на круглой \(\mathbb{RP}^3\) имеет комплексную кратность \(2\), то есть вещественный ранг \(4\). Поэтому совместный канонический projector имеет ранг \(1\cdot4\cdot1=4\), а не \(1\).

Уже quaternion subgroup \(Q_8\subset Spin(3)\), действующая на \(\mathbb R^4_{spin}\) левыми quaternion-множителями, запрещает инвариантную вещественную прямую: оператор умножения на \(i\) имеет квадрат \(-I_4\) и не имеет вещественных собственных прямых. Reynolds averaging любого rank-one spinor projector по \(Q_8\) даёт \(I_4/4\) с support rank \(4\). Следовательно, существующие \(S_3\)-, Dirac/spin- и cycle-симметрии не могут вывести rank-one selector в \(\mathbb R^{24}\).
\end{nogocriterion}

\begin{corollary}[Ковариантные альтернативы счёту 23]
Если удаляется generation-singlet, self-dual cycle и полный lowest-spinor блок, то удаляемый ранг равен \(4\), а дополнительный ранг равен \(24-4=20\). Если удаляется полный generation-singlet внутренний модуль, остаётся \(24-8=16\). Эти числа являются representation-consistent кандидатами, но ещё не физическими знаменателями. Сохранение \(23\) требует нового symmetry-breaking datum: выделенной spinor polarization, defect boundary condition, condensate или выбранного Killing spinor.
\end{corollary}

\begin{proposition}[Условный twisted-Majorana defect selector]
Пусть systolic geodesic \(\gamma=\mathbb{RP}^1\subset\mathbb{RP}^3\) является core нечётного class-D vortex дефекта двухкомпонентного вещественного Majorana mass order parameter. В радиальном канале минимальный оператор
\[
A_\perp=\frac{d}{dx}+\tanh x
\]
имеет единственную нормируемую вещественную нулевую моду
\[
\psi_0(x)=\frac{1}{\sqrt2}\operatorname{sech}x,
\qquad
A_\perp\psi_0=0,
\]
тогда как нулевая мода сопряжённого оператора ненормируема. Поэтому transverse mod-two index равен единице.

Если эта Majorana-мода связана с нетривиальной плоской \(\mathbb Z_2\)-линией на \(\mathbb{RP}^3\), то на \(\gamma\) её знак голономии равен \(-1\). Он сокращает анти-периодический spin-знак \(-1\), так что полный longitudinal holonomy равен
\[
h_{tot}=h_{spin}h_{\mathbb Z_2}=(-1)(-1)=+1.
\]
Одноканальный вещественный оператор на цикле тогда имеет одну постоянную нулевую моду. Вместе с generation-singlet и self-dual cycle line получается совместный kernel rank
\[
1\cdot1\cdot1\cdot1=1.
\]
\end{proposition}

\begin{remark}[Статус defect-обхода no-go]
Этот механизм не противоречит предыдущему no-go: defect-free spinor module не содержит канонической прямой, а vortex mass texture физически нарушает spin-isometry multiplet и выбирает её mod-two индексом. Тем самым rank \(23=24-1\) вновь становится возможным, но только условно. Необходимо вывести mass order parameter из уже существующего S2T-поля, доказать нечётность winding, построить ambient-to-core restriction map и показать, что Majorana/BdG описание не добавляет физических состояний сверх исходного \(\mathbb R^{24}\)-модуля.
\end{remark}

\begin{corollary}[Совместимость defect selector с \(Q_{cycle}\)]
Defect selector действует в ambient spinor factor, тогда как \(Q_{cycle}\) задаёт reciprocal metric в двухканальном cycle factor. Поэтому unique Majorana kernel и self-dual coupling \(v_\nu=(1,1)\) совместимы, а норма
\[
v_\nu^TQ_{cycle}v_\nu=\pi+\pi^{-1}
\]
сохраняется. Defect-механизм может восстановить rank-one count, не разрушая ранее построенную seesaw-свёртку.
\end{corollary}

\begin{proposition}[Square-root torsion origin нечётного vortex]
Пусть \(C=\mathbb{RP}^3\setminus N(\gamma)\) --- complement трубчатой окрестности generator-core. Это solid torus с
\[
H_1(C;\mathbb Z)=\mathbb Z\langle y\rangle,
\qquad
[\mu_\gamma]=2[y],
\]
где \(\mu_\gamma\) --- meridian удалённого core. Нетривиальная ambient torsion line имеет character \(\chi_L(y)=-1\), но \(\chi_L(\mu_\gamma)=+1\), поэтому сама гладкая \(\mathbb Z_2\)-линия не создаёт локальный odd vortex.

На complement существует square-root character
\[
\chi_S(y)^2=\chi_L(y)=-1,
\qquad
\chi_S(y)=\pm i.
\]
Тогда \(\chi_S(\mu_\gamma)=\chi_S(2y)=-1\). Следовательно, \(S\) не продолжается через core как гладкая flat line и задаёт дискретный quarter-holonomy defect с \(\pi\)-flux на meridian. Для charge-two Majorana pair field условие конечной энергии даёт
\[
\Delta\arg\Phi=2\oint_{\mu_\gamma}A=2\pi,
\]
то есть vortex winding равен единице и mod-two index нечётен.
\end{proposition}

\begin{remark}[Связь с уже вычисленной gauge-ветвью]
В gauge-holonomy аудите значение \(\beta=1/4\) даёт conjugate phases \(e^{i\theta_+}=-i\) и \(e^{i\theta_-}=+i\). Их квадраты равны \(-1\). Поэтому square-root defect не вводит непрерывную подгоночную фазу: он совпадает с уже присутствующей quarter-holonomy doublet и выбирается алгебраическим условием \(z^2=-1\).
\end{remark}

\begin{nogocriterion}[Оставшийся core-gluing gate]
Для окончательного восстановления числа \(23\) необходимо вывести singular square-root line \(S\) из S2T-действия, задать её boundary condition на core и доказать, что единственная transverse Majorana-мода наследует longitudinal line с полным holonomy \(+1\). До этого шага odd winding имеет непараметрический топологический источник, но rank-one kernel остаётся условным.
\end{nogocriterion}

\begin{proposition}[Локальное замыкание Majorana core line]
В минимальном Nambu-базисе \(\Psi=(\psi,\psi^\dagger)^T\) particle--hole involution задаётся как \(C=\tau_xK\). Локальный zero-mode basis при phase \(\vartheta\) pair field можно выбрать в виде
\[
v(\vartheta)=\frac1{\sqrt2}
\begin{pmatrix}
e^{i\vartheta/2}\\
e^{-i\vartheta/2}
\end{pmatrix},
\qquad
Cv(\vartheta)=v(\vartheta).
\]
При square-root переходе pair phase меняется на \(\pi\), а Nambu transition равен
\[
G_{1/4}=\operatorname{diag}(i,-i),
\qquad
v(\vartheta+\pi)=G_{1/4}v(\vartheta).
\]
Следовательно, quarter-holonomy переносит сам локальный basis нулевой моды и не создаёт дополнительный coefficient phase.

Оставшийся вещественный coefficient line получает spin transition \(-1\) и ambient torsion transition \(-1\). Их произведение равно \(+1\), поэтому полный transport
\[
T_{core}=(-1)_{spin}(-1)_{L}G_{1/4}=G_{1/4}
\]
переводит начальный zero basis точно в конечный. Coefficient holonomy равен \(+1\), и longitudinal Majorana operator имеет одну постоянную вещественную нулевую моду. Без ambient torsion line coefficient holonomy равен \(-1\), и эта мода исчезает.
\end{proposition}

\begin{corollary}[Отсутствие Nambu-удвоения и rank \(23\)]
Поскольку \(Cv=v\), физические коэффициенты в \(av\) должны быть вещественными: \(C(av)=\bar a v=av\) только при \(a\in\mathbb R\). Поэтому zero eigenspace даёт одну вещественную Majorana-прямую, а не дополнительную комплексную частицу. В условной defect-модели совместный rank равен
\[
1_{gen}\cdot1_{transverse}\cdot1_{longitudinal}\cdot1_{cycle}=1,
\]
и его дополнение в \(\mathbb R^{24}\) имеет ранг \(23\).
\end{corollary}

\begin{nogocriterion}[Глобальный action-embedding gate]
Локальный transition theorem не заменяет вывода динамики. Для статуса S2T-теоремы требуется построить глобальный tubular-neighborhood BdG/Majorana operator, получить его pairing и square-root transition из S2T-действия, вложить единственный kernel в ранее объявленный \(\mathbb R^{24}\)-модуль и доказать, что quotient этого kernel действительно создаёт heavy denominator, а не новую наблюдаемую light state.
\end{nogocriterion}

\begin{definition}[Глобальный tubular EFT-кандидат]
На трубчатой окрестности \(N(\gamma)\) минимальный квадратичный функционал записывается как
\[
S_{tube}=\frac12\int_{N(\gamma)}
\langle\Psi,\mathcal B_\Phi\Psi\rangle,
\qquad
\mathcal B_\Phi=\not\!D_S+\Phi_1\Gamma_1+\Phi_2\Gamma_2,
\]
где \(S\) --- square-root torsion line на complement, \(\Phi=\Phi_1+i\Phi_2\) --- charge-two pairing section с unit winding, а его asymptotic magnitude отождествляется с уже разрешённым heavy scale \(M_*\). Coupling к light-neutrino sector локализован на core:
\[
S_Y=\int_\gamma
\overline\nu_L\,\frac{m_e^2}{m_\mu}
v_\nu^TQ_{cycle}^{1/2}\mathcal R_\gamma\Psi+\mathrm{h.c.},
\]
где \(\mathcal R_\gamma\) проектирует на нормированную Majorana zero mode. Топологический профиль не вводит нового непрерывного fit-параметра.
\end{definition}

\begin{proposition}[Heavy quotient глобального defect operator]
Пусть \(P_{ker}=|u_0\rangle\langle u_0|\) --- projector на единственную вещественную нулевую моду, а \(P_H=I_{24}-P_{ker}\). Тогда
\[
\operatorname{rank}P_H=23.
\]
Минимальный symmetry-covariant mass operator на quotient имеет вид
\[
M_H=m\,P_H.
\]
Он self-adjoint и инвариантен относительно \(O(23)\)-замен basis в \(\operatorname{im}P_H\), но symmetry не фиксирует scalar \(m\).
\end{proposition}

\begin{nogocriterion}[Rank не является tree-level denominator]
Если \(M_H=M_*P_H\), то для канонически нормированного coupling vector \(w\in\operatorname{im}P_H\)
\[
w^TM_H^+w=\frac1{M_*},
\]
независимо от rank. Для ненормированного democratic vector с единичной компонентой в каждом из \(23\) каналов получается, напротив,
\[
w_{dem}^TM_H^+w_{dem}=\frac{23}{M_*},
\]
то есть число \(23\) входит в числитель. В Gaussian functional rank проявляется как
\[
\log\det{}'M_H=23\log M_*.
\]
Следовательно, один факт \(\operatorname{rank}P_H=23\) не выводит mass eigenvalue \(23M_*\) и не доказывает denominator \(23+\pi^{-1}\).
\end{nogocriterion}

\begin{protocol}[Оставшийся collective-normalization gate]
Желаемый оператор
\[
M_H^{target}=M_*\left(\Tr P_H+\pi^{-1}\right)P_H
\]
алгебраически даёт нужный denominator, но не может быть принят как определение. Для закрытия необходимо вывести coefficient \(\Tr P_H+\pi^{-1}\) из того же действия: как loop self-energy, spectral trace или collective stiffness с заранее фиксированной нормировкой. До такого вывода defect-модель доказывает rank-one kernel и rank-23 quotient, но не численную тяжёлую массу \(D_\nu M_*\).
\end{protocol}

\begin{definition}[Collective pairing tangent]
Рассмотрим не mass eigenvalue отдельного fermion, а единственную deformation amplitude \(a\), одновременно меняющую heavy pairing operator и dual cycle component. Её tangent задаётся прямой суммой
\[
\Xi=(P_H,e_1)
\in \operatorname{End}(\mathbb R^{24})\oplus\Omega^1(\gamma).
\]
При канонической Hilbert--Schmidt/Hodge product-метрике
\[
\|\Xi\|^2
=\Tr(P_H^TP_H)+\|e_1\|^2
=\Tr P_H+\pi^{-1}
=23+\pi^{-1}.
\]
Здесь использованы projector identity \(P_H^2=P_H\) и уже выведенная dual-cycle норма \(\|e_1\|^2=\pi^{-1}\).
\end{definition}

\begin{proposition}[Collective-stiffness вывод denominator]
Пусть \(a\) --- auxiliary amplitude этой pairing-деформации с квадратичным функционалом
\[
S_{coll}[a,J_\nu]
=\frac{M_*}{2}\|\Xi\|^2a^2+y_{cycle}aJ_\nu.
\]
Исключение \(a\) по алгебраическому уравнению движения даёт
\[
\Delta S_{eff}
=-\frac{y_{cycle}^2}{2M_*\|\Xi\|^2}J_\nu^2
=-\frac{y_{cycle}^2}{2M_*(23+\pi^{-1})}J_\nu^2.
\]
Таким образом, rank входит через каноническую норму одной collective deformation, а не как mass eigenvalue и не как ненормированная сумма \(23\) propagator-вкладов. Это согласуется с предыдущим tree-level no-go.
\end{proposition}

\begin{corollary}[Восстановление нейтринной шкалы]
При
\[
y_{cycle}=\sqrt{\pi+\pi^{-1}}\,\frac{m_e^2}{m_\mu},
\qquad
M_*=m_\mu,
\]
получаем
\[
\mu_\nu
=\frac{y_{cycle}^2}{M_*\|\Xi\|^2}
=\left(\pi+\pi^{-1}\right)
\frac{m_e^4}{m_\mu^3(23+\pi^{-1})}.
\]
Результат инвариантен относительно rescaling \(a\mapsto ca\), поскольку stiffness и vertex преобразуются совместно.
\end{corollary}

\begin{nogocriterion}[Relative product-metric gate]
Наиболее общий product metric имеет вид
\[
\|\Xi\|_{\alpha,\beta}^2
=\alpha\Tr(P_H^2)+\beta\|e_1\|^2.
\]
Точное \(D_\nu=23+\pi^{-1}\) следует только при общей unweighted-нормировке \(\alpha=\beta=1\). Поэтому равенство весов должно быть выведено из одного spectral-action или superconnection trace. Если относительный вес выбирается по нейтринным данным, collective-stiffness механизм становится скрытой подгонкой.
\end{nogocriterion}

\begin{proposition}[Единый graded superconnection trace]
Relative-weight gate закрывается, если tangent записать как единую неоднородную superconnection-форму
\[
\boldsymbol\Xi
=P_H\otimes\widehat e_0
+P_{ker}\otimes e_1,
\qquad
\widehat e_0=\pi^{-1/2},
\qquad
e_1=\frac{ds}{\pi}.
\]
Здесь \(\widehat e_0\) является канонически нормированной dynamical zero-mode wavefunction,
\[
\|\widehat e_0\|^2=1,
\]
а \(e_1\) сохраняет integral-period normalization,
\[
\oint_\gamma e_1=1,
\qquad
\|e_1\|^2=\pi^{-1}.
\]
Один graded trace--Hodge norm даёт
\[
\|\boldsymbol\Xi\|^2
=\Tr(P_H^2)\|\widehat e_0\|^2
+\Tr(P_{ker}^2)\|e_1\|^2
=23+\pi^{-1}.
\]
Смешанный член равен нулю одновременно по form degree и по projector identity \(P_HP_{ker}=0\).
\end{proposition}

\begin{remark}[Почему нормировки различаются]
Различие между \(\widehat e_0\) и integral generator \(e_1\) не является относительным fit-весом. Degree-zero component описывает propagating collective amplitude и поэтому L2-нормируется. Degree-one component описывает quantized holonomy/charge-lattice element и поэтому фиксируется условием единичного периода. Это то же различение между wavefunction normalization и integral coupling vector, которое уже использовалось при построении \(Q_{cycle}\).
\end{remark}

\begin{corollary}[Статус знаменателя внутри minimal superconnection model]
В минимальной graded-superconnection модели denominator
\[
D_\nu=23+\pi^{-1}
\]
выводится единым trace--Hodge functional без нового непрерывного коэффициента. После явного tubular embedding оставшаяся обязанность относится уже не к геометрическому restriction и не к relative weight, а к статусу самого functional: необходимо либо принять canonical configuration metric как первичный кинематический объект S2T, либо вывести специальную spectral-kernel/background identity, воспроизводящую тот же Hessian.
\end{corollary}

\begin{proposition}[Ambient superconnection и tubular restriction]
На parent carrier \(K=\mathbb{RP}^3\times S^1\) берём
\[
\mathcal H_{par}=L^2(K,S_K)\otimes\mathbb R^{24}\otimes\mathcal H_{Nambu},
\qquad
D_{par}=D_K\otimes I+\Gamma_K\otimes D_F.
\]
В defect sector одна collective variation имеет вид
\[
\delta\mathbb A_a
=a\,\rho(r)\left(
P_H\otimes\widehat e_0
+P_{ker}\otimes e_1
\right),
\]
где radial profile \(\rho\) и spectator lowest mode на дополнительной \(S^1\)-ветви L2-нормированы. Интегрирование этих двух факторов даёт единицу, поэтому restriction на core точно возвращает \(a\boldsymbol\Xi\), а ambient trace--Hodge norm равен
\[
\|\delta\mathbb A_a\|^2
=a^2\left(23+\pi^{-1}\right).
\]
Таким образом, геометрический parent-to-tube-to-core embedding закрывается без дополнительного коэффициента.
\end{proposition}

\begin{nogocriterion}[Generic spectral-kernel Hessian gate]
Предыдущее предложение фиксирует canonical configuration metric
\[
\langle\delta\mathbb A,\delta\mathbb A\rangle
=\int\Str(\delta\mathbb A^\dagger*\delta\mathbb A).
\]
Однако для общего bosonic functional
\[
S_f(a)=\Tr f\!\left((D_0+a\delta\mathbb A)^2\right)
\]
квадратичные веса massive heavy sector и zero-mode connection sector зависят от \(f'(0)\), \(f'(M_*^2)\) и \(f''(M_*^2)\). Контроль на linear, exponential и rational kernels показывает, что exact equality автоматически удерживается для linear/quadratic configuration functional, но не для generic heat-kernel Hessian около ненулевого heavy background.

Следовательно, denominator является теоремой canonical superconnection metric, но не универсальным следствием произвольного \(\Tr f(D^2)\). Parent S2T должен либо объявить configuration trace metric первичным кинематическим functional, либо вывести специальную kernel/background identity, восстанавливающую те же веса.
\end{nogocriterion}

\begin{theorem}[Единственность гладкой background-independent метрики]
В редуцированном commuting-контроле квадратичные веса kernel-line и massive heavy sector равны соответственно
\[
w_0=2f'(0),
\qquad
w_M=2f'(M_*^2)+4M_*^2f''(M_*^2).
\]
Если их равенство требуется не в одной заранее выбранной точке, а гладко для любого \(x=M_*^2>0\), то
\[
f'(0)=f'(x)+2xf''(x).
\]
Полагая \(y=f'\), получаем \(y+2xy'=y(0)\). Общее решение на \(x>0\) имеет вид
\[
y(x)=y(0)+C x^{-1/2}.
\]
Гладкость в нуле требует \(C=0\), поэтому единственный допустимый класс есть
\[
f(x)=A+Bx.
\]
Следовательно, background-independent equality не выводит специальный нелинейный spectral cutoff: она единственным образом возвращает affine kernel, то есть canonical quadratic configuration metric.
\end{theorem}

\begin{nogocriterion}[Positive heat-kernel no-go]
Для стандартного положительного Laplace-представления
\[
f(x)=\int_0^\infty e^{-tx}\,d\mu(t),
\qquad d\mu(t)\ge0,
\]
разность весов при \(x>0\) равна
\[
w_M-w_0
=2\int_0^\infty t\left[1-e^{-tx}+2txe^{-tx}\right]d\mu(t)>0
\]
для любой нетривиальной меры с поддержкой при \(t>0\). Поэтому generic positive heat-kernel spectral action не может закрыть relative-metric gate на ненулевом heavy background. Равенство в одной выбранной точке является лишь одним скалярным ограничением на бесконечномерный класс kernels и без независимого принципа представляет hidden kernel tuning.
\end{nogocriterion}

\begin{corollary}[Окончательный статус нейтринной метрики]
Знаменатель \(23+\pi^{-1}\) является строгим результатом внутри явно объявленной affine configuration-metric модели с defect sector. Он не является универсальным следствием произвольного bosonic spectral action. Поэтому дальнейшее продвижение требует не нового поиска удобного kernel, а явного выбора: принять canonical metric как первичный кинематический постулат S2T либо сохранить абсолютную нейтринную шкалу в условном статусе.
\end{corollary}

\chapter{Электрослабый и QCD-блоки}

\section{Минимальные условия закрытия}

Электрослабый и сильный сектора нельзя закрывать тем же численным приёмом, что и электромагнитный якорь. Для них требуется не одна скалярная нормировка, а полный EFT-пакет:
\begin{equation}
    G_{SM}=SU(3)_c\times SU(2)_L\times U(1)_Y,
\end{equation}
спектральное происхождение трёх калибровочных констант, правило смешивания \(U(1)_Y\) и \(SU(2)_L\), механизм задания вакуумного среднего Хиггса и нормировка сильной бегущей константы. Без этих данных строки \(\sin^2\theta_W\), \(\alpha_s(M_Z)\), \(M_W\), \(M_Z\) и \(M_H\) не являются предсказаниями S2T; ниже хиггсовская строка закрывается отдельно через \(v_{S2T}\) и \(\lambda_H^{S2T}\).

\begin{protocol}[Definition of Done электрослабого блока]
Электрослабый блок считается закрытым только если без новых blind-якорей выведены величины
\begin{equation}
    g_2,\qquad g_Y,\qquad v,\qquad \lambda_H,
\end{equation}
после чего стандартные EFT-соотношения
\begin{align}
    e&=g_2\sin\theta_W=g_Y\cos\theta_W,\\
    \sin^2\theta_W&=\frac{g_Y^2}{g_2^2+g_Y^2},\\
    M_W&=\frac12 g_2v,\\
    M_Z&=\frac12\sqrt{g_2^2+g_Y^2}\,v,\\
    M_H&=\sqrt{2\lambda_H}\,v
\end{align}
дают blind-значения без подстановки \(M_W\), \(M_Z\), \(M_H\) или \(\sin^2\theta_W\) в качестве входов.
\end{protocol}

\begin{nogocriterion}[Провал электрослабого блока]
Если \(v\) или \(\lambda_H\) берутся из наблюдаемых \(G_F\), \(M_W\), \(M_Z\) или \(M_H\), то электрослабый блок становится реконструкцией стандартной модели, а не S2T-предсказанием. В версии II.A этот запрет снимается только для хиггсовской строки после вывода \(v_{S2T}\) и \(\lambda_H^{S2T}\); калибровочные низкоэнергетические величины остаются контрольными целями до полного RG-переноса.
\end{nogocriterion}

\section{Спектральные следы калибровочных представлений}

Первый замыкаемый слой EW/QCD не требует знания \(v\) или \(\lambda_H\). Он касается только относительной нормировки калибровочных кинетических членов. Для одного поколения стандартного фермионного модуля берём левые дублеты
\begin{equation}
    Q_L:(3,2)_{1/6},
    \qquad
    L_L:(1,2)_{-1/2},
\end{equation}
и правые синглеты
\begin{equation}
    u_R:(3,1)_{2/3},
    \qquad
    d_R:(3,1)_{-1/3},
    \qquad
    e_R:(1,1)_{-1}.
\end{equation}
Следы генераторов в фермионном модуле дают коэффициенты калибровочных кинетических членов
\begin{equation}
    C_a=\Tr_{\mathcal{H}_F}(T_a^2).
\end{equation}
В стандартной нормировке индекса фундаментального представления \(SU(N)\), \(T(\mathbf{N})=1/2\), для одного поколения получаем
\begin{align}
    C_2&=4\cdot\frac12=2,\\
    C_3&=(2+1+1)\cdot\frac12=2,\\
    C_Y&=6\left(\frac16\right)^2+3\left(\frac23\right)^2
    +3\left(\frac13\right)^2+2\left(\frac12\right)^2+1
    =\frac{10}{3}.
\end{align}
Множитель трёх поколений сокращается в отношениях.

\begin{proposition}[Спектральное электрослабое смешивание на S2T-масштабе]
Если кинетическая нормировка задаётся спектральными следами \(C_i\), то
\begin{equation}
    \frac{1}{g_i^2(\Lambda_{S2T})}\propto C_i.
\end{equation}
Отсюда
\begin{equation}
    \sin^2\theta_W(\Lambda_{S2T})
    =\frac{g_Y^2}{g_2^2+g_Y^2}
    =\frac{C_2}{C_2+C_Y}
    =\frac{2}{2+10/3}=\frac38.
\end{equation}
\end{proposition}

\begin{proof}
Спектральное действие нормирует калибровочный член через след квадрата соответствующего генератора по внутреннему фермионному модулю. Поэтому \(g_i^2\) обратно пропорционален \(C_i\). Подстановка \(C_2=2\) и \(C_Y=10/3\) даёт \(3/8\). Это значение относится к спектральному S2T-масштабу, а не напрямую к \(M_Z\).
\end{proof}

\begin{corollary}[Граничное условие сильной и слабой нормировок]
Так как \(C_3=C_2=2\), на спектральном масштабе выполняется
\begin{equation}
    g_3(\Lambda_{S2T})=g_2(\Lambda_{S2T}).
\end{equation}
Это закрывает относительную \(SU(3)_c/SU(2)_L\)-нормировку, но не заменяет RG-перенос к \(M_Z\).
\end{corollary}

\begin{nogocriterion}[Провал следовой нормировки]
Если внутренний фермионный модуль отличается от стандартного набора \((Q_L,L_L,u_R,d_R,e_R)\) или использует другую гиперзарядовую нормировку, то значения \(C_2=2\), \(C_3=2\), \(C_Y=10/3\) должны быть пересчитаны. В этом случае утверждения \(\sin^2\theta_W(\Lambda_{S2T})=3/8\) и \(g_3=g_2\) теряют статус вывода.
\end{nogocriterion}

\section{Хиггсовская шкала и самосвязь}

После закрытия заряженно-лептонной ветви доступна размерная величина
\begin{equation}
    m_\tau^{S2T}=m_\mu\left(\pi^2+2\pi+\frac23-\frac{\alpha}{3}\right).
\end{equation}
Электрослабая шкала должна быть построена из неё и уже выведенного электромагнитного вакуумного скаляра \(S_{vac}\), без использования \(G_F\), \(M_W\), \(M_Z\) или \(M_H\). Минимальная спектральная нормировка хиггсовского вакуума задаётся как
\begin{equation}
    v_{S2T}=m_\tau^{S2T}S_{vac}\left(1+\frac{1}{\pi^4}\right).
\end{equation}
Множитель \(S_{vac}\) переносит лептонную массу в электромагнитную вакуумную шкалу, а \(1+\pi^{-4}\) является тем же четвёртым дзета-остатком, который уже разрешён в электромагнитном секторе.

\begin{proposition}[Численная электрослабая шкала]
При \(m_\mu=105.6583755\,\mathrm{MeV}\) и \(S_{vac}=137.035999173522\) получаем
\begin{equation}
    v_{S2T}=245.993409\,\mathrm{GeV}.
\end{equation}
Это значение не использует \(G_F\) как вход.
\end{proposition}

\begin{proposition}[Blind-постдикция константы Ферми]
Если \(v_{S2T}\) отождествить с tree-level electroweak vacuum scale стандартного EFT, то без нового входа следует
\[
G_F^{S2T}=\frac{1}{\sqrt2\,v_{S2T}^2}
=1.1685251368\cdot10^{-5}\,\mathrm{GeV}^{-2}.
\]
Контрольное значение
\[
G_F^{exp}=1.1663788(6)\cdot10^{-5}\,\mathrm{GeV}^{-2}
\]
ниже на относительную величину
\[
\frac{G_F^{S2T}-G_F^{exp}}{G_F^{exp}}
=1.84017\cdot10^{-3}.
\]
Эквивалентная экспериментальная шкала равна \(v_F=246.219640\,\mathrm{GeV}\), то есть превышает \(v_{S2T}\) на \(0.226231\,\mathrm{GeV}\).
\end{proposition}

\begin{nogocriterion}[Fermi matching gate]
Tree-level отождествление \(v_{S2T}=v_F\) не проходит precision-проверку. Для согласования требуется вывести конечный matching-factor
\[
C_F=\frac{G_F^{exp}}{G_F^{S2T}}=0.998163209,
\qquad
\frac{v_F}{v_{S2T}}=1.000919663.
\]
Универсальное умножение всех масс, пропорциональных \(v\), на второй множитель даёт лишь
\[
M_W=79.9965\,\mathrm{GeV},
\qquad
M_Z=90.4151\,\mathrm{GeV},
\]
что не закрывает ранее полученные невязки относительно измеренных \(M_W\) и \(M_Z\). Поэтому один fitted scale-factor недопустим. Matching должен быть заранее выведен из полного frozen EW/KK/defect spectrum и совместно проверить \(G_F\), \(M_W\), \(M_Z\) и \(\sin^2\theta_W\).
\end{nogocriterion}

\begin{definition}[Спектральная самосвязь Хиггса]
Минимальная хиггсовская самосвязь задаётся средним между четырёхкомпонентной скалярной нормировкой \(1/8\) и первой кривизной поколенческой плоскости:
\begin{equation}
    \lambda_H^{S2T}=\frac18\left(1+\frac{1}{3\pi^2}\right).
\end{equation}
Здесь \(3\) --- число поколений, а \(\pi^{-2}\) --- первая окружная спектральная кривизна.
\end{definition}

\begin{proposition}[Хиггсовская масса]
Стандартное EFT-соотношение
\begin{equation}
    M_H=\sqrt{2\lambda_H}\,v
\end{equation}
даёт
\begin{equation}
    M_H^{S2T}=125.056486\,\mathrm{GeV}.
\end{equation}
Следовательно, хиггсовская строка получает статус S2T-следствия версии II.A после вывода минимума эффективного потенциала \(U(H,\chi)\).
\end{proposition}

\begin{nogocriterion}[Провал хиггсовской нормировки]
Если множитель \(1+\pi^{-4}\) в \(v_{S2T}\) или поправка \((3\pi^2)^{-1}\) в \(\lambda_H\) не удерживаются как спектральные остатки уже зафиксированного носителя, то \(v\), \(\lambda_H\) и \(M_H\) возвращаются в статус open. Запрещено заменять эти множители числом, подобранным по наблюдаемой массе Хиггса.
\end{nogocriterion}

\begin{definition}[Эффективный потенциал хиггсовского моста]
Пусть \(\chi\) --- геометрический модуль электромагнитно-лептонной вакуумной шкалы с массовой размерностью один, а
\begin{equation}
    \chi_0=m_\tau^{S2T}S_{vac},
    \qquad
    a_H=1+\pi^{-4}.
\end{equation}
Минимальный S2T-потенциал задаётся как
\begin{equation}
    U(H,\chi)= -\mu_H^2(\chi)H^\dagger H+\lambda_H^{S2T}(H^\dagger H)^2
    +\frac{\kappa_\chi}{2}\left(\chi-\chi_0\right)^2,
\end{equation}
где \(\kappa_\chi>0\) имеет массовую размерность два,
\begin{equation}
    \lambda_H^{S2T}=\frac18\left(1+\frac{1}{3\pi^2}\right),
    \qquad
    \mu_H^2(\chi)=\lambda_H^{S2T}a_H^2\left(2\chi_0\chi-\chi^2\right).
\end{equation}
Такая форма фиксирует не только значение \(\mu_H^2(\chi_0)\), но и условие \(\partial_\chi\mu_H^2|_{\chi_0}=0\), поэтому модульный минимум не сдвигается хиггсовским конденсатом. Параметр \(\kappa_\chi\) задаёт только жёсткость модульного направления и не является феноменологической ручкой для \(v\), \(\lambda_H\) или \(M_H\).
\end{definition}

\begin{proposition}[Минимум и устойчивость хиггсовского потенциала]
Для \(H=(0,v/\sqrt2)^T\) стационарные условия
\begin{equation}
    \frac{\partial U}{\partial (H^\dagger H)}=0,
    \qquad
    \frac{\partial U}{\partial \chi}=0
\end{equation}
дают
\begin{equation}
    \chi=\chi_0=m_\tau^{S2T}S_{vac},
    \qquad
    v=\frac{\mu_H(\chi_0)}{\sqrt{\lambda_H^{S2T}}}
    =m_\tau^{S2T}S_{vac}(1+\pi^{-4}).
\end{equation}
В этом минимуме
\begin{equation}
    \frac{\partial^2 U}{\partial \Phi^2}=2\lambda_H^{S2T}>0,
    \qquad
    \frac{\partial^2 U}{\partial \chi^2}=\kappa_\chi+2\lambda_H^{S2T}a_H^2\Phi_0>0,
    \qquad \Phi_0=\frac{v^2}{2},
\end{equation}
так что хиггсовская шкала \(v_{S2T}\) является устойчивым минимумом эффективного потенциала, а не отдельной численной вставкой.
\end{proposition}

\begin{proof}
Положим \(\Phi=H^\dagger H\). При фиксированном \(\chi\) потенциал по \(\Phi\) имеет стандартную форму
\begin{equation}
    U_H=-\mu_H^2(\chi)\Phi+\lambda_H^{S2T}\Phi^2.
\end{equation}
Его нетривиальный минимум удовлетворяет \(\Phi_0=\mu_H^2(\chi)/(2\lambda_H^{S2T})\), то есть \(v^2=2\Phi_0=\mu_H^2(\chi)/\lambda_H^{S2T}\). В точке \(\chi=\chi_0\) имеем \(\mu_H^2(\chi_0)=\lambda_H^{S2T}a_H^2\chi_0^2\) и \(\partial_\chi\mu_H^2|_{\chi_0}=0\), поэтому второе стационарное условие не получает скрытого сдвига от \(\Phi_0\). Отсюда \(v=a_H\chi_0=m_\tau^{S2T}S_{vac}(1+\pi^{-4})\). Положительность второго вариационного оператора следует из \(\lambda_H^{S2T}>0\) и \(\kappa_\chi>0\); значение \(\kappa_\chi\) не входит в \(v\) или \(M_H\).
\end{proof}

\begin{theorem}[Замыкание хиггсовского моста]
В версии II.A хиггсовский мост внутренне закрыт внутри минимальной EFT \(U(H,\chi)\): величины
\begin{equation}
    v_{S2T}=m_\tau^{S2T}S_{vac}(1+\pi^{-4}),
    \qquad
    \lambda_H^{S2T}=\frac18\left(1+\frac{1}{3\pi^2}\right)
\end{equation}
получаются из минимума потенциала без использования \(G_F\), \(M_W\), \(M_Z\) или \(M_H\) как входов. Однако blind-аудит \(G_F\) показывает, что физическое отождествление \(v_{S2T}\) с charged-current weak scale требует ещё не выведенного finite matching. Поэтому \(M_H\) остаётся выведенным внутри Higgs EFT, но полный электрослабый статус моста является условным до совместного matching-gate.
\end{theorem}

\section{Сильный сектор}

Для QCD недостаточно указать группу \(SU(3)_c\). Нужно вывести либо спектральную начальную нормировку \(g_3(\Lambda_{S2T})\), либо безразмерный аналог \(\Lambda_{QCD}\), после чего разрешено использовать стандартный RG-бег
\begin{equation}
    \mu\frac{dg_3}{d\mu}=\beta_{QCD}(g_3,n_f)
\end{equation}
до масштаба \(M_Z\). Только тогда величина
\begin{equation}
    \alpha_s(M_Z)=\frac{g_3^2(M_Z)}{4\pi}
\end{equation}
может считаться вычисленной.

\begin{protocol}[Definition of Done QCD-блока]
QCD-блок считается закрытым только если выполнены три условия:
\begin{enumerate}
    \item из спектра внутреннего носителя получена нормировка \(SU(3)_c\)-кинетического члена;
    \item пороги кварков и число активных ароматов \(n_f\) заданы EFT-лагранжианом, а не подобраны по \(\alpha_s(M_Z)\);
    \item RG-перенос к \(M_Z\) воспроизводим без использования \(\alpha_s(M_Z)\) как входа.
\end{enumerate}
\end{protocol}

\begin{nogocriterion}[Провал QCD-блока]
Если начальная сильная константа, \(\Lambda_{QCD}\) или пороговая схема выбираются после просмотра \(\alpha_s(M_Z)\), то QCD-блок не проходит blind-валидацию. В версии II.A сильный сектор имеет статус локализованной незакрытой задачи, а не отрицательного численного предсказания.
\end{nogocriterion}

\section{Однопетлевой перенос к \texorpdfstring{\(M_Z\)}{MZ}}

Для отделения строгого вывода от пороговой феноменологии фиксируем минимальный однопетлевой перенос. Пусть
\begin{equation}
    T_{EW}=\log\frac{\Lambda_{S2T}}{M_Z}
    =D_\nu+2\pi+\pi^{-1}=29.919805079547\ldots.
\end{equation}
Коэффициенты стандартной модели в нормировке \(g_Y\), \(g_2\), \(g_3\) равны
\begin{equation}
    b_Y=\frac{41}{6},
    \qquad
    b_2=-\frac{19}{6},
    \qquad
    b_3=-7.
\end{equation}
Однопетлевое уравнение имеет вид
\begin{equation}
    \frac{1}{g_i^2(\mu)}=\frac{1}{g_i^2(\Lambda_{S2T})}
    -\frac{b_i}{8\pi^2}\log\frac{\mu}{\Lambda_{S2T}}.
\end{equation}

\begin{proposition}[Однопетлевые EW-оценки на масштабе \(M_Z\)]
Используя \(\alpha^{-1}\) как train-якорь для \(e^2=4\pi\alpha\), граничное условие \(\sin^2\theta_W(\Lambda_{S2T})=3/8\) и шкалу \(v_{S2T}=245.993409\,\mathrm{GeV}\), получаем
\begin{align}
    \sin^2\theta_W(M_Z)_{1\ell}&=0.217181,\\
    M_W^{1\ell}&=79.923\,\mathrm{GeV},\\
    M_Z^{1\ell}&=90.332\,\mathrm{GeV}.
\end{align}
Эти значения являются однопетлевым S2T-прогоном без пороговых и двухпетлевых поправок.
\end{proposition}

\begin{proposition}[QCD-прогон без порогов]
При \(g_3(\Lambda_{S2T})=g_2(\Lambda_{S2T})\) тот же однопетлевой перенос даёт
\begin{equation}
    \alpha_s(M_Z)_{1\ell,no\;thr}=0.086898.
\end{equation}
Следовательно, сильный сектор не может быть объявлен закрытым без пороговой схемы и, вероятно, двухпетлевого RG-переноса.
\end{proposition}

\begin{nogocriterion}[Провал RG-замыкания]
Если пороговые поправки к \(M_W\), \(M_Z\) или \(\alpha_s(M_Z)\) выбираются численно после просмотра blind-значений, то RG-слой считается подгонкой. Допустимы только пороги, выведенные из EFT-лагранжиана и масс уже закрытых секторов.
\end{nogocriterion}

\section{Спектральные пороги и модуль \texorpdfstring{\(\rho\)}{rho}}

Оставшийся калибровочный разрыв можно локализовать в конечных порогах тяжёлых GUT-состояний. На компактном носителе
\begin{equation}
    K=\mathbb{RP}^3_{R_3}\times S^1_{R_1}
\end{equation}
естественный безразмерный модуль формы равен
\begin{equation}
    \rho=\frac{R_1}{R_3}.
\end{equation}
Если тяжёлые \(X,Y\)-бозоны и хиггсовские триплеты чувствуют разные спектральные ветви \(\mathbb{RP}^3\) и \(S^1\), их массы получают конечное геометрическое расщепление
\begin{align}
    M_{XY}^2&=\Lambda_{S2T}^2\left(1+\xi_{XY}\log\rho\right),\\
    M_{H_3}^2&=\Lambda_{S2T}^2\left(1+\xi_H\log\rho\right),
\end{align}
где \(\xi_{XY}\) и \(\xi_H\) должны быть вычислены из собственных значений лапласиана на соответствующих представлениях.

В минимальной SU(5)-подобной пороговой записи конечные добавки имеют вид
\begin{align}
    \Delta_3^{GUT}&=\frac{1}{12\pi}\left[-21\log\frac{M_{XY}}{\Lambda_{S2T}}+
    \log\frac{M_{H_3}}{\Lambda_{S2T}}\right],\\
    \Delta_2^{GUT}&=\frac{1}{12\pi}\left[-12\log\frac{M_{XY}}{\Lambda_{S2T}}\right],\\
    \Delta_1^{GUT}&=\frac{1}{12\pi}\left[-22\log\frac{M_{XY}}{\Lambda_{S2T}}-\frac25\log\frac{M_{H_3}}{\Lambda_{S2T}}\right].
\end{align}
Тогда низкоэнергетическая нормировка записывается как
\begin{equation}
    \alpha_i^{-1}(M_Z)=\alpha_i^{-1}(\Lambda_{S2T})-\frac{b_i}{2\pi}\log\frac{M_Z}{\Lambda_{S2T}}+\Delta_i^{GUT}.
\end{equation}

\begin{proposition}[Резонансный кандидат для модуля]
Без обращения к blind-значениям естественный дискретный кандидат формы задаётся отношением окружной и трёхмерной ветвей с уже выведенным лептонным QED-сдвигом:
\begin{equation}
    \rho_{S2T}=\frac34\left(1+\frac{\alpha}{3}\right)=0.751824338\ldots.
\end{equation}
Число \(3/4\) фиксирует минимальный резонанс между трёхмерной \(\mathbb{RP}^3\)-ветвью и четырёхмерной спинорной нормировкой, а \(\alpha/3\) является уже закрытым компактным однопетлевым сдвигом.
\end{proposition}

\begin{protocol}[Допустимое закрытие порогов]
EW/QCD-пороги получают статус ``выведено'' только если выполнены три условия:
\begin{enumerate}
    \item \(\xi_{XY}\) и \(\xi_H\) вычислены из спектра лапласиана на \(\mathbb{RP}^3\times S^1\), а не подогнаны по \(\alpha_s(M_Z)\) или \(\sin^2\theta_W(M_Z)\);
    \item модуль \(\rho\) фиксирован как \(\rho_{S2T}\) до просмотра EW/QCD blind-строк;
    \item пороговые формулы с этим \(\rho_{S2T}\) воспроизводят \(\sin^2\theta_W(M_Z)\), \(M_W\), \(M_Z\) и \(\alpha_s(M_Z)\) в заранее заданной точности.
\end{enumerate}
\end{protocol}

\begin{nogocriterion}[Обратная реконструкция не является blind-предсказанием]
Если \(\rho\), \(\xi_{XY}\) или \(\xi_H\) определяются из наблюдаемых \(\alpha_s(M_Z)\), \(\sin^2\theta_W(M_Z)\), \(M_W\) или \(M_Z\), то это допустимая реконструкция геометрии, но не blind-замыкание. В таком режиме EW/QCD-блок остаётся в статусе ``частичный успех'', даже если восстановленное значение близко к \(\rho_{S2T}=\frac34(1+\alpha/3)\).
\end{nogocriterion}

\begin{proposition}[Аудит минимального порогового ansatz]
Используем однопетлевые значения без порогов
\begin{equation}
    \sin^2\theta_W(M_Z)_{1\ell}=0.217181,
    \qquad
    \alpha_s(M_Z)_{1\ell,no\;thr}=0.086898.
\end{equation}
Для контрольных значений \(\sin^2\theta_W(M_Z)=0.23122\) и \(\alpha_s(M_Z)=0.11810\) требуемые сдвиги обратных калибровочных констант равны
\begin{equation}
    \delta\alpha_1^{-1}=-1.923858,
    \qquad
    \delta\alpha_2^{-1}=+1.923858,
    \qquad
    \delta\alpha_3^{-1}=-3.040313.
\end{equation}
Если пытаться воспроизвести первые два сдвига минимальными SU(5)-порогами выше, получаются
\begin{equation}
    \log\frac{M_{XY}}{\Lambda_{S2T}}=-6.043978,
    \qquad
    \log\frac{M_{H_3}}{\Lambda_{S2T}}=513.738102,
\end{equation}
что одновременно нефизично и даёт неверный знак/размер для \(\delta\alpha_3^{-1}\).
\end{proposition}

\begin{proof}
Из \(\delta\alpha_2^{-1}=(-12/12\pi)\log(M_{XY}/\Lambda_{S2T})\) следует \(\log(M_{XY}/\Lambda_{S2T})=-\pi\delta\alpha_2^{-1}=-6.043978\). Подстановка этого значения в уравнение для \(\delta\alpha_1^{-1}\) требует \(\log(M_{H_3}/\Lambda_{S2T})=513.738102\). Такое значение не может быть получено из малой геометрической деформации \(\rho_{S2T}=0.751824338\). Более того, линейная масса \(M^2=\Lambda^2(1+\xi\log\rho)\) при \(\log\rho<0\) требует \(\xi<-1/\log\rho=3.505665\) для положительности, тогда как восстановленная \(H_3\)-ветвь выходит далеко за этот диапазон. Следовательно, минимальная двухпороговая SU(5)-схема не закрывает EW/QCD-разрыв.
\end{proof}

\begin{theorem}[No-go минимального EW/QCD-порогового закрытия]
Пороговый ansatz с одним модулем \(\rho\) и двумя жёсткостями \(\xi_{XY},\xi_H\) в форме
\begin{equation}
    M_J^2=\Lambda_{S2T}^2(1+\xi_J\log\rho)
\end{equation}
не переводит EW/QCD-блок версии II.A в статус ``успех''. Он локализует место разрыва: нужна либо более полная спектральная башня Калуцы--Клейна, либо иной групповой breaking-блок, либо двухпетлевые EFT-пороги, выведенные до сравнения с blind-данными.
\end{theorem}

\section{Рабочая программа версии II.B}

Честный итог EW/QCD-аудита состоит не в замене проваленного ansatz новым подгоночным коэффициентом, а в расширении спектральной задачи. Следующая версия должна включить три обязательных узла:
\begin{enumerate}
    \item полную KK-башню тяжёлых калибровочных и скалярных состояний на \(\mathbb{RP}^3\times S^1\), а не только два эффективных порога \(M_{XY}\) и \(M_{H_3}\);
    \item адъюнктный breaking-сектор \(\Sigma\in\mathbf{24}\) с состояниями \(\Sigma_8:(8,1,0)\) и \(\Sigma_3:(1,3,0)\);
    \item двухпетлевый RG-перенос с порогами, выведенными из EFT-лагранжиана до сравнения с \(\alpha_s(M_Z)\), \(\sin^2\theta_W(M_Z)\), \(M_W\) и \(M_Z\).
\end{enumerate}

Минимальная расширенная пороговая структура должна иметь форму
\begin{align}
    \Delta_3^{GUT}&=\frac{1}{12\pi}\left[-21\log\frac{M_{XY}}{\Lambda}+
    \log\frac{M_{H_3}}{\Lambda}+\frac12\log\frac{M_{\Sigma_8}}{\Lambda}\right],\\
    \Delta_2^{GUT}&=\frac{1}{12\pi}\left[-12\log\frac{M_{XY}}{\Lambda}+\frac13\log\frac{M_{\Sigma_3}}{\Lambda}\right],\\
    \Delta_1^{GUT}&=\frac{1}{12\pi}\left[-22\log\frac{M_{XY}}{\Lambda}-\frac25\log\frac{M_{H_3}}{\Lambda}\right].
\end{align}
Эти формулы являются задачей II.B, а не завершённым выводом II.A: массы \(M_{\Sigma_8}\) и \(M_{\Sigma_3}\) должны быть получены из спектра, а не выбраны для компенсации провала минимальной модели.

\begin{protocol}[Два незакрытых узла II.A]
К концу версии II.A остаются две честно локализованные задачи:
\begin{enumerate}
    \item EW/QCD-пороги: минимальный двухпороговый SU(5)-ansatz провален; требуется KK-башня, \(\Sigma_8/\Sigma_3\) и двухпетлевой RG;
    \item нейтринная дираковская вставка: формула \(m_D^{(\nu)}=\sqrt{\pi+\pi^{-1}}m_e^2/m_\mu\) должна быть выведена из \(\mathcal{L}_\nu\).
\end{enumerate}
\end{protocol}

\begin{theorem}[Статус электрослабого и QCD-блоков в версии II.A]
В версии II.A электрослабый и QCD-блоки закрыты только на уровне спектральных граничных нормировок
\begin{equation}
    \sin^2\theta_W(\Lambda_{S2T})=\frac38,
    \qquad
    g_3(\Lambda_{S2T})=g_2(\Lambda_{S2T}).
\end{equation}
Они ещё не закрыты полностью как низкоэнергетические blind-предсказания \(\sin^2\theta_W(M_Z)\), \(\alpha_s(M_Z)\), \(M_W\) и \(M_Z\), потому что пороговые условия и полный двухпетлевой RG-перенос от \(\Lambda_{S2T}\) к \(M_Z\) пока не выведены. Хиггсовская строка \(M_H\) закрыта отдельной спектральной нормировкой \(v\) и \(\lambda_H\).
\end{theorem}

\chapter{Численный аудит и blind-валидация}

\section{Train/blind-разделение}

Train-блок фиксируется как
\begin{equation}
    \mathcal{D}_{train}=\{\alpha^{-1},m_e,m_\mu\}.
\end{equation}
Blind-блок включает величины, не использованные для калибровки:
\begin{equation}
    \mathcal{D}_{blind}=\{m_\tau,\sin^2\theta_W,\alpha_s(M_Z),M_W,M_Z,M_H,\Delta m^2_{21},\Delta m^2_{31}\}.
\end{equation}

\section{Машинно-воспроизводимый аудит закрытых строк}

Численные строки, заявленные как закрытые в версии II.A, воспроизводятся скриптом \texttt{s2t\_tome2\_audit.py}. Скрипт принимает только train-якоря
\begin{equation}
    \alpha^{-1}=137.035999177,
    \qquad
    m_e=0.51099895069\,\mathrm{MeV},
    \qquad
    m_\mu=105.6583755\,\mathrm{MeV},
\end{equation}
и записывает результат в \texttt{s2t\_tome2\_results.json}. Закрытый численный блок получается как
\begin{align}
    S_{geo}&=137.036303775878,\\
    S_{vac}&=137.035999173522,\\
    m_\tau^{S2T}&=1776.859428563\,\mathrm{MeV},\\
    v_{S2T}&=245.993409261\,\mathrm{GeV},\\
    \lambda_H^{S2T}&=0.129221715985,\\
    M_H^{S2T}&=125.056486039\,\mathrm{GeV}.
\end{align}
Контрольные невязки равны
\begin{equation}
    S_{vac}-\alpha^{-1}=-3.48\cdot10^{-9},
    \qquad
    \frac{m_\tau^{S2T}-m_\tau^{ctrl}}{m_\tau^{ctrl}}=-3.22\cdot10^{-7}.
\end{equation}
Эта проверка не закрывает EW/QCD-пороги и нейтринную дираковскую вставку; она только фиксирует воспроизводимость уже заявленных закрытых строк.

\section{Сводный blind-scorecard и локализация недостающего сектора}

После замораживания всех строк версии II.A выполнено сравнение с контрольными значениями PDG 2024. Экспериментальные ошибки используются только как диагностическая шкала; поскольку theory matching error ещё не вычислена, это не глобальный \(\chi^2\)-fit.

Близкий кластер образуют
\begin{align}
m_\tau &: -0.78\sigma,\\
M_H &: -1.30\sigma,\\
\lambda_H-\frac{(M_H^{exp})^2}{2v_F^2} &: -0.26\sigma,\\
\Delta m_{21}^2 &: -0.26\sigma,\\
\Delta m_{32}^2 &: +1.95\sigma.
\end{align}
Последние две строки сохраняют низкий evidential weight, поскольку нейтринное action-замыкание условно. Напротив, zero-matching gauge cluster проваливается совместно:
\begin{align}
\frac{\delta G_F}{G_F}&=+0.1840\%,\\
\frac{\delta \sin^2\theta_W^{OS}}{\sin^2\theta_W^{OS}}&=-2.70\%,\\
\frac{\delta M_W}{M_W}&=-0.555\%,\\
\frac{\delta M_Z}{M_Z}&=-0.939\%,\\
\frac{\delta\alpha_s(M_Z)}{\alpha_s(M_Z)}&=-26.36\%.
\end{align}

\begin{proposition}[Кластеризация невязок по секторам]
После исправления только общей weak scale по \(G_F\) остаются необходимые множители
\[
\frac{g_2^{req}}{g_2^{S2T}}=1.004659,
\qquad
\frac{g_Z^{req}}{g_Z^{S2T}}=1.008549,
\qquad
\frac{g_3^{req}}{g_3^{S2T}}=1.165296.
\]
Следовательно, невязка не является одной общей ошибкой нормировки \(v\). Она локализуется в representation-dependent gauge matching и RG-переносе; особенно сильный сектор требует изменения \(g_3\), несводимого к малой Fermi-scale поправке. В то же время scalar quartic proxy остаётся близким к контролю.
\end{proposition}

\begin{protocol}[Следующий совместный falsification gate]
Новый KK/breaking-sector spectrum должен быть зафиксирован до сравнения и одним расчётом дать отдельные threshold corrections для \(U(1)_Y\), \(SU(2)_L\) и \(SU(3)_c\). Он проходит gate только при совместном улучшении \(G_F\), \(\sin^2\theta_W\), \(M_W\), \(M_Z\) и \(\alpha_s\). Исправление одной строки отдельным коэффициентом считается fit.
\end{protocol}

\section{Первый representation-cone gate для KK-сектора}

После фиксации \(v\) по \(G_F\) требуемые модули сдвигов обратных связей в tree-level physical proxy равны
\[
|\delta\alpha_i^{-1}|_{Y,2,3}
=(4.65824,\,0.275386,\,3.03317),
\]
то есть после нормировки на \(SU(2)\)
\[
(16.91,\,1,\,11.01).
\]
Полная SM generation имеет one-loop matter direction
\[
b^{gen}_{Y,2,3}=\left(\frac{20}{9},\frac43,\frac43\right).
\]
Оптимизация только одного общего threshold amplitude не спасает этот луч: он несёт более чем в восемь раз слишком большой \(SU(2)\)-вклад относительно требуемого направления. Следовательно, обычная KK-репликация полных поколений проваливает первый no-fit direction gate.

\begin{proposition}[Split-representation hint]
Низкосложный расщеплённый набор из одного vectorlike up-type singlet, двух down-type singlets и одного scalar doublet имеет направление
\[
b^{U+2D+H}_{Y,2,3}
=\left(\frac{17}{6},\frac16,2\right)
\sim(17,1,12),
\]
которое отличается от требуемого отношения не более чем на \(5.6\%\) после выбора одного общего amplitude. Это не blind-предсказание, поскольку набор найден после построения residual vector. Но он локализует необходимое качество missing sector: гиперзарядные и цветные \(SU(2)\)-синглеты должны переживать compactification значительно чаще, чем их doublet partners.
\end{proposition}

\begin{nogocriterion}[Magnitude и reverse-engineering gate]
Фиксированный модуль \(\rho_{S2T}=0.751824338\) даёт одну геометрическую единицу
\[
\frac{|\log\rho_{S2T}|}{2\pi}=0.0453994.
\]
Для найденных representation rays требуются примерно \(7\)--\(35\) когерентных единиц в зависимости от multiplicity. Поэтому один low KK level недостаточен, а суммарный эффект башни нельзя постулировать. Необходимо вывести anomaly-free parent representation, holonomy/boundary projection и regulated finite tower sum до сравнения. Простое добавление \(U+2D+H\) по результату inverse-аудита считается fit.
\end{nogocriterion}

\section{Anomaly-free parent и совместная \texorpdfstring{\(Z_2/Z_4\)}{Z2/Z4}-проекция}

Split-ray допускает anomaly-free реализацию до вычисления threshold magnitude. Возьмём vectorlike parent content
\[
(\mathbf{10}+\overline{\mathbf{10}})
+2(\mathbf{5}+\overline{\mathbf{5}})
+\mathbf{5}_H.
\]
Фермионная \(SU(5)^3\)-аномалия сокращается между conjugate representations. В фундаментальном представлении введём order-four элемент
\[
h=e^{i3\pi Y}=\operatorname{diag}(-1,-1,-1,-i,-i),
\qquad
P_5=h^2=\operatorname{diag}(1,1,1,-1,-1).
\]
Оба имеют determinant one. Интерпретируем \(P_5\) как \(\mathbb{RP}^3\) \(Z_2\)-голономию, а \(h\) как quarter-ветвь дополнительной окружности.

Назначим conjugate flat characters: \(+1\) для \(\mathbf{10}+\overline{\mathbf{10}}\), \(-1\) для двух \(\mathbf{5}+\overline{\mathbf{5}}\) и \(-i\) для \(\mathbf{5}_H\). Полная фаза компоненты равна
\[
\Phi=P_5\,i^{6Y}\,\chi_{flat}.
\]
Прямая таблица даёт
\[
\Phi_U=1,
\quad \Phi_Q=-i,
\quad \Phi_E=-1,
\quad \Phi_D=1,
\quad \Phi_L=i,
\quad \Phi_H=1,
\quad \Phi_{T_H}=i.
\]

\begin{proposition}[Точная проекция split-ray]
Периодические zero modes существуют ровно для одного vectorlike \(U\), двух vectorlike \(D\) и одного complex scalar doublet \(H\). Их one-loop matter direction равен
\[
b_{Y,2,3}=\left(\frac{17}{6},\frac16,2\right).
\]
Поля \(Q,L,E\) и цветной triplet \(T_H\) переходят соответственно в quarter, quarter, half и quarter shifted branches. Низкоэнергетические gauge anomalies равны нулю, поскольку surviving fermions vectorlike.
\end{proposition}

\begin{protocol}[Следующий determinant gate]
Конструкция закрывает compatibility representation content с anomaly cancellation, но не threshold magnitude. Следующий расчёт фиксирован: вычислить regulated разность \(\mathbb{RP}^3\times S^1\) KK determinants между periodic \(U,D,H\)-ветвями и quarter/half-shifted партнёрами \(Q,L,E,T_H\). Flat-character assignment должен быть независимо выведен из sector attribution; если он выбирается ради сохранения \(U+2D+H\), конструкция остаётся reverse-engineered completion.
\end{protocol}

\section{Determinant gate проецированных KK-партнёров}

Для каждого общего положительного \(\mathbb{RP}^3\)-mass parameter \(\rho\) конечная shifted/periodic circle-разность задаётся функцией
\[
L_\beta(\rho)=
\log\frac{\cosh(2\pi\rho)-\cos(2\pi\beta)}{\cosh(2\pi\rho)-1}.
\]
Локальные heat-kernel коэффициенты в этом отношении сокращаются. Quarter-shifted partners имеют суммарное beta-direction
\[
b^{(1/4)}=b_Q+2b_L+b_{T_H}
=\left(\frac53,\frac{10}{3},\frac32\right),
\]
а half-shifted \(E\)-branch даёт
\[
b^{(1/2)}=\left(\frac43,0,0\right).
\]
Следовательно, любой общий shell и любая положительная сумма одинаковых spectra имеют вид
\[
\Delta=A\,b^{(1/4)}+B\,b^{(1/2)},
\qquad A,B>0.
\]

\begin{theorem}[Direction no-go shifted-partner determinant]
После нормировки на \(SU(2)\) получаем
\[
\frac{\Delta}{\Delta_2}
=\left(\frac{5+4r}{10},1,\frac9{20}\right),
\qquad r=\frac BA>0.
\]
Таким образом, color-to-weak ratio тождественно равен \(9/20\), тогда как scorecard требует \(11.014\). Поэтому finite holonomy determinant проецированных \(Q,L,E,T_H\)-партнёров не может дать нужное направление при общем \(\mathbb{RP}^3\)-спектре независимо от overall normalization, cutoff и числа shells.
\end{theorem}

\begin{nogocriterion}[Sign no-go intermediate-running маршрута]
Inverse residual действительно почти коллинеарен beta-вектору periodic split-sector \(U+2D+H\). Формальная амплитуда противоположного знака соответствует
\[
\log\frac{\Lambda_{S2T}}{M_{split}}=10.065\ldots,
\]
и при уже используемом \(T_{EW}=29.919805\) задаёт диагностическую шкалу
\[
M_{split}\simeq3.8\cdot10^{10}\,\mathrm{GeV}.
\]
Однако физический RG-перенос вниз удовлетворяет
\[
\alpha_i^{-1}(\mu)=\alpha_i^{-1}(\Lambda)
+\frac{b_i}{2\pi}\log\frac{\Lambda}{\mu}.
\]
Все компоненты \(\Delta b=(17/6,1/6,2)\) положительны, поэтому обычный промежуточный running сдвигает обратные связи в противоположную сторону. Реконструированная шкала является sign-reversed inverse diagnostic, а не допустимой vectorlike mass prediction.
\end{nogocriterion}

\section{Минимальный геометрический action для split-scale}

Из уже существующих до gauge-scorecard инвариантов можно составить безразмерный кандидат
\[
S_{split}^{geom}
=\Vol(\mathbb{RP}^3)+\frac12\|e_1\|^2
=\pi^2+\frac{1}{2\pi}
=10.0287593442.
\]
Здесь первый член является фиксированным бозонным объёмом носителя, а второй --- стандартным quadratic half-weight integral cycle mode с \(\|e_1\|^2=\pi^{-1}\). Тогда
\[
M_{split}^{geom}=\Lambda_{S2T}e^{-S_{split}^{geom}}
=3.9674\cdot10^{10}\,\mathrm{GeV}.
\]
Это на \(3.63\%\) выше прежней inverse-реконструкции, но сама реконструкция ниже исключается как RG-механизм из-за неверного знака.

\begin{nogocriterion}[Провал legacy gauge-scorecard]
Ранее использованная shortcut-формула вычитала положительный matter-вклад,
\[
\alpha_{i,legacy}^{-1}(M_Z)
=\alpha_{i,no\;thr}^{-1}(M_Z)
-\frac{\Delta b_i}{2\pi}S_{split}^{geom},
\]
и тем самым формально давала
\[
M_W=80.2826\,\mathrm{GeV},
\qquad
M_Z=91.0795\,\mathrm{GeV},
\]
\[
\sin^2\theta_W^{OS}=0.223035,
\qquad
\alpha_s(M_Z)=0.120257.
\]
Но этот знак противоположен RG-уравнению. Кроме того, сумма \(\alpha_Y^{-1}+\alpha_2^{-1}\), то есть train-якорь \(\alpha_{em}^{-1}\), изменяется с \(137.035999\) до \(132.245398\), или на \(-3.50\%\). Поэтому указанное одновременное улучшение недействительно и исключается из evidence ledger.
\end{nogocriterion}

\begin{nogocriterion}[Post-scorecard status]
Несмотря на отсутствие fitted coefficient, комбинация \(\pi^2+(2\pi)^{-1}\) предложена после знания reconstructed exponent и потому не является blind evidence. Ещё более точный кандидат
\[
\pi^2+\frac{1}{2\pi}+\frac{\alpha\pi^2}{2}
\]
исключается из evidence ledger, поскольку coefficient \(\alpha/2\) не был независимо выведен из действия. Сам geometric action может исследоваться как самостоятельная defect-конструкция, но ordinary running сектора \(U+2D+H\) больше не является допустимым gauge-repair механизмом.
\end{nogocriterion}

\section{Defect-saddle gate для промежуточной шкалы}

Пусть \(X:\Sigma\to\mathbb{RP}^3\) представляет degree-one wrapping фундаментального класса, а на систолическом цикле \(\gamma\) задана вещественная форма \(a\) единичного периода. Рассмотрим минимальный функционал
\[
S_{wrap}[X,a]
=\Vol(X)+\frac12\int_\gamma a\wedge *_\gamma a,
\qquad
\deg X=1,
\qquad
\oint_\gamma a=1.
\]

\begin{theorem}[Условное constrained saddle]
На фиксированном единичном носителе глобальный минимум функционала равен
\[
S_{wrap}^{on\text{-}shell}
=\pi^2+\frac{1}{2\pi}.
\]
Объёмная форма калибрует degree-one wrapping, поэтому \(\Vol(X)\geq\Vol(\mathbb{RP}^3)=\pi^2\). Если \(a=f(s)ds\) и \(\ell_\gamma=\pi\), то
\[
1=\left(\int_\gamma f\,ds\right)^2
\leq \ell_\gamma\int_\gamma f^2ds,
\]
причём равенство достигается единственно при \(f=\pi^{-1}\). Следовательно,
\[
a=e_1=\frac{ds}{\pi},
\qquad
\frac12\|a\|^2=\frac{1}{2\pi}.
\]
Для любой zero-mean деформации \(f=\pi^{-1}+u\) выполняется точное разложение
\[
S_\gamma[f]-S_\gamma[e_1]
=\frac12\int_\gamma u^2ds>0,
\]
так что cycle saddle строго устойчив.
\end{theorem}

\begin{nogocriterion}[Dynamical-radius obstruction]
Предыдущая теорема использует уже объявленную unit-radius нормировку носителя. Если общий радиус \(R\) также варьируется, on-shell функционал принимает вид
\[
S_{wrap}(R)=\pi^2R^3+\frac{1}{2\pi R}.
\]
Но
\[
S_{wrap}'(1)=3\pi^2-\frac{1}{2\pi}=29.449658\ldots\neq0,
\]
а собственная stationary point находится при
\[
R_*=(6\pi^3)^{-1/4}=0.270770\ldots.
\]
Поэтому двухчленный action выводит split exponent только как constrained saddle на независимо фиксированной геометрии; он не стабилизирует единичный радиус и не выводит общий relative normalization wrapping- и cycle-вкладов из parent dynamics.
\end{nogocriterion}

Таким образом, численная комбинация повышается до точного вариационного результата внутри минимальной fixed-background defect-модели. Её статус в полной S2T остаётся условным до вывода unit tension и cycle kinetic normalization из одного superconnection trace либо до появления независимо обязательного radius-stabilizing сектора. Этот математический результат больше не интерпретируется как успешный gauge threshold.

\section{Двухпетлевой stress-test split-сектора}

Общие product-group формулы были сначала проверены на стандартной модели. В нормировке \((g_Y,g_2,g_3)\) они точно воспроизводят
\[
b^{SM}=\left(\frac{41}{6},-\frac{19}{6},-7\right),
\qquad
B^{SM}=\begin{pmatrix}
199/18&9/2&44/3\\
3/2&35/6&12\\
11/6&9/2&-26
\end{pmatrix}.
\]
Для одного vectorlike \(U\), двух vectorlike \(D\) и дополнительного complex scalar doublet \(H\) получаются
\[
\Delta b=\left(\frac{17}{6},\frac16,2\right),
\qquad
\Delta B=\begin{pmatrix}
19/6&3/2&32/3\\
1/2&13/6&0\\
4/3&0&38
\end{pmatrix}.
\]

\begin{theorem}[Anchor-preserving two-loop verdict]
Решим coupled boundary problem с условиями
\[
\sin^2\theta_W(\Lambda_{S2T})=\frac38,
\qquad
g_3(\Lambda_{S2T})=g_2(\Lambda_{S2T}),
\]
и сохраним train-якорь \(\alpha_{em}^{-1}(M_Z)=137.035999177\). Gauge-only двухпетлевой перенос с split-sector выше \(M_{split}=3.9674\cdot10^{10}\,\mathrm{GeV}\) даёт
\[
M_W=82.0226\,\mathrm{GeV},
\qquad
M_Z=92.0616\,\mathrm{GeV},
\]
\[
\sin^2\theta_W^{OS}=0.206203,
\qquad
\alpha_s(M_Z)=0.080177.
\]
Относительные невязки равны \(+2.06\%\), \(+0.96\%\), \(-7.62\%\) и \(-32.05\%\). Включение one-loop running top Yukawa в sensitivity-диапазоне \(y_t(M_Z)=0.94\ldots1.00\) меняет \(\alpha_s(M_Z)\) лишь от \(0.080213\) до \(0.080222\) и не влияет на verdict.
\end{theorem}

\begin{nogocriterion}[Закрытие ordinary split-running ветви]
Если вместо повторной калибровки сохранить UV-нормировку no-threshold theory, split-sector предсказывает \(\alpha_{em}^{-1}(M_Z)=141.8168\), также нарушая train-якорь. Следовательно, обе допустимые интерпретации отрицательны: при сохранении \(\alpha_{em}\) blind observables резко ухудшаются, а при сохранении UV coupling сама \(\alpha_{em}\) перестаёт воспроизводиться. Ordinary intermediate running \(U+2D+H\) закрывается отрицательно. Ветвь может быть открыта вновь только конечным threshold-вкладом требуемого отрицательного знака либо независимо выведенной high-scale нормировкой, не откалиброванной по low-energy \(\alpha_{em}\).
\end{nogocriterion}

\section{Finite-threshold sign-cone gate}

После закрытия intermediate running построим corrected finite-matching target относительно alpha-anchored gauge-only two-loop SM baseline. Фиксируя \(M_W\) при исходном \(v_{S2T}\), сохраняя \(\alpha_{em}\) и используя контроль \(\alpha_s\), получаем требуемый сдвиг
\[
\delta\boldsymbol\alpha^{-1}_{Y,2,3}
=(0.368015,-0.368015,-2.459153)
\sim(1,-1,-6.68221).
\]
Уже на этом уровне остаётся отдельный electroweak gap: полученные \(g_Y,g_2\) дают
\[
M_Z=90.6969\,\mathrm{GeV},
\]
что на \(0.54\%\) ниже контроля. Следовательно, одни gauge-coupling thresholds не могут одновременно закрыть \(M_W\), \(M_Z\) и \(\alpha_{em}\) при замороженном \(v_{S2T}\).

Для расширенного \(SU(5)\)-подобного basis введём
\[
x_J=-\log\frac{M_J}{\Lambda_{S2T}}\geq0.
\]
Тогда columns threshold cone в порядке \((X,Y),H_3,\Sigma_8,\Sigma_3\) равны
\[
\frac1{12\pi}
\begin{pmatrix}
22&2/5&0&0\\
12&0&0&-1/3\\
21&-1&-1/2&0
\end{pmatrix}.
\]

\begin{theorem}[Алгебраическая достижимость и физическая невозможность]
Неограниченный nonnegative cone содержит corrected target. Однако решение, минимизирующее максимальный mass logarithm, имеет вид
\[
(x_{XY},x_{H_3},x_{\Sigma_8},x_{\Sigma_3})
=(0,34.6846,116.0466,41.6215).
\]
Доступный интервал между \(M_Z\) и \(\Lambda_{S2T}\) равен только
\[
0\leq x_J\leq T_{EW}=29.9198.
\]
Поэтому exact threshold system в физическом окне несовместна. Даже если разрешить массы по обе стороны matching scale, минимально возможное
\[
\max_J\left|\log\frac{M_J}{\Lambda}\right|=56.3025
\]
почти вдвое превосходит весь доступный RG-интервал.
\end{theorem}

\begin{nogocriterion}[Закрытие минимального logarithmic threshold basis]
Best constrained approximation насыщает \(H_3,\Sigma_8,\Sigma_3\) на нижней границе \(M_Z\), но сохраняет \(50.7\%\) target vector по относительной \(L^2\)-норме. Следовательно, basis \((X,Y),H_3,\Sigma_8,\Sigma_3\) закрывается отрицательно как simultaneous EW/QCD repair. Допустимое продолжение требует genuinely nonlogarithmic finite matching с независимо фиксированными коэффициентами. Добавление новых свободных threshold masses после просмотра target считается fit.
\end{nogocriterion}

\section{Шаблон итоговой таблицы}

{\small
\setlength{\tabcolsep}{3pt}
\begin{longtable}{>{\raggedright\arraybackslash}p{0.15\textwidth}
>{\raggedright\arraybackslash}p{0.25\textwidth}
>{\raggedright\arraybackslash}p{0.18\textwidth}
>{\raggedright\arraybackslash}p{0.15\textwidth}
>{\raggedright\arraybackslash}p{0.11\textwidth}}
\toprule
Величина & Формула S2T & Теория & Контр. & Откл. \\
\midrule
\endhead
\(\alpha^{-1}\) & \(S_{vac}\) & \(137.035999174\) & \(137.035999177\) & \(<4\cdot10^{-9}\) \\
\(m_\tau\) & \(m_\mu(\pi^2+2\pi+2/3-\alpha/3)\) & \(1776.859\,\mathrm{MeV}\) & \(1776.86\,\mathrm{MeV}\) & \(-3.2\cdot10^{-7}\) \\
\(G_F\) & \((\sqrt2v^2)^{-1}\) & \(1.1685251\cdot10^{-5}\) & \(1.1663788\cdot10^{-5}\) & \(+1.84\cdot10^{-3}\) \\
\(\sin^2\theta_W(M_Z)\) & 1-loop RG & \(0.217181\) & blind & threshold open \\
\(\alpha_s(M_Z)\) & 1-loop no-threshold & \(0.086898\) & blind & threshold open \\
\(M_W\) & 1-loop RG + \(v\) & \(79.923\,\mathrm{GeV}\) & blind & threshold open \\
\(M_Z\) & 1-loop RG + \(v\) & \(90.332\,\mathrm{GeV}\) & blind & threshold open \\
\(M_H\) & \(\sqrt{2\lambda_H}v\) & \(125.056\,\mathrm{GeV}\) & blind & закрыто \\
\(\Delta m^2_{21}\) & \(\mu_\nu^2\) & \(7.35648\cdot10^{-5}\,\mathrm{eV}^2\) & blind & условно \\
\(\Delta m^2_{31}\) & \(R_\nu\mu_\nu^2\) & \(2.49051\cdot10^{-3}\,\mathrm{eV}^2\) & blind & условно \\
\bottomrule
\end{longtable}
}

\begin{remark}
Строки со статусом open не считаются проваленными предсказаниями: они исключены из набора закрытых утверждений версии II.A. Их включение в будущую blind-таблицу допускается только после вывода соответствующих формул без новых непрерывных параметров.
\end{remark}

\section{Permutation-control и ablation}

Permutation-control и ablation теперь выполняются тем же скриптом \texttt{s2t\_tome2\_audit.py}. Для электромагнитного скаляра используется функционал ошибки
\begin{equation}
    E=|S_{vac}-\alpha^{-1}|.
\end{equation}
Физическая сборка даёт
\begin{equation}
    E_{phys}=3.48\cdot10^{-9}.
\end{equation}

Ablation-тест удаляет по одному структурному слагаемому из
\begin{equation}
    4\pi^3,
    \qquad
    \pi^2,
    \qquad
    \pi,
    \qquad
    \frac{1}{24S_{geo}},
    \qquad
    \frac{1}{\pi^4S_{geo}^2}.
\end{equation}
Полученные отношения ошибок к физической ошибке равны
\begin{center}
\begin{tabular}{>{\raggedright\arraybackslash}p{0.34\textwidth}r}
\toprule
Удалённый член & \(E_{abl}/E_{phys}\) \\
\midrule
\(4\pi^3\) & \(3.57\cdot10^{10}\) \\
\(\pi^2\) & \(2.84\cdot10^{9}\) \\
\(\pi\) & \(9.03\cdot10^{8}\) \\
\(1/(24S_{geo})\) & \(8.74\cdot10^{4}\) \\
\(1/(\pi^4S_{geo}^2)\) & \(1.56\cdot10^{2}\) \\
\bottomrule
\end{tabular}
\end{center}
Следовательно, все пять членов являются идентифицируемыми: удаление любого из них ухудшает электромагнитную нормировку.

Основной permutation-control переставляет коэффициенты \((4,1,1)\) между ролями \(\pi^3,\pi^2,\pi\), сохраняя форму спектральных вычитаний. При \(N_{total}=3\) не найдено ни одной перестановки лучше физической, поэтому
\begin{equation}
    p_{perm}^{geo}=\frac{1+N_{better}}{1+N_{total}}=\frac14.
\end{equation}
Дополнительный sign-role stress-test переставляет все пять структурных токенов между ролями ``плюс, плюс, плюс, минус, минус''. Он даёт \(N_{total}=120\), \(N_{better}=8\),
\begin{equation}
    p_{perm}^{sign}=0.0744.
\end{equation}
Этот стресс-тест не является доказательством уникальности, потому что малые вычитательные остатки почти вырождены при обмене ролями; его статус --- диагностический. Основной вывод ablation/permutation-аудита: закрытая строка \(S_{vac}\) численно воспроизводима и чувствительна ко всем структурным слагаемым, но строгая уникальность регуляризационной схемы остаётся предметом no-go критерия выбора альтернативных членов того же порядка.

\chapter{Заключение: программа спектрального замыкания}

Заключение является не риторическим итогом, а программой закрытия S2T-версии. Каждый ключевой коэффициент получает окончательный статус:
\begin{equation}
    \text{успех},\qquad \text{частичный успех},\qquad \text{неудача}.
\end{equation}
Статус ``успех'' означает отсутствие скрытого непрерывного параметра и наличие явного спектрального, топологического или EFT-вывода. Статус ``частичный успех'' означает численно и структурно проверенную формулу с локализованным недостающим диаграммным или индексным выводом. Статус ``неудача'' означает, что блок не закрыт в версии II.A и не может использоваться как доказанное blind-предсказание.



\begin{protocol}[Обновление C6/C11 на 2026-07-15]
После последних операторных аудитов узел C6 больше не находится на уровне только ``вывести mixed-trace''. Для диагонального блока
\begin{equation}
    C_{\mathrm{conn2}}[1,1]
\end{equation}
в шести quotient-нормированных Killing-состояниях на \(L(2,1)\) уже вычислена вся second-connection часть грубого one-form Laplacian. Она разложена на три подблока:
\begin{equation}
    C_{\mathrm{conn2}}[1,1]
    =C_{\Gamma_A\Gamma_B}+C_{\Gamma_{AB}}+C_{m\Gamma}.
\end{equation}
Здесь \(C_{\Gamma_A\Gamma_B}\) --- product slot, \(C_{\Gamma_{AB}}\) --- single mixed-connection slot, а \(C_{m\Gamma}\) --- inverse-metric/first-connection cross slot. Все три таблицы посчитаны в raw-базисе \(\mathrm{Sym}^2(\mathbb R^4)\) и затем сложены.

Итог сборки:
\begin{itemize}
    \item product slot ненулевой для всех \(55\) симметричных пар деформаций;
    \item single \(\Gamma_{AB}\) slot ненулевой для \(52/55\) пар;
    \item metric-cross slot ненулевой для всех \(55\) пар;
    \item полный connection block \(C_{\mathrm{conn2}}[1,1]\) ненулевой для всех \(55\) пар;
    \item ранги полного connection block распределены как \(2:6\), \(4:39\), \(6:10\).
\end{itemize}
Следовательно, внутренняя отмена second-connection блока не происходит. Если полная вторая вариация \(\delta^2\Delta_1\) всё же сокращается в determinant trace, это сокращение должно прийти из других машин: principal second-symbol, Ricci/curvature, coexact projector, Hilbert/basis transport или same-scheme local/compensation bookkeeping. Это является повышением технической конкретности C6, но не повышает статус \(S_{vac}\) до mature theorem.
\end{protocol}

\section{Шкала научной зрелости и честный путь повышения}

Для внешнего обзора вводим отдельную шкалу зрелости, не совпадающую с внутренней уверенностью автора:
\begin{equation}
    0=\text{не наука},
    \qquad
    10=\text{неоспоримое открытие}.
\end{equation}
Текущий статус версии II.A оценивается как
\begin{equation}
    R_{sci}=4/10.
\end{equation}
Это не маркетинговая оценка, а контрольная метрика зрелости. Она означает: структура математизирована, численно воспроизводима и имеет no-go критерии, но massive-аудит I--III выявил positivity no-go для ordinary coexact Bessel-tail; модель ещё не учитывает полную coexact KK-башню без paired-сокращения, EW/QCD-порогового закрытия, нейтринной вставки и внешней рецензии.

После вывода first ambient strain, coupling \(B\to P_{0,2}\), полного coexact-аудита, paired-sector search, C11-сборки и внешнего winding-determinant gate версия II.A остаётся на уровне \(4/10\), но C6 больше не является открытым вычислительным обещанием. Он закрыт отрицательно для объявленной Maxwell--ghost absorption-модели: ordinary transverse coexact tower ненулева, standard covariant FP оставляет scalar half-determinant residual, principal+connection+Ricci geometry не сокращается, обязательный paired-сектор не найден, а Casimir Bessel kernel не совпадает с Euclidean winding log-determinant. Поэтому переход \(4\to5\) нельзя получить простым завершением прежней C6-таблицы. Он требует либо новой заранее выведенной физической структуры, либо ревизии версии II.B, где exact \(\pi^{-4}\) исключён из доказанных утверждений и оставшиеся результаты проходят независимое воспроизведение.

\begin{center}
\begin{tabular}{>{\raggedright\arraybackslash}p{0.16\textwidth}>{\raggedright\arraybackslash}p{0.32\textwidth}>{\raggedright\arraybackslash}p{0.36\textwidth}}
\toprule
Порог & Что нужно закрыть & Что это меняет \\
\midrule
\(4\to5\) & вывести новую обязательную EM-структуру или выпустить ревизию без exact \(\pi^{-4}\)-теоремы и пройти независимый audit & отрицательно закрытый C6 больше не маскируется как незавершённый расчёт \\
\(5\to6\) & независимое воспроизведение численного аудита и blind-таблицы & модель становится проверяемой внешним исследователем \\
\(6\to7\) & EW/QCD-пороги без новых параметров & open строки \(\sin^2\theta_W,\alpha_s,M_W,M_Z\) переходят из roadmap в расчёт \\
\(7\to8\) & вывод нейтринной дираковской вставки и phenomenology-check полной матрицы & нейтринный блок получает абсолютную шкалу \\
\(8\to9\) & рецензируемая публикация или независимый preprint-аудит & утверждения выходят за пределы внутреннего документа \\
\(9\to10\) & новая проверяемая blind-предсказательная удача & появляется основание говорить об открытии, а не только о сильной программе \\
\bottomrule
\end{tabular}
\end{center}

\begin{nogocriterion}[Запрет искусственного повышения шкалы]
Оценка \(R_{sci}\) не повышается за счёт новых совпадений с уже просмотренными числами, переименования статусов или добавления свободных коэффициентов. Каждый шаг вверх требует либо нового вывода без параметров, либо независимого воспроизведения, либо blind-проверки, заранее отделённой от train-набора.
\end{nogocriterion}


Для текущей версии II.A финальная программа даёт следующий вердикт:
\begin{center}
\begin{tabular}{>{\raggedright\arraybackslash}p{0.34\textwidth}>{\raggedright\arraybackslash}p{0.24\textwidth}>{\raggedright\arraybackslash}p{0.30\textwidth}}
\toprule
Блок & Вердикт & Основание \\
\midrule
\(S_{geo}\) & успех & объёмы, накрытие, \(\mathbb{Z}_2\)-голономия \\
\(S_{vac}\) & структурная спектральная компрессия & periodic \(1/24\) сохраняется как рабочая ветвь; C6 закрыт отрицательно, поскольку full same-scheme compensation отсутствует и Casimir kernel не равен winding determinant \\
\(m_\tau\) & сильная условная связь & formula уникальна в frozen grammar; \(\rho_0\) и projection normalization \(\mathcal J_{RP^3}\) open \\
Нейтринная матрица & частичный успех & \(M_R\) и отношение выведены; абсолютная шкала зависит от вывода \(m_D^{(\nu)}\) \\
Higgs EFT & условный bridge & форма potential и \(\lambda_H\) выведены; absolute \(v,M_H\) наследуют \(m_\tau/S_{vac}\) \\
EW/QCD & отрицательно для минимальной ветви & ordinary split running и logarithmic threshold cone закрыты; nonlogarithmic matching не выведен \\
Единое parent action & отрицательно в minimal metric & canonical trace--Hodge norm сохраняет только neutrino stiffness; tau seed/loop и exact EM closure не воспроизводятся \\
\bottomrule
\end{tabular}
\end{center}

Итак, электромагнитный сектор тома II имеет окончательно разделённый статус: \(S_{geo}\) получает статус ``успех'', periodic \(1/24\) сохраняется как условная Maxwell--ghost ветвь, а exact \(\pi^{-4}\)-absorption понижается до strong structural compression. C6 закрыт отрицательно для версии II.A: внешне воспроизведённый спектр подтверждает coexact tower и scalar half-residual, полный same-scheme audit не находит компенсации, а Casimir Bessel kernel не является Euclidean winding determinant. В нейтринном секторе canonical graded-superconnection metric сохраняет norm \(23+\pi^{-1}\), однако новый parent-action gate показывает, что та же нормировка не выводит charged-lepton seed и coefficient \(1/3\). Поэтому найден не единый action всей программы, а один успешно нормированный нейтринный submodel. Электрослабый и QCD-блоки также закрыты отрицательно для минимальных repair-ветвей. Следующая честная задача II.B --- не новый поиск совпадений, а вывод prior noncanonical measure или stiffness operator из симметрии, boundary principle либо квантовой меры. Если такой вывод не существует, S2T следует сохранять как геометрическую исследовательскую гипотезу, а не как готовую единую физическую теорию.

\begin{nogocriterion}[Финальный no-go]
No-go сработал и для charged-lepton proof-chain: coefficient \(1/3\) требует невыведенного projection normalization, а \(\rho_0\) пока постулируется. Поэтому \(m_\tau\) и absolute Higgs-scale bridge сохраняются как сильные conditional relations, но не как закрытые блоки. Для \(S_{vac}\) exact \(\pi^{-4}\)-determinant claim также запрещено возвращать без одного из явно перечисленных условий переоткрытия C6.

Дополнительно сработал unified-action no-go: canonical parent metric проходит только один normalization-sensitive sector из требуемых двух. Запрещено объявлять существующие EM, lepton, Higgs и neutrino formulas редукциями одной сущности, пока общий action не зафиксирует их меры и Hessian-веса до сравнения с данными.
\end{nogocriterion}

\section{Глобальный falsification-аудит остаточных утверждений}

После закрытия минимального parent action необходимо проверить, не сохраняется ли физическая теория за счёт остальных положительных результатов по отдельности. Для этого вводится более строгий критерий, чем внутренняя математическая корректность.

\begin{protocol}[Критерий выжившего физического предсказания]
Утверждение считается закрытым физическим предсказанием S2T только если одновременно выполнены условия:
\begin{enumerate}
    \item оно относится к эмпирически измеримой величине;
    \item оператор и мера фиксированы до сравнения с этой величиной;
    \item имеется демонстрируемый prospective blind-статус, а не только последующая заморозка формулы;
    \item результат порождается тем же parent action, который проходит хотя бы один независимый второй сектор;
    \item вычисление воспроизводимо и не требует закрытого отрицательно механизма.
\end{enumerate}
\end{protocol}

Машинная матрица \texttt{s2t\_global\_falsification\_closure\_results.json} проверяет двенадцать остаточных групп утверждений. Результат имеет вид
\begin{equation}
N_{claims}=12,
\qquad
N_{emp}=5,
\qquad
N_{closed\ physical}=0.
\end{equation}
При этом десять из двенадцати групп содержат корректные тождества, известные спектральные результаты или внутренне согласованные конструкции в объявленных моделях. Это различие принципиально:
\begin{equation}
\text{корректная математика внутри модели}
\quad\not\Rightarrow\quad
\text{подтверждённая физическая теория}.
\end{equation}

\begin{center}
\small
\begin{tabular}{>{\raggedright\arraybackslash}p{0.29\textwidth}>{\raggedright\arraybackslash}p{0.26\textwidth}>{\raggedright\arraybackslash}p{0.33\textwidth}}
\toprule
Группа утверждений & Остаточный статус & Почему это не закрытое физическое предсказание \\
\midrule
Выбор \(\mathbb{RP}^3\times S^1\) & условная теорема в узком классе & критерии \(\Vol=\pi^2\) и фазовый шаг \(\pi\) заданы до выбора, но не выведены как требование природы \\
Lens-space spectra & воспроизведённая известная математика & подтверждает реализацию операторов, но не является новым наблюдаемым эффектом S2T \\
\(S_{geo}\) и \(1/24\) & геометрическая комбинация и локальная determinant-ветвь & отсутствует единый независимый map к наблюдаемой; full EM determinant не закрыт \\
\(S_{vac}\) & train-anchor reproduction & \(\alpha^{-1}\) входит в разрешённые train-данные, а exact \(\pi^{-4}\)-механизм закрыт отрицательно \\
\(m_\tau\) & сильный postdictive pattern & seed и projection measure не выведены; prospective preregistration отсутствует \\
Higgs bridge & зависимый ansatz & наследует \(m_\tau\) и \(S_{vac}\); независимый \(G_F\)-gate провален \\
\(Q_{cycle}\), defect и \(23+\pi^{-1}\) & теоремы условного нейтринного submodel & физический vertex и общий parent action не выведены \\
Нейтринные splittings & условная phenomenology & absolute Dirac insertion и общая мера остаются открытыми \\
EW/QCD & отрицательно для minimal realization & blind scorecard, correct-sign running и finite threshold cone не проходят \\
\bottomrule
\end{tabular}
\end{center}

\begin{proposition}[Отсутствие закрытых независимых предсказаний версии II.A]
В версии II.A нет ни одного остаточного эмпирического утверждения, которое одновременно имеет заранее фиксированные оператор и меру, prospective blind-статус и происхождение из выжившего двухсекторного parent action.
\end{proposition}

\begin{proof}
Эмпирическими группами являются \(S_{vac}\), charged-lepton relation, Higgs bridge, neutrino splittings и EW/QCD observables. Первая использует train-якорь и отрицательно закрытый determinant mechanism. Вторая не имеет выведенных seed и loop measure. Третья зависит от первых двух и независимо не проходит \(G_F\)-matching. Четвёртая не имеет закрытой Dirac insertion и unified action. Пятая провалена прямым blind-аудитом и двумя repair-аудитами. Следовательно, каждая эмпирическая группа нарушает по меньшей мере одно обязательное условие протокола, а число закрытых независимых предсказаний равно нулю.
\end{proof}

\begin{nogocriterion}[Окончательный физический вердикт версии II.A]
Версия S2T-II.A отвергается как замкнутая единая предсказательная физическая теория. Её minimal realization закрыта отрицательно в gauge-секторе и в unified-normalization gate; ни одно численное совпадение не может использоваться для обратного повышения статуса.

Этот вердикт не эквивалентен утверждению, что весь корпус является бессодержательным. Выбор носителя в ограниченном классе, воспроизведение lens-space spectra, Gram identities, условные defect-теоремы и полученные no-go результаты сохраняют математическую корректность. Не опровергнута и широкая онтологическая гипотеза, поскольку она не задаёт единственной количественной теории и потому в текущем виде не имеет полного falsification-критерия.

Физический статус может быть переоткрыт только новой, заранее зарегистрированной версией с единым action, фиксированной мерой и минимум двумя независимыми observables. Продолжение феноменологической нумерологии внутри II.A прекращается.
\end{nogocriterion}

\part{Часть II.B: переоткрытие через новые операторные структуры}

\chapter{Протокол перехода от II.A к II.B}

Часть II.B не отменяет отрицательный вердикт II.A и не переименовывает старые
совпадения в доказательства. Её задача --- исследовать только те новые
механизмы, которые были сформулированы после закрытия II.A как самостоятельные
операторные модели с собственными условиями провала. Поэтому результаты II.B
не могут ретроспективно повышать статус Maxwell--ghost C6, charged-lepton seed,
минимального parent action или прежних blind-наблюдаемых.

\begin{protocol}[Правило допуска механизма в II.B]
Новый механизм допускается к исследованию только если одновременно выполнены
условия:
\begin{enumerate}
    \item объявлены новое поле, Hilbert space, оператор, boundary conditions и
    статистика всех степеней свободы;
    \item мера и относительные determinant-веса фиксированы до сравнения с
    \(\alpha\), массами и mixing-данными;
    \item указано, является ли конструкция продолжением существующего action
    или новой моделью с повторным аудитом всех зависимых коэффициентов;
    \item сформулирован отрицательный gate, после которого механизм не может
    сохраняться как ``почти выведенный'';
    \item положительный результат обязан пройти минимум один независимый
    второй сектор.
\end{enumerate}
\end{protocol}

Два направления II.B возникли из полного перебора ближайших продолжений II.A:
семейная incidence-алгебра и электромагнитная inverse susceptibility. Первое
направление отвечает на вопрос, может ли четырёхсостоянийное spin-меню породить
полное flavour mixing. Второе проверяет, может ли точный коэффициент
\(\pi^{-4}\) возникать как обратная жёсткость нового spectral Hessian, а не как
невыведенный остаток старого Maxwell determinant.

\chapter{Семейный selector на четырёхсостоянийном меню}

\section{Два факторных оператора и предел commuting-алгебры}

Четыре spin-состояния меню маркируются точками
\begin{equation}
    F_2^2=\{(p,q)\},
\end{equation}
а физический семейный carrier определяется тройкой нетривиальных характеров.
Два фактора \(\mathbb{RP}^3\) и \(S^1\) задают независимые трансляции
\begin{equation}
    T_3(p,q)=(p+1,q),
    \qquad
    T_1(p,q)=(p,q+1).
\end{equation}
В характерном базисе они совместно диагональны и различают все три семейных
направления. Общий факторный лапласиан
\begin{equation}
    L(w_3,w_1)=w_3(I-T_3)+w_1(I-T_1)
\end{equation}
имеет собственные значения
\begin{equation}
    \{2w_3,2w_1,2(w_3+w_1)\}.
\end{equation}
Таким образом, topology и product metric способны различить поколения, но
любые два Yukawa-оператора вида
\begin{equation}
    Y_u=f_u(T_3,T_1),
    \qquad
    Y_d=f_d(T_3,T_1)
\end{equation}
имеют общий собственный базис и дают
\begin{equation}
    V_{CKM}=I.
\end{equation}

Среди проверенных parameter-free kernels наиболее сильную иерархию даёт
аддитивное winding-правило
\begin{equation}
    (m_1,m_2,m_3)\propto(e^{-2\pi},e^{-\pi},1),
\end{equation}
но оно не является универсальным: log-RMS равен \(0.591\) для объявленных
down-quark ratios, \(1.342\) для charged leptons и \(3.754\) для up-quarks.
Следовательно, factor geometry даёт различимость и грубую hierarchy, но не
выводит единый Yukawa kernel.

\section{Полный affine-incidence перебор}

Трансляции вместе с геометрическим shear порождают подгруппу порядка \(8\) в
\(AGL(2,2)\simeq S_4\). Её операторная алгебра на семейном триплете имеет
размерность \(5\) и коммутант размерности \(2\), то есть сохраняет разложение
\(1+2\).

Полный машинный перебор шестнадцати affine-элементов вне этой подгруппы даёт:
\begin{center}
\begin{tabular}{ccc}
\toprule
Тип элемента & Число & Алгебра на триплете \\
\midrule
transposition \((2+1+1)\) & 4 & полная \(M_3(\mathbb C)\) \\
three-cycle \((3+1)\) & 8 & полная \(M_3(\mathbb C)\) \\
four-cycle \((4)\) & 4 & редуцируемая размерности 5 \\
\bottomrule
\end{tabular}
\end{center}
Следовательно, двенадцать одиночных incidence-операторов способны породить
полное трёхсемейное mixing. Алгебраического препятствия больше нет. Но
\(\mathbb{RP}^3\) и \(S^1\) являются неэквивалентными факторами, и текущая
геометрия не выбирает один из двенадцати элементов. Усреднение любого кандидата
по существующей группе порядка \(8\) возвращает его в \(1+2\)-коммутант.

\begin{proposition}[Existence без selector-а]
Минимальное четырёхсостоянийное меню уже содержит операторы, достаточные для
полного \(M_3\)-mixing. Открытой является не задача существования, а вывод
конкретного incidence-направления, его веса относительно факторного лапласиана
и различного, но связанного назначения up/down-секторам.
\end{proposition}

\section{Projective flux и intrinsic-selector no-go}

Нетривиальный projective \(\mathbb Z_2\)-flux задаёт
\begin{equation}
    U^2=V^2=I,
    \qquad
    UV=-VU.
\end{equation}
Он является точным coefficient-free источником некоммутативности. Однако для
обратимых \(U,V\) в размерности \(d\)
\begin{equation}
    \det(UV)=(-1)^d\det(VU),
\end{equation}
поэтому при нечётном \(d\) projective representation не существует. Flux не
может действовать непосредственно на трёхмерном family carrier. В четырёх
измерениях он смешивает uniform reference state с триплетом, а в шестимерной
реализации \(\mathbb C^2_{flux}\otimes\mathbb C^3_{mult}\) вся flavour-задача
остаётся в свободном коммутанте \(M_3(\mathbb C)\).

Когомологическое кольцо
\begin{equation}
    H^*(\mathbb{RP}^3\times S^1;\mathbb F_2)
    =\mathbb F_2[a,b]/(a^4,b^2)
\end{equation}
и product metric различают три ненулевых класса по квадрату, свободной
\(S^1\)-компоненте и длинам
\begin{equation}
    \ell(a)=\pi,
    \qquad
    \ell(b)=2\pi,
    \qquad
    \ell(a+b)=\sqrt5\,\pi.
\end{equation}
Но все эти scalar readouts диагональны. Minimal cellular Dirac на квадрате
spin-структур также возвращает commuting spectrum
\begin{equation}
    \left\{0,\pi^{-1},2\pi^{-1},3\pi^{-1}\right\}.
\end{equation}
Determinant-line connection имеет только trivial holonomy или уже рассмотренный
\(\pi\)-flux. Ни cohomology, ни metric, ни canonical determinant line не
выбирают outside-\(D_8\) incidence-оператор.

\begin{nogocriterion}[Закрытие intrinsic family-selector]
В пределах данных \(\mathbb{RP}^3\times S^1\), четырёх spin-структур,
product metric, cohomology ring и canonical determinant line полный
family-selector не выводится. Переоткрытие требует новой prior boundary
connection или дискретного Dirac-комплекса, заранее фиксирующего один из
двенадцати рабочих incidence-операторов и его sector map.
\end{nogocriterion}

{\raggedright
Машинные основания этого раздела сохранены в
\texttt{s2t\_family\_factor\_operator\_results.json},
\texttt{s2t\_family\_affine\_incidence\_exhaustive\_results.json},
\texttt{s2t\_even\_projective\_carrier\_exhaustive\_results.json},
\texttt{s2t\_topological\_family\_selector\_results.json} и
\texttt{s2t\_spin\_menu\_dirac\_determinant\_line\_results.json}.
\par}

\chapter{Hurwitz--Hessian и inverse susceptibility для
\texorpdfstring{\(\pi^{-4}\)}{pi(-4)}}

\section{Провал прямого scale-Hessian}

Тождество
\begin{equation}
    \zeta\left(4,\frac12\right)
    =(2^4-1)\zeta(4)=\frac{\pi^4}{6}
\end{equation}
уже использовалось в II.A как нормировочный шаблон четвёртого остатка. Для
half-shifted tower
\begin{equation}
    \lambda_n(R)=\frac{(n+1/2)^2}{R^2},
    \qquad
    Z_A(s;R)=R^{2s}\zeta(2s,1/2).
\end{equation}
Поскольку \(\zeta(0,1/2)=0\), дзета-регуляризованный \(\log\det A_R\) не имеет
scale-Hessian по \(\log R\). Первый стандартный functional, содержащий
четвёртый момент,
\begin{equation}
    F(R)=\Tr A_R^{-2}=R^4\zeta(4,1/2),
\end{equation}
даёт
\begin{equation}
    \frac{d^2F}{d(\log R)^2}\bigg|_{R=1}
    =16\zeta(4,1/2)=\frac{8\pi^4}{3}.
\end{equation}
Следовательно, direct moment/Hessian даёт \(\pi^4\), а не \(\pi^{-4}\).

\begin{nogocriterion}[Провал прямого Hessian-объяснения]
Само появление \(\zeta(4,1/2)\) во второй вариации не выполняет обращение
\(\pi^4\mapsto\pi^{-4}\). Для такого обращения необходим inverse-response
functional с независимо фиксированной нормой deformation coordinate.
\end{nogocriterion}

\section{Точная шесть-канальная Gaussian identity}

Пусть шесть равновесных каналов имеют действие
\begin{equation}
    \Gamma(x)=\frac12\zeta\left(4,\frac12\right)
    \sum_{a=1}^{6}x_a^2,
    \qquad
    \bar x=\frac16\sum_{a=1}^{6}x_a.
\end{equation}
Тогда ordinary component-sum Gaussian measure даёт
\begin{equation}
    \operatorname{Var}(\bar x)
    =\frac1{6\zeta(4,1/2)}
    =\frac1{\pi^4}.
\end{equation}
Это является точной моделью inverse susceptibility, а не только численным
тождеством. Но normalized trace
\(\tau_6(X)=\Tr(X)/6\) сокращает rank suppression и даёт
\begin{equation}
    \operatorname{Var}_{\tau_6}(\bar x)=\frac6{\pi^4}.
\end{equation}
Поэтому trace normalization является физическим gate.

Ближайшие естественные шестёрки II.A не проходят полный тест:
\begin{itemize}
    \item \(\Lambda^2\) в 4D имеет raw rank \(6\), но Maxwell curvature
    constrained by \(F=dA\), а gauge/ghost modes periodic;
    \item \(\operatorname{Sym}^2T^*\mathbb{RP}^3\) имеет размерность \(6\), но
    constant conformal variation является одной trace-direction;
    \item mixed \(SU(5)\)-блок имеет complex dimension \(6\), но gauge Hessian
    считает индексы \((5/2,3/2,1)\), а reality требует conjugate block;
    \item chiral \(Q=(3,2)\) имеет шесть half-shifted Weyl-компонент, но
    Grassmann statistics и electromagnetic weight \(5/3\), а не \(6\).
\end{itemize}

\section{Self-dual full tower и determinant seed}

На ориентированном евклидовом четырёхмерном носителе
\begin{equation}
    \Lambda^2=\Lambda^2_+\oplus\Lambda^2_- ,
    \qquad
    \operatorname{rank}\Lambda^2_+=3.
\end{equation}
Полный antiperiodic spectrum удовлетворяет
\begin{equation}
    \sum_{n\in\mathbb Z}(n+1/2)^{-4}
    =2\zeta(4,1/2)=\frac{\pi^4}{3},
\end{equation}
поэтому
\begin{equation}
    3\sum_{n\in\mathbb Z}(n+1/2)^{-4}=\pi^4.
\end{equation}
Шесть каналов тем самым получают предварительное геометрическое чтение: три
самодуальные компоненты, умноженные на две стороны полного спектра Фурье.

Та же сумма возникает как Hessian mass-деформации
\begin{equation}
    \Gamma(q)=c\Tr\log(K+qI),
    \qquad
    K_n=(n+1/2)^2,
\end{equation}
поскольку
\begin{equation}
    \Gamma''(0)=-c\Tr K^{-2}.
\end{equation}
Если ошибочно считать компоненты независимыми, естественные ранги дают
\(|\Gamma''|=\pi^4\), а знак зависит от статистики: бозонные двухформенные
блоки дают \(-\pi^4\), тогда как грассмановы --- \(+\pi^4\).

Положительная жёсткость, следовательно, возникает только в духовом блоке.
Стандартная приводимая калибровочная теория двухформенного поля, однако,
имеет полную чередующуюся детерминантную башню
\begin{equation}
    Z_2\propto
    \det{}'(\Delta_2^T)^{-1/2}
    \det{}'(\Delta_1^T)^{+1/2}
    \det{}'(\Delta_0)^{-1/2}.
\end{equation}
Полная форма записи использует степени
\begin{equation}
 Z_2\propto
 \det{}'(\Delta_2)^{-1/2}
 \det{}'(\Delta_1)^{+1}
 \det{}'(\Delta_0)^{-3/2}.
\end{equation}
Сочетание поперечных степеней с полными рангами, ранее давшее
\(-\pi^4/2\), было ошибкой нормировки. Полный счёт с рангами \((6,4,1)\) и
поперечный счёт с рангами \((3,3,1)\) совпадают. Для пространственно
постоянной полусдвинутой башни оба дают
\begin{equation}
 H_{\mathrm{BV}}^{(\ell=0)}=-\frac{\pi^4}{6}.
\end{equation}
Полный четырёхмерный след дополнительно содержит ненулевые оболочки
\(\mathbb{RP}^3\) и является логарифмически расходящимся до выбора локальной
схемы вычитаний. БРСТ-комплекс тем самым не сохраняет изолированное
\(+\pi^4\).

\begin{nogocriterion}[Провал автоматического тензорно-духового спасения]
Точный положительный гессиан \(+\pi^4\) существует только при подсчёте шести
несвязанных грассмановых компонент. После полного калибровочного факторизования
один комплекс с обращённой статистикой даёт лишь \(+\pi^4/6\). Изоляция
выгодного детерминанта без полного БВ-комплекса повторяет запрещённый выбор
духового вклада.
\end{nogocriterion}

\section{Source normalization и точная поправка}

Даже если новый parent action создаёт коллективную stiffness
\begin{equation}
    H=\pi^4,
\end{equation}
интегрирование координаты
\begin{equation}
    \Gamma(q)=\frac12Hq^2+\frac{Jq}{S_{geo}}
\end{equation}
даёт
\begin{equation}
    \Delta S=-\frac{J^2}{2HS_{geo}^2}.
\end{equation}
Для точной формулы II.A требуется \(J=\sqrt2\), либо независимо выведенный
двухисточниковый или variance-readout mechanism. Следовательно, знак,
\(S_{geo}^{-2}\) и единичный коэффициент ещё не следуют из spectral identity.

\begin{hypothesis}[Скрученная градуированная конструкция II.B]
Существует полное по БРСТ расширение исходного действия с шестью независимо
выведенными комплексами двухформенных полей с обращённой статистикой, для
которых после полного учёта детерминантов остаётся коллективный гессиан
\(H=\pi^4\), а отображение источника фиксирует
\(J^2/2=1\). Тогда
\begin{equation}
    \Delta S=-\frac1{\pi^4S_{geo}^2}
\end{equation}
получается без подгонки.
\end{hypothesis}

Эта гипотеза является новой моделью II.B, а не восстановлением C6. Число
шесть запрещено вводить только ради требуемого коэффициента. Гипотеза требует
повторного вычисления периодической ветви \(1/24\), локальных коэффициентов
теплового ядра, нулевых мод, сокращений духовых вкладов и хотя бы одной
независимой второй наблюдаемой величины.

\chapter{Общая голономия: первый прямой двухсекторный критерий}

\section{Одна граничная структура для пространства поколений и тензорного триплета}

Четыре состояния спинового меню несут перестановочное представление \(S_4\),
а пространство поколений является его трёхмерным подпространством с нулевой
суммой компонент. Это позволяет проверить наиболее экономную объединяющую
гипотезу: одна постоянная голономия \(g\in S_4\), возможно умноженная на
центральный знак, одновременно выбирает направление смежности поколений и
задаёт три граничные фазы самодуального тензорного триплета на \(S^1\).

Для собственной фазы \(a\ne0\) тензорный отклик равен
\begin{equation}
 \sum_{n\in\mathbb Z}\frac1{(n+a)^4}
 =\pi^4\frac{2+\cos(2\pi a)}{3\sin^4(\pi a)},
\end{equation}
а для периодического канала после удаления нулевой моды
\begin{equation}
 \sum_{n\in\mathbb Z\setminus\{0\}}\frac1{n^4}=\frac{\pi^4}{45}.
\end{equation}
Тем самым все \(24\) элемента и оба центральных знака проверяются без
непрерывных коэффициентов.

\section{Полный перебор \texorpdfstring{\(S_4\)}{S4} с дополнительными знаками}

Результат зависит только от структуры циклов:
\begin{center}
\begin{tabular}{cccc}
\toprule
Тип & Знак & Фазы на триплете & \(H/\pi^4\) \\
\midrule
\(1+1+1+1\) & \(+\) & \(0,0,0\) & \(1/15\) \\
\(1+1+1+1\) & \(-\) & \(1/2,1/2,1/2\) & \(1\) \\
\(2+1+1\) & \(+\) & \(0,0,1/2\) & \(17/45\) \\
\(2+1+1\) & \(-\) & \(0,1/2,1/2\) & \(31/45\) \\
\(2+2\) & \(+\) & \(0,1/2,1/2\) & \(31/45\) \\
\(2+2\) & \(-\) & \(0,0,1/2\) & \(17/45\) \\
\(3+1\) & \(+\) & \(0,1/3,2/3\) & \(9/5\) \\
\(3+1\) & \(-\) & \(1/4,1/2,3/4\) & \(17/3\) \\
\(4\) & \(+\) & \(1/4,1/2,3/4\) & \(17/3\) \\
\(4\) & \(-\) & \(0,1/4,3/4\) & \(241/45\) \\
\bottomrule
\end{tabular}
\end{center}

Ровно один из сорока восьми вариантов даёт \(H=\pi^4\): центральная
антипериодическая голономия \(-I\). Однако на триплете поколений она действует
скалярно и не расширяет исходную редуцируемую алгебру. Все двенадцать
направлений смежности вне \(D_8\), соответствующие транспозициям и трёхциклам,
которые порождают полную \(M_3\)-алгебру, не дают точного тензорного отклика.
Ближайший случай с полной \(M_3\)-алгеброй имеет
\begin{equation}
 H=\frac{31}{45}\pi^4,
 \qquad
 \frac{|H-\pi^4|}{\pi^4}=\frac{14}{45}.
\end{equation}

\begin{nogocriterion}[Провал простой общей голономии]
Одна постоянная \(S_4\)-голономия с дополнительным знаком не может
одновременно породить точный тензорный гессиан \(H=\pi^4\) и выбрать
направление смежности, дающее полную \(M_3\)-алгебру поколений. Точная
антипериодичность сохраняет вырождение поколений, а все направления смешивания
меняют спектральный отклик.
\end{nogocriterion}

Это первая прямая проверка без подстраиваемых параметров, в которой две
исходные операторные конструкции II.B помещены в одну схему. Отрицательный
ответ не возвращает программу к
исходной неопределённости, а закрывает самый короткий мост между ними.
Следующий допустимый класс должен содержать неабелеву граничную связность с
упорядочением вдоль пути, вильсоновский оператор на нескольких рёбрах или
единый исходный принцип симметрии, который порождает разные, но связанные
представления в секторе поколений и тензорном секторе. Простого назначения
одного элемента группы обоим секторам уже недостаточно.

Машинное основание сохранено в
\texttt{s2t\_shared\_holonomy\_two\_sector\_results.json}.

\chapter{Непрерывная вильсоновская голономия и уравнение щели}

\section{Восемь совместных операторных решений}

Дискретный критерий для \(S_4\) с дополнительными знаками закрывает только
конечный набор голономий. Для непрерывной вильсоновской петли в \(SO(3)\) с фазами
\begin{equation}
 (0,a,1-a)
\end{equation}
полный тензорный отклик в единицах \(\pi^4\) равен
\begin{equation}
 R(c)
 =\frac1{45}+\frac{8(2+c)}{3(1-c)^2},
 \qquad
 c=\cos(2\pi a).
\end{equation}
Условие \(R(c)=1\) редуцируется к
\begin{equation}
 11c^2-52c-49=0.
\end{equation}
Единственный допустимый корень имеет вид
\begin{equation}
 c_*=\frac{26-9\sqrt{15}}{11},
 \qquad
 a_*=\frac{\arccos c_*}{2\pi}
 \approx0.3989625047.
\end{equation}

Для каждой канонической оси \(\hat n\) соответствующий оператор смежности
поколений равен
\begin{equation}
 A_{\hat n}
 =\frac{R_{\hat n}(\theta_*)+R_{\hat n}(\theta_*)^T}{2}
 =c_*I+(1-c_*)\hat n\hat n^T.
\end{equation}
Полный перебор тринадцати октаэдрических \(S_4\)-осей даёт восемь
направлений вне \(D_8\), для которых
\begin{equation}
 \dim\operatorname{Alg}
 \bigl(T_{\mathbb{RP}^3},T_{S^1},S,A_{\hat n}\bigr)=9.
\end{equation}
Все восемь одновременно дают \(R(c_*)=1\), то есть точный
\(H=\pi^4\).

\begin{proposition}[Существование непрерывного двухсекторного решения]
В классе непрерывных \(SO(3)\)-связностей с вильсоновской голономией
существуют восемь общих операторов поколений и тензорного сектора, которые
одновременно порождают полную \(M_3\)-алгебру и точный тензорный отклик
\(H=\pi^4\).
\end{proposition}

\section{Устойчивый спектральный функционал с уравнением щели}

Найденный угол допускает более сильное чтение, чем простая обратная
реконструкция. Введём
\begin{equation}
 V_{\mathrm{gap}}(c)
 =
 \frac83\left(
 \frac3{1-c}+\log(1-c)
 \right)
 -\frac{44}{45}c.
\end{equation}
Прямое дифференцирование даёт
\begin{equation}
 V_{\mathrm{gap}}'(c)=R(c)-1.
\end{equation}
Следовательно, \(c_*\) является стационарной точкой этого спектрального
функционала. Его кривизна положительна:
\begin{equation}
 V_{\mathrm{gap}}''(c_*)
 =R'(c_*)
 =\frac{8(5+c_*)}{3(1-c_*)^3}
 \approx1.901648596>0.
\end{equation}
Тем самым точный вильсоновский угол является устойчивым локальным минимумом,
а не
изолированной точкой численного совпадения.

Функционал содержит резольвентный член \((1-c)^{-1}\) и логарифм
детерминанта \(\log(1-c)\). Он является точной первообразной найденной
восприимчивости, но пока не выведен интегрированием конкретного локального
полного состава полей БВ/БРСТ.

\section{Новый критерий нормировки}

Уравнение стационарности фактически сравнивает спектральный отклик с
жёсткостью древесного уровня, нормированной к единице. Если заменить её на
\(\kappa\), то
\begin{equation}
 R(c_\kappa)=\kappa,
\end{equation}
и локальная чувствительность равна
\begin{equation}
 \frac{dc_\kappa}{d\kappa}\bigg|_{\kappa=1}
 =\frac1{R'(c_*)}
 \approx0.525859511.
\end{equation}
Поэтому единица не может считаться безобидным выбором единиц: она должна
следовать из канонической нормы поля или из исходного действия.

\begin{nogocriterion}[Критерий единичной жёсткости]
Если коэффициент вильсоновской жёсткости древесного уровня остаётся
свободным, то \(c_*\) и точный \(H=\pi^4\) являются настроенной стационарной
точкой. Положительное
закрытие требует вывести \(\kappa=1\) до сравнения с электромагнитным
остатком.
\end{nogocriterion}

\section{Сокращение числа осей прежней факторной геометрией}

Для восьми оставшихся осей проверен уже существующий факторный оператор
\begin{equation}
 L(w_3,w_1)
 =w_3(I-T_{\mathbb{RP}^3})+w_1(I-T_{S^1}).
\end{equation}
Три ранее объявленных ядра --- обратная длина, обратный квадрат длины и
туннельное подавление --- независимо дают один и тот же качественный результат:
минимизация
\begin{equation}
 E_L(\hat n)=\hat n^TL\hat n
\end{equation}
сокращает
\begin{equation}
 8\longrightarrow2
\end{equation}
и оставляет два направления транспозиций. Это первый случай, когда
существующая геометрия произведения не только допускает оператор с полной
\(M_3\)-алгеброй, но и уменьшает вырождение выбора.

Однако взаимодействие \(E_L(\hat n)\) ещё не получено из того же действия,
что и \(V_{\mathrm{gap}}\). Оставшаяся двукратность может быть вырождением
вакуумов, связанных симметрией, но такой статус требует явной проверки
действия остаточной группы на обеих осях.

\begin{proposition}[Промежуточный статус вильсоновской ветви]
На текущем этапе установлены существование совместного оператора, устойчивая
стационарная точка спектрального функционала и частичный геометрический
селектор. Невыведенными остаются единое локальное действие, каноническая
единичная жёсткость, взаимодействие оси с факторным оператором, знаки
детерминантов БВ/БРСТ и нормировка источника.
\end{proposition}

{\raggedright
Машинное основание сохранено в
\texttt{s2t\_continuous\_wilson\_gap\_action\_results.json}.
\par}

\chapter{Граничное роторное действие для вильсоновской ветви}

\section{Локальное происхождение положительного обратного члена}

Функционал с уравнением щели имеет вид
\begin{equation}
 V_{\mathrm{gap}}(c)
 =
 \frac8{1-c}
 +\frac83\log(1-c)
 -\frac{44}{45}c.
\end{equation}
Обычное исключение устойчивого вещественного поля с фиксированным источником
даёт отрицательный член \(-J^2/(2m^2)\). Положительный обратный член
возникает у циклической граничной степени свободы в ансамбле фиксированного
заряда. Для ротора с
\begin{equation}
 m^2(c)=1-c
\end{equation}
рутьева функция равна
\begin{equation}
 H_{\mathrm{rot}}=\frac{q^2}{2(1-c)}>0.
\end{equation}
Это даёт устойчивый локальный механизм требуемого знака.

\section{Минимальная арифметическая архитектура}

Пусть имеется \(N\) вещественных граничных роторов, действие усредняется
нормированным следом ранга \(d\), а средняя квадратичная норма заряда равна
\(q^2\). Однопетлевой детерминант и рутьева функция дают
\begin{equation}
 \Gamma_{\det}
 =\frac{N}{2d}\log(1-c),
 \qquad
 \Gamma_{\mathrm{rot}}
 =\frac{Nq^2}{2d(1-c)}.
\end{equation}
Точные коэффициенты требуют
\begin{equation}
 \frac{N}{2d}=\frac83,
 \qquad
 \frac{Nq^2}{2d}=8.
\end{equation}
Полный целочисленный перебор даёт минимальное решение
\begin{equation}
 N=16,
 \qquad
 d=3,
 \qquad
 q^2=3.
\end{equation}
Размерность \(16\) согласуется с алгеброй операторов четырёхсостоянийного
меню, а ранг \(3\) --- с пространством поколений и самодуальным тензорным
триплетом.

\section{Приводимое представление \(SU(2)\)}

Для неприводимого представления спина \(j\)
\begin{equation}
 C_2(j)=j(j+1),
\end{equation}
поэтому отдельное неприводимое представление с \(C_2=3\) отсутствует.
Однако перебор приводимых представлений общей размерности \(16\) и среднего
казимира \(3\) даёт три решения. Единственное решение без синглетов, в
котором все спины полуцелые, имеет вид
\begin{equation}
 \boxed{
 \mathcal R_{\partial}
 =
 2V_{1/2}\oplus3V_{3/2}.
 }
\end{equation}
Для него
\begin{equation}
 \dim\mathcal R_{\partial}
 =2\cdot2+3\cdot4=16,
\end{equation}
\begin{equation}
 \Tr_{\mathcal R_{\partial}}C_2
 =
 2\cdot2\cdot\frac34
 +3\cdot4\cdot\frac{15}{4}
 =48.
\end{equation}
После нормировки на \(d=3\)
\begin{equation}
 \frac{\dim\mathcal R_{\partial}}{2d}
 =\frac83,
 \qquad
 \frac{\Tr_{\mathcal R_{\partial}}C_2}{2d}
 =8.
\end{equation}
Следовательно,
\begin{equation}
 \Gamma_{\partial}(c)
 =
 \frac83\log(1-c)
 +\frac8{1-c}.
\end{equation}

Все блоки \(\mathcal R_{\partial}\) имеют полуцелый спин. Центральный элемент
\(-I\in SU(2)\) действует на каждом блоке знаком минус, поэтому то же
представление автоматически задаёт антипериодическую граничную ветвь и
полусдвинутый спектр.

\begin{proposition}[Кандидат на уровне среднего казимира]
Приводимое представление
\(2V_{1/2}\oplus3V_{3/2}\) является минимальным кандидатом, который
одновременно воспроизводит размерность \(16\), средний квадратичный казимир
\(3\), коэффициенты \(8\) и \(8/3\), а также антипериодическое действие
центра \(SU(2)\). Это утверждение относится к усреднённым квадратичным
инвариантам и само по себе ещё не гарантирует совпадения полного весового
детерминанта.
\end{proposition}

\section{Незакрытые условия}

Найденная конструкция снимает возражение о неканонической норме заряда, но
ещё не является полным исходным действием. Требуется:
\begin{enumerate}
    \item вывести кратности \(2\) и \(3\) из геометрии или симметрии;
    \item выписать явный граничный кинетический оператор на
    \(\mathbb{RP}^3\times S^1\);
    \item получить линейный член \(-44c/45\) из древесного действия и
    периодической ветви;
    \item построить полный комплекс БВ/БРСТ и проверить знаки детерминантов;
    \item связать представление с осью поколений и нормировкой источника.
\end{enumerate}

\begin{nogocriterion}[Критерий граничного ротора]
Если полный детерминант представления
\(2V_{1/2}\oplus3V_{3/2}\), его рутьева функция фиксированного заряда или
комплекс БВ/БРСТ изменяют коэффициенты \(8/3\), \(8\) либо
антипериодический знак, то роторная ветвь закрывается.
\end{nogocriterion}

Машинное основание сохранено в
\texttt{s2t\_wilson\_boundary\_rotor\_action\_results.json}.

\chapter{Точный весовой детерминант граничного ротора}

\section{Почему среднего казимира недостаточно}

Для вещественной двумерной плоскости, на которой голономия действует
поворотом на \(q\theta\), дзета-детерминант ковариантного лапласиана на
окружности содержит множитель
\begin{equation}
 \lambda_q(\theta)=2-2\cos(q\theta).
\end{equation}
Следовательно, точный детерминант определяется не только суммой квадратов
весов, но и полным набором самих весов.

В представлении
\begin{equation}
 2V_{1/2}\oplus3V_{3/2}
\end{equation}
после выбора нормировки, в которой фундаментальные веса имеют
\(|q|=1\), положительные зарядовые плоскости распределены как
\begin{equation}
 q=1:\ 5,
 \qquad
 q=3:\ 3.
\end{equation}
Поэтому нормированный логарифмический вклад равен
\begin{equation}
 \Gamma_{\det}
 =
 \frac13\left(
 5\log\lambda_1
 +3\log\lambda_3
 \right),
\end{equation}
тогда как функционал щели требует
\begin{equation}
 \Gamma_{\mathrm{req}}
 =\frac83\log\lambda_1.
\end{equation}
Разность не является константой:
\begin{equation}
 \Gamma_{\det}-\Gamma_{\mathrm{req}}
 =
 \log\frac{\lambda_3}{\lambda_1}
 =
 2\log|1+2\cos\theta|.
\end{equation}
При целевом угле она равна приблизительно
\begin{equation}
 -0.987490047.
\end{equation}

\begin{nogocriterion}[Провал казимирного завершения]
Представление \(2V_{1/2}\oplus3V_{3/2}\) совпадает с требуемой схемой по
размерности и среднему казимиру, но не по полному детерминанту. Поэтому оно
не является точным локальным завершением функционала щели.
\end{nogocriterion}

\section{Точный логарифмический рисунок}

Точный член
\begin{equation}
 \frac83\log\lambda_1
\end{equation}
возникает для восьми одинаковых вещественных зарядовых плоскостей с
\(|q|=1\). Такой весовой рисунок соответствует четырём комплексным
фундаментальным дублетам \(SU(2)\), имеющим суммарную вещественную
размерность \(16\).

Таким образом, кандидат
\begin{equation}
 \mathcal R_{\log}=4V_{1/2}
\end{equation}
точно воспроизводит логарифмическую часть функционала и сохраняет
антипериодическое действие центра.

\section{Исправление нормировочной ошибки}

При том же нормированном следе ранга \(3\) восемь зарядовых плоскостей дают
фиксированный зарядовый вклад
\begin{equation}
 \frac{8q^2}{6(1-c)}.
\end{equation}
Для требуемого коэффициента \(8/(1-c)\) необходимо
\begin{equation}
 q^2_{\mathrm{req}}=6.
\end{equation}
Каноническая осевая норма равна \(1\), поэтому согласованный с тем же
генератором вклад составляет
\begin{equation}
 \frac{4}{3(1-c)}.
\end{equation}
Ранее сюда был подставлен полный удвоенный казимир фундаментального
представления, равный \(3\). Такая подстановка неверна: детерминант на
окружности определяется выбранным осевым генератором, тогда как полный
казимир суммирует три генератора \(SU(2)\). Их смешение искусственно уменьшало
расхождение. При единой осевой нормировке недостаток равен не двум, а шести.

\begin{proposition}[Текущий точный статус роторной ветви]
Четыре фундаментальных дублета \(SU(2)\) точно воспроизводят весовой
детерминант и антипериодическую ветвь, но отождествление голономного веса с
фиксированным роторным импульсом оставляет шестикратный недостаток обратного
вклада.
\end{proposition}

Этот результат не закрывает роторную ветвь, поскольку голономный вес и
собственное значение фиксированного импульса являются различными квантовыми
числами. Он закрывает только их преждевременное отождествление.

Машинное основание сохранено в
\texttt{s2t\_wilson\_rotor\_exact\_determinant\_results.json}.

\chapter{Квантованный импульсный сектор граничного ротора}

\section{Разделение двух квантовых чисел}

Единичный вес \(q=1\) задаёт действие вильсоновской голономии на флуктуации
и тем самым определяет множитель
\begin{equation}
 \lambda_1=2-2\cos\theta.
\end{equation}
Фиксированный роторный импульс \(p\in\mathbb Z\), напротив, является
собственным значением канонического импульса циклической координаты. В
ансамбле фиксированного импульса он даёт рутьев вклад
\begin{equation}
 H_p(c)=\frac{p^2}{2(1-c)}.
\end{equation}
Следовательно, равенство \(p^2=q^2\) не является кинематическим тождеством.
Именно это неявное отождествление породило нормировочную ошибку предыдущей
проверки.

\section{Точный целочисленный перебор}

Для восьми вещественных вращательных плоскостей с единичным голономным весом
логарифмическая часть остаётся точной:
\begin{equation}
 \Gamma_{\det}
 =\frac13\sum_{i=1}^{8}\log\lambda_1
 =\frac83\log\lambda_1.
\end{equation}
Чтобы одновременно получить требуемый обратный член, фиксированные импульсы
должны удовлетворять
\begin{equation}
 \frac1{2\cdot3}\sum_{i=1}^{8}p_i^2=8,
 \qquad
 \sum_{i=1}^{8}p_i^2=48.
\end{equation}

Полный перебор неотрицательных целых импульсов даёт четырнадцать
неупорядоченных решений в диапазоне \(0\leq p_i\leq7\). Требование
положительности и нечётности каждого импульса, необходимое для одинакового
знака минус при полуповороте, оставляет две ветви:
\begin{equation}
 (1,1,1,1,1,3,3,5),
 \qquad
 (1,1,1,3,3,3,3,3).
\end{equation}
Если дополнительно применить принцип минимального максимального импульса,
то \(p_{\max}=3\), и остаётся единственный сектор
\begin{equation}
 \boxed{
 (p_1,\ldots,p_8)=(1,1,1,3,3,3,3,3).
 }
\end{equation}
Для него
\begin{equation}
 3\cdot1^2+5\cdot3^2=48,
\end{equation}
а потому
\begin{equation}
 \Gamma_{\mathrm{rot}}
 =\frac{48}{6(1-c)}
 =\frac8{1-c}.
\end{equation}

\begin{proposition}[Операторное совпадение двух коэффициентов]
Восемь единичных голономных весов и минимальный положительный нечётный
импульсный сектор \(1^3 3^5\) одновременно воспроизводят точные коэффициенты
\(8/3\) логарифмического члена и \(8\) обратного члена, не подменяя осевой
заряд полным казимиром.
\end{proposition}

\section{Новый критерий замыкания}

Полученное совпадение пока является операторным, а не динамическим. Не
выведены:
\begin{enumerate}
    \item проектор канонического интеграла по траекториям на сектор
    \(1^3 3^5\);
    \item принцип, запрещающий вторую нечётную ветвь с \(p_{\max}=5\);
    \item сохранение логарифмического детерминанта после фиксации импульсов;
    \item мера нулевых мод и полный комплекс БВ/БРСТ;
    \item связь кратностей \(3+5\) с геометрией произведения и осью поколений.
\end{enumerate}

\begin{nogocriterion}[Провал импульсного завершения]
Если проекция на фиксированный импульс удаляет требуемый детерминант,
порождает дополнительный множитель меры, не допускает единого нечётного
сектора или не выводит принцип минимального \(p_{\max}\), то совпадение
\(1^3 3^5\) остаётся целочисленной реконструкцией и не замыкает локальное
действие.
\end{nogocriterion}

Машинное основание сохранено в
\texttt{s2t\_wilson\_rotor\_momentum\_sector\_results.json}.

\chapter{Каноническая проекция: запрет однослойного ротора}

\section{Точный след по импульсным состояниям}

Для свободного квантового ротора с инерцией \(I(c)=1-c\) гамильтониан имеет
вид
\begin{equation}
 \widehat H=\frac{\widehat p^2}{2I(c)},
 \qquad
 p\in\mathbb Z.
\end{equation}
Канонический след равен
\begin{equation}
 Z(\beta,I)
 =\sum_{p\in\mathbb Z}
 \exp\left(-\frac{\beta p^2}{2I}\right).
\end{equation}
После проекции на конкретный импульсный сектор остаётся
\begin{equation}
 Z_p(\beta,I)
 =\exp\left(-\frac{\beta p^2}{2I}\right).
\end{equation}
Следовательно, его эффективное действие содержит обратную инерцию, но не
содержит независимого члена \(\log I\).

\section{Куда исчезает гауссов множитель}

Пуассоновское суммирование даёт эквивалентную запись полного следа
\begin{equation}
 \sum_{p\in\mathbb Z}e^{-\beta p^2/(2I)}
 =
 \sqrt{\frac{2\pi I}{\beta}}
 \sum_{w\in\mathbb Z}
 e^{-2\pi^2 Iw^2/\beta}.
\end{equation}
Корневой множитель, который можно было бы принять за источник
логарифмического члена, относится к полной сумме по намоткам. При разложении
по отдельным импульсным секторам он не сохраняется как самостоятельный
детерминант. Численная проверка двух представлений проведена с остатком менее
\(10^{-13}\).

Для восьми роторов в секторе \(1^3 3^5\) поэтому получается
\begin{equation}
 \Gamma_{p}
 =\frac8{1-c},
 \qquad
 \Gamma_{\log}^{\mathrm{same\ carrier}}=0,
\end{equation}
а не требуемая сумма двух членов.

\begin{nogocriterion}[Провал однослойного роторного завершения]
Один и тот же канонический роторный носитель не может одновременно дать
фиксированный импульсный вклад \(8/(1-c)\) и логарифмический детерминант
\((8/3)\log(1-c)\). Проекция, необходимая для первого члена, удаляет второй
как независимый секторный множитель.
\end{nogocriterion}

\section{Оставшийся допустимый класс}

После этой проверки сохраняется только двухслойная конструкция:
\begin{enumerate}
    \item восемь единично заряженных гауссовых граничных плоскостей создают
    точный логарифмический детерминант;
    \item отдельный топологический роторный слой в секторе \(1^3 3^5\)
    создаёт точный положительный обратный член;
    \item одно исходное симметрийное действие должно связать их меры,
    кратности и относительную нормировку до сравнения с наблюдаемыми.
\end{enumerate}
Без третьего условия конструкция является суммой двух специально подобранных
носителей и нарушает запрет скрытой подгонки.

{\raggedright
Машинное основание сохранено в
\texttt{s2t\_rotor\_fixed\_momentum\_projection\_results.json}.
\par}

\chapter{Полный БВ-комплекс скрученного двухформенного поля}

\section{Минимальный комплекс и единая скрутка}

Пусть \(L_-\to S^1\) --- плоское вещественное расслоение с голономией
\(-1\). Для бозонного двухформенного поля
\begin{equation}
 B_2\in\Omega^2(K,L_-)
\end{equation}
калибровочное преобразование и его приводимость имеют вид
\begin{equation}
 \delta B_2=d_L\Lambda_1,
 \qquad
 \Lambda_1\sim\Lambda_1+d_L\phi_0.
\end{equation}
Минимальный БВ-комплекс поэтому содержит грассманов однородный дух
\(c_1\in\Omega^1(K,L_-)\) и бозонный дух для духа
\(c_0\in\Omega^0(K,L_-)\). Все три ступени обязаны быть сечениями одного
расслоения \(L_-\), поскольку ковариантный дифференциал не меняет
коэффициентное расслоение. Полусдвиг нельзя назначить только выгодному
детерминанту.

\section{Две эквивалентные формы детерминанта}

После калибровочного фиксирования полная форма записи имеет вид
\begin{equation}
 Z_2\propto
 \det{}'(\Delta_2)^{-1/2}
 \det{}'(\Delta_1)^{+1}
 \det{}'(\Delta_0)^{-3/2},
\end{equation}
а после разложения Ходжа ---
\begin{equation}
 Z_2\propto
 \det{}'(\Delta_2^T)^{-1/2}
 \det{}'(\Delta_1^T)^{+1/2}
 \det{}'(\Delta_0)^{-1/2}.
\end{equation}
Для эффективного действия \(\Gamma=-\log Z\) первая запись с полными рангами
\((6,4,1)\) даёт
\begin{equation}
 \frac12\cdot6-1\cdot4+\frac32\cdot1=\frac12.
\end{equation}
Вторая запись с поперечными рангами \((3,3,1)\) независимо даёт
\begin{equation}
 \frac12\cdot3-\frac12\cdot3+\frac12\cdot1=\frac12.
\end{equation}
Совпадение двух представлений является внутренней проверкой: приводимая
калибровочная симметрия сокращает шесть компонент двухформенного потенциала
до одной физической степени свободы, как и требует четырёхмерная дуальность
со скаляром.

\section{Точный полусдвинутый гессиан}

Для общей массовой деформации каждой ступени
\begin{equation}
 \sum_{n\in\mathbb Z}\frac1{(n+1/2)^4}=\frac{\pi^4}{3}.
\end{equation}
Пространственно постоянная оболочка полного бозонного комплекса даёт
\begin{equation}
 \boxed{H_{B_2}^{(\ell=0)}=-\frac{\pi^4}{6}}.
\end{equation}
Обращение статистики всего комплекса меняет знак, но не число физических
степеней свободы:
\begin{equation}
 \boxed{H_{\Pi B_2}^{(\ell=0)}=+\frac{\pi^4}{6}}.
\end{equation}
На уровне только этой оболочки точное значение \(+\pi^4\) потребовало бы
шести независимых комплексов с обращённой статистикой. Ранг шесть одного
двухформенного поля не заменяет шесть полных калибровочных комплексов.

\section{Пространственно-гармонический контроль}

Вещественные числа Бетти \(\mathbb{RP}^3\) равны
\begin{equation}
 (b_0,b_1,b_2,b_3)=(1,0,0,1).
\end{equation}
Следовательно, на \(\mathbb{RP}^3\) нет гармонического триплета одно- или
двухформенных полей. Единственная чистая башня окружности в стандартном
комплексе происходит из пространственно постоянного скалярного духа для
духа и снова даёт \(-\pi^4/6\) в бозонной теории.

Для ненулевого пространственного собственного значения \(x^2\) вклад имеет
вид
\begin{equation}
 \sum_{n\in\mathbb Z}
 \frac1{\bigl((n+1/2)^2+x^2\bigr)^2}
 =
 \frac{\pi\tanh(\pi x)}{2x^3}
 -\frac{\pi^2}{2x^2\cosh^2(\pi x)}.
\end{equation}
Он зависит от полного спектра форм на \(\mathbb{RP}^3\) и не является
универсальным рациональным кратным \(\pi^4\). При больших чётных \(\ell\)
с учётом кратности \((\ell+1)^2\) вклад одной оболочки ведёт себя как
\begin{equation}
 \frac{\pi}{4\ell}.
\end{equation}
Поэтому сумма по оболочкам расходится логарифмически. Полный локальный
четырёхмерный гессиан требует контрчлена и схемы вычитаний; конечное число
\(\pi^4\) из одной башни окружности не является полным БВ-результатом.

\begin{theorem}[Вердикт стандартного БВ-замыкания]
Стандартный локальный комплекс одного скрученного двухформенного поля на
\(\mathbb{RP}^3\times S^1\) не выводит конечный положительный гессиан
\(H=\pi^4\). Пространственно постоянная оболочка даёт \(-\pi^4/6\) для
бозонной статистики и \(+\pi^4/6\) после полного обращения статистики, а
полный четырёхмерный след логарифмически расходится. Поэтому стандартная
БВ-ветвь закрыта как беспараметрический вывод электромагнитного остатка.
\end{theorem}

\begin{nogocriterion}[Запрет шести постфактум введённых комплексов]
Шесть комплексов с обращённой статистикой арифметически дают \(+\pi^4\)
только в выделенной пространственно постоянной оболочке. Они не считаются
выводом, пока их кратность, скрутка, мера, схема вычитаний и источник не
следуют из одного заранее объявленного градуированного действия и не проходят
второй независимый сектор.
\end{nogocriterion}

Машинное основание сохранено в
\texttt{s2t\_twisted\_twoform\_bv\_complete\_results.json}.

\chapter{Остаточная симметрия двух осей поколений}

\section{Точный централизатор факторных операторов}

Два направления, оставшиеся после минимизации всех трёх факторных ядер,
соответствуют в четырёхсостоянийном носителе векторам
\begin{equation}
 n_-\sim e_1-e_2,
 \qquad
 n_+\sim e_0-e_3.
\end{equation}
Они ортогональны. Для проверки их физической различимости был перебран весь
\(S_4\) и найден точный централизатор двух факторных трансляций
\(T_{\mathbb{RP}^3}\) и \(T_{S^1}\). Он имеет порядок четыре:
\begin{equation}
 G_{\mathrm{res}}
 =\{e,(01)(23),(02)(13),(03)(12)\}
 \simeq \mathbb Z_2\times\mathbb Z_2.
\end{equation}
Обмен самих факторов не включается в эту группу, поскольку
\(\mathbb{RP}^3\) и \(S^1\) неэквивалентны и несут разные веса.

Элемент \((01)(23)\) переводит
\begin{equation}
 e_1-e_2\longmapsto e_0-e_3
\end{equation}
с точностью до знака и одновременно сопрягает транспозицию \((12)\) в
\((03)\). Следовательно, соответствующие проекторы удовлетворяют
\begin{equation}
 P_{n_+}=Q P_{n_-}Q^{-1},
 \qquad
 Q=(01)(23),
\end{equation}
а вся вильсоновская конструкция с фиксированным углом переходит сама в себя.

\begin{theorem}[Единственность осевой орбиты]
Два минимума факторного функционала образуют одну проективную орбиту
остаточной группы \(G_{\mathrm{res}}\). Размер орбиты равен двум, а стабилизатор
каждого проективного направления имеет порядок два. Поэтому равенство энергий
двух осей является точным следствием симметрии, а не случайным вырождением.
\end{theorem}

Любой скалярный инвариант нескошенной факторной геометрии постоянен на этой
орбите. Поэтому искать дополнительный внутренний функционал, различающий
\(n_+\) и \(n_-\), было бы ошибкой: выбор конкретного представителя возможен
только при спонтанном или явном нарушении остаточной симметрии.

\begin{proposition}[Уточнение селекторного статуса]
Проблема выбора одной оси в отдельном секторе сокращается
\(8\to2\to1\) после факторизации по остаточной симметрии. Незакрытым остаётся
не выбор представителя, а относительная ориентация осевых орбит в верхнем и
нижнем кварковых секторах.
\end{proposition}

Машинное основание сохранено в
\texttt{s2t\_family\_residual\_axis\_orbit\_results.json}.

\chapter{Относительная осевая орбита и минимальный критерий CKM}

\section{Два относительных класса}

В базисе трёх нетривиальных характеров факторный оператор диагонален:
\begin{equation}
 L=\diag\bigl(2w_3,2w_1,2(w_3+w_1)\bigr).
\end{equation}
Два представителя единственной осевой орбиты имеют вид
\begin{equation}
 n_\pm=\frac1{\sqrt2}(1,\pm1,0)^T.
\end{equation}
Для пары верхнего и нижнего секторов диагональное действие остаточной группы
оставляет два относительных класса:
\begin{enumerate}
    \item совпадающий класс \(n_u=n_d\);
    \item перекрёстный класс \(n_u\perp n_d\).
\end{enumerate}
Их дискретный относительный инвариант равен соответственно
\begin{equation}
 |n_u^Tn_d|=1
 \quad\text{или}\quad
 |n_u^Tn_d|=0.
\end{equation}

\section{Минимальный массовый оператор}

Рассмотрим наиболее экономный оператор сектора
\begin{equation}
 M_s=L(w_3,w_1)+gA_{n_s},
 \qquad
 A_n=c_*I+(1-c_*)nn^T.
\end{equation}
Совпадающий класс имеет общий собственный базис и даёт единичную матрицу
смешивания. В перекрёстном классе первые два характера образуют блок
\(2\times2\), тогда как третий характер полностью отделён. Относительная
матрица имеет вид
\begin{equation}
 V_{ud}
 =
 \begin{pmatrix}
 \cos\theta_{\mathrm{mix}}&\sin\theta_{\mathrm{mix}}&0\\
 -\sin\theta_{\mathrm{mix}}&\cos\theta_{\mathrm{mix}}&0\\
 0&0&1
 \end{pmatrix},
\end{equation}
где
\begin{equation}
 \tan\theta_{\mathrm{mix}}
 =\frac{g(1-c_*)}{2(w_3-w_1)}.
\end{equation}

Таким образом, дискретный выбор относительного класса не фиксирует угол:
он непрерывно зависит от невыведенного отношения коэффициента смежности к
разности факторных весов. Кроме того, третье поколение не смешивается, а
матрица вещественна и имеет нулевую CP-фазу.

\begin{nogocriterion}[Провал минимального осевого CKM]
Операторы вида \(L+gA_n\) не дают беспараметрическую полную матрицу CKM.
Совпадающий класс даёт единицу, а перекрёстный --- только одно вещественное
двухсемейное вращение со свободным углом. Подбор \(g\) по углу Кабиббо
запрещён.
\end{nogocriterion}

Ранее доказанное равенство порождённой алгебры полной \(M_3\) не противоречит
этому запрету. Полная алгебра возникает из набора
\begin{equation}
 \{T_{\mathbb{RP}^3},T_{S^1},S,A_n\},
\end{equation}
но конкретный массовый оператор \(L+gA_n\) не содержит геометрический сдвиг
\(S\) и остаётся блочным. Следующий допустимый шаг должен вывести из одного
действия дополнительный некоммутирующий член, его коэффициент и комплексную
ориентацию до использования кварковых данных.

Машинное основание сохранено в
\texttt{s2t\_family\_relative\_orbit\_ckm\_results.json}.

\chapter{Полный перебор минимального операторного меню поколений}

\section{Заранее объявленный конечный класс}

Чтобы проверить, достаточно ли уже найденной полной алгебры для физического
вывода смешивания, рассмотрен конечный класс эрмитовых операторов
\begin{equation}
 M_s
 =L+aS+bA_n+d\,i[S,A_n],
 \qquad
 a,b,d\in\{-1,0,1\}.
\end{equation}
Здесь \(S\) --- ранее объявленный геометрический сдвиг, \(A_n\) --- оператор
вильсоновской оси, а \(i[S,A_n]\) --- минимальный эрмитов член, способный
создать комплексную ориентацию. Тройки коэффициентов верхнего и нижнего
секторов перебирались независимо. Экспериментальные массы кварков, элементы
CKM и CP-фаза в расчёт не загружались.

Для контроля произвольности один и тот же перебор выполнен для трёх
факторных ядер и трёх нормировок операторов: исходной, по спектральной норме
и по норме Фробениуса.

\section{Результат полного перебора}

В основной нормировке для каждого факторного ядра имеется \(54\) формальных
оператора сектора и \(2916\) упорядоченных верхне-нижних пар. После удаления
точных повторов остаётся \(51\) оператор сектора. Для ядра обратной длины:
\begin{center}
\begin{tabular}{lr}
\toprule
Свойство & Число пар \\
\midrule
Невырожденные спектры & 2916 \\
Полное ненулевое смешивание трёх поколений & 2748 \\
Ненулевой инвариант Ярлскога & 2352 \\
Различные абсолютные CP-рисунки & 533 \\
\bottomrule
\end{tabular}
\end{center}
Две другие нормировки и два других факторных ядра сохраняют число
CP-ненулевых пар \(2352\), хотя сами матричные элементы и значения инварианта
Ярлскога меняются.

Следовательно, добавление минимального коммутатора действительно снимает
структурный запрет: конечная алгебра легко создаёт невырожденные спектры,
смешивание всех трёх поколений и CP-нарушение. Однако именно лёгкость
получения результата становится новой проблемой: вместо одного вывода
возникают сотни неэквивалентных кандидатов.

\section{Абляционные проверки}

Перебор даёт три строгих вывода.

Во-первых, при
\begin{equation}
 d_u=d_d=0
\end{equation}
оба массовых оператора вещественны и инвариант Ярлскога тождественно равен
нулю во всех проверенных нормировках.

Во-вторых, если верхний и нижний секторы получают одинаковую тройку
\((a,b,d)\), то CP также равна нулю даже при выборе разных представителей
осевой орбиты. Следовательно, остаточная осевая геометрия сама по себе не
создаёт физическую комплексную фазу: требуется секторно-асимметричное правило
назначения коэффициентов.

В-третьих, одного ненулевого коммутаторного члена только в одном секторе уже
достаточно для полного смешивания и ненулевой CP. Это доказывает
существование механизма, но одновременно показывает отсутствие жёсткости
предсказания.

Знак \(d\) также не выбран. Преобразование
\begin{equation}
 d\longmapsto-d
\end{equation}
комплексно сопрягает массовый оператор. Без ориентированного или хирального
принципа две ветви являются CP-сопряжёнными и даже знак инварианта Ярлскога
остаётся неопределённым.

\begin{proposition}[Существование полного смешивания]
Конечное меню \(L+aS+bA_n+d\,i[S,A_n]\) содержит многочисленные операторы с
невырожденными трёхсемейными спектрами, полным смешиванием и ненулевой
CP-фазой. Алгебраического запрета на CKM больше нет.
\end{proposition}

\begin{nogocriterion}[Провал селектора конечного меню]
Конечный перебор не является предсказанием CKM, поскольку текущая геометрия
не выбирает одну тройку \((a,b,d)\) для каждого сектора, их относительную
нормировку и знак комплексной ориентации. Выбор кандидата после просмотра
кварковых данных является запрещённой дискретной подгонкой.
\end{nogocriterion}

Следующий допустимый шаг должен вывести обе тройки коэффициентов из
представлений исходного действия, граничного оператора Дирака или единого
правила путевого упорядочения. Если такого правила нет, операторное меню
закрывается со статусом «существование доказано, физический селектор не
найден».

Машинное основание сохранено в
\texttt{s2t\_family\_minimal\_operator\_menu\_exhaustive\_results.json}.

\chapter{Ориентированные вильсоновские слова и путевой селектор}

\section{Примитивные операторы и геометрическое ограничение}

В качестве примитивных шагов рассмотрены правильное вращение
\begin{equation}
 Q=-S,
 \qquad
 Q^2=I,
 \qquad
 \det Q=1,
\end{equation}
и вильсоновское вращение \(R_n(\theta_*)\) вокруг выбранной осевой орбиты,
а также его обратное \(R_n^{-1}\). Знак в определении \(Q\) переводит
отражение стандартного представления транспозиции в правильное
октаэдрическое вращение. Этот выбор сам является частью знакового расширения
и должен быть выведен граничной теорией.

Слова строились из алфавита
\begin{equation}
 \{Q,R,R^{-1}\}
\end{equation}
с удалением немедленных сокращений \(QQ\), \(RR^{-1}\) и \(R^{-1}R\).
Важно, что
\begin{equation}
 \pi_1(\mathbb{RP}^3\times S^1)
 \simeq\mathbb Z_2\times\mathbb Z
\end{equation}
абелева. Поэтому некоммутирующие \(Q\) и \(R\) не могут быть двумя плоскими
голономиями фундаментальных петель. Такая конструкция требует новой
неплоской связности, многорёберного комплекса или дефекта с заранее заданным
путевым упорядочением.

\section{Три чтения путевого слова}

Для каждого слова \(U\) проверены три минимальных отображения в массовый
оператор.

Прямое чтение
\begin{equation}
 Y=L+U
\end{equation}
оставляет \(Y\) вещественным. Левая диагонализация \(YY^T\) поэтому даёт
ортогональную матрицу смешивания и тождественно нулевую CP-фазу.

Симметрическое чтение
\begin{equation}
 M=L+\frac{U+U^T}{2}
\end{equation}
также является вещественным и даёт нулевой инвариант Ярлскога для всех слов.

Ненулевая CP возникает только в логарифмическом чтении
\begin{equation}
 M=L-i\operatorname{Log}U,
\end{equation}
где главный логарифм правильного вращения является эрмитовым чисто мнимым
антисимметрическим генератором. Это чтение вводит дополнительные условия:
ветвь матричного логарифма, масштаб генератора относительно \(L\) и правило
для собственных углов, равных \(\pi\).

\section{Полный перебор коротких слов}

Экспериментальные данные в перебор не загружались. Кумулятивные результаты
главного логарифмического чтения имеют вид:
\begin{center}
\small
\begin{tabular}{rrrr}
\toprule
Длина & Допустимые операторы & CP-пары & CP-рисунки \\
\midrule
1 & 4 & 0 & 0 \\
2 & 16 & 160 & 40 \\
3 & 36 & 1048 & 252 \\
4 & 84 & 6520 & 1598 \\
5 & 172 & 28552 & 7080 \\
\bottomrule
\end{tabular}
\end{center}
Первый CP-ненулевой пример возникает при общей длине двух секторных слов,
равной трём, например для пары \(R\) и \(QR\). Следовательно, путевое
упорядочение действительно способно породить комплексную ориентацию без
непрерывного коэффициента.

Однако увеличение допустимой длины не выбирает решение, а быстро увеличивает
число кандидатов. При длине не более четырёх изменение нормировки логарифма
также существенно меняет результат:
\begin{center}
\begin{tabular}{lrr}
\toprule
Нормировка & CP-ненулевые пары & Различные CP-рисунки \\
\midrule
Исходный главный логарифм & 6520 & 1598 \\
Единичная спектральная норма & 6360 & 844 \\
Деление на длину слова & 6552 & 1474 \\
\bottomrule
\end{tabular}
\end{center}

Слова с собственным углом \(\pi\) исключались из логарифмического перебора,
поскольку ось главного генератора в этой точке не определяется единственно.
Обращение ориентации слова меняет знак генератора и создаёт CP-сопряжённую
ветвь. Без ориентированного исходного комплекса этот знак также не выбран.

\begin{proposition}[Путевое происхождение CP возможно]
Некоммутирующее путевое произведение и эрмитов логарифм его голономии могут
создать полное трёхсемейное смешивание и ненулевую CP уже для слов общей
длины три. Тем самым существует операторный механизм возникновения
коммутаторной ориентации без непрерывной подгонки коэффициента.
\end{proposition}

\begin{nogocriterion}[Провал свободного путевого меню]
Перебор слов не является селектором, пока не заданы граф допустимых рёбер,
начальная и конечная вершины, правило назначения слов верхнему и нижнему
секторам, ветвь логарифма, нормировка генератора и ориентация пути. Выбор
слова после просмотра CKM является запрещённой дискретной подгонкой.
\end{nogocriterion}

Таким образом, путевое упорядочение повышает статус коммутаторного члена от
формальной вставки до реализуемого механизма, но не выбирает единственный
массовый оператор. Следующий положительный шаг возможен только после
построения конкретного ориентированного рёберного комплекса или неплоской
граничной связности, допускающей ровно одно минимальное слово в каждом
кварковом секторе.

Машинное основание сохранено в
\texttt{s2t\_family\_oriented\_wilson\_word\_results.json}.

\chapter{Модульная цепь, градуировка и минимальный CP-коцикл}

\section{Спектральное происхождение цепи}

Свободное меню вильсоновских слов оставило открытым граф допустимых путей.
Проверим, задаётся ли он уже выведенным факторным оператором
\begin{equation}
 L=\frac1\pi(I-T_{\mathbb{RP}^3})
   +\frac1{2\pi}(I-T_{S^1}).
\end{equation}
В совместном характерном базисе двух коммутирующих инволюций получаем точный
спектр
\begin{equation}
 \Spec L=\left\{\frac1\pi,\frac2\pi,\frac3\pi\right\}.
\end{equation}
Для невырожденного состояния
\begin{equation}
 \rho_\beta=\frac{e^{-\beta L}}{\Tr e^{-\beta L}},
 \qquad \beta>0,
\end{equation}
модульный оператор на \(M_3(\mathbb C)\) имеет вид
\begin{equation}
 \Delta_\rho(X)=\rho_\beta X\rho_\beta^{-1},
 \qquad
 K_{\mathrm{mod}}=-\log\Delta_\rho.
\end{equation}
Положительность \(\Delta_\rho\) однозначно фиксирует его самосопряжённый
логарифм. При этом \(K_{\mathrm{mod}}\) является супероператором на \(M_3\),
а не выбранной ветвью \(\operatorname{Log}U\) вильсоновской голономии;
отождествлять их без отдельного отображения нельзя.

Если уровни нумеровать по возрастанию, то
\begin{equation}
 K_{\mathrm{mod}}E_{ij}
 =\beta(\lambda_i-\lambda_j)E_{ij}.
\end{equation}
Минимальной положительной частоте \(\beta/\pi\) соответствуют ровно две
стрелки
\begin{equation}
 E_{10}:0\longrightarrow1,
 \qquad
 E_{21}:1\longrightarrow2,
\end{equation}
и ровно одна композиция длины два,
\begin{equation}
 E_{21}E_{10}=E_{20}.
\end{equation}
Следовательно, факторный спектр выводит ориентированную цепь \(A_3\), а
реальная структура \(J(X)=X^*\) обращает её направление.

\section{Единственная градуировка}

Пусть
\begin{equation}
 \Gamma=\diag(\gamma_0,\gamma_1,\gamma_2),
 \qquad \gamma_i\in\{+1,-1\}.
\end{equation}
Условие нечётности конечного оператора Дирака,
\begin{equation}
 \Gamma D+D\Gamma=0,
\end{equation}
требует противоположных знаков на концах каждого ребра. Полный перебор
восьми назначений оставляет только
\begin{equation}
 (-,+,-),\qquad (+,-,+).
\end{equation}
Они различаются общим знаком, поэтому цепь имеет единственную проективную
градуировку
\begin{equation}
 \Gamma\simeq\diag(1,-1,1).
\end{equation}

\section{Когомологический запрет CP на дереве}

Пусть двум рёбрам общей \(U(1)\)-связности назначены произвольные фазы.
Вершинная калибровка действует как
\begin{equation}
 E_{ij}\longmapsto e^{i(\chi_i-\chi_j)}E_{ij}.
\end{equation}
Инцидентная матрица \(A_3\) имеет ранг два, откуда
\begin{equation}
 \dim H^1(A_3;U(1))=|E|-\operatorname{rank}B=2-2=0.
\end{equation}
Все фазы одной общей древесной связности являются границами и снимаются
одной заменой базиса.

Прямой слепой тест использовал \(17\) значений каждой фазы, оба примитивных
ребра и чтение ``примитив против полной цепи''. Все
\begin{equation}
 2\times17^2=578
\end{equation}
случаев дали полное смешивание, но
\begin{equation}
 |J_{CP}|<5.1\times10^{-17}.
\end{equation}
Независимые фазы верхнего и нижнего секторов способны дать CP, однако это
уже две разные связности и невыведенный относительный коцикл.

\begin{theorem}[Tree CP no-go]
Одна общая \(U(1)\)-связность на модульной цепи \(A_3\) не порождает
физическую CP-фазу. Ненулевая CP требует либо секторно-различных связностей,
либо добавления цикла.
\end{theorem}

\section{Минимальный цикл и точный CP-инвариант}

Добавление хорды \(E_{20}\) замыкает цепь в треугольник. Теперь
\begin{equation}
 \dim H^1=3-2=1
\end{equation}
и существует единственный калибровочно-инвариантный поток
\begin{equation}
 \Phi=\phi_{10}+\phi_{21}-\phi_{20}.
\end{equation}
Для минимального разбиения
\begin{equation}
 M_u=L+H_{20}(\phi_{20}),
 \qquad
 M_d=L+H_{10}(\phi_{10})+H_{21}(\phi_{21})
\end{equation}
символьное вычисление для самих эрмитовых матриц даёт вспомогательный
инвариант
\begin{equation}
 \Tr[M_u,M_d]^3
 =\frac{12i(1+\pi^2)}{\pi^3}\sin\Phi.
\end{equation}
Стандартный кварковый CP-тест должен строиться из
\begin{equation}
 H_u=M_uM_u^\dagger,
 \qquad
 H_d=M_dM_d^\dagger.
\end{equation}
Независимый red-team пересчёт даёт
\begin{equation}
 \Tr[H_u,H_d]^3
 =\frac{192i(\pi^2-15)(2\pi^2-15)(1+\pi^2)}{\pi^9}\sin\Phi.
\end{equation}
Коэффициент ненулевой. Поэтому один треугольный поток достаточен для
физической CP при \(\sin\Phi\neq0\), но именно вторая формула является
стандартным инвариантом квадратов масс.

Однако хорда соединяет вершины одинаковой \(\Gamma\)-чётности. Треугольник
не двудолен, и полный перебор показывает, что не существует градуировки,
при которой все три ребра одновременно нечётны:
\begin{equation}
 \#\{\Gamma:\Gamma D_\triangle+D_\triangle\Gamma=0\}=0.
\end{equation}
Значит, хорда не может быть третьим ребром того же нечётного оператора
Дирака. Она должна принадлежать отдельному чётному кривизностному или
композитному слою. Без расширения градуированного пространства её нельзя
называть стандартным нулеформенным Higgs-компонентом нечётной
Quillen-суперсвязности.

Минимальное положительное действие потока
\begin{equation}
 V(\Phi)=\kappa(1-\cos\Phi),\qquad\kappa>0,
\end{equation}
выбирает \(\Phi=0\) и сохраняет CP. При противоположном знаке выбирается
\(\Phi=\pi\), где CP также равна нулю. Нетривиальный поток требует
фрустрации, дискретного закона квантования или ориентированного члена, пока
не выведенных из исходной геометрии.

\begin{proposition}[Точная развилка семейного сектора]
Модульный спектр выводит цепь \(A_3\) и её единственную градуировку, но
дерево имеет строгий CP-no-go. Минимальный CP-порождающий цикл существует,
однако требует двухслойного parent action: нечётной Dirac-цепи и отдельной
чётной хорды с динамически выбранным ненулевым потоком.
\end{proposition}

\begin{nogocriterion}[Провал свободного треугольного потока]
Если хорда, поток \(\Phi\), его ориентация или коэффициенты потенциала
выбираются после просмотра CKM, конструкция является параметризацией, а не
выводом CP. Положительный статус возможен только если одно градуированное
действие выводит хорду, \(\Phi\notin\{0,\pi\}\) и проходит независимый
normalization-sensitive сектор.
\end{nogocriterion}

Машинное основание сохранено в
\texttt{s2t\_modular\_grading\_cocycle\_results.json}.

\chapter{Red-team двухслойной CKM-текстуры}

\section{Физический инвариант и градуировка}

Для CKM диагонализуются левые квадраты
\begin{equation}
 H_f=M_fM_f^\dagger,
 \qquad f=u,d,
\end{equation}
а стандартным CP-нечётным тестом является \(\Tr[H_u,H_d]^3\). В
безразмерной нормировке \(\pi L=\diag(1,2,3)\) рассмотрен ansatz
\begin{equation}
 M_u=\pi L+qH_{20}(-\Phi),
 \qquad
 M_d=\pi L+pH_{10}(0)+pH_{21}(0).
\end{equation}
Точный red-team результат имеет вид
\begin{equation}
 \Tr[H_u,H_d]^3
 =192ip^2q(2p^2-15)(q^2-15)(q^2+1)\sin\Phi.
\end{equation}
Нули при \(p^2=15/2\) и \(q^2=15\) соответствуют вырождениям квадратов
масс. Для модульных весов \(p=1,q=2\) вырождений нет.

Численно проверены тождество Ярлскога
\begin{equation}
 \Tr[H_u,H_d]^3
 =6iJ_{CP}
 \prod_{i<j}(m_{u,i}^2-m_{u,j}^2)
 \prod_{i<j}(m_{d,i}^2-m_{d,j}^2),
\end{equation}
инвариантность относительно общей диагональной репараметризации и обращение
знака \(J_{CP}\) при \(\Phi\mapsto-\Phi\).

Градуировочная проверка также уточняет постановку. Цепь антикоммутирует с
\(\Gamma=\diag(1,-1,1)\), тогда как хорда \(E_{20}\) с \(\Gamma\)
коммутирует. Поэтому текущая матрица является двухслойным градуированным
ansatz, но ещё не построенной Quillen-суперсвязностью.

\section{Blind-расчёт и PDG-контроль}

Условно отождествляя \(\cos\Phi\) с устойчивым седлом ранее найденного
gap-функционала и используя модульные веса \(p=1,q=2\), до загрузки CKM
получаем
\begin{equation}
 |V|=
 \begin{pmatrix}
 0.76319&0.30542&0.56944\\
 0.57735&0.57735&0.57735\\
 0.29018&0.75722&0.58515
 \end{pmatrix},
\end{equation}
\begin{equation}
 (\theta_{12},\theta_{23},\theta_{13})
 =(21.81^\circ,44.62^\circ,34.71^\circ),
 \qquad
 J_{CP}=0.05104.
\end{equation}
Перенос gap-функционала на \(\Phi\) остаётся гипотезой, поскольку общего
parent action для двух переменных нет.

После фиксации результата выполнено сравнение с Particle Data Group 2024:
\begin{equation}
 \sin\theta_{12}=0.22501,
 \quad
 \sin\theta_{23}=0.04183,
 \quad
 \sin\theta_{13}=0.003732,
 \quad
 J_{CP}=3.12\times10^{-5}.
\end{equation}
Минимальная текстура предсказывает неиерархические углы и завышает
\(J_{CP}\) более чем на три порядка.

Диагностическая оптимизация общего нечётного веса \(p\) и веса хорды \(q\),
проведённая только после blind-сравнения, в лучшей точке оставляет
\begin{equation}
 (\theta_{12},\theta_{23},\theta_{13})
 =(7.73^\circ,7.73^\circ,0.53^\circ).
\end{equation}
Следовательно, один общий вес двух нечётных рёбер структурно не способен
создать наблюдаемое разделение \(\theta_{12}\gg\theta_{23}\).

\begin{proposition}[Исправленный вердикт двухслойной ветви]
Физический CP-инвариант двухслойной матрицы ненулевой и проходит проверки
репараметризации. Однако конструкция не является выведенной
суперсвязностью, перенос gap-функционала условен, а симметричная нечётная
метрика проваливает CKM-иерархию.
\end{proposition}

\begin{nogocriterion}[Запрет неверного повышения статуса]
Запрещено использовать \(\Tr[M_u,M_d]^3\) вместо физического
\(\Tr[H_u,H_d]^3\) или называть чётную хорду готовым Higgs-компонентом
нечётной суперсвязности. Переоткрытие возможно только после расширения
градуированного пространства либо построения отдельного чётного
кривизностного слоя и вывода двух неравных нечётных весов до просмотра CKM.
\end{nogocriterion}

Машинное основание сохранено в
\texttt{s2t\_two\_layer\_physical\_ckm\_redteam\_results.json}.

\chapter{Минимальный двудольный \texorpdfstring{\(C_4\)}{C4}-лифт}

\section{Исчерпывающая классификация графа}

Предыдущий red-team оставил точную структурную задачу. Координатная цепь
\(A_3=0-1-2\) имеет градуировку
\begin{equation}
 \Gamma_3=\diag(1,-1,1),
\end{equation}
но замыкающая хорда \(0-2\) является чётной. Чтобы получить ненулевой
калибровочный поток внутри одного полностью нечётного графа, требуется
двудольный цикл. Любой цикл двудольного графа имеет чётную длину, поэтому
треугольник запрещён, а минимальная возможная длина равна четырём.

Добавим одну вершину \(3\) противоположной чётности и потребуем:
\begin{enumerate}
 \item сохранения рёбер \(0-1\) и \(1-2\);
 \item связности графа;
 \item нечётности каждого ребра относительно
 \(\Gamma_4=\diag(1,-1,1,-1)\);
 \item наличия хотя бы одного независимого цикла.
\end{enumerate}
Полный перебор всех простых графов на четырёх помеченных вершинах оставляет
единственный минимальный вариант
\begin{equation}
 C_4:\qquad 0-1-2-3-0.
\end{equation}
Его incidence-матрица имеет ранг \(3\), поэтому
\begin{equation}
 b_1=|E|-\operatorname{rank}B=4-3=1.
\end{equation}
Тем самым квадрат действительно несёт ровно один калибровочно-инвариантный
\(U(1)\)-поток. Это не выбор из меню: при фиксированных \(A_3\) и
градуировке минимальный lift уникален с точностью до переименования
вспомогательной вершины.

\section{Несовместимость с каноническим разложением \(1+3\)}

Предыдущая классификация относится к четырём координатным вершинам. Однако
канонические три поколения в четырёхмодовом affine-меню задаются не
координатным подпространством \(\operatorname{span}\{e_0,e_1,e_2\}\), а
ортогональным дополнением к равномерному синглету
\begin{equation}
 u=\frac12(1,1,1,1)^T,
 \qquad
 P_1=uu^\dagger,
 \qquad
 P_3=I-P_1.
\end{equation}
Поэтому добавление четвёртой координатной вершины ещё не является lift
канонического семейного триплета.

Для unit-edge adjacency-оператора квадрата с потоком \(\Phi\) точный размер
утечки равномерного синглета в триплет равен
\begin{equation}
 \left\|P_3A_{C_4}(\Phi)u\right\|^2
 =\frac{(1-\cos\Phi)(3+\cos\Phi)}4.
\end{equation}
Он обращается в нуль только при \(\Phi=0\pmod{2\pi}\), то есть именно в
CP-сохраняющей точке. Для условно перенесённого Wilson-потока квадрат нормы
утечки равен \(0.9905101417\). Следовательно, ненулевой поток почти
максимально смешивает канонический синглет с семейным триплетом.

Кроме того,
\begin{equation}
 \left\|[\Gamma_4,P_1]\right\|_F^2=2.
\end{equation}
Градуировка \(\Gamma_4=\diag(1,-1,1,-1)\) переводит равномерный синглет в
sum-zero триплет и потому не спускается на каноническое разложение
\(\mathbb C^4=\mathbf1\oplus\mathbf3\).

Наконец, координатная вершина \(e_3\), исключаемая ниже, имеет
\begin{equation}
 \|P_1e_3\|^2=\frac14,
 \qquad
 \|P_3e_3\|^2=\frac34.
\end{equation}
Поэтому её Schur-исключение удаляет не affine-синглет, а преимущественно
компоненту самого семейного триплета. Это корневое ограничение конструкции,
предшествующее любому сравнению с CKM.

\section{Точный Schur-комплемент}

Пусть в диагностической координатной модели уровень вершины \(e_3\) равен
\(\mu\), а связи с
вершинами \(0\) и \(2\) собраны в вектор
\begin{equation}
 v=\begin{pmatrix}1&0&e^{i\Phi}\end{pmatrix}^{T}.
\end{equation}
Для unit-edge квадрата родительская матрица имеет блочный вид
\begin{equation}
 M_4=
 \begin{pmatrix}
  L_3+D_{A_3} & v\\
  v^\dagger & \mu
 \end{pmatrix}.
\end{equation}
Её внедиагональная часть строго антикоммутирует с \(\Gamma_4\). Исключение
вспомогательной компоненты даёт не постоянную матрицу, а точный
энергозависимый Schur-комплемент
\begin{equation}
 M_{eff}(\lambda)
 =L_3+D_{A_3}-\frac{vv^\dagger}{\mu-\lambda},
\end{equation}
для которого символьно проверено тождество
\begin{equation}
 \det(M_4-\lambda I)
 =(\mu-\lambda)\det(M_{eff}(\lambda)-\lambda I_3).
\end{equation}

Это уточняет прежнюю картину. При \(\lambda=0\) координатный обходной путь
порождает не только хорду \(0-2\), но одновременно диагональные сдвиги
\begin{equation}
 -\frac{vv^\dagger}{\mu}
 =-\frac1\mu
 \begin{pmatrix}
 1&0&e^{-i\Phi}\\
 0&0&0\\
 e^{i\Phi}&0&1
 \end{pmatrix}.
\end{equation}
Следовательно, вес хорды и сдвиги крайних уровней уже не независимы. Это
положительный алгебраический результат: квадрат объясняет, как чётная
эффективная хорда может возникнуть из полностью нечётного координатного
оператора. Он не доказывает сохранение канонических трёх поколений.

Канонический четырёхмодовый факторный лапласиан имеет спектр
\begin{equation}
 \left(0,\frac1\pi,\frac2\pi,\frac3\pi\right),
\end{equation}
причём равномерный синглет является нулевой, а не тяжёлой модой. Поэтому его
статическое исключение при \(\lambda=0\) вообще не определено без нового
массового оператора и его нормировки.

Даже для координатной вершины статическая подстановка \(\lambda=0\) является
точной физической
редукцией только если вспомогательная мода алгебраическая либо если из
действия выведен соответствующий низкоэнергетический предел. Для обычной
динамической тяжёлой моды точный оператор остаётся зависимым от
\(\lambda\).

\section{Проблема двух секторных чтений}

Один общий \(M_4\) и одна общая Schur-редукция дают один и тот же левый
собственный базис. Если этот оператор одинаково читать в \(u\)- и
\(d\)-секторах, то
\begin{equation}
 V_{CKM}=I,
 \qquad J_{CP}=0.
\end{equation}
Значит, сам квадрат ещё не выводит смешивание. Для продолжения требуется
объяснить, почему два сектора используют разные представления, проекторы
или разные пути одного parent-оператора.

Минимальный диагностический split назначает нижнему сектору исходный путь
\(0-1-2\), а верхнему --- путь через вспомогательную вершину:
\begin{equation}
 M_d=L_3+H_{10}+H_{21},
 \qquad
 M_u=L_3-\frac{vv^\dagger}{\mu}.
\end{equation}
Это уже не следствие одной Schur-редукции, а дополнительная гипотеза о
секторных projectors. Для неё точный физический инвариант равен
\begin{equation}
 \Tr[H_u,H_d]^3
 =\frac{1248i(2\mu-1)(15\mu-8)}{\mu^3}\sin\Phi.
\end{equation}
Нули \(\mu=1/2\) и \(\mu=8/15\) совпадают с вырождениями квадратов масс.
Формула проверена символьно, через тождество Ярлскога и на \(25\) случайных
точках.

\section{Blind-меню и отрицательный результат текстуры}

До загрузки CKM были объявлены unit-edge связи, условный поток из прежнего
continuous Wilson saddle и диагностическое, но не выведенное из спектра,
меню безразмерных координатных уровней
\begin{equation}
 \mu\in\{1,2,3,4,\pi,2\pi,\pi^2\}.
\end{equation}
Во всём меню получены диапазоны
\begin{equation}
 \theta_{12}=37.62^\circ\text{--}47.39^\circ,
 \qquad
 \theta_{23}=35.57^\circ\text{--}35.98^\circ,
\end{equation}
\begin{equation}
 \theta_{13}=7.06^\circ\text{--}10.68^\circ,
 \qquad
 |J_{CP}|=0.00575\text{--}0.04035.
\end{equation}
Таким образом, устранение градуировочного противоречия не устраняет
неиерархичность минимальной текстуры.

После blind-фиксации проведён диагностический scan
\(e^{-8}\leq\mu\leq e^{12}\). Лучшая sampled-точка уходит к почти
декаплированной вспомогательной моде и даёт
\begin{equation}
 (\theta_{12},\theta_{23},\theta_{13})
 \simeq(36.21^\circ,36.21^\circ,12.20^\circ).
\end{equation}
Минимальное sampled-значение \(\theta_{13}\) остаётся около
\(7.03^\circ\). Этот scan является сильным диагностическим провалом, но не
объявляется аналитическим доказательством для всех \(\mu>0\).

\begin{proposition}[Точный статус минимального \(C_4\)-лифта]
Квадрат является единственным минимальным двудольным расширением
координатной цепи \(A_3\), сохраняет нечётность всех рёбер и выводит
эффективную комплексную хорду вместе с фиксированными диагональными
сдвигами. Однако эта координатная градуировка несовместима с каноническим
affine-разложением \(1+3\), а ненулевой поток смешивает синглет и триплет.
Дополнительно один общий parent-оператор даёт единичное смешивание, тогда как
complementary-path split требует невыведенных секторных проекторов и уровня
\(\mu\). Поэтому ветвь закрывается раньше феноменологического CKM-теста.
\end{proposition}

\begin{nogocriterion}[Провал свободного \(C_4\)-расширения]
Если координатные вершины без доказательства отождествляются с тремя
поколениями, секторные представления, уровень \(\mu\), статическая точка
Schur-редукции или разные рёберные веса вводятся после просмотра CKM, то
\(C_4\)-лифт остаётся новой параметризацией. Положительный статус возможен
только в большем градуированном пространстве, где affine-синглет и триплет
инвариантны, либо после отказа от прежнего механизма счёта поколений и его
независимого повторного вывода.
\end{nogocriterion}

Машинное основание сохранено в
\texttt{s2t\_family\_bipartite\_c4\_lift\_results.json}.

\chapter{Хирально удвоенный семейный триплет}

\section{Корректный нечётный конечный оператор}

Координатный \(C_4\)-лифт провалился потому, что его градуировка не сохраняла
канонический affine-триплет. Минимальная постановка, сохраняющая сам
трёхмерный семейный носитель, состоит не в добавлении четвёртого поколения,
а в хиральном удвоении:
\begin{equation}
 \mathcal H_F=V_{3,L}\oplus V_{3,R},
 \qquad
 \Gamma_F=\diag(-I_3,+I_3).
\end{equation}
Для любого комплексного блока \(Y\in M_3(\mathbb C)\) оператор
\begin{equation}
 D_F(Y)=
 \begin{pmatrix}
 0&Y\\
 Y^\dagger&0
 \end{pmatrix}
\end{equation}
самосопряжён и удовлетворяет
\begin{equation}
 \{\Gamma_F,D_F\}=0,
 \qquad
 D_F^2=
 \diag(YY^\dagger,Y^\dagger Y).
\end{equation}
Тем самым физический левый оператор квадратов масс \(H=YY^\dagger\)
возникает автоматически. Максимальные численные ошибки самосопряжённости и
антикоммутатора во всём дискретном меню равны нулю.

Это исправляет постановку, но одновременно показывает предел одной
градуировки: любой комплексный \(3\times3\) блок \(Y\), содержащий \(18\)
вещественных параметров, совместим с теми же условиями. Для двух секторов до
факторизации по базисным преобразованиям имеется \(36\) вещественных
направлений. Следовательно, нечётность \(D_F\) сама по себе не является
правилом выбора Yukawa-матриц.

\section{Физический пересчёт дискретного меню}

Повторно проверено ранее объявленное меню
\begin{equation}
 Y_s=L+aS+bA_n+d\,i[S,A_n],
 \qquad
 (a,b,d)\in\{-1,0,1\}^3,
\end{equation}
для двух остаточных осей, трёх нормировок и трёх факторных ядер. В отличие
от прежнего предварительного перебора теперь во всех случаях
диагонализовался именно \(YY^\dagger\), а не непосредственно эрмитов блок
\(Y\).

Для raw inverse-length ветви получено:
\begin{center}
\begin{tabular}{lr}
\toprule
Величина & Результат \\
\midrule
Секторных кандидатов до удаления повторов & \(54\) \\
Различных Yukawa-блоков & \(51\) \\
Невырожденных по квадратам масс кандидатов & \(54\) \\
Упорядоченных двухсекторных пар & \(2916\) \\
Пар с полным смешиванием & \(2748\) \\
Пар с ненулевым \(J_{CP}\) & \(2352\) \\
Различных абсолютных CP-сигнатур & \(534\) \\
\bottomrule
\end{tabular}
\end{center}
Ни одна из объявленных матриц не приобрела вырождения квадратов масс.
Независимая SVD-проверка совпала с диагонализацией \(YY^\dagger\) с ошибкой
спектра \(5.33\times10^{-15}\) и ошибкой левых сингулярных направлений
\(6.66\times10^{-16}\).

Качественные ablation-результаты сохранились. Если в обоих секторах удалён
коммутаторный член, CP равен нулю. Одинаковые тройки коэффициентов в
\(u\)- и \(d\)-секторах также не дают CP. Но разрешение независимых троек
сразу порождает \(2352\) CP-ненулевые пары. Значит, удвоение не уменьшает
неоднозначность: оно лишь помещает прежнее меню в корректную
лево--правую архитектуру.

\begin{proposition}[Статус хирально удвоенного триплета]
Хиральное удвоение сохраняет канонические три поколения, даёт строго
нечётный самосопряжённый конечный оператор и автоматически порождает
физические квадраты масс. Однако градуировка допускает произвольный
\(Y\in M_3(\mathbb C)\), а существующее дискретное меню оставляет сотни
неэквивалентных CP-сигнатур. Поэтому получена корректная кинематическая
оболочка, но не вывод Yukawa-сектора.
\end{proposition}

\begin{nogocriterion}[Провал свободного Yukawa-блока]
Если конечная алгебра, её представление, вещественная структура,
условие первого порядка и секторное различие \(Y_u,Y_d\) не фиксированы до
просмотра масс и CKM, то выбор двух Yukawa-блоков является подгонкой внутри
разрешённого \(M_3(\mathbb C)\). Следующий положительный статус возможен
только после классификации блоков, переживающих все аксиомы конечной
геометрии.
\end{nogocriterion}

Машинное основание сохранено в
\texttt{s2t\_chiral\_doubled\_triplet\_yukawa\_results.json}.

\chapter{Что фактически достигнуто: от наивной модели к строгой программе}

\section{Начальные принципы наивной модели}

Исходная популярная модель строилась не от спектрального оператора, а от
последовательности простых геометрических образов. Её отправной тезис состоял
в смене точки зрения: не частицы являются первичными объектами на фоне
пустого пространства, а само пространство является единственной основной
сущностью, тогда как материя представляет собой его устойчивые дефекты.
Частица сравнивалась с узлом на верёвке: узел выглядит отдельным объектом, но
состоит из того же материала, что и сама верёвка.

Второй принцип был принципом компактного резонатора. Обычная волна в
бесконечном пространстве рассеивается, тогда как стабильная частица должна
сохранять энергию. Поэтому фундаментальная геометрия представлялась замкнутой
«банкой», в которой волна многократно отражается и образует стоячий режим.
Стабильными объявлялись только те конфигурации, которые после полного обхода
возвращаются в фазе и усиливают сами себя. Этот механизм назывался принципом
самосогласованности или Уробороса.

Третьим шагом был выбор циклической геометрии \(S^3\times S^1\). В исходном
рассуждении переход от сферы к тороидальной структуре мотивировался
необходимостью непрерывного потока без «мёртвой зоны», а дополнительное
скручивание объяснялось аналогиями с лентой Мёбиуса и расслоением Хопфа.
Предполагалось, что после склейки исчезает разделение на внутреннюю и внешнюю
стороны, а полный обход создаёт нетривиальный фазовый переворот.

Эта часть была именно эвристической. В частности, аналогия с теоремой о
причёсывании ежа не является строгим доказательством невозможности
\(S^3\): трёхмерная сфера параллелизуема. В строгой программе выбор
\(\mathbb{RP}^3\times S^1\) поэтому был заново поставлен как отдельная
задача геометрического и спектрального отбора.

Четвёртый принцип описывал вакуум как ткань, сплетённую из одной
фундаментальной «нити света». Многомерный объём интерпретировался как
результат плотной намотки одномерной траектории, подобно клубку нити или
кривой, заполняющей пространство. Отсюда возникала идея измерять не локальные координаты,
а полную геометрическую ёмкость замкнутого пути.

Пятый принцип вводил три вклада единой формы:
\begin{enumerate}
    \item объёмное ядро;
    \item поверхность склейки;
    \item голономию полного обхода.
\end{enumerate}
В единичной нормировке им сопоставлялись
\begin{equation}
 4\pi^3,\qquad \pi^2,\qquad \pi,
\end{equation}
а их сумма трактовалась как идеальный геометрический импеданс вакуума:
\begin{equation}
 S_{\mathrm{geo}}=4\pi^3+\pi^2+\pi,
\end{equation}
\begin{equation}
 4\pi^3
 =\Vol(S^3)\,\operatorname{Length}(S^1),
 \qquad
 \pi^2=\Vol(\mathbb{RP}^3).
\end{equation}
Сложение величин разной геометрической природы предполагалось допустимым в
планковских единицах, но ещё не выводилось из единого действия.

Шестой принцип утверждал, что идеальная гладкая геометрия должна получить
малую поправку из-за дискретной структуры вакуума. В популярной версии она
связывалась с 24-мерной решёткой Лича и записывалась как
\begin{equation}
 \delta_{24}
 =\frac1{24S_{\mathrm{geo}}}.
\end{equation}
Число \(24\) читалось как размерность максимально плотной структуры, а
поправка --- как потеря геометрической ёмкости из-за зернистости.

Седьмой принцип добавлял ещё меньшую поправку. Вакуумный «кристалл»
предполагался не абсолютно холодным, а обладающим нулевыми колебаниями и
термодинамической потерей. Четвёртая степень \(\pi\) связывалась с
четырёхмерным тепловым интегралом, а эффект второго порядка --- с квадратом
основного геометрического масштаба:
\begin{equation}
 \delta_{4}
 =\frac1{\pi^4S_{\mathrm{geo}}^2}.
\end{equation}

Наконец, восьмым и наиболее сильным постулатом было прямое отождествление
полученного вакуумного скаляра с обратной постоянной тонкой структуры:
\begin{equation}
 S_{\mathrm{vac}}
 =
 S_{\mathrm{geo}}
 -\frac1{24S_{\mathrm{geo}}}
 -\frac1{\pi^4S_{\mathrm{geo}}^2},
 \qquad
 \alpha^{-1}\equiv S_{\mathrm{vac}}.
\end{equation}
Таким образом, исходная цепочка имела вид
\begin{equation}
\boxed{
\begin{gathered}
\text{пространство первично}
\longrightarrow
\text{материя как топологический дефект}
\\
\longrightarrow
\text{компактный резонатор}
\longrightarrow
\text{скрученная циклическая геометрия}
\\
\longrightarrow
\text{геометрический импеданс}
\longrightarrow
\text{решёточная и тепловая поправки}
\\
\longrightarrow
\alpha^{-1}.
\end{gathered}
}
\end{equation}

Численная точность этой формулы и стала причиной дальнейшего исследования.
Но каждый переход в исходной цепочке ещё требовал отдельной проверки:
топологические метафоры не фиксировали оператор, решётка Лича не выводила
автоматически коэффициент детерминанта, а аналогия с чёрным телом не
доказывала точную обратную \(\pi^{-4}\)-нормировку.

\section{Как трактат перевёл эти принципы в проверяемую форму}

Главное достижение трактата состоит в том, что общая философия была переведена
в технические требования:
\begin{center}
\small
\begin{tabular}{>{\raggedright\arraybackslash}p{0.27\textwidth}>{\raggedright\arraybackslash}p{0.31\textwidth}>{\raggedright\arraybackslash}p{0.30\textwidth}}
\toprule
Первоначальная идея & Технический перевод & Результат \\
\midrule
Мир есть структура согласованности &
пространство-носитель, гильбертовы пространства, операторы и допустимые редукции &
гипотезу можно опровергать по частям \\
Частицы являются устойчивыми режимами &
спектральные башни, ядра, дефектные моды и секторы эффективной теории поля &
появился язык существования и устойчивости состояний \\
Геометрия не является фоном &
\(\mathbb{RP}^3\times S^1\), накрытие, объёмы и циклы &
получен конкретный источник нормировок и граничных условий \\
Топология хранит физическую информацию &
\(\mathbb Z_2\)-голономия, спиновые структуры и секторы фактор-пространства &
дискретные фазы перестали быть свободными украшениями \\
Сектора имеют общий источник &
единая мера и обязательный двухсекторный критерий &
набор независимых формул больше нельзя называть объединением \\
Совпадения должны что-то означать &
слепая проверка, перестановочный контроль, анализ чувствительности и проверка исключением &
численная близость отделена от операторного вывода \\
\bottomrule
\end{tabular}
\end{center}

Таким образом, фраза «всё может быть проявлением одной структуры» превратилась
в список условий, после которых разрешено или запрещено говорить, что общий
механизм действительно найден.

\section{Математическое ядро}

К настоящему моменту получены следующие устойчивые результаты.
\begin{enumerate}
    \item Рабочий носитель зафиксирован как
    \(K=\mathbb{RP}^3\times S^1\), а фермионное накрытие --- как
    \(S^3\times S^1\). Ошибочное смешение фактор-пространства и накрытия
    теперь контролируется явно.

    \item Объём \(\mathbb{RP}^3\), длины циклов и нетривиальная
    \(\mathbb Z_2\)-голономия получили определённое место в формулах.
    Слагаемые \(\pi^2\) и \(\pi\) связаны с геометрией и глобальной фазой,
    а не просто выбраны из набора красивых чисел.

    \item Воспроизведены скалярный и козамкнутый спектры на
    \(L(2,1)\simeq\mathbb{RP}^3\), их кратности, отбор по чётности и
    детерминантные продолжения.

    \item Разделены локальные коэффициенты теплового ядра, глобальная голономия,
    казимировские ветви, детерминанты обмоточных мод, нулевые моды и
    детерминанты духовых полей. Наивная модель
    смешивала эти источники; трактат показал, что они не взаимозаменяемы.

    \item Построены модели дефектов, в которых топология действительно
    ограничивает ранг ядра и число допустимых ветвей. Это прямой технический
    пример тезиса о первичности глобальной допустимости.
\end{enumerate}

\section{Статус исходных численных формул}

Формулы не были удалены или объявлены бессмысленными. Их научный статус был
разделён.
\begin{enumerate}
    \item \(S_{\mathrm{geo}}=4\pi^3+\pi^2+\pi\) сохраняется как сильный
    исходная геометрическая конструкция.

    \item Поправка \(1/(24S_{\mathrm{geo}})\) связана с локальной
    детерминантной ветвью, но полная нормировка схемы Максвелла с духовыми
    полями не позволяет автоматически объявить её окончательным
    электромагнитным выводом.

    \item Старый C6-механизм не вывел
    \(1/(\pi^4S_{\mathrm{geo}}^2)\). Вместо этого обнаружено новое чтение:
    \(\pi^{-4}\) может быть обратной восприимчивостью, то есть обратной
    жёсткостью коллективной спектральной моды.

    \item Соотношение масс \(\tau\)-лептона и мюона прошло строгий
    аудит малосложных формул и сохранилось как необычно сильная структурная
    связь, но его исходный объёмный множитель и петлевая нормировка пока не
    получены из общего действия.

    \item Хиггсовский масштаб сохраняется как зависимая условная конструкция:
    он наследует силу и незакрытые нормировки лептонного сектора.

    \item Нейтринная подмодель и норма \(23+\pi^{-1}\) сохранились как
    содержательная операторная модель. Минимальное следово-ходжево действие,
    однако, не воспроизвело независимо второй сектор, чувствительный к
    нормировке.
\end{enumerate}

Трактат тем самым разделил четыре разных понятия: точное тождество, численно
сильную связь, модельный вывод и физическое предсказание.

\section{Трёхсемейность и смешивание}

Наивная версия говорила о трёх поколениях преимущественно на уровне аналогии.
Теперь имеется конечномерная операторная конструкция.
\begin{enumerate}
    \item Четыре спиновых состояния образуют меню \(F_2^2\), а три
    нетривиальных характера дают естественное трёхмерное пространство
    поколений.

    \item Два факторных оператора различают все три семейных направления без
    непрерывного угла смешивания.

    \item Одновременно доказан критерий невозможности для коммутирующих
    операторов: функции только от факторных операторов имеют общий
    собственный базис и дают
    \(V_{\mathrm{CKM}}=I\).

    \item Полный перебор аффинных операторов смежности нашёл двенадцать
    операторов,
    расширяющих семейную алгебру до полной \(M_3(\mathbb C)\). Значит,
    алгебраического запрета на полное смешивание больше нет.

    \item Проективный поток, когомологии, клеточный оператор Дирака и
    детерминантное расслоение проверены и закрыты там, где они не способны
    выбрать требуемый оператор.
\end{enumerate}

\section{Самый сильный новый результат II.B}

Два ранее независимых затруднения --- точный коэффициент \(\pi^{-4}\) и
полная алгебра поколений --- впервые помещены в одну конструкцию.

Перебор всех \(48\) \(S_4\)-голономий с дополнительным знаком показал, что
простая дискретная версия не работает: точный тензорный отклик даёт только
\(-I\), который не смешивает семейства. После перехода к непрерывной
вильсоновской петле в \(SO(3)\)
получено уравнение
\begin{equation}
 11c^2-52c-49=0,
 \qquad
 c_*=\frac{26-9\sqrt{15}}{11}.
\end{equation}
Для этого угла существуют восемь осей, которые одновременно:
\begin{enumerate}
    \item дают точный тензорный отклик \(H=\pi^4\);
    \item расширяют алгебру поколений до \(M_3\);
    \item происходят из классифицированных направлений вне \(D_8\).
\end{enumerate}

Более того, угол является устойчивым минимумом функционала
\begin{equation}
 V_{\mathrm{gap}}(c)
 =
 \frac83\left(\frac3{1-c}+\log(1-c)\right)
 -\frac{44}{45}c,
\end{equation}
поскольку
\begin{equation}
 V_{\mathrm{gap}}'(c)=R(c)-1,
 \qquad
 V_{\mathrm{gap}}''(c_*)\approx1.90165>0.
\end{equation}
Следовательно, нужный угол уже можно читать не как произвольно вставленное
число, а как устойчивую стационарную точку спектрального функционала.

Наконец, три ранее введённых ядра факторных операторов независимо сокращают восемь
совместных осей до двух:
\begin{equation}
 8\longrightarrow2.
\end{equation}
Исходная геометрия произведения впервые не только допускает совместное решение, но
и частично выбирает его.

\section{Отрицательные результаты тоже являются достижениями}

Большая часть продвижения состоит в том, что программа научилась запрещать
собственные неверные объяснения.
\begin{itemize}
    \item Чистый эффект Казимира и простая дуальная сумма не выводят всю
    электромагнитную формулу.
    \item Полная схема Максвелла с духовыми полями не разрешает оставить
    только численно выгодный вклад духов.
    \item Обычная ренормгрупповая эволюция дополнительного сектора расщепления
    имеет неверное
    направление поправок.
    \item Минимальный конус пороговых поправок не достигает требуемых
    электрослабых и КХД-сдвигов.
    \item Каноническая нормировка заряженных лептонов не создаёт нужные объёмные
    коэффициенты.
    \item Минимальное единое следово-ходжево действие не проходит два
    независимых сектора, чувствительных к нормировке.
    \item Внутренний селектор поколений не возникает автоматически из
    когомологий, проективного потока или канонического детерминантного
    расслоения; две последние оси, однако, оказались одной остаточной орбитой.
    \item Минимальный коммутаторный член допускает полное смешивание и CP, но
    конечный перебор порождает сотни кандидатов без правила выбора.
    \item Путевое упорядочение реализует комплексную ориентацию, но свободное
    меню слов, ветвей логарифма и нормировок увеличивает неоднозначность.
    \item Модульный спектр выводит цепь \(A_3\) и единственную градуировку.
    Первая когомология дерева тривиальна; минимальный CP-поток требует чётной
    хорды, несовместимой с полностью нечётным Dirac-графом.
    \item Red-team заменил вспомогательный \(\Tr[M_u,M_d]^3\) физическим
    \(\Tr[H_u,H_d]^3\) и запретил трактовать чётную хорду как готовый
    Higgs-компонент нечётной суперсвязности.
    \item Минимальный двудольный \(C_4\)-лифт устраняет конфликт
    координатной градуировки и выводит эффективную хорду через
    Schur-комплемент, но не сохраняет каноническое affine-разложение
    \(1+3\): ненулевой поток смешивает синглет и триплет.
    \item Одна дискретная общая голономия не связывает сектор поколений и
    тензорный сектор.
\end{itemize}

В наивной модели каждый провал можно было скрыть новым коэффициентом. Теперь
подобное продолжение прямо запрещено.

\section{Было и стало}

Эволюцию проекта можно записать так:
\begin{equation}
\begin{gathered}
\text{философская интуиция}
\longrightarrow
\text{геометрическая нумерология}
\longrightarrow
\text{операторная программа}
\\
\longrightarrow
\text{система воспроизводимых аудиторских проверок}
\\
\longrightarrow
\text{совместные двухсекторные решения}.
\end{gathered}
\end{equation}

Наивная модель спрашивала: «Можно ли составить из \(\pi\) красивое физическое
число?» Текущая программа спрашивает иначе:
\begin{quote}
Какое пространство-носитель, оператор, спектр, мера, статистика и граничное
условие
неизбежно порождают это число, одновременно объясняя хотя бы ещё один
независимый сектор?
\end{quote}

\section{Честный текущий статус}

Нельзя утверждать, что единая теория мира уже построена. Из одного локального
действия пока не выведены:
\begin{itemize}
    \item каноническая единичная жёсткость \(\kappa=1\);
    \item полный состав полей БВ/БРСТ вильсоновско-тензорного сектора;
    \item взаимодействие оси с факторным оператором;
    \item представленная конечная алгебра, вещественная структура и условие
    первого порядка, выделяющие \(Y_u,Y_d\) внутри хирально удвоенного
    affine-триплета;
    \item нормировка источника \(J^2/2=1\);
    \item общая редукция к электромагнитному сектору, сектору поколений и
    третьей независимой наблюдаемой величине.
\end{itemize}

Но возвращаться к исходной точке уже неверно. Получены:
\begin{equation}
\boxed{
\begin{gathered}
\text{фиксированное геометрическое ядро},\\
\text{воспроизводимые спектральные результаты},\\
\text{полная алгебра поколений},\\
\text{точная совместная тензорно-семейственная операторная конструкция},\\
\text{устойчивая стационарная точка спектрального функционала},\\
\text{единственная осевая орбита остаточной симметрии},\\
\text{строгая карта закрытых и открытых механизмов}.
\end{gathered}
}
\end{equation}

Наиболее точный итог таков: первоначальная философская концепция ещё не стала
завершённой физической теорией, но уже перестала быть только философией. Она
превратилась в ограниченную, воспроизводимую и способную к
самофальсификации математико-физическую программу, внутри которой впервые
найден конкретный мост между двумя ранее независимыми секторами.

\chapter{Аудит покрытия наблюдаемого мира}

Предыдущие главы проверяли отдельные механизмы. Теперь требуется обратный
аудит: перечислить свойства нашего мира, которые надёжно установлены, и
проверить, какие из них текущая версия S2T воспроизводит независимо, какие
получает только условно, а какие не выводит вообще. Входная величина,
postdiction или численно близкая формула с незамкнутым операторным механизмом
не считается физическим предсказанием.

\section{Численный scorecard}

Наиболее жёсткий независимый тест даёт electroweak/QCD-ветвь. При сохранении
электромагнитного train-якоря двухпетлевой split-running предсказывает
\begin{equation}
 M_W=82.0226\,\mathrm{GeV},\qquad
 M_Z=92.0616\,\mathrm{GeV},
\end{equation}
\begin{equation}
 \sin^2\theta_W=0.206203,\qquad
 \alpha_s(M_Z)=0.080177.
\end{equation}
Относительные невязки равны соответственно \(+2.06\%\), \(+0.96\%\),
\(-7.62\%\) и \(-32.05\%\). Варьирование top-Yukawa не исправляет результат,
а минимальная logarithmic-threshold ветвь уже закрыта отрицательно.

Физический CKM red-team даёт
\begin{equation}
 J_{S2T}=5.1041\cdot10^{-2},
 \qquad
 J_{CKM}=3.12\cdot10^{-5},
\end{equation}
и
\begin{equation}
 \sin\theta_{13}^{S2T}=0.5694,
 \qquad
 \sin\theta_{13}^{CKM}=0.003732.
\end{equation}
Следовательно, найденный CP-инвариант не воспроизводит наблюдаемую
иерархическую матрицу смешивания. Хиральное удвоение исправляет градуировку,
но допускает произвольный комплексный \(3\times3\) Yukawa-блок и потому не
создаёт предсказания.

Сверхточные численные строки имеют иной тип недостатка:
\begin{itemize}
 \item \(S_{vac}\) воспроизводит \(\alpha^{-1}\), но \(\alpha\) является
 train-якорем, а exact-\(\pi^{-4}\) determinant route закрыт отрицательно;
 \item \(m_\tau/m_\mu\) совпадает с контролем до \(3.2\cdot10^{-7}\), но
 seed \(\rho_0\) и projection normalization не выведены;
 \item нейтринное отношение отличается от контрольного примерно на
 \(0.90\%\), но integer \(23\) не получен representation-consistent способом,
 абсолютная шкала и PMNS-матрица отсутствуют;
 \item \(M_H\) численно близка к наблюдаемой массе, но наследует условную
 абсолютную шкалу и пока не сопровождается выводом couplings и ширины.
\end{itemize}

\section{Что уже покрывает коллективный атлас}

Coverage-аудит нельзя читать так, будто за пределами нескольких строк модель
вообще не содержит численных образов наблюдаемого мира. Ранее построенный
коллективный \(\pi\)-atlas включает одиннадцать формул для
\begin{equation}
 \alpha_s,\quad \sin^2\theta_W,\quad
 \frac{m_b}{m_p},\quad\frac{m_s}{m_p},\quad
 \frac{m_d}{m_e},\quad\frac{m_u}{m_e},\quad
 \frac{m_\tau}{m_\mu},\quad V_{cb},\quad
 \Omega_\Lambda,\quad\Omega_{\mathrm{dm}},\quad\Omega_b.
\end{equation}
При общей замене \(\pi\mapsto x\) его RMS logarithmic score равен
\begin{equation}
 E(\pi)=0.01827893,
 \qquad
 x_{\mathrm{best}}=3.16422,
 \qquad
 E(x_{\mathrm{best}})=0.01587965.
\end{equation}
На равномерной сетке \(x\in[2,4]\) значение \(\pi\) входит в лучшие
\(2.27\%\) баз, а leave-one-out optimum остаётся около \(3.16\).
Следовательно, атлас действительно даёт устойчивое многосекторное численное
сжатие: кварковые отношения, \(V_{cb}\) и космологические доли нельзя
называть полностью отсутствующими в программе.

Но этот положительный факт имеет строго ограниченный статус. Формулы и
назначения observables собраны ретроспективно; данные не выделяют именно
\(\pi\), поскольку \(\sqrt{10}\) и некоторые другие базы дают меньшую общую
ошибку. Не выведены правило адресации observable к формуле, единый оператор,
renormalization-scale dependence и динамика, порождающая соответствующие
массы, mixing и космологические доли.

\begin{proposition}[Различие атласного покрытия и физического вывода]
Атлас показывает, что первичный арифметический образ уже охватывает больше
наблюдаемых, чем отражал первоначальный coverage-реестр. Поэтому корректный
статус этих строк --- «ретроспективно воспроизведены, но операторно не
выведены», а не «отсутствуют». Это повышает связанность программы, но не
делает строки blind-предсказаниями.
\end{proposition}

\section{Что в мире есть, а модель пока не выводит}

С учётом атласа заголовок этой секции означает не отсутствие численной
формулы, а отсутствие единого выведенного механизма. Минимальный реестр
обязательных физических результатов имеет вид:
\begin{enumerate}
 \item единый mass/Yukawa-оператор для электрона, мюона, шести кварков и трёх
 нейтрино, включая абсолютные шкалы и running, несмотря на наличие части
 атласных массовых отношений;
 \item четыре совместных параметра CKM и полная PMNS-матрица, несмотря на
 отдельную атласную формулу для \(V_{cb}\);
 \item одновременное воспроизведение \(G_F,M_W,M_Z,\sin^2\theta_W\) и
 \(\alpha_s(M_Z)\) из gauge dynamics без threshold-fit, а не только из
 атласных формул;
 \item конфайнмент, адронный спектр, ширины распадов и scattering amplitudes;
 \item единый low-energy action с наблюдаемым набором полей, корректным
 BRST/BV-комплексом, квантовой мерой и anomaly cancellation;
 \item динамическое пространство-время, ньютоновская нормировка и
 гравитационные уравнения;
 \item космологические уравнения, объясняющие расширение, наблюдаемую
 недостающую гравитирующую массу, ускорение и барионную асимметрию, несмотря
 на атласное численное покрытие
 \(\Omega_\Lambda,\Omega_{\mathrm{dm}},\Omega_b\).
\end{enumerate}

Это не означает, что первичный образ опровергнут. Установлен более точный
результат: имеющаяся редукция первичного образа ещё не идентифицирует наш мир
среди эмпирически различимых продолжений. В частности, даже коллективное
совпадение одиннадцати атласных строк не заменяет вывод взаимодействий,
scale dependence и динамических наблюдаемых.

\begin{protocol}[Definition of Done покрытия наблюдаемого мира]
Минимальное физическое покрытие считается достигнутым только если один
заранее фиксированный parent action и один readout-map без новых коэффициентов:
\begin{enumerate}
 \item задают полный light-field content и его калибровочные представления;
 \item выводят массы и mixing обоих фермионных секторов;
 \item проходят совместный electroweak/QCD precision scorecard;
 \item дают хотя бы один адронный или scattering observable;
 \item содержат динамический гравитационный предел;
 \item после заморозки всех входов дают новую preregistered blind-величину.
\end{enumerate}
\end{protocol}

\begin{nogocriterion}[Запрет компенсации неполноты точностью одной константы]
Никакая точность \(\alpha^{-1}\), \(m_\tau/m_\mu\), \(M_H\) или иной отдельной
строки не компенсирует отсутствие CKM, PMNS, strong dynamics, gravity или
общего action. Если высокоточная строка использована как train-вход либо её
операторная нормировка остаётся свободной, она учитывается как указатель на
возможную структуру, но не как доказательство идентификации нашего мира.
\end{nogocriterion}

Текущий вердикт gate отрицателен: закрытых независимых многосекторных
физических предсказаний нет. Приоритетом становится не добавление следующей
константы, а вывод общего действия и карты чтения, способных одновременно
закрыть хотя бы electroweak/QCD scorecard и полные матрицы масс и смешивания.
Машинный реестр сохранён в
\texttt{s2t\_observed\_world\_coverage\_results.json}.

\chapter{Карта чтения и наблюдатель как недостающий слой}

Отрицательный coverage gate допускает две интерпретации. Либо первичный
источник недостаточен, либо текущая карта перехода от него к наблюдаемым
выбрана неверно. Проверим вторую возможность как конечномерную операторную
гипотезу.

\section{Допустимый класс карт}

Пусть \(\mathcal A_{\mathrm{source}}\) является алгеброй первичных операторов,
а \(\mathcal B_{\mathrm{obs}}\) --- подалгеброй физически доступных
наблюдаемых. Минимально допустимая карта чтения
\begin{equation}
 \mathcal R:\mathcal A_{\mathrm{source}}
 \longrightarrow\mathcal B_{\mathrm{obs}}
\end{equation}
является унитальной полностью положительной картой. Сильная версия требует
идемпотентности, сохранения выделенного состояния \(\omega\) и свойства
условного ожидания.

\section{Почему известные данные не выбирают карту}

Пусть \(\mathcal R_0\) совпадает с наблюдениями в конечном наборе
\(x_1,\ldots,x_n\). Тогда
\begin{equation}
 \mathcal R_\lambda(x)
 =\mathcal R_0(x)
 +\lambda\prod_{j=1}^{n}(x-x_j)
\end{equation}
совпадает с ней во всех известных точках, но меняет значение в новой точке.
Для набора \(\{-2,-1,0,1,2\}\) поправка равна нулю на train-точках и равна
\(120\) при \(x=3\). Конечное число известных констант не определяет
единственную карту чтения.

\begin{nogocriterion}[Запрет эмпирического выбора readout]
Если \(\mathcal R\), наблюдаемая подалгебра или её проектор выбираются после
просмотра физических констант, масс или mixing, то карта является
интерполяцией данных, а не редукцией первичного образа.
\end{nogocriterion}

\section{Центральные скаляры не исправляются наблюдением}

Для любой унитальной линейной карты
\begin{equation}
 \mathcal R(cI)=c\mathcal R(I)=cI.
\end{equation}
Это проверено для random-unitary CP-карты с пятью Kraus-ветвями: ошибка
сохранения \(137.035999177I\) меньше \(8\cdot10^{-14}\). Поэтому наблюдение
не исправляет неверную константу, если она кодирована центральным скаляром.
Для изменения числа требуется нецентральный оператор или новая динамика.

\section{Канонический симметрийный readout}

На четырёхсостоянийном affine-меню естественная полностью
\(AGL(2,2)\simeq S_4\)-ковариантная карта равна
\begin{equation}
 E_{S_4}(A)=\frac1{24}\sum_{g\in S_4}U_gAU_g^\dagger.
\end{equation}
Машинный расчёт даёт
\begin{equation}
 \dim\operatorname{Fix}(E_{S_4})=2,
 \qquad
 \operatorname{Fix}(E_{S_4})
 =\operatorname{span}\{P_{\mathrm{sing}},P_{\mathrm{trip}}\}.
\end{equation}
На каноническом триплете любой усреднённый оператор пропорционален \(I_3\) с
ошибкой меньше \(6\cdot10^{-16}\). Такая карта сохраняет разложение \(1+3\),
но уничтожает семейную иерархию, ось и CKM.

Даже fixed-point states не уникальны:
\begin{equation}
 \rho(p)=pP_{\mathrm{sing}}
 +(1-p)\frac{P_{\mathrm{trip}}}{3},
 \qquad 0\leq p\leq1.
\end{equation}
Существование неподвижной точки не эквивалентно неизбежности одного мира.

\section{Неординарный положительный результат}

Для изометрии \(V:\mathbb C^3\to\mathbb C^4\) компрессия
\begin{equation}
 \Phi_V(A)=V^\dagger A V
\end{equation}
является унитальной полностью положительной картой. Каноническая
affine-триплетная изометрия сохраняет
\begin{equation}
 [\Phi_V(T_{\mathbb{RP}^3}),\Phi_V(T_{S^1})]=0,
\end{equation}
поэтому прежний \(V_{\mathrm{CKM}}=I\) остаётся.

Но для всех \(256\) детерминированных случайных rank-three компрессий получен
ненулевой коммутатор. Его Frobenius-норма лежит между \(0.277\) и \(1.399\),
а медиана равна \(1.049\). Наблюдательная компрессия действительно может
создать эффективную некоммутативность и mixing из коммутирующей первичной
записи.

Этот результат не является решением: множество \(V\) непрерывно, и
произвольный выбор из грассманиана является новым selector-ом. Уже найденная
единственная осевая орбита также оставляет два относительных up/down-класса:
aligned и orthogonal.

\begin{proposition}[Исправленный статус гипотезы readout]
Недостающая физика может находиться в карте наблюдения: допустимая компрессия
способна породить эффективные структуры, отсутствующие в коммутирующем
представлении источника. Однако полностью симметрийная карта слишком груба,
а общая CP-карта неидентифицируема. Поэтому readout должен быть выведен из
состояния, модульного потока и динамики до сравнения с данными.
\end{proposition}

\begin{protocol}[Definition of Done карты чтения]
Карта \(\mathcal R\) получает физический статус только если:
\begin{enumerate}
 \item первичная алгебра \(\mathcal A\) и состояние \(\omega\) фиксированы;
 \item подалгебра \(\mathcal B_{\mathrm{obs}}\) выведена без загрузки target;
 \item доказана её инвариантность относительно модульного потока \(\omega\);
 \item построено единственное \(\omega\)-сохраняющее условное ожидание
 \(\mathcal E_\omega:\mathcal A\to\mathcal B_{\mathrm{obs}}\);
 \item неподвижный физический сектор уникален modulo объявленной симметрии;
 \item та же карта одновременно проходит scalar, family и gauge gates;
 \item после заморозки конструкции получено новое blind-предсказание.
\end{enumerate}
\end{protocol}

Текущий gate имеет смешанный вердикт. Гипотеза карты чтения подтверждена как
реальная математическая возможность и новое направление программы. Но
канонической карты, выделяющей наш мир, пока нет. Машинные результаты
сохранены в
\texttt{s2t\_observer\_readout\_fixed\_point\_results.json}.

\chapter{Wilson-state и относительный модульный коцикл}

Предыдущий gate показал, что произвольная карта чтения может создать
эффективное mixing, но не объясняет собственный выбор. Проверим более сильное
предположение: задаёт ли уже найденная Wilson-стационарная точка выделенное
состояние \(\omega\), из которого readout следует модульно.

\section{Минимальный класс Gibbs-state}

Wilson saddle фиксирует
\begin{equation}
 c_*=\frac{26-9\sqrt{15}}{11}
\end{equation}
и одну projective axis-orbit, но не задаёт Hamiltonian состояния. Поэтому
проверен весь минимальный класс
\begin{equation}
 \rho_{k,s,\hat n}
 =\frac{\exp(-H_{k,s,\hat n})}
 {\Tr\exp(-H_{k,s,\hat n})},
 \qquad
 H_{k,s,\hat n}
 =B\bigl(L_k+sA_{\hat n}\bigr)B^\dagger.
\end{equation}
Здесь \(k\) пробегает inverse-length, inverse-square и tunneling kernel,
\(s=\pm1\), \(B:\mathbb C^3\to\mathbb C^4\) является канонической
affine-триплетной изометрией, \(\beta=1\), а singlet energy положена равной
нулю. Последние два условия являются частью test-class, а не выводом.

Две минимальные оси дают одинаковые spectra, mutual information и
PPT-negativity с точностью \(10^{-12}\), как и должно быть для одной
остаточной симметрийной орбиты. После факторизации по ней остаются шесть
различных классов состояния: три kernel и два знака.

\section{Нефакторизуемость не равна неизбежной запутанности}

Все двенадцать состояний нефакторизуемы относительно двух двоичных факторов:
\begin{equation}
 \|\rho-\rho_{S^1}\otimes\rho_{\mathbb{RP}^3}\|_F
 \in[0.231,0.395].
\end{equation}
Их mutual information лежит в интервале
\begin{equation}
 I(\mathbb{RP}^3:S^1)\in[0.131,0.274].
\end{equation}
Но PPT-критерий разделяет знаки. Все шесть состояний с \(s=+1\) сепарабельны,
тогда как все шесть состояний с \(s=-1\) запутаны и имеют negativity от
\(0.0740\) до \(0.0949\).

\begin{proposition}[Wilson-корреляция без inevitability запутанности]
Wilson-угол и factor geometry действительно создают межфакторную корреляцию.
Однако квантовая запутанность зависит от знака вставки, который gap-functional
не фиксирует. Следовательно, тезис о первичной запутанности не закрывается до
вывода знака из parent action.
\end{proposition}

\section{Модульно-инвариантные блоки}

Все состояния faithful и имеют невырожденный спектр. Канонический
triplet-projector удовлетворяет
\begin{equation}
 [\rho,P_3]=0
\end{equation}
с ошибкой меньше \(2\cdot10^{-16}\). Поэтому block algebra
\begin{equation}
 \mathcal B_{\mathrm{block}}
 =P_3M_4P_3\oplus\mathbb CP_1
\end{equation}
модульно инвариантна, а block-dephasing сохраняет состояние.

Но это не редукция к чистому \(M_3\): singlet остаётся наблюдаемым центральным
исходом. Conditioning
\begin{equation}
 \rho\longmapsto\frac{P_3\rho P_3}{\Tr(P_3\rho)}
\end{equation}
нелинейно и отбрасывает от \(18.1\%\) до \(31.7\%\) полной вероятности.
Кроме того, одно невырожденное \(\rho\) допускает как минимум пятнадцать
канонических спектральных block-partition, включая четыре разложения
\(3+1\). Состояние не выбирает уникальную наблюдаемую подалгебру.

\section{Спектральный readout не создаёт CKM}

Для каждого состояния проверены четыре rank-three complement его
собственных векторов. Всего исследовано \(48\) модульно выделенных
компрессий. Во всех случаях
\begin{equation}
 [\Phi(T_{\mathbb{RP}^3}),\Phi(T_{S^1})]=0
\end{equation}
с максимальной Frobenius-нормой \(5.7\cdot10^{-14}\). Следовательно,
нефакторизуемое состояние само по себе не превращает факторные трансляции в
источник mixing.

Если две остаточные оси назначить up/down-секторам, overlap собственных
базисов содержит одну вещественную двухсемейную ротацию и один spectator.
Для всех шести классов
\begin{equation}
 J=0.
\end{equation}

\section{Возврат ориентированной Wilson-фазы}

Предыдущий Hamiltonian использует симметричную часть Wilson-оператора и
теряет ориентацию. Поэтому проверены два комплексных эрмитовых продолжения:
\begin{equation}
 H_{\mathrm{gen}}
 =L_k+s\theta_*\,i[\hat n]_\times
\end{equation}
и
\begin{equation}
 H_{\mathrm{full}}
 =L_k+A_{\hat n}
 +s\sin\theta_*\,i[\hat n]_\times.
\end{equation}
Для трёх kernel и четырёх относительных ориентаций в каждом классе получено
\(24\) sector-pair. Во всех случаях overlap имеет девять ненулевых модулей:
ориентированная фаза устраняет spectator и создаёт full three-family mixing.

Однако Jarlskog invariant во всех \(24\) случаях остаётся нулевым с ошибкой
меньше \(3\cdot10^{-17}\). Значит, простая комплексность Wilson-generator не
нарушает общую антиунитарную симметрию двух секторов.

\begin{nogocriterion}[Провал одного ориентированного Wilson-state]
Если up/down Hamiltonian различаются только осью одной Wilson-орбиты и знаком
общего ориентированного генератора, то они могут давать full mixing, но не
дают CP violation. Такая ветвь не является CKM-механизмом независимо от
красоты её модулей.
\end{nogocriterion}

\begin{hypothesis}[Следующий обязательный объект]
Недостающим механизмом является не отдельное запутанное состояние, а
относительный модульный коцикл двух секторов, который нельзя устранить общей
вещественной структурой. Он должен одновременно фиксировать sign, \(\beta\),
singlet energy и относительную комплексную ориентацию up/down-readout.
\end{hypothesis}

\begin{protocol}[Definition of Done relative modular cocycle]
Ветвь получает физический статус только если:
\begin{enumerate}
 \item оба состояния \(\omega_u,\omega_d\) выведены из одного parent action;
 \item их relative cocycle фиксирован без CKM-данных;
 \item общая antiunitary symmetry действительно нарушена;
 \item результирующий overlap имеет девять допустимых модулей и \(J\neq0\);
 \item ни sign, ни \(\beta\), ни singlet conditioning не остаются свободными;
 \item та же пара состояний проходит независимый mass-hierarchy gate.
\end{enumerate}
\end{protocol}

Машинные результаты сохранены в
\texttt{s2t\_wilson\_modular\_state\_readout\_results.json}.

\chapter{Прямой модульный коцикл и BCH-backreaction}

Предыдущий gate выделил относительный модульный коцикл как следующий
кандидат на источник CP violation. Теперь необходимо разделить два разных
механизма: непосредственное действие коцикла на overlap и динамическую
обратную реакцию его некоммутативного генератора на sector Hamiltonian.

\section{Точный no-go непосредственного действия}

Пусть
\begin{equation}
 H_u=U_u\diag(h_u)U_u^\dagger,
 \qquad
 H_d=U_d\diag(h_d)U_d^\dagger,
 \qquad
 V=U_u^\dagger U_d.
\end{equation}
Тогда прямое конечномерное действие относительного модульного коцикла имеет
вид
\begin{equation}
 V(t)=\diag(e^{-it h_u})V\diag(e^{it h_d}).
\end{equation}
Следовательно, оно только перефазирует строки и столбцы mixing matrix.
Любой rephasing invariant, в частности Jarlskog invariant, сохраняется
точно:
\begin{equation}
 J(V(t))=J(V).
\end{equation}

Для двух ориентированных Wilson-классов, трёх frozen factor-kernel, четырёх
относительных ориентаций и трёх значений modular time выполнены \(72\)
численных контроля. Максимальное изменение \(J\) равно
\(2.1\cdot10^{-17}\).

\begin{nogocriterion}[No-go прямого относительного коцикла]
Если relative cocycle действует на mixing matrix только как
\(\rho_u^{it}V\rho_d^{-it}\), он не может превратить CP-сохраняющий overlap
в CKM-матрицу с \(J\ne0\). Изменение modular time не устраняет этот запрет.
\end{nogocriterion}

\section{Первый оператор BCH-backreaction}

Первый эрмитов оператор, содержащий относительную некоммутативность двух
секторов и не сводящийся к перефазировке, равен
\begin{equation}
 Q=\frac{i}{2}[H_u,H_d].
\end{equation}
Коэффициент \(1/2\) фиксирован первым BCH-членом до сравнения с CKM.

Симметричная двухсторонняя реакция
\begin{equation}
 H_u\longmapsto H_u+Q,
 \qquad
 H_d\longmapsto H_d-Q
\end{equation}
проверена для всех \(24\) frozen sector-pair. Во всех случаях
\begin{equation}
 J=0
\end{equation}
с численной точностью. Общая антиунитарная связь между секторами снова
сохраняется.

Односторонние реакции
\begin{equation}
 (H_u,H_d)\longmapsto(H_u+Q,H_d)
\end{equation}
и
\begin{equation}
 (H_u,H_d)\longmapsto(H_u,H_d-Q)
\end{equation}
дают ненулевой Jarlskog invariant во всех \(48\) тестах. Up-only и down-only
ветви дают противоположные знаки \(J\). Это первый coefficient-free оператор
текущей программы, который действительно создаёт трёхсемейное CP violation.

\section{Post-blind проверка масштаба}

CKM-величина не использовалась при построении \(Q\). После заморозки полного
перебора введён контроль
\begin{equation}
 J_{\mathrm{CKM}}=3.12\cdot10^{-5}.
\end{equation}
Канонический BCH-результат лежит в диапазоне
\begin{equation}
 |J|\in[1.21\cdot10^{-3},3.43\cdot10^{-2}].
\end{equation}
Даже минимальный кандидат превышает контроль в \(38.7\) раза. Поэтому
ненулевой CP является структурным положительным результатом, но численное
совпадение отсутствует.

Нельзя исправить расхождение произвольным коэффициентом перед \(Q\): это
создало бы новый непрерывный selector. Кроме того, mass-hierarchy gate пока
не определён, поскольку отсутствует выведенная карта от собственных значений
\(H_u,H_d\) к физическим Yukawa-массам.

\begin{proposition}[Уточнённый следующий механизм]
Недостающим объектом является не прямой Connes-cocycle, а выведенная
хирально- или зарядово-асимметричная обратная реакция относительного
модульного коммутатора. Именно асимметрия действия \(Q\), а не его наличие,
нарушает общую antiunitary symmetry.
\end{proposition}

\begin{protocol}[Definition of Done chiral modular backreaction]
Ветвь получает физический статус только если:
\begin{enumerate}
 \item конечная алгебра, её представление, \(J\) и градуировка фиксированы;
 \item неравные up/down-связи с \(Q\) следуют из одного parent action;
 \item коэффициенты этих связей заморожены до использования CKM;
 \item получены девять допустимых модулей и ненулевой Jarlskog invariant;
 \item масштаб \(J\) проходит preregistered tolerance без подгонки;
 \item та же конструкция задаёт Yukawa-readout и проходит независимый
 mass-hierarchy gate.
\end{enumerate}
\end{protocol}

\begin{nogocriterion}[Запрет CP через свободное подавление]
Если односторонность выбирается вручную или множитель при \(Q\) настраивается
по \(J_{\mathrm{CKM}}\), результат классифицируется как феноменологическая
текстура, а не как вывод первичного образа.
\end{nogocriterion}

Текущий gate имеет смешанный, но направляющий вердикт. Прямой модульный
коцикл и симметричный BCH-механизм закрыты отрицательно. Односторонний
BCH-backreaction впервые создаёт CP без вставки CKM-данных, но требует
выведенной секторной асимметрии и более сильного подавления. Машинные
результаты сохранены в
\texttt{s2t\_relative\_modular\_cocycle\_bch\_results.json}.

\chapter{Алгебраические нормировки BCH и полный CKM gate}

Предыдущий gate показал, что канонический односторонний BCH-коммутатор
создаёт CP violation, но без подавления даёт слишком большой \(J\).
Следующий вопрос состоит в том, может ли нужный порядок величины возникнуть
из уже имеющихся конечных размерностей, а не из коэффициента, подобранного по
CKM.

\section{Замороженное меню нормировок}

До загрузки CKM-control зафиксировано меню
\begin{equation}
 1,\quad\frac12,\quad\frac13,\quad\frac14,\quad
 \frac16,\quad\frac18,\quad\frac1{12},\quad
 \frac1{24},\quad\frac1{48}.
\end{equation}
Его элементы соответствуют raw insertion, хиральному проектору,
нормированным следам семейного и четырёхточечного пространств, chiral
doubling, двенадцати incidence-channel, \(S_4\)-orbit average и signed-chiral
удвоению этой орбиты.

Для двух ориентированных Wilson-Hamiltonian, трёх frozen kernel, четырёх
относительных ориентаций, девяти нормировок и двух односторонних реакций
построены
\begin{equation}
 2\cdot3\cdot4\cdot9\cdot2=432
\end{equation}
кандидата. Их mixing matrices вычислены до сравнения с CKM.

\section{Скалярный успех Jarlskog scale}

После заморозки меню выполнен post-blind контроль. У \(52\) кандидатов
\(|J|\) находится в пределах factor two от наблюдаемой величины. Лучший
кандидат даёт
\begin{equation}
 \frac{|J|}{J_{\mathrm{CKM}}}=0.9888283.
\end{equation}
Таким образом, конечные алгебраические размерности естественно достигают
требуемого порядка подавления. Ранее обнаруженный overshoot не требует
обязательно нового произвольного малого числа.

Но совпадение одного CP-скаляра ещё не является CKM-выводом. Близкие
значения возникают для нескольких разных нормировок, kernel и ориентаций;
сама величина \(J\) не выбирает единственный механизм.

\section{Полный CKM red-team}

Для каждого из \(432\) кандидатов перебраны все \(3!\times3!=36\)
независимых перенумераций up/down-поколений. Jarlskog magnitude при этом не
меняется, но проверяется возможность неверного назначения собственных
векторов поколениям.

Ни один из \(52\) J-compatible кандидатов не воспроизводит одновременно
\(\sin\theta_{12},\sin\theta_{23},\sin\theta_{13}\) даже в грубом интервале
factor two. Лучший logarithmic RMS трёх углов равен
\begin{equation}
 1.7283,
\end{equation}
а лучший Frobenius score для матрицы модулей равен
\begin{equation}
 \bigl\lVert |V|-|V_{\mathrm{CKM}}|\bigr\rVert_F=0.8175.
\end{equation}
Следовательно, алгебраическая нормировка может исправить масштаб \(J\), но
не создаёт иерархию mixing angles.

\begin{nogocriterion}[Запрет CKM по одному Jarlskog]
Совпадение \(|J|\) не засчитывается как CKM-предсказание, если та же
конструкция не воспроизводит три угла и полную матрицу модулей при заранее
фиксированном назначении поколений.
\end{nogocriterion}

\section{Центральная неидентифицируемость масс}

Для любых \(c_u,c_d\in\mathbb R\) преобразования
\begin{equation}
 H_u\longmapsto H_u+c_uI,
 \qquad
 H_d\longmapsto H_d+c_dI
\end{equation}
оставляют неизменными собственные векторы, mixing matrix и Jarlskog
invariant. Одновременно все отношения собственных значений непрерывно
меняются с \(c_u,c_d\).

\begin{proposition}[Mass-shift no-go]
Wilson/BCH-Hamiltonian сами по себе не определяют массовую иерархию.
Интерпретация их собственных значений как масс или squared masses требует
parent action, фиксирующего центральные offsets, знак и положительный
Yukawa-readout. Без этого один и тот же CKM-сектор совместим с континуумом
различных mass spectra.
\end{proposition}

\begin{protocol}[Definition of Done полного семейного readout]
Следующая ветвь считается закрытой только если:
\begin{enumerate}
 \item finite algebra выводит неравные family-edge metrics;
 \item центральные части up/down-операторов фиксированы до mass-data;
 \item один положительный Yukawa-readout задаёт массы и mixing;
 \item три CKM-угла и \(J\) проходят один совместный blind-test;
 \item та же конструкция даёт хотя бы два независимых mass ratios;
 \item ни generation permutation, ни normalization не выбираются по данным.
\end{enumerate}
\end{protocol}

Текущий результат разделён. Положительно установлено, что размерности
конечной алгебры дают естественный CP-scale. Отрицательно установлено, что
этого недостаточно для полной CKM-иерархии и масс. Следующим обязательным
объектом является не ещё один общий suppressing coefficient, а
нецентральный Yukawa-readout с неравными рёберными метриками. Машинные
результаты сохранены в
\texttt{s2t\_chiral\_bch\_normalization\_ckm\_results.json}.

\chapter{Экспоненциальный Yukawa-readout: mass-trained и CKM-blind gate}

Mass-shift no-go показывает, что линейные eigenvalues Hamiltonian не
определяют иерархию без центральной нормировки. Минимальный способ исключить
additive offset без нового непрерывного параметра задаётся нормированным
экспоненциальным functional calculus:
\begin{equation}
 m_i(H;\beta)
 =\frac{\exp[-\beta(E_i-E_{\min})]}
 {\max_j\exp[-\beta(E_j-E_{\min})]}.
\end{equation}
Он положителен, зависит только от spectral gaps и точно инвариантен
относительно \(H\mapsto H+cI\).

\section{Mass-trained протокол}

До quark mass comparison фиксируется дискретное меню
\begin{equation}
 \beta\in\{1,\pi,2\pi\},
\end{equation}
соответствующее безразмерной единице, длине кратчайшего
\(\mathbb{RP}^3\)-цикла и длине \(S^1\). Проверены все девять независимых
up/down-назначений \((\beta_u,\beta_d)\) для \(432\) предыдущих
Wilson/BCH-кандидатов, всего
\begin{equation}
 432\cdot9=3888
\end{equation}
readout. Выбор кандидата проводится только по четырём лёгким quark mass
ratios. CKM-данные до выбора не используются.

\section{Результат mass selector}

Глобальный mass-only optimum сам выбирает
\begin{equation}
 \beta_u=\beta_d=\pi.
\end{equation}
Для нормированных patterns multiplicative errors равны
\begin{equation}
 (10.242,\ 0.841,\ 0.511,\ 1.076),
\end{equation}
где первые две строки относятся к up-сектору, вторые две --- к down-сектору.
Combined logarithmic RMS равен \(1.2145\). Ни один из \(3888\) кандидатов не
помещает все четыре отношения даже в factor-ten interval.

Положительный частичный результат всё же имеется: middle-generation ratios
и оба down-sector ratios получаются правильного порядка, а центральная
неопределённость устранена точно. Прямой численный shift-control с
\(c_u=7.3\), \(c_d=-4.1\) сохраняет mass patterns и mixing с ошибкой меньше
\(10^{-12}\).

\section{Blind CKM-провал}

После заморозки mass-selected кандидата открыт CKM-control. Получено
\begin{equation}
 \frac{|J|}{J_{\mathrm{CKM}}}=357.9
\end{equation}
и отношения предсказанных углов к наблюдаемым
\begin{equation}
 \left(
 \frac{\sin\theta_{12}}{\sin\theta_{12}^{CKM}},
 \frac{\sin\theta_{23}}{\sin\theta_{23}^{CKM}},
 \frac{\sin\theta_{13}}{\sin\theta_{13}^{CKM}}
 \right)
 =(4.34,\ 13.96,\ 199.06).
\end{equation}
Следовательно, mass-trained selector не переносится на mixing и не
идентифицирует наш мир.

\begin{theorem}[Functional-calculus mixing no-go]
Для любой скалярной функции \(f\), определённой на спектре эрмитова \(H\),
операторы \(H\) и \(f(H)\) имеют одни и те же spectral projectors. Поэтому
functional calculus может менять eigenvalues, создавать hierarchy и
устранять central offset, но не может изменить mixing basis.
\end{theorem}

\begin{nogocriterion}[Провал eigenvalue-only Yukawa-readout]
Если Yukawa-map имеет вид \(Y=f(H)\) отдельно в каждом секторе, она не может
исправить CKM-провал исходных собственных базисов независимо от выбора
экспоненты, степени или другого scalar spectral kernel.
\end{nogocriterion}

\begin{protocol}[Definition of Done некоммутирующей Yukawa-вставки]
Следующая ветвь получает физический статус только если:
\begin{enumerate}
 \item новый оператор не является scalar function исходного \(H\);
 \item его неравные edge weights следуют из finite algebra или action;
 \item центральные offsets и positivity фиксированы до mass-data;
 \item mass-trained выбор проходит blind CKM;
 \item один и тот же оператор закрывает минимум четыре mass ratios и четыре
 CKM-параметра;
 \item ablation новой некоммутирующей части уничтожает совместный успех.
\end{enumerate}
\end{protocol}

Экспоненциальная ветвь закрыта как полезный, но недостаточный readout. Она
решает central-shift problem и частично создаёт hierarchy, однако общий
functional-calculus no-go доказывает, что следующий объект обязан менять
eigenvectors. Машинные результаты сохранены в
\texttt{s2t\_exponential\_yukawa\_readout\_results.json}.

\chapter{Gap-derived weighted \texorpdfstring{\(A_3\)}{A3} Yukawa gate}

Functional-calculus no-go требует оператора, меняющего eigenvectors.
Factor Laplacian уже содержит каноническую неравномерность: после аффинной
нормировки его spectrum в \([0,1]\) соседние gaps задают два edge weight
\(p_{10},p_{21}\). Композиция
\begin{equation}
 E_{21}E_{10}=E_{20}
\end{equation}
фиксирует chord weight
\begin{equation}
 q=p_{10}p_{21}.
\end{equation}
Непрерывные коэффициенты не вводятся.

Проверены три factor kernel, две ориентации цепи, два назначения
chain/chord к up/down-секторам и девять пар
\(\beta_{u,d}\in\{1,\pi,2\pi\}\), всего \(108\) mass-trained строк.

Единственный кандидат, помещающий все четыре лёгких quark ratios в factor
five, выбирает inverse-square spectrum:
\begin{equation}
 (p_{10},p_{21},q)
 =\left(\frac34,\frac14,\frac3{16}\right),
 \qquad
 \beta_u=\beta_d=2\pi.
\end{equation}
Его multiplicative errors равны
\begin{equation}
 (3.675,\ 4.426,\ 1.090,\ 1.315),
\end{equation}
а combined logarithmic RMS равен \(0.9986\). По сравнению с equal-gap
ablation это первый реальный hierarchy gain от некоммутирующей edge metric.

После mass-only выбора открыт CKM-control:
\begin{equation}
 \frac{|J|}{J_{\mathrm{CKM}}}=540.36,
 \qquad
 \frac{\sin\theta_{ij}}{\sin\theta_{ij}^{CKM}}
 =(3.76,\ 6.59,\ 129.92).
\end{equation}
Следовательно, hierarchy success не переносится на mixing.

\begin{proposition}[Частичный успех weighted edges]
Gap-derived неравные рёбра являются первым source-based механизмом текущей
цепи, который совместно улучшает четыре quark mass ratios без непрерывной
подгонки. Однако назначение одного сектора цепи, а другого хорде создаёт
слишком сильное относительное вращение и не является CKM-механизмом.
\end{proposition}

\begin{protocol}[Следующий sector-relative graph gate]
Требуется вывести из одного finite action две коррелированные
некоммутирующие Yukawa-матрицы, каждая из которых содержит chain и chord, но
с sector-relative weights. Mass-data могут служить train, однако вся CKM
матрица должна оставаться blind.
\end{protocol}

Машинные результаты сохранены в
\texttt{s2t\_weighted\_gap\_yukawa\_readout\_results.json}.

\chapter{APS, орбифолд и anomaly inflow: red-team смены языка}

Предложенная смена рамки содержательна: она переносит внимание с локального
bookkeeping на глобальные спектральные асимметрии, flat bundles и
bulk--boundary consistency. Однако смена языка не может автоматически
перевести прежние условные коэффициенты в статус вывода. Ниже отделены
точные тождества от новых операторных гипотез.

\section{Арифметика \texorpdfstring{\(1/3\)}{1/3} и APS-gate}

Числовое равенство
\begin{equation}
 \frac14+\left|\zeta_R(-1)\right|
 =\frac14+\frac1{12}=\frac13
\end{equation}
точно. Но его два слагаемых принадлежат разным конструкциям. APS-терм имеет
вид \(-\frac12(\eta+h)\) для конкретного граничного оператора, ориентации,
метрики, spin-структуры и twisting bundle. В уже проверенной конвенции для
двух spin-структур \(\mathbb{RP}^3\) получены
\begin{equation}
 \eta_D=\pm\frac14,
\end{equation}
а не универсальное значение \(-1/2\). Поэтому число \(+1/4\) не следует из
имеющегося eta-аудита без дополнительного удвоения или иного twisted
оператора.

Кроме того, \(\zeta_R(-1)\) регуляризует линейную сумму уровней, тогда как
\(\det' D\) определяется производной spectral zeta-функции. Переход к
абсолютному значению и аддитивное включение в \(\delta m/m\) требуют одного
явного relative operator; это не является следствием одного удаления
нулевой моды.

Топологически
\begin{equation}
 \pi_1(\mathbb{RP}^3\times S^1)=\mathbb Z_2\times\mathbb Z.
\end{equation}
Torsion holonomy на первом генераторе и spin sign окружности на втором
генераторе являются независимой парой. Их произведение описывает специально
выбранный диагональный цикл, а не автоматически определённую ``полную
голономию компактного направления''. Наконец, назначение нетривиальной
\(\mathbb Z_2\)-линии мюону, но не тау, само является family-selector.

\begin{nogocriterion}[APS-разложение коэффициента]
Равенство \(1/3=1/4+1/12\) получает физический статус только после задания
одного twisted Dirac/relative-mass operator, для которого вычислены eta,
kernel, determinant/Casimir prescription и знак обоих вкладов, а
mu/tau-twist следует из представления конечной алгебры до использования их
масс.
\end{nogocriterion}

\section{Почему \texorpdfstring{\(24-1\)}{24-1} ещё не выводит 23}

Арифметика \(24-1=23\) не эквивалентна gauge reduction. Конформный фактор
метрики, нулевая мода диффеоморфизма вдоль \(S^1\) и направление в
конечномерном Majorana-модуле принадлежат разным комплексам. Факторизация
функционального интеграла по gauge orbit требует явных constraints, ghost
операторов и kernel counting; она не является вычитанием одного заранее
выбранного вектора из \(\mathbb R^{24}\).

Коэффициент \(1/24\) в разложении \(\widehat A\)-рода фиксирует
характеристический класс, но сам по себе не устанавливает физический rank
\(24-1\). Тем самым прежний \(Q_8\)-no-go действительно не применим к
настоящей gauge-факторизации, но вместо него возникает более сильная
обязанность построить полный quadratic/BV или constrained Majorana complex.

\section{SU(5)-орбифолд: что уже выведено}

Матрица из предложенной таблицы
\begin{equation}
 P_{bad}=\diag(-1,-1,-1,+1,+1)
\end{equation}
имеет determinant \(-1\) и не принадлежит \(SU(5)\). Корректный order-two
элемент имеет вид
\begin{equation}
 P_5=\diag(+1,+1,+1,-1,-1)\in SU(5).
\end{equation}
Он уже проверен в предыдущем gate и оставляет
\(S(U(3)\times U(2))\). Однако один \(P_5\) фиксирует adjoint breaking, но
не определяет все intrinsic parities matter multiplets.

Более сильная существующая конструкция использует
\begin{equation}
 h=\diag(-1,-1,-1,-i,-i),\qquad h^2=P_5,
\end{equation}
то есть совместную \(Z_2/Z_4\)-проекцию и conjugate flat characters для
vectorlike parent
\begin{equation}
 (\mathbf{10}+\overline{\mathbf{10}})
 +2(\mathbf5+\overline{\mathbf5})+\mathbf5_H.
\end{equation}
Она действительно оставляет ровно \(U+2D+H\) и даёт
\begin{equation}
 \Delta b_{Y,2,3}=\left(\frac{17}{6},\frac16,2\right).
\end{equation}
Следовательно, representation direction уже не является голой подгонкой.
Но flat characters пока не выведены из parent action, finite shifted-partner
determinant проваливает direction gate, а ordinary running имеет неверный
знак для требуемого legacy-сдвига. Абсолютный threshold не закрыт.

\section{Inverse susceptibility и предел anomaly inflow}

Шесть-канальное тождество
\begin{equation}
 \operatorname{Var}(\bar x)
 =\frac1{6\zeta(4,1/2)}=\frac1{\pi^4}
\end{equation}
остаётся точным положительным результатом. Геометрическое чтение
\(3\times2\) через self-dual two-forms и двусторонний half-integer spectrum
также согласовано с уже выполненным аудитом. Но оно требует нового
chiral, antiperiodic, fourth-order auxiliary sector и ordinary trace.

Anomaly inflow фиксирует аномальную вариацию или фазу boundary partition
function. Он не фиксирует автоматически вещественный finite response
\(-1/(\pi^4S_{geo}^2)\), unit source coefficient и относительную нормировку
с periodic \(1/(24S_{geo})\)-ветвью. Более того, существующие surviving
\(U,D\)-фермионы vectorlike, поэтому их gauge anomaly сокращается и сама по
себе не задаёт ненулевой inflow level.

\section{Размерностная поправка parent action}

Предложенная форма
\begin{equation}
 \int_{M^5}\Tr\left(A\wedge dA+\frac23A^3\right)
\end{equation}
является интегралом трёхформы и не определяет пяти-мерное действие. Gauge
Chern--Simons five-form имеет схематический вид
\begin{equation}
 \omega_5(A)=\Tr\left(AF^2-\frac12A^3F+\frac1{10}A^5\right),
 \qquad d\omega_5=\Tr F^3,
\end{equation}
с convention-dependent overall normalization.

Есть и APS-dimensional gate. Если filling имеет boundary
\(\mathbb{RP}^3\), можно использовать четырёхмерную APS-задачу с
трёхмерным eta-оператором. Если boundary объявлена равной
\(\mathbb{RP}^3\times S^1\), необходимо заново задать соответствующий
пятимерный оператор и его четырёхмерную boundary spectral asymmetry;
\(\eta_D(\mathbb{RP}^3)\) нельзя просто перенести без product theorem.

\begin{protocol}[Definition of Done APS/orbifold/inflow ветви]
Новая рамка считается физически переоткрывшей II.A только если:
\begin{enumerate}
 \item заданы bulk manifold, boundary, bundles и корректная five-form;
 \item вычислен anomaly polynomial полного проецированного спектра;
 \item inflow level квантован и фиксирован cancellation, а не наблюдаемым;
 \item twisted APS operator выводит mu/tau branches и общий mass term;
 \item полный constrained complex даёт ровно физический rank 23;
 \item та же конструкция фиксирует sign и magnitude EW/QCD thresholds;
 \item tensor sector выводит half-shift, fourth-order stiffness, ordinary
 trace и unit response map;
 \item хотя бы два normalization-sensitive сектора проходят blind-test.
\end{enumerate}
\end{protocol}

Текущий вердикт: это конструктивная смена пространства поиска, но не
устранение блокировок. Сохраняются три сильных кандидата --- APS eta как
sector discriminator, точная anomaly-free \(Z_2/Z_4\) split-проекция и
self-dual inverse susceptibility. Не выведены общий mass operator,
family-selector, rank 23, inflow level, threshold magnitude, seed
\(\rho_0\), CKM/PMNS и strong dynamics.

Машинный red-team сохранён в
\texttt{s2t\_aps\_orbifold\_inflow\_redteam\_results.json}.

\chapter{Фаза детерминанта и вещественная масса: тест критерия Ф.2}

После выбора anomaly inflow как вспомогательной invertible/topological
теории остаётся один решающий вопрос: может ли eta-фаза фермионного
детерминанта породить вещественный аддитивный член \(-\alpha/3\) в
относительном операторе заряженных лептонов?

\section{Полный product Dirac operator}

Для product geometry квадрат оператора имеет вид
\begin{equation}
 D_K^2=D_{\mathbb{RP}^3}^2+D_{S^1}^2,
\end{equation}
поэтому каждому трёхмерному собственному значению \(\lambda_j\) и окружной
моде \(n+\beta\) соответствует положительная энергия
\begin{equation}
 E_{j,n}^2=\lambda_j^2+\frac{(n+\beta)^2}{R_1^2}.
\end{equation}
Физический charged lepton является vectorlike Dirac field. Его полный
Clifford-блок содержит пару \(\pm iE_{j,n}\), и при вещественной массе
\begin{equation}
 (m+iE_{j,n})(m-iE_{j,n})=m^2+E_{j,n}^2>0.
\end{equation}
Следовательно,
\begin{equation}
 \log(m+iE)+\log(m-iE)=\log(m^2+E^2),
\end{equation}
а mass response равен
\begin{equation}
 \frac{\partial}{\partial m}\log\det(D_K+m)
 =\sum_{j,n}\frac{2m}{m^2+E_{j,n}^2}.
\end{equation}
Обе формулы вещественны blockwise. Спектральная асимметрия нечётномерного
фактора не создаёт eta-члена в производной вещественного vectorlike
effective action.

\section{Численный asymmetric-spectrum control}

Машинный тест намеренно использует несимметричный набор положительных и
отрицательных трёхмерных eigenvalues, после чего строит полный
четырёхмерный Clifford pair для periodic и antiperiodic окружностей. Для
каждого блока, полной суммы и antiperiodic-minus-periodic difference мнимая
часть равна нулю с машинной точностью. Производная по массе также полностью
вещественна. Тем самым cancellation не опирается на случайную симметрию
выбранного cutoff spectrum.

\section{Почему \texorpdfstring{\(\zeta(-1)\)}{zeta(-1)} не является этим
детерминантом}

Для одного massive product-level конечное circle determinant ratio задаётся
\begin{equation}
 \mathcal I_\rho(\beta)=
 \log\frac{\cosh(2\pi\rho)-\cos(2\pi\beta)}
 {\cosh(2\pi\rho)-1}.
\end{equation}
При \(\rho=1\) и \(\beta=1/2\) получаем
\begin{equation}
 \mathcal I_1(1/2)=0.0074697796\ldots,
\end{equation}
что точно воспроизводит прежний projected-KK audit, но не равно
\begin{equation}
 |\zeta_R(-1)|=\frac1{12}=0.0833333\ldots.
\end{equation}
Первое число является mass-dependent Euclidean log-determinant; второе ---
Casimir-регуляризацией линейной суммы уровней. Их нельзя складывать как
части одного determinant без нового functional bridge.

\section{Статус трёхмерной eta-фазы}

В текущем spin-аудите для двух структур \(\mathbb{RP}^3\) зафиксировано
\begin{equation}
 \eta_D=\pm\frac14.
\end{equation}
Значение \(-1/2\) отсутствует в замороженном operator menu и потребовало бы
явного twisted или doubled Dirac operator. Но даже после такого изменения
eta-фаза относится к chiral determinant line. В vectorlike charged-lepton
determinant две хиральности образуют положительный модуль и фазовый вклад
сокращается.

\begin{theorem}[Минимальный phase-to-mass no-go]
Для вещественной положительной Dirac-массы и vectorlike charged-lepton
operator на \(\mathbb{RP}^3\times S^1\) eta-инвариант не даёт аддитивного
вклада в \(\partial_m\Re\log\det(D_K+m)\). Поэтому разложение
\begin{equation}
 \frac13=\frac14+\frac1{12}
\end{equation}
не выводит вещественный mass shift \(-\alpha/3\) в минимальной Gaussian
реализации.
\end{theorem}

\begin{proof}
Каждый product-level входит парой \(m\pm iE\). Их произведение положительно,
мнимая часть логарифма сокращается, а mass derivative зависит только от
\(m^2+E^2\). Отдельно, окружной determinant является функцией
\(\mathcal I_\rho(\beta)\), а не константой \(|\zeta_R(-1)|\). Следовательно,
ни eta-фаза, ни Casimir-остаток не поставляют требуемый вещественный
коэффициент в объявленном operator class.
\end{proof}

\begin{protocol}[Условие переоткрытия Ф.2]
Phase-to-mass ветвь может быть переоткрыта только если из parent action
следуют:
\begin{enumerate}
 \item complex chiral Yukawa mass с физически неустранимой фазой;
 \item CP-odd background или сумма topological sectors, превращающая
 фазовую интерференцию в вещественный potential;
 \item заранее фиксированная связь этого potential с lepton mass modulus;
 \item линейный коэффициент \(\alpha/3\) без использования \(m_\tau\);
 \item ablation eta/topological sector уничтожает именно вещественный mass
 shift, а не только фазу partition function.
\end{enumerate}
\end{protocol}

Текущий вердикт критерия Ф.2 отрицателен. 4D topological interpretation
остаётся допустимой для global anomaly и determinant-phase bookkeeping, но
не закрывает \(\tau/\mu\)-формулу. Единственным существующим вещественным
кандидатом остаётся прежняя QED winding self-energy, для которой unit
normalization даёт коэффициент \(0.0616969\ldots\), а требуемый projection
weight \(5.4027533\ldots\) не выведен.

Машинный протокол сохранён в
\texttt{s2t\_eta\_phase\_mass\_gate\_results.json}.

\chapter{Конформная мода и Majorana-rank: тест критерия Ф.3}

После отрицательного phase-to-mass gate остаётся проверить вторую часть
предложенной 4D-интерпретации: может ли удаление одной conformal/gauge mode
превратить полный Majorana-модуль размерности \(24\) в физический rank \(23\)?

\section{Diffeomorphism gauge complex}

В четырёх измерениях metric fluctuation имеет десять локальных компонент,
а линейное действие diffeomorphism задаётся
\begin{equation}
 \delta_\xi h_{\mu\nu}=\nabla_\mu\xi_\nu+\nabla_\nu\xi_\mu.
\end{equation}
На Fourier-level с momentum \(k\) его symbol равен
\begin{equation}
 G_k(\xi)_{\mu\nu}=k_\mu\xi_\nu+k_\nu\xi_\mu.
\end{equation}
Для \(k\neq0\) этот map имеет rank \(4\), но pure trace tensor
\(h_{\mu\nu}=c g_{\mu\nu}\) не принадлежит его image. Для constant mode
\(k=0\) symbol вообще равен нулю. Глобально тот же факт следует из
\begin{equation}
 \delta_\xi\Vol(K)
 =\int_K\nabla_\mu\xi^\mu\,d\Vol=0
\end{equation}
на компактном manifold без границы. Следовательно, constant volume mode не
является diffeomorphism gauge direction.

\section{Почему gauge volume не действует на \texorpdfstring{\(\mathbb R^{24}\)}{R24}}

Минимальное поле конфигураций имеет direct-sum структуру
\begin{equation}
 \operatorname{Sym}^2(T^*K)\oplus\mathcal M_{Maj},
 \qquad \dim_{\mathbb R}\mathcal M_{Maj}=24.
\end{equation}
В замороженной постановке BRST/diffeomorphism map действует в первом блоке:
\begin{equation}
 \xi\longmapsto(G\xi,0).
\end{equation}
Поэтому quotient по \(\operatorname{Diff}(K)\), ghost determinant и
gauge-volume factor не меняют rank второго summand:
\begin{equation}
 \dim\mathcal M_{Maj}^{physical}=24,
\end{equation}
если отсутствует отдельный constraint или projector внутри самого
Majorana-модуля.

\section{Weyl и wrong-sign не спасают подсчёт}

Pure trace становится gauge direction только при дополнительной локальной
Weyl symmetry
\begin{equation}
 \delta_\sigma g_{\mu\nu}=2\sigma g_{\mu\nu}.
\end{equation}
Такая symmetry отсутствует в текущем massive parent action: фиксированные
радиусы, Majorana/Yukawa scales и dimensionful defect couplings нарушают её.
Даже если формально добавить Weyl quotient, он rescale-ит поля по их Weyl
weights, но не выбирает одну вещественную polarization внутри
\(\mathbb R^{24}\). На zero-field background линейный gauge map в
Majorana-блок равен нулю.

Также «неправильный знак» Euclidean conformal factor означает выбор или
rotation integration contour. Это не gauge quotient и не уменьшение
конечномерного rank.

Наконец, коэффициент \(1/24\) в
\begin{equation}
 \widehat A=1-\frac{p_1}{24}+\cdots
\end{equation}
является нормировкой characteristic class, а не числом Majorana degrees of
freedom. Из него не следует операция \(24-1\).

\begin{theorem}[Conformal-rank no-go]
В 4D parent architecture без точной локальной Weyl gauge symmetry
обработка conformal factor и деление на gravitational gauge volume не
превращают \(\mathcal M_{Maj}\simeq\mathbb R^{24}\) в rank-\(23\) модуль.
Для такого сокращения необходим rank-one operator, действующий внутри
Majorana fiber.
\end{theorem}

\section{Сравнение с defect-route}

Предыдущий dimension gate уже доказал, что удаление generation singlet
вычитает полный внутренний spinor module, а не одну вещественную компоненту.
Но class-D defect на систолической геодезической даёт условный обход этого
no-go: transverse mod-two kernel, longitudinal periodic real mode,
generation singlet и self-dual cycle line имеют combined rank \(1\). Внутри
этой defect-модели дополнение в \(\mathbb R^{24}\) действительно имеет rank
\(23\).

Различие принципиально. Conformal factor находится в metric field space и
не действует на Majorana fiber. Defect zero mode является собственным
вектором Majorana/BdG operator и потому способен определить настоящий
rank-one projector. Его оставшийся gap --- вывод odd-winding mass order
parameter и restriction map из общего S2T action.

\begin{protocol}[Условие переоткрытия Ф.3]
Конформная ветвь может быть переоткрыта только если:
\begin{enumerate}
 \item полное massive действие обладает точной Weyl gauge symmetry;
 \item построен её BRST-complex вместе с diffeomorphism ghosts;
 \item mixed gauge map имеет rank-one image внутри \(\mathcal M_{Maj}\);
 \item quotient удаляет ровно одну реальную Majorana direction, а не общий
 scale или полный internal multiplet;
 \item тот же complex независимо фиксирует relative trace--Hodge weight.
\end{enumerate}
Альтернативный допустимый путь --- вывести уже построенный class-D defect из
parent action без нового выбранного order parameter.
\end{protocol}

Текущий вердикт Ф.3 отрицателен для conformal/gauge-volume объяснения.
Число \(23\) не выводится из \(\widehat A\)-знаменателя и не возникает при
удалении gravitational conformal mode. Оно сохраняет условный статус только
через явный rank-one Majorana defect либо иной operator в \(\mathbb R^{24}\).

Машинный протокол сохранён в
\texttt{s2t\_conformal\_majorana\_rank\_gate\_results.json}.

\chapter{Динамическое происхождение Majorana-дефекта}

После conformal-rank no-go единственным явным маршрутом к rank \(23\)
остаётся class-D defect на систолической геодезической. Его topology,
mod-two kernel, core gluing и tubular embedding уже согласованы. Теперь
проверяется более сильное условие: следует ли сам defect как saddle из
текущего parent action?

\section{Минимальное tubular action}

На tubular neighborhood \(N(\gamma)\) минимальный локальный кандидат имеет
вид
\begin{align}
 S_B={}&\int_{N(\gamma)}
 \left[
 Z_\Phi |(d-2ia)\Phi|^2
 +m_\Phi^2|\Phi|^2+\lambda_\Phi|\Phi|^4
 +\frac1{4e_{root}^2}|f_a|^2
 \right],\\
 S_F={}&\frac12\int_{N(\gamma)}
 \Psi^T C\left(
 D_a+\Re\Phi\,\Gamma_1+\Im\Phi\,\Gamma_2
 \right)\Psi.
\end{align}
Здесь \(a\) является square-root connection, а \(\Phi\) --- charge-two
Majorana pairing field.

\section{Что топология действительно вынуждает}

Для square-root branch meridian holonomy равна
\begin{equation}
 \oint_{\mu_\gamma}a=\pi\pmod{2\pi}.
\end{equation}
Если \(|\Phi|\neq0\) вне core, finite-energy condition даёт
\begin{equation}
 \oint_{\mu_\gamma}(d\arg\Phi-2a)=0,
\end{equation}
и потому
\begin{equation}
 \Delta\arg\Phi=2\pi,
 \qquad
 \operatorname{wind}(\Phi)=1.
\end{equation}
Следовательно, odd winding не является fitted integer. При наличии root
flux и pairing condensate он вынужден topology.

\section{Что топология не вынуждает}

Для минимального potential
\begin{equation}
 V(\Phi)=m_\Phi^2|\Phi|^2+\lambda_\Phi|\Phi|^4
\end{equation}
при \(\lambda_\Phi>0\) существуют два разных режима:
\begin{equation}
 m_\Phi^2\geq0
 \quad\Longrightarrow\quad
 \Phi=0,
\end{equation}
и
\begin{equation}
 m_\Phi^2<0
 \quad\Longrightarrow\quad
 |\Phi|^2=-\frac{m_\Phi^2}{2\lambda_\Phi}.
\end{equation}
В первом режиме vortex отсутствует: winding фазы не определён, Majorana
gap не открыт и rank-one index не реализуется физически. Поэтому topology
не выбирает знак \(m_\Phi^2\), величину condensate или существование
attractive gap equation.

Есть и независимый connection gap. Smooth ambient \(\mathbb Z_2\)-линия
имеет trivial meridian holonomy \(+1\). Её square root на complement имеет
holonomy \(-1\) и не продолжается через core. Значит, локальное действие
должно содержать либо dynamical \(Z_4/U(1)\) root connection, либо явный
disorder/defect sector. Совпадение с уже найденной quarter-holonomy ветвью
показывает совместимость, но не выводит её назначение Majorana-паре.

\begin{theorem}[Conditional-defect theorem]
Square-root \(\pi\)-flux и ненулевой charge-two pairing condensate строго
вынуждают odd vortex и сохраняют единственный реальный Majorana kernel.
Но ambient \(\mathbb Z_2\)-topology сама по себе не вынуждает ни singular
root sector, ни pairing condensation. Поэтому defect является условным
saddle расширенного tubular action, а не следствием замороженного II.A
field content.
\end{theorem}

\section{Что при этом сохраняется}

Отрицательный dynamical gate не разрушает локальный результат. Если root
sector и condensed saddle получены из будущего parent action, то уже
проверены:
\begin{itemize}
 \item unit winding и нечётный mod-two index;
 \item один transverse real Majorana zero mode;
 \item periodic longitudinal coefficient line;
 \item отсутствие физического Nambu doubling;
 \item complement rank \(23\) в \(\mathbb R^{24}\);
 \item точный ambient-to-tube superconnection restriction.
\end{itemize}
Таким образом, отсутствует не математическая конструкция defect, а механизм
его обязательного возникновения.

\begin{protocol}[Definition of Done defect parent action]
Defect-route получает статус вывода только если:
\begin{enumerate}
 \item root \(Z_4/U(1)\)-connection следует из существующей finite algebra
 или обязательной суммы topological sectors;
 \item pairing Hessian имеет выведенное отрицательное направление либо
 получена parameter-free attractive gap equation;
 \item nonzero vortex saddle энергетически предпочтительнее \(\Phi=0\);
 \item quarter-holonomy sector map к Majorana pair фиксирован до neutrino
 data;
 \item full fluctuation/BV complex сохраняет один mod-two kernel;
 \item та же action normalization воспроизводит второй независимый sector.
\end{enumerate}
\end{protocol}

Текущий вердикт смешанный. Topological implication
\[
 \pi\text{-flux}+\Phi\neq0\Longrightarrow\operatorname{wind}\Phi=1
\]
закрыта. Динамическое утверждение
\[
 \text{S2T parent action}\Longrightarrow \pi\text{-flux и }\Phi\neq0
\]
не закрыто. Добавление этих структур вручную означало бы новую версию
модели.

Машинный протокол сохранён в
\texttt{s2t\_majorana\_defect\_parent\_action\_gate\_results.json}.

\chapter{Источник square-root connection: исчерпание существующего меню}

Предыдущий gate показал, что class-D defect требует singular root
connection и pairing condensate. Теперь проверяется, содержится ли нужный
order-four action уже в SU(5), spin, flat-character или Nambu-структурах
проекта.

\section{SU(5)-quarter holonomy}

Используемый в split-sector элемент
\begin{equation}
 h=\exp(i3\pi Y)
\end{equation}
действительно имеет order \(4\) на некоторых Standard Model components.
Но sterile right-handed mode имеет
\begin{equation}
 N^c\sim(\mathbf1,\mathbf1)_0,
 \qquad Y(N^c)=0.
\end{equation}
Она является дополнительным \(SU(5)\)-singlet, отсутствующим в минимальном
\(\mathbf{10}+\overline{\mathbf5}\). Поэтому
\begin{equation}
 h|_{N^c}=1,
 \qquad
 P_5|_{N^c}=h^2|_{N^c}=1.
\end{equation}
SU(5)-quarter branch, разделяющая \(Q,L,T_H\), не создаёт sterile
square-root connection.

\section{Абстрактная gauge-holonomy doublet}

Ранний gauge audit использует U(1)-like charges \(q=\pm1\). При
\(\beta=1/4\) он даёт conjugate phases
\begin{equation}
 (-i,+i),
 \qquad
 (-i)^2=(+i)^2=-1.
\end{equation}
Алгебраически это именно требуемый root pair. Но audit не задаёт
representation map этого U(1)-like charge на neutral sterile sector.
Назначить \(q=\pm1\) постфактум означало бы ввести новый charge assignment.

То же относится к multiplet flat characters. В SU(5) split-table characters
назначены \(\mathbf{10}\), \(\mathbf5\) и \(\mathbf5_H\); sterile singlet в
таблице отсутствует. Character \(\pm i\) можно добавить, но он не выводится
из текущей finite algebra.

\section{Spin и Nambu не являются root-source}

Spin-структуры дают coefficient holonomies \(\pm1\), а не независимую линию
с квадратом \(-1\). Nambu transition
\begin{equation}
 \diag(i,-i)
\end{equation}
переносит basis нулевой моды после того, как pair phase уже изменилась на
\(\pi\). Поэтому он является следствием vortex texture и не может служить
причиной её возникновения. Использовать его как source означало бы
логический круг.

\section{Единственный настоящий root}

Требуемый объект уже классифицирован топологически: это square root
ambient torsion line на
\(\mathbb{RP}^3\setminus\gamma\). Его character на generator равен
\(\pm i\), а meridian holonomy равна \(-1\). Именно поэтому line не
продолжается через core.

Следовательно, root существует математически, но не является smooth field
замороженного action. Он должен войти как один из трёх новых объектов:
\begin{enumerate}
 \item sterile-sector \(Z_4\) flat character;
 \item dynamical \(U(1)\) или \(Z_4\) connection с charge-two pairing field;
 \item disorder operator, включающий nonextendable sector в path integral.
\end{enumerate}

\begin{theorem}[Existing-menu root no-go]
Ни SU(5)-holonomy, ни spin structure, ни Nambu transition не задают
обязательный независимый square-root connection на sterile Majorana mode.
Абстрактная quarter-holonomy doublet содержит правильные phases, но её
sterile charge assignment не выведен. Поэтому defect-route требует явного
расширения finite algebra или topological-sector measure.
\end{theorem}

\begin{protocol}[Definition of Done root-source extension]
Новая root-ветвь допустима только если:
\begin{enumerate}
 \item sterile \(Z_4\)-charge следует из finite algebra и real structure;
 \item root sector включён в path-integral measure до neutrino data;
 \item проверены discrete и global anomaly constraints;
 \item та же representation выводит charge-two pairing field;
 \item pairing Hessian или gap equation создаёт condensate;
 \item extension проходит второй независимый sector test.
\end{enumerate}
\end{protocol}

Текущий итог: нужная фаза в математическом меню существует, но не назначена
sterile-моде никакой обязательной структурой. Поэтому переиспользование
SU(5)-quarter branch не закрывает defect parent action.

Машинный протокол сохранён в
\texttt{s2t\_majorana\_root\_source\_menu\_results.json}.

\chapter{Минимальное расширение \texorpdfstring{\(U(1)_{B-L}\)}{U(1) B-L}
как источник Majorana-root}

Исчерпание существующего меню показывает, что root connection нельзя
получить внутри минимального \(SU(5)\). Следующий честный шаг --- не скрывать
новую структуру, а проверить минимальное зарегистрированное расширение
\[
 G_{ext}=G_{SM}\times U(1)_{B-L}.
\]

\section{Charge assignment и anomaly gate}

Используем left-handed Weyl convention. Для одного поколения charges равны
\[
\begin{array}{c|rrrrrr}
 &Q&u^c&d^c&L&e^c&N^c\\
\hline
 B-L&1/3&-1/3&-1/3&-1&1&1
\end{array}
\]
с обычными multiplicities. Прямой rational audit даёт
\begin{align}
 \mathcal A_{(B-L)^3}&=0,&
 \mathcal A_{\mathrm{grav}^2(B-L)}&=0,\\
 \mathcal A_{SU(3)^2(B-L)}&=0,&
 \mathcal A_{SU(2)^2(B-L)}&=0,\\
 \mathcal A_{Y^2(B-L)}&=0,&
 \mathcal A_{Y(B-L)^2}&=0.
\end{align}
Таким образом, одна sterile mode на поколение, уже требуемая seesaw-сектором,
закрывает continuous anomaly gate без новых chiral fermions.

\section{Root holonomy и pairing field}

Пусть holonomy на generator complement равна
\begin{equation}
 \oint_y a_{BL}=\frac{\pi}{2}.
\end{equation}
Тогда \(N^c\) с charge \(+1\) получает phase
\begin{equation}
 \chi_{N^c}(y)=i,
 \qquad
 \chi_{N^c}(\mu_\gamma)=i^2=-1.
\end{equation}
Это ровно missing square-root connection.

Введём один complex scalar
\begin{equation}
 \Phi_{BL},\qquad q_{BL}(\Phi_{BL})=-2.
\end{equation}
Тогда Majorana vertex
\begin{equation}
 y_N\Phi_{BL}N^cN^c+\mathrm{h.c.}
\end{equation}
gauge invariant, поскольку \(-2+1+1=0\). На meridian finite-energy
condition вынуждает winding \(\pm1\), знак которого соответствует orientation.
Следовательно, одно расширение одновременно поставляет root charge и
ранее отсутствующий pairing order parameter.

\section{Почему это ещё не замыкание}

Расширение требует новых continuous данных:
\begin{equation}
 g_{BL},\quad v_{BL},\quad\lambda_{BL},\quad
 \lambda_{H\Phi},\quad\epsilon_{YB},\quad y_N.
\end{equation}
Topology фиксирует quarter phase и winding, но не знак pairing Hessian,
масштаб \(v_{BL}\) и Majorana Yukawa matrix.

Поля \(N^c\) и \(\Phi_{BL}\) являются Standard Model singlets и не дают
прямого one-loop вклада в \((b_Y,b_2,b_3)\). Однако обычные SM fermions
заряжены одновременно по \(Y\) и \(B-L\), причём
\begin{equation}
 \sum_{\mathrm{Weyl}}Y(B-L)=\frac83
\end{equation}
на поколение. Поэтому kinetic mixing двух abelian factors генерируется
радиативно. Его boundary value и running становятся новым
normalization-sensitive gauge-sector test.

Кроме того, \(B-L\), действующий нетривиально на \(N^c\), не содержится в
существующем minimal \(SU(5)\)-fiber, где \(N^c\) является external singlet.
Нужно либо добавить abelian factor, либо перейти к более крупной finite
algebra/unification structure.

\begin{proposition}[Первый coherent root extension]
\(U(1)_{B-L}\) с одним \(N^c\) на поколение и scalar charge \(-2\) является
первым проверенным расширением, которое одновременно:
\begin{enumerate}
 \item anomaly-free во всех continuous mixed channels;
 \item даёт sterile order-four holonomy;
 \item содержит gauge-invariant charge-two Majorana pairing field;
 \item воспроизводит conditional odd-vortex/rank-one defect route.
\end{enumerate}
Но оно является новой моделью с невыведенными couplings и scale, а не
замыканием II.A.
\end{proposition}

\begin{protocol}[Definition of Done \(B-L\)-ветви]
Расширение получает S2T-статус только если:
\begin{enumerate}
 \item \(U(1)_{B-L}\) или его \(Z_4\)-subsector следует из finite algebra;
 \item pairing Hessian и \(v_{BL}\) выводятся из frozen spectral scales;
 \item kinetic mixing фиксирован до gauge controls;
 \item discrete/global anomalies остаточного \(Z_4\) проверены;
 \item heavy neutrino scale и rank-one defect возникают из одного saddle;
 \item тот же extension проходит независимый EW/QCD или tensor gate.
\end{enumerate}
\end{protocol}

Текущий вердикт положителен на representation/anomaly уровне и условен на
dynamics уровне. Дополнительный global consistency gate, приведённый далее,
уточняет этот вывод: ordinary charge-two scalar несовместим с root holonomy
на generator complement, и требуется ambient-torsion-twisted pairing
section. Пространство поиска сузилось до конкретного расширения, но цена
расширения теперь полностью видима.

Машинный протокол сохранён в
\texttt{s2t\_bl\_root\_extension\_gate\_results.json}.

\chapter{Consistency gate для
\texorpdfstring{\(B-L\)}{B-L}-root и pairing condensate}

Предыдущий gate показал, что \(U(1)_{B-L}\) даёт правильную sterile phase
\(+i\) и сокращает continuous anomalies. Однако этого недостаточно:
необходимо проверить, может ли charge-two pairing field иметь ненулевой
condensate на всём complement носителя.

\section{Обструкция обычного charge-two scalar}

Для complement
\[
 H_1(\mathbb{RP}^3\setminus\gamma;\mathbb Z)=\mathbb Z\langle y\rangle,
 \qquad \mu_\gamma=2y.
\]
Root connection имеет
\[
 \operatorname{Hol}_{N^c}(y)=i,
 \qquad
 \oint_y a_{BL}=\frac{\pi}{2}.
\]
Поэтому обычный scalar \(\Phi_{BL}\) с charge \(-2\) видит
\[
 \operatorname{Hol}_{\Phi_{BL}}(y)
 =\exp\left(-2i\oint_y a_{BL}\right)=-1.
\]
Ненулевой covariantly constant section на замкнутом цикле требует
holonomy \(+1\). Следовательно, обычный nowhere-zero \(B-L\) condensate
несовместим с order-four root на generator \(y\). На meridian обструкция
исчезает, поскольку
\[
 \operatorname{Hol}_{\Phi_{BL}}(\mu_\gamma)=(-1)^2=+1,
\]
но этого недостаточно для глобального condensate на complement.

Та же проблема видна из symmetry breaking. Charge-two VEV оставляет
\(Z_2\), а элемент с parameter \(\pi/2\) действует на VEV как \(-1\) и не
принадлежит unbroken subgroup. Charge-four VEV оставил бы \(Z_4\), но scalar
charge \(-4\) не допускает линейный Majorana vertex
\(\Phi N^cN^c\).

\section{Torsion-twisted rescue}

На носителе уже имеется ambient torsion line \(L_{\mathrm{tor}}\) с
\[
 \operatorname{Hol}_{L_{\mathrm{tor}}}(y)=-1,
 \qquad
 \operatorname{Hol}_{L_{\mathrm{tor}}}(\mu_\gamma)=+1.
\]
Поэтому pairing field следует рассматривать не как обычный scalar, а как
section
\begin{equation}
 \Phi_{\mathrm{pair}}\in
 \Gamma\!\left(
 L_{\mathrm{tor}}\otimes L_{B-L}^{-2}
 \right).
\end{equation}
Его полный holonomy равен
\[
 (-1)\times(-1)=+1
\]
на \(y\) и также \(+1\) на meridian. Это устраняет обструкцию для
parallel condensate. Однако следующий полный vertex gate показывает, что
bilinear \(N^cN^c\) всё ещё имеет holonomy \(i^2=-1\): torsion twist
переносит знак в Yukawa map и потому не является завершённым rescue.

\begin{proposition}[Уточнённый статус \(B-L\)-ветви]
Обычный \(U(1)_{B-L}\) Higgs charge \(-2\) не может одновременно
поддерживать nonzero condensate на всём complement и сохранять sterile
root holonomy \(i\). Минимальный consistency rescue существует, если
pairing field скручен ambient \(\mathbb Z_2\)-torsion line, но этот rescue
совместим только с condensate и не закрывает scalar Yukawa vertex. Полный
статус устанавливается следующим trilemma gate.
\end{proposition}

\begin{protocol}[Definition of Done twisted \(B-L\)-ветви]
\begin{enumerate}
 \item вывести \(L_{\mathrm{tor}}\otimes L_{B-L}^{-2}\) как representation
 finite algebra, а не назначить после обнаружения обструкции;
 \item построить global pairing potential для twisted section;
 \item проверить core boundary condition и fluctuation determinant;
 \item показать, что twisted condensate сохраняет rank-one kernel;
 \item связать breaking scale и kinetic mixing с независимым сектором.
\end{enumerate}
\end{protocol}

Машинный протокол сохранён в
\texttt{s2t\_bl\_root\_condensate\_consistency\_results.json}.

\chapter{Root--mass--condensate trilemma}

Предыдущий consistency gate показал, что ambient torsion twist устраняет
holonomy-обструкцию для самого pairing field. Теперь необходимо проверить
всю Majorana vertex. Эта проверка даёт более сильный результат.

\section{Три несовместимых требования}

На generator \(y\) sterile root line имеет
\[
 r=\operatorname{Hol}_{N^c}(y)=i,
 \qquad r^2=-1.
\]
Для scalar Yukawa vertex
\[
 \Phi N^cN^c
\]
bundle invariance требует
\begin{equation}
 \phi r^2=1,
 \qquad\Longrightarrow\qquad
 \phi=-1,
\end{equation}
где \(\phi\) --- holonomy pairing field. Но uniform covariantly constant
nonzero condensate требует \(\phi=+1\). Поэтому order-four root, ordinary
linear Majorana vertex и uniform parallel condensate не могут выполняться
одновременно.

\section{Почему torsion twist переносит обструкцию}

Для twisted field
\[
 \Phi_{\mathrm{tw}}\in
 \Gamma(L_{\mathrm{tor}}\otimes L_{B-L}^{-2})
\]
полный holonomy равен
\[
 \phi_{\mathrm{tw}}=(-1)\times(-1)=+1.
\]
Это действительно допускает parallel condensate. Однако vertex теперь
имеет holonomy
\[
 \phi_{\mathrm{tw}}r^2=(+1)(-1)=-1
\]
и не является scalar functional на всём complement.

Можно формально добавить torsion-odd Yukawa spurion \(Y_{\mathrm{tor}}\)
с holonomy \(-1\). Тогда
\[
 Y_{\mathrm{tor}}\Phi_{\mathrm{tw}}N^cN^c
\]
инвариантен. Но spurion сам не имеет nowhere-zero parallel trivialization.
Следовательно, знак не исчезает, а переносится в новый defect, zero locus
или вручную выбранную trivialization.

\begin{nogocriterion}[Root--mass--uniform-vacuum no-go]
При \(r=i\) не существует одного scalar pairing field и scalar constant
Yukawa coefficient, которые одновременно дают invariant linear Majorana
vertex и uniform parallel nonzero condensate на всём core complement.
\end{nogocriterion}

\section{Выжившая физическая ветвь}

No-go не уничтожает class-D defect route. Ordinary charge-two field с
holonomy \(-1\) может существовать как nonuniform section с forced gradient
или zero locus. Тогда Majorana mass является defect texture, а не
homogeneous bulk parameter. Это возвращает программу к исходному
тубулярному BdG-механизму, но теперь с более строгим условием: texture
должна быть получена как lowest-action section, а не объявлена.

\begin{protocol}[Следующий nonuniform-section gate]
\begin{enumerate}
 \item построить минимальную energy functional на solid-torus complement;
 \item найти lowest-action section с holonomy \(-1\);
 \item определить обязательный zero locus и winding;
 \item вычислить BdG kernel и fluctuation determinant;
 \item проверить, остаётся ли rank-one результат без torsion Yukawa spurion.
\end{enumerate}
\end{protocol}

Таким образом, torsion twist не является полным rescue homogeneous
\(B-L\)-модели. Он точно локализует место обструкции и показывает, что
выжившая ветвь обязана быть defect-like.

Машинный протокол сохранён в
\texttt{s2t\_root\_mass\_condensate\_trilemma\_results.json}.

\chapter{Неоднородная конфигурация наименьшего действия}

Трилемма «корневая голономия --- масса --- конденсат» оставляет
единственную локальную возможность: поле майорановского спаривания должно
быть неоднородным. Проверим, возникает ли такая стационарная конфигурация
автоматически и что именно в ней фиксирует геометрия.

\section{Минимальный функционал Гинзбурга--Ландау}

На generator \(y\) длины \(L\) рассмотрим
\begin{equation}
 F[\Phi]=\int_0^L
 \left[
 |D_y\Phi|^2+
 \frac{\lambda}{2}\left(|\Phi|^2-v^2\right)^2
 \right]dy,
\end{equation}
где \(\oint_y a=\pi/2\), а заряд поля спаривания равен \(-2\).
Для периодического скалярного поля с числом намотки
\(n\in\mathbb Z\) ковариантный импульс равен
\begin{equation}
 k_n=\frac{2\pi n+\pi}{L}.
\end{equation}
Следовательно, геометрия даёт полуцелую решётку импульсов. Две нижние ветви
\[
 n=0,\qquad n=-1
\]
имеют \(k=\pm\pi/L\) и точно вырождены.

\section{Точный порог конденсации}

После минимизации по постоянной амплитуде \(\rho=|\Phi|\) получаем
\begin{equation}
 \rho^2=v^2-\frac{k_n^2}{\lambda}.
\end{equation}
Ненулевая стационарная конфигурация существует только при
\begin{equation}
 \boxed{\lambda v^2>k_n^2}.
\end{equation}
Для нижней ветви
\[
 k_{\min}^2=\frac{\pi^2}{L^2}.
\]
Для единичной \(\mathbb{RP}^3\) систолическая длина равна \(L=\pi\),
поэтому порог
принимает особенно простой вид
\begin{equation}
 \lambda v^2>1.
\end{equation}

Плотность энергии конденсированной ветви и выигрыш относительно
\(\Phi=0\) равны
\begin{align}
 f_{\mathrm{cond}}
 &=k_{\min}^2v^2-\frac{k_{\min}^4}{2\lambda},\\
 f_0-f_{\mathrm{cond}}
 &=\frac{(\lambda v^2-k_{\min}^2)^2}{2\lambda}.
\end{align}
Радиальный гессиан
\begin{equation}
 H_\rho=4(\lambda v^2-k_{\min}^2)
\end{equation}
положителен строго выше порога.

\section{Что фиксировано и что остаётся свободным}

Топология фиксирует полуцелый сдвиг, сопряжённую пару \(n=0,-1\) и
единичную поперечную намотку вокруг сердцевины при наличии конденсата.
Но она не фиксирует \(\lambda v^2\), поэтому не гарантирует сам факт
конденсации. Она также не выбирает знак продольного импульса.

\begin{proposition}[Условная неоднородная конфигурация]
Корневая голономия порождает устойчивую неоднородную конфигурацию
спаривания, если \(\lambda v^2>\pi^2/L^2\). Она возникает в виде двух
вырожденных сопряжённых ветвей. Следовательно, топология определяет
спектральный порог и структуру ветвей, но не масштаб вакуума и не
ориентацию.
\end{proposition}

Поперечный индекс по модулю два условно остаётся равным единице: если
стационарная конфигурация существует, меридиональный поток \(\pi\)
по-прежнему заставляет намотку поля спаривания быть нечётной. Однако
одномерное ядро оператора Боголюбова--де Жена должно быть пересчитано на
действительном фоне наименьшего действия.

\begin{protocol}[Следующая проверка коэффициентов исходного действия]
\begin{enumerate}
 \item вывести \(\lambda v^2\) из спектрального действия до обращения к
 нейтринным данным;
 \item проверить знак \(\lambda v^2-\pi^2/L^2\);
 \item вывести член, расщепляющий ориентации \(n=0,-1\);
 \item решить поперечную задачу Боголюбова--де Жена на выбранном фоне;
 \item вычислить отношение определителей к нормальной ветви.
\end{enumerate}
\end{protocol}

Машинный протокол сохранён в
\texttt{s2t\_nonuniform\_pairing\_saddle\_results.json}.

\chapter{Спектральная проверка жёсткости майорановского спаривания}

Предыдущий раздел установил геометрический порог
\[
 \lambda v^2>\frac{\pi^2}{L^2}.
\]
Теперь требуется выяснить, следует ли левая часть из уже принятого
спектрального действия.

\section{Разложение спектральной функции}

Пусть \(a\) --- амплитуда майорановской деформации нулевой моды. В
одномерном контрольном секторе
\begin{equation}
 f(a^2)
 =f(0)+f'(0)a^2+\frac12f''(0)a^4+\cdots .
\end{equation}
Сравнение с потенциалом Гинзбурга--Ландау даёт
\begin{equation}
 c_2=f'(0),\qquad
 c_4=\frac12f''(0),\qquad
 \lambda=f''(0),\qquad
 v^2=-\frac{f'(0)}{f''(0)}.
\end{equation}
Поэтому пороговая комбинация имеет особенно простой вид
\begin{equation}
 \boxed{\lambda v^2=-f'(0)}.
\end{equation}
Это равенство относится к приведённому одномодовому сектору с канонической
нормировкой производного члена. Полное действие обязано независимо вывести
эту нормировку.

\section{Проверка принятого набора функций}

Для уже использованных в программе спектральных функций получаем:
\[
\begin{array}{c|c|c}
 f(x) & \lambda v^2 & \text{состояние при }L=\pi\\
\hline
 x & \text{устойчивый минимум не возникает} & \text{нормальная фаза}\\
 e^{-x} & 1 & \text{критическая точка}\\
 e^{-2x} & 2 & \text{конденсат}\\
 (1+x)^{-3} & 3 & \text{конденсат}\\
 (1+x)^{-4} & 4 & \text{конденсат}
\end{array}
\]
Следовательно, различные допустимые спектральные функции дают различные
фазы. Сам носитель не выбирает между ними.

Каноническая норма пространства конфигураций
\[
 23+\pi^{-1}
\]
задаёт положительную кинематическую метрику коллективной деформации, но не
определяет знак \(f'(0)\). Поэтому она не может сама по себе вызвать
неустойчивость нормальной фазы.

\begin{nogocriterion}[Неопределённость спектральной функции]
Пока функция \(f\), её масштаб и относительная нормировка производного и
потенциального членов не выведены из единого действия, наличие
майорановского конденсата не является предсказанием программы.
\end{nogocriterion}

\section{Запрет выбора ориентации чётным действием}

Две нижние ветви имеют ковариантные импульсы
\[
 k_{0}=+1,\qquad k_{-1}=-1
\]
в единичной геометрии. Для них совпадают все чётные спектральные
инварианты:
\[
 k_0^2=k_{-1}^2,\qquad
 k_0^4=k_{-1}^4,\qquad
 |k_0|=|k_{-1}|.
\]
Значит, всякое вещественное действие, зависящее только от \(D^2\) и чётных
степеней ковариантного импульса, сохраняет точное вырождение ветвей.

Ранее проверенная \(\eta\)-фаза в минимальной вектороподобной модели
сокращается между сопряжёнными блоками и не создаёт вещественного
расщепления энергии. Поэтому для выбора ориентации требуется либо
выведенный киральный граничный член, либо другой вещественный вклад,
нечётный относительно \(k\mapsto-k\).

\begin{proposition}[Двойная незамкнутость майорановской ветви]
Текущее спектральное ядро не определяет ни сам факт конденсации, ни выбор
одной из двух ориентаций. Первое зависит от невыведенной спектральной
функции и её масштаба, второе запрещено для вещественного чётного
спектрального действия.
\end{proposition}

\begin{protocol}[Условия продолжения]
\begin{enumerate}
 \item вывести спектральную функцию и её масштаб до обращения к
 нейтринным данным;
 \item получить коэффициенты производного, квадратичного и четвертичного
 членов из одного следа;
 \item доказать превышение геометрического порога;
 \item вывести вещественный ориентационно-нечётный член;
 \item после этого повторно вычислить ядро оператора Боголюбова--де Жена.
\end{enumerate}
\end{protocol}

Машинный протокол сохранён в
\texttt{s2t\_spectral\_pairing\_stiffness\_gate\_results.json}.

\chapter{Рабочий пакет ветви \texorpdfstring{\(B-L\)}{B-L}-root и неоднородного дефекта}

Настоящий раздел сводит результаты отдельных гейтов в одну статусную
формулу. Подтверждение относится к явно проверенному минимальному меню и не
является доказательством абсолютной единственности среди всех возможных
расширений.

\section{Точно подтверждённые узлы}

В левокиральной вейлевской конвенции добавление одного стерильного поля
\(N^c\) на поколение с зарядом \(B-L=1\) зануляет все шесть локальных
непрерывных аномалий:
\begin{equation}
 \mathcal A_{(B-L)^3}=\mathcal A_{\mathrm{grav}^2(B-L)}
 =\mathcal A_{SU(3)^2(B-L)}=\mathcal A_{SU(2)^2(B-L)}
 =\mathcal A_{Y^2(B-L)}=\mathcal A_{Y(B-L)^2}=0.
\end{equation}
При
\begin{equation}
 \oint_y a_{B-L}=\frac{\pi}{2}
\end{equation}
стерильная линия получает корневую голономию
\begin{equation}
 \operatorname{Hol}_{N^c}(y)=i.
\end{equation}
Поле спаривания \(\Phi_{B-L}\) с зарядом \(-2\) допускает
калибровочно-инвариантную вершину
\begin{equation}
 y_N\Phi_{B-L}N^cN^c,
 \qquad -2+1+1=0.
\end{equation}

На меридиане корневой поток и ненулевой конденсат дают точное условное
следствие
\begin{equation}
 \oint_{\mu_\gamma}a_{B-L}=\pi,\quad |\Phi_{B-L}|\neq0
 \quad\Longrightarrow\quad
 \left|\operatorname{wind}(\Phi_{B-L})\right|=1.
\end{equation}
Слово ``условное'' существенно: топология фиксирует намотку при наличии
конденсата, но не вынуждает сам конденсат и не выбирает знак после смены
ориентации меридиана.

Для минимального функционала Гинзбурга--Ландау ковариантные импульсы равны
\begin{equation}
 k_n=\frac{2\pi n+\pi}{L}.
\end{equation}
В проверенном седловом анзаце с постоянным модулем ненулевая стационарная
конфигурация существует тогда и только тогда, когда
\begin{equation}
 \lambda v^2>\frac{\pi^2}{L^2}.
\end{equation}
При \(L=\pi\) это даёт точный порог \(\lambda v^2>1\). Нижние ветви
\(n=0\) и \(n=-1\) имеют импульсы \(k=+1\) и \(k=-1\), одинаковую энергию
и образуют точно вырожденную сопряжённую пару.

\begin{proposition}[Выжившая конструктивная ветвь минимального меню]
После проведённых гейтов единственной не закрытой конструктивной ветвью
проверенного минимального меню является
\begin{equation}
 U(1)_{B-L}+N^c+\Phi_{-2}
 \quad\longrightarrow\quad
 \text{неоднородная дефектная текстура при }\lambda v^2>1.
\end{equation}
Она подтверждена на уровне представлений, локальных аномалий, корневой
голономии и условного стационарного решения, но ещё не выведена из единого
спектрального действия.
\end{proposition}

\section{Уточнение роли торсионного скручивания}

Торсионное скручивание не является полным спасением однородной модели.
Оно меняет голономию поля спаривания с \(-1\) на \(+1\), но тогда линейная
вершина Юкавы получает остаточный торсионный знак \(-1\). Поэтому полностью
скрученный однородный конденсат не удовлетворяет одновременно условиям
корневой голономии, скалярной вершины и ненулевого параллельного вакуума.
Выжившую конфигурацию корректно называть неоднородной \(B-L\)-дефектной
текстурой. Термин ``торсионно-скрученная'' допустим только после вывода
соответствующего торсионно-нечётного отображения Юкавы, которого в текущем
действии нет.

\section{Две независимые точные конструкции}

Шестиканальная гауссова модель с обычной суммой по компонентам даёт точную
идентичность
\begin{equation}
 \operatorname{Var}(\bar x)
 =\frac{1}{6\zeta(4,1/2)}
 =\frac{1}{\pi^4}.
\end{equation}
Независимо от неё совместная \(\mathbb Z_2/\mathbb Z_4\)-проекция
вектороподобного \(SU(5)\)-родителя оставляет нулевые моды
\(U+2D+H\) и точно даёт направление
\begin{equation}
 \Delta b_{Y,2,3}=\left(\frac{17}{6},\frac16,2\right).
\end{equation}
Обе формулы воспроизводимы, но пока не следуют из одного оператора, одной
меры или одного действия. Поэтому запрещено записывать их композицию как
единый физический вывод. Первая конструкция остаётся кандидатом на
обратную восприимчивость, вторая --- точной проекцией представлений с
невыведенной пороговой величиной и неверным знаком обычного промежуточного
ренормгруппового пробега.

\section{Невыведенные динамические данные}

Топология и голономия не фиксируют:
\begin{itemize}
 \item величину \(\lambda v^2\) и сам факт превышения порога;
 \item выбор одной из ориентаций \(n=0,-1\);
 \item масштаб нарушения \(v_{B-L}\) и матрицу \(y_N\);
 \item граничное значение и пробег кинетического смешивания двух
 абелевых факторов.
\end{itemize}
Последний пункт особенно важен, поскольку
\begin{equation}
 \sum_{\mathrm{Weyl}}Y(B-L)=\frac83
\end{equation}
на поколение, и смешивание в общем случае радиативно порождается, но его
нормировка не определяется топологией.

\section{Закрытые маршруты}

Без новой динамики повторно не открываются следующие гипотезы:
\begin{enumerate}
 \item разложение \(1/3=1/4+1/12\) как вывод вещественного массового
 сдвига \(-\alpha/3\) из APS-фазы;
 \item универсальное значение \(\eta_D=-1/2\);
 \item отождествление \(|\zeta_R(-1)|=1/12\) с окружным
 логарифмическим определителем;
 \item получение \(23=24-1\) удалением калибровочной или конформной моды;
 \item использование трёхформы Черна--Саймонса
 \(\operatorname{Tr}(A\wedge dA+\tfrac23A^3)\) как пятимерного объёмного
 действия;
 \item получение стерильного корня из минимального \(SU(5)\), спиновой
 структуры или намбу-удвоения;
 \item однородный конденсат обычного поля заряда \(-2\);
 \item торсионное скручивание как одновременное спасение конденсата и
 вершины Юкавы;
 \item электрослабое и КХД-замыкание через обычный split-running или
 минимальный логарифмический конус порогов;
 \item формула \(23=k+1\) с \(k=22\) как выведенный уровень
 Черна--Саймонса или Весса--Зумино--Виттена.
\end{enumerate}

\begin{protocol}[Условие следующего продвижения]
Ветвь повышает статус только после вывода из одного заранее фиксированного
действия величины \(\lambda v^2\), ориентационно-нечётного члена,
\(v_{B-L}\), нормировки кинетического смешивания и оператора
Боголюбова--де Жена на выбранной конфигурации. Новые численные совпадения
без такого вывода не считаются продолжением ветви.
\end{protocol}

\chapter{Пакетный аудит гипотез: четыре леммы и новые прунеры}

Десять гипотез были повторно проверены в замороженных конвенциях
Протокола~0.1: train-якоря \(\{\alpha^{-1},m_e,m_\mu\}\), единичные
радиусы, спектральное меню
\(\eta_D(\mathbb{RP}^3)=\pm\tfrac14\) и окружной функционал
\begin{equation}
 \mathcal I_\rho(\beta)=
 \log\frac{\cosh(2\pi\rho)-\cos(2\pi\beta)}
 {\cosh(2\pi\rho)-1}.
\end{equation}
Аудит не добавил закрытого физического предсказания, но дал несколько
точных уточнений и сузил пространство повторных поисков.

\section{Воспроизведённые отрицательные результаты}

Для APS-кандидата необходимо различать две допустимые ветви. При
\(\eta=-1/4\)
\begin{equation}
 -\frac{\eta}{2}+|\zeta_R(-1)|
 =\frac18+\frac1{12}=\frac5{24},
\end{equation}
а при \(\eta=+1/4\) та же комбинация равна \(-1/24\). Ни одна ветвь не
даёт \(1/3\), а значение \(\eta=-1/2\) отсутствует в замороженном меню.
Независимо от этой арифметики вектороподобный блок
\begin{equation}
 (m+iE)(m-iE)=m^2+E^2
\end{equation}
сокращает фазу по каждому уровню, поэтому \(\eta\)-часть не создаёт
аддитивного вещественного массового сдвига.

Конформный и диффеоморфный фактор действует только в метрическом блоке.
Прямая сумма с майорановским пространством оставляет его ранг равным
\(24\), так что маршрут \(24-1=23\) вновь закрывается отрицательно.

Матрица
\begin{equation}
 \operatorname{diag}(-1,-1,-1,+1,+1)
\end{equation}
имеет определитель \(-1\) и не принадлежит \(SU(5)\). Исправленная
\(\mathbb Z_2/\mathbb Z_4\)-проекция с
\begin{equation}
 P_5=\operatorname{diag}(1,1,1,-1,-1),\qquad
 h=\operatorname{diag}(-1,-1,-1,-i,-i)
\end{equation}
точно оставляет \(U+2D+H\) и направление
\((17/6,1/6,2)\), но не выводит величину порога и не исправляет знак
обычного промежуточного ренормгруппового пробега.

Предложенный интегранд
\(\operatorname{Tr}(A\wedge dA+\tfrac23A^3)\) является трёхформой, а не
пятиформой. Корректный шаблон имеет вид
\begin{equation}
 \omega_5(A)=\operatorname{Tr}
 \left(AF^2-\frac12A^3F+\frac1{10}A^5\right),
 \qquad d\omega_5=\operatorname{Tr}F^3,
\end{equation}
с точностью до конвенции и точных форм. Это исправляет размерность, но не
создаёт отсутствующий операторный мост от
\(\eta_D(\mathbb{RP}^3)\) к границе
\(\mathbb{RP}^3\times S^1\).

\section{Сохранённые точные леммы}

Обычная шестиканальная гауссова сумма даёт
\begin{equation}
 \operatorname{Var}(\bar x)
 =\frac1{6\zeta(4,1/2)}=\frac1{\pi^4},
\end{equation}
но нормированный след даёт \(6/\pi^4\), а полный физический носитель и
нормировка источника не выведены.

Минимальный \(U(1)_{B-L}\)-контент с одним \(N^c\) на поколение проходит
все шесть непрерывных аномальных сумм, даёт корневую фазу \(i\) и
калибровочно-инвариантную вершину поля заряда \(-2\). Однородный конденсат
закрыт трилеммой, поэтому сохраняется только неоднородная текстура.
Смешанный след
\begin{equation}
 \sum_{\mathrm{Weyl}}Y(B-L)=\frac83
\end{equation}
задаёт будущий нормировочно-чувствительный второй гейт, но ещё не является
вторым пройденным сектором.

Для приведённого седлового анзаца точно сохраняются
\begin{equation}
 k_n=\frac{2\pi n+\pi}{L},\qquad
 \lambda v^2>\frac{\pi^2}{L^2},
\end{equation}
а при \(L=\pi\) --- порог \(\lambda v^2>1\) и вырожденная пара
\(n=0,-1\). Само превышение порога и выбор ориентации действием не
зафиксированы.

\section{Новые точные прунеры}

Окружной функционал чётен по голономной ветви:
\begin{equation}
 \mathcal I_\rho(\beta)=\mathcal I_\rho(-\beta)
 =\mathcal I_\rho(1-\beta).
\end{equation}
Следовательно, он тождественно не расщепляет сопряжённые ориентации.
Численное сравнение при \(\rho=1\),
\begin{align}
 \mathcal I_1(1/2)&=0.0074697796100663\ldots,\\
 2\mathcal I_1(1/4)&=0.0074837289794912\ldots,
\end{align}
является отдельным неравенством, а не доказательством вырождения; их
отношение равно \(0.9981360403\ldots\).

Попытка заменить член \(-\alpha/3\) двухпетлевой формой
\begin{equation}
 \left(\frac{\alpha}{\pi}\right)^2X=\frac\alpha3
\end{equation}
требует
\begin{equation}
 X=\frac{\pi^2}{3\alpha}=450.8303668617\ldots.
\end{equation}
Строгий дефект этой гипотезы состоит не в субъективной
``неестественности'' числа, а в циркулярности: коэффициент получен из
train-якоря \(\alpha\), который он должен был независимо объяснить.

Кандидат проекционного веса
\begin{equation}
 \frac{12\pi}{7}=5.3855874062\ldots
\end{equation}
не равен требуемому
\(5.4027533072\ldots\); относительный дефицит равен
\(3.17725\times10^{-3}\). Без независимого оператора это приближение
отклоняется Протоколом~0.1.

\section{Незарегистрированное утверждение}

Фраза ``осевой разностный сдвиг тождественно равен нулю'' не сопровождается
в пакете определением соответствующего оператора. Существующий аудит осей
доказывает одну остаточную проективную орбиту, а относительный двухсекторный
аудит оставляет свободный вещественный угол смешивания. Поэтому новый
нулевой операторный прунер не регистрируется до предъявления явной формулы.

\begin{proposition}[Итог пакетного аудита]
Сохраняются четыре точных или условно точных конструктивных результата:
шестиканальная дисперсия, \(B-L\)-аномальная свобода и корневая фаза,
неоднородный порог с сопряжённой парой и исправленная
\(\mathbb Z_2/\mathbb Z_4\)-проекция представлений. Нового закрытого
физического предсказания не возникает:
\begin{equation}
 N_{\mathrm{closed\ physical}}=0.
\end{equation}
Следующий допустимый шаг требует заранее фиксированного родительского
действия, а не новой вставки, восстановленной из целевого числа.
\end{proposition}

Машинный протокол сохранён в
\texttt{s2t\_hypothesis\_batch\_pruner\_results.json}.

\chapter{Первый гейт ярусного родительского действия}

Предложенная архитектура различает интегрально нормированные компактные
координаты и \(L^2\)-нормированные распространяющиеся амплитуды. Это
различение уже имеет точное содержание в нейтринной суперсвязности:
нулевая форма \(\widehat e_0\) является нормированной волновой функцией, а
одна форма \(e_1\) имеет единичный период. Необходимо проверить, переносится
ли то же правило на заряженный лептонный сектор без второго скрытого
считывания меры.

\section{Алгебраический тест трёх чисел}

Доступные нормировки действительно воспроизводят три требуемые величины:
\begin{align}
 \rho_\tau^{\mathrm{raw}}
 &=\|1_{\mathbb{RP}^3}\|^2+
   \|1_{S^1}\|^2+
   \left\langle\|P_\perp n\|^2\right\rangle_{S^2}
 =\pi^2+2\pi+\frac23,\\
 \rho_\tau^{\mathrm{can}}
 &=1+1+\frac23=\frac83,\\
 \|\boldsymbol\Xi_\nu\|^2
 &=23\|\widehat e_0\|^2+\|e_1\|^2
 =23+\pi^{-1}.
\end{align}
Следовательно, числовой P1-гейт проходит как тождество норм. Однако это
ещё не означает, что все три значения являются выводами одного действия.

\section{Гейт типовой классификации}

Квантованный период имеет только нейтринная одна форма:
\begin{equation}
 \oint_\gamma e_1=1.
\end{equation}
Напротив, \(1_{\mathbb{RP}^3}\) и \(1_{S^1}\) являются постоянными
нуль-формами. Их амплитуды не имеют интегрального периода и не являются
вильсоновскими координатами только потому, что носитель компактен. При
канонической нормировке
\begin{equation}
 1_{\mathbb{RP}^3}\mapsto\pi^{-1}1_{\mathbb{RP}^3},
 \qquad
 1_{S^1}\mapsto(2\pi)^{-1/2}1_{S^1},
\end{equation}
и два объёмных вклада превращаются в \(1+1\).

Тем самым детерминированное правило
\begin{quote}
интегральный период \(\Rightarrow\) компактная нормировка;
обычная амплитуда нуль-формы \(\Rightarrow L^2\)-нормировка
\end{quote}
даёт одновременно \(23+\pi^{-1}\) и \(8/3\), но не даёт сырой
лептонный seed. Правило, сохраняющее ненормированные постоянные функции,
даёт seed, но не объясняет, почему те же направления в вершине должны
читаться канонически.

\begin{nogocriterion}[Один ярус --- два лептонных считывания]
Одна фиксированная типовая классификация не может назначить двум
постоянным лептонным нуль-формам одновременно сырые нормы
\((\pi^2,2\pi)\) и канонические нормы \((1,1)\). Для получения обоих
результатов требуется второе отображение считывания или перенормировка
вершины. Пока это отображение не выведено из симметрии, граничного действия
или квантовой меры, оно является тем же скрытым относительным весом, который
запрещён двухсекторным гейтом.
\end{nogocriterion}

Число \(8/3\) в этом тесте не является новым наблюдаемым сектором. Оно
остаётся контрольным результатом канонической нормировки, которая ранее
уничтожила сырой seed. Воспроизведение двух альтернативных координат одного
направления не считается двумя физическими предсказаниями.

\section{Статус архитектур}

Стратифицированная суперсвязность сохраняет нейтринный результат, но в
текущем виде не выводит лептонный seed. Слой БФ/граничного ротора может
оставаться независимым кандидатом на петлевой функционал, однако он не
устраняет древесную неоднозначность без явного отображения между фоновыми
модулями и нормированными лептонными вершинами. Поэтому гибрид A+B не
закрыт отрицательно как будущая модель, но ещё не является одним
родительским действием.

\begin{proposition}[Вердикт P1]
P1 проходит алгебраический гейт трёх чисел и проваливает гейт единственной
меры и типовой классификации:
\begin{equation}
 N_{\mathrm{predictive\ sectors}}=1<2.
\end{equation}
Единственным прошедшим нормировочно-чувствительным сектором остаётся
нейтринная коллективная жёсткость. Нового закрытого физического
предсказания не возникает.
\end{proposition}

\begin{protocol}[Условие переоткрытия A+B]
Необходимо выписать локальное или граничное действие, которое до сравнения
с \(m_\tau\):
\begin{enumerate}
 \item выводит отображение фоновых компактных модулей в канонически
 нормированную лептонную вершину;
 \item совместно преобразует stiffness и coupling при смене координат;
 \item фиксирует, почему древесная норма сохраняет
 \(\pi^2+2\pi+2/3\), а петлевая вершина получает единственный заранее
 определённый вес;
 \item тем же правилом проходит независимый электромагнитный или
 роторный гейт.
\end{enumerate}
До этого P2--P4 являются тестами отдельных слоёв, а не тестами одного
родительского действия.
\end{protocol}

Машинный протокол сохранён в
\texttt{s2t\_tiered\_parent\_action\_p1\_results.json}.

\chapter{Гейт единой карты весов}

После провала простого ярусного правила полезно отделить три различных
утверждения: существование числовой таблицы весов, происхождение этой
таблицы из одного действия и логическую независимость оставшихся
динамических разрывов. Эти утверждения не эквивалентны.

\section{Фиксация опорной метрики}

Карта весов \(W\) определена только относительно заранее фиксированного
скалярного произведения. Поэтому варианты ``\(W=1\) с каноническими
нормами'' и ``\(W=1\) с сырыми нормами'' не являются двумя правилами одного
класса: второй вариант уже заменяет базовую метрику. В каноническом базисе
целевые значения равны
\begin{align}
 \rho_\tau&=\pi^2+2\pi+\frac23,\label{eq:WgateTau20260807}\\
 \rho_\nu&=23+\pi^{-1}.\label{eq:WgateNu20260807}
\end{align}

Для постоянного веса \(W=1\) получаются
\begin{equation}
 \rho_\tau^{C0}=\frac83,
 \qquad
 \rho_\nu^{C0}=24,
\end{equation}
так что класс \(C0\) проваливает оба целевых значения.

\section{Классы \(C1\) и \(C2\)}

Правило, использующее только степень формы и наличие квантованного периода,
назначает обычным нуль-формам единичный вес, а интегральной одной форме ---
вес \(\pi^{-1}\). Оно даёт
\begin{equation}
 \rho_\tau^{C1}=\frac83,
 \qquad
 \rho_\nu^{C1}=23+\pi^{-1}.
\end{equation}
Следовательно, \(C1\) проходит нейтринный сектор и проваливает лептонный.

Формулировка ``все нуль-формы имеют сырую норму'' не принадлежит \(C1\):
сырая норма равна \(\pi^2\) на \(\mathbb{RP}^3\), \(2\pi\) на \(S^1\) и
другому числу на ином факторе. Это уже зависимость от фактора, то есть
класс \(C2\).

В \(C2\) действительно существует таблица
\begin{equation}
 W_{(\mathbb{RP}^3,0)}=\pi^2,
 \quad
 W_{(S^1,0)}=2\pi,
 \quad
 W_{(\mathrm{fin},0)}=1,
 \quad
 W_{(S^1,1)_{\mathrm{int}}}=\pi^{-1},
\end{equation}
которая одновременно воспроизводит \eqref{eq:WgateTau20260807} и
\eqref{eq:WgateNu20260807}.
Она не использует названия ``лептон'' и ``нейтрино'', поэтому формально не
является запрещённым правилом \(C3\). Однако разбиение на четыре ячейки уже
содержит все требуемые относительные нормировки. Без вариационного
функционала, симметрии или граничной симплектической формы это согласованная
таблица метрики, но не предсказание.

\begin{nogocriterion}[Табличное прохождение не равно единой мере]
Прохождение двух агрегатных норм таблицей \(C2\) не считается прохождением
двухсекторного физического гейта, пока:
\begin{enumerate}
 \item веса ячеек не получены как гессиан одного заранее заданного действия;
 \item не задан закон их совместного преобразования с вершинами при смене
 координат поля;
 \item тем же правилом не вычислена независимая наблюдаемая, не
 использованная при построении таблицы.
\end{enumerate}
\end{nogocriterion}

\section{Гейт ковариантности и логический статус}

При замене координаты поля \(x'=a x\) числовой коэффициент квадратичной
формы меняется обратно квадрату масштаба, а коэффициенты вершин должны
преобразоваться согласованно. Поэтому физическим объектом является не
список норм сам по себе, а ковариантная пара ``метрика плюс карта
наблюдаемой''. Предложенная таблица \(C2\) такого закона пока не задаёт.

Провал \(C0\)--\(C1\) и недовыведенность \(C2\) показывают, что единая мера
в текущем формализме не построена. Они не доказывают, что более глубокое
действие не может вывести \(C2\), и не доказывают логическую независимость
разрывов меры, CP-источника, выбора вакуума и порогов.

\begin{proposition}[Вердикт \(W\)-гейта]
Алгебраический статус равен
\begin{equation}
 C0:\ \mathrm{fail},
 \qquad
 C1:\ \mathrm{one\ sector},
 \qquad
 C2:\ \mathrm{lookup\ pass}.
\end{equation}
Операторный и двухсекторный предсказательный гейты для \(C2\) провалены,
поскольку родительское происхождение и ковариантность весов отсутствуют.
Число новых закрытых физических предсказаний остаётся нулевым.
\end{proposition}

\begin{protocol}[Условие переоткрытия \(W\)-гейта]
Следующая модель должна вывести указанную \(C2\)-метрику как гессиан или
граничную форму одного действия до обращения к массам, задать совместное
преобразование stiffness и couplings и затем пройти независимый
электромагнитный либо роторный тест. Только после этого допустимо проверять,
являются ли остальные динамические разрывы проекциями одного принципа.
\end{protocol}

Машинный протокол сохранён в
\texttt{s2t\_weight\_map\_gate\_results.json}.

\chapter{Каноническая мера и локализация лептонного разрыва}

Гипотеза H-M предлагает принять одну каноническую кинетическую меру,
сохранить нейтринную норму \(23+\pi^{-1}\), читать лептонный тангент с
нормой \(8/3\), а сырой seed \(\pi^2+2\pi+2/3\) перенести из кинетического
члена в вершину. Эта схема согласована как разложение ролей, но необходимо
различать совместимость меры, вывод меры и физический двухсекторный тест.

\section{Тест совместимости H-M}

Единичный вес \(W=I\) следует понимать относительно уже фиксированной
градуированной trace--Hodge метрики. Он не означает евклидову норму один
для каждого компонента: интегральная одна форма сохраняет геометрическую
норму \(\pi^{-1}\). При таком понимании
\begin{align}
 \|\boldsymbol\Xi_\nu\|^2
 &=23\|\widehat e_0\|^2+\|e_1\|^2
 =23+\pi^{-1},\\
 \|\boldsymbol\Xi_\tau\|^2
 &=1+1+\frac23=\frac83.
\end{align}
Следовательно, H-M проходит алгебраический тест совместимости.

Однако каноническая метрика в формулировке H-M зафиксирована как исходное
условие. Ни симметрия, ни гессиан родительского действия, ни квантовая мера,
единственным образом выбирающие её, не предъявлены. Поэтому корректный
статус равен «каноническая мера допустима», а не «каноническая мера
выведена».

\section{Физический двухсекторный гейт}

Число \(8/3\) ранее возникло как контроль канонической нормировки, которая
удаляет объёмные множители из лептонного seed. Оно не является независимой
наблюдаемой заряженного лептонного сектора. Поэтому пара
\begin{equation}
 23+\pi^{-1},
 \qquad
 \frac83
\end{equation}
задаёт два согласованных математических чтения одной метрики, но не два
закрытых физических предсказания. Число прошедших независимых
normalization-sensitive секторов остаётся равным одному.

\begin{nogocriterion}[Совместимость не является выводом меры]
Назначение \(W=I\) проходит H-M-гейт только как фиксированная кинематическая
конвенция. Для статуса выведенной меры требуется получить ту же
градуированную метрику из вариационного функционала и вычислить с её помощью
независимую наблюдаемую, не использованную при выборе нормировки.
\end{nogocriterion}

\section{Перенос сырого seed в вершину}

Исключение
\begin{equation}
 \rho_\tau^{\mathrm{raw}}
 =\pi^2+2\pi+\frac23
\end{equation}
из кинетической нормы устраняет внутреннее противоречие H-M, но не выводит
сам seed. Теперь необходимо независимо построить объёмную вершину и
петлевой проекционный вес
\begin{equation}
 \mathcal J_{\mathrm{req}}
 =\frac{1/3}{|I_\tau|/\pi}
 =5.4027533071\ldots.
\end{equation}
Следовательно, разрыв действительно локализован в вершинной карте, но не
закрыт.

\section{Статус single-vertex утверждения}

Проверенный циклический оператор даёт
\begin{equation}
 \Tr Q_{\mathrm{cycle}}=\pi+\pi^{-1},
\end{equation}
а не \(\pi^2+2\pi\), поэтому именно этот минимальный кандидат проваливается.
Из этого не следует no-go для любой одной градуированной вершины.

Для общего утверждения необходимо заранее определить:
\begin{enumerate}
 \item допустимую алгебру вершин;
 \item разрешённые блоки и градуировки;
 \item единственный ковариантный readout;
 \item запрет на коэффициенты, содержащие целевые числа.
\end{enumerate}
Без этих ограничений один блочно-диагональный оператор может формально
содержать произвольные объёмный и winding-коэффициенты; это target-loaded
запись, но она опровергает возможность универсального no-go без объявления
класса.

\begin{proposition}[Исправленный вердикт H-M]
Гипотеза H-M проходит тест кинематической совместимости:
\begin{equation}
 W=I
 \quad\Longrightarrow\quad
 \left(23+\pi^{-1},\,\frac83\right).
\end{equation}
Она не проходит гейт происхождения меры и не добавляет второго физического
сектора. Перенос сырого seed в вершину является допустимой локализацией
разрыва, но объёмная вершина и \(\mathcal J_{\mathrm{req}}\) остаются
невыведенными. Для \(Q_{\mathrm{cycle}}\) получен частный отрицательный
результат; общий single-vertex no-go не доказан.
\end{proposition}

\begin{protocol}[Следующий вершинный гейт]
Следует заморозить конечномерную алгебру допустимых объёмных и winding-вершин
до сравнения с \(m_\tau\), вывести её из одного действия и проверить, даёт ли
один ковариантный след оба коэффициента без sector-specific проекторов.
Только провал полного объявленного класса позволит сформулировать строгий
single-vertex no-go.
\end{protocol}

Машинный протокол сохранён в
\texttt{s2t\_canonical\_measure\_vertex\_localization\_results.json}.

\chapter{Ревизия R1 и триаж ходов H0--H2}

После отрицательного закрытия C6 предложены три маршрута: честная ревизия
без exact \(\pi^{-4}\)-теоремы, повторное открытие канала
\(1\leftrightarrow3\) и объединение reciprocal-функционала с defect-saddle.
Их статусы существенно различаются.

\section{H0: независимое воспроизведение ревизии R1}

Из замороженных train-входов
\begin{equation}
 \alpha^{-1}=137.035999177,
 \qquad
 m_e=0.51099895069\ {\rm MeV},
 \qquad
 m_\mu=105.6583755\ {\rm MeV}
\end{equation}
новая прямая реализация, не импортирующая прежние численные скрипты,
воспроизводит
\begin{align}
 S_{\rm geo}&=137.0363037758784,\\
 S_{\rm vac}
 &=S_{\rm geo}-\frac1{24S_{\rm geo}}
   -\frac1{\pi^4S_{\rm geo}^2}
 =137.0359991735224,\\
 m_\tau&=1776.8594285630\ {\rm MeV},\\
 v_{\rm S2T}&=245.9934092611\ {\rm GeV},\\
 \lambda_H&=0.1292217159851,\\
 M_H&=125.0564860390\ {\rm GeV}.
\end{align}
Остаток \(S_{\rm vac}-\alpha^{-1}\) равен
\(-3.47765\cdot10^{-9}\), а относительный остаток относительно
\(m_\tau=1776.86\ {\rm MeV}\) равен
\(-3.21599\cdot10^{-7}\).

Это подтверждает машинную воспроизводимость формул, но не меняет их
доказательный статус. \(S_{\rm geo}\) остаётся геометрическим успехом,
periodic \(1/24\) --- условной determinant-ветвью,
\(\pi^{-4}\) --- strong structural compression, \(m_\tau\) --- сильной
условной связью, а абсолютный Higgs-мост наследует эти условия.

\begin{proposition}[Внутренний переход зрелости R1]
Согласно заранее записанной шкале тома, ревизия без exact
\(\pi^{-4}\)-теоремы вместе с независимой прямой реализацией выполняет
внутренний порог
\begin{equation}
 R_{\rm sci}:4/10\longrightarrow5/10.
\end{equation}
Это повышение качества фиксации и воспроизводимости, а не появление
закрытого физического предсказания:
\begin{equation}
 N_{\rm closed\ physical}=0.
\end{equation}
Порог \(5\to6\) требует внешнего воспроизведения аудита и blind-таблицы.
\end{proposition}

\section{H1: канал \(1\leftrightarrow3\)}

Предложенный тест уже выполнен. В явном quotient-нормированном
тридцатимерном coexact-базисе оболочки \(n=3\) проекция шести leaked
Killing-образов имеет
\begin{equation}
 \Tr G_{1\to3}=80,
 \qquad
 \operatorname{rank}G_{1\to3}=6,
\end{equation}
а не ноль. Консервативный denominator-proxy равен \(80/12^2=5/9\), а не
малому остатку C6. Последующий полный same-scheme аудит не нашёл
обязательной компенсации.

\begin{nogocriterion}[Запрет повторного открытия H1]
H1 закрыт отрицательно в текущей Maxwell--ghost модели. Его допустимо
переоткрыть только при выводе нового обязательного same-spectrum сектора,
точного spectral identity либо внешне доказанной эквивалентности нового
primary quotient. Повторная проекция уже вычисленного канала не является
новым маршрутом.
\end{nogocriterion}

\section{H2: reciprocal плюс defect}

Точка \(r=1\) является минимумом
\(G(r)+G(r^{-1})\) вследствие добавленной симметризации. Ambient-аудит
доказал, что текущий Maxwell--FP спектр не несёт требуемой
\(r\leftrightarrow r^{-1}\)-дуальности. Defect-saddle также остаётся
условным на parent-derived stiffness и condensation.

Для суммы этих объектов не заданы общая размерность действия, относительная
нормировка, единый вариационный принцип и карта к
\((M_W,M_Z,\sin^2\theta_W)\). Поэтому H2 пока не является вычислимым
multi-blind gate. Назначение единичного относительного коэффициента было бы
новым скрытым весом.

Кроме того, шкала тома требует сначала пройти \(5\to6\) внешним
воспроизведением, а \(6\to7\) определён как parameter-free EW/QCD matching.
H2 не может перескочить эти ворота одним стационарным пунктом.

\section{Повторная свёртка полной coexact-башни}

Повторный расчёт
\begin{equation}
 \Delta_{\rm abs}
 =\frac{\pi^2}{2}\frac{T_{\rm coex}}{S_{\rm geo}}
  \left(1-\frac{10}{24S_{\rm geo}}\right)
 =5.4667669181\cdot10^{-7}
\end{equation}
совпадает с прежним absorption-аудитом. Относительное расхождение с
\(1/(\pi^4S_{\rm geo}^2)\) равно \(3.0407\cdot10^{-6}\), а абсолютное
расхождение равно
\begin{equation}
 1.6623\cdot10^{-12},
\end{equation}
не \(1.66\cdot10^{-11}\). Оно действительно меньше train-остатка
\(|S_{\rm vac}-\alpha^{-1}|=3.4777\cdot10^{-9}\), но сравнение двух
train-остатков не создаёт нового доказательного статуса.

Необходимый коэффициент \(N_{\rm need}=10.0099700224\) показывает
несовпадение с projector-rank \(10\) в данной интерпретации. Однако
нецелое эффективное значение само по себе не запрещено для общего
взвешенного следа; запрет относится именно к отождествлению этого числа с
\(\Tr P_{0,2}\). Более фундаментально, положительная Bessel-сумма
\(T_{\rm coex}\) является Casimir kernel, а не Euclidean winding
log-determinant. Поэтому даже точное числовое устранение рангового остатка
не закрыло бы determinant theorem без функционального моста.

\begin{remark}[Статус повторной свёртки]
Формулировка «absorption с малым ранговым остатком» является уточнённым
описанием уже существующей strong structural compression, а не повышением
её статуса. Кроме того, после ревизии R1 внутренний уровень уже равен
\(R_{\rm sci}=5/10\); путь \(4\to5\) не заблокирован C6, поскольку он пройден
через честное понижение exact-утверждения и независимое воспроизведение.
\end{remark}

\begin{protocol}[Рекомендуемая последовательность]
\begin{enumerate}
 \item заморозить и выпустить пакет R1 со статусом \(R_{\rm sci}=5/10\);
 \item подготовить claim--evidence--assumption--failure-mode таблицу для
 внешнего воспроизведения \(5\to6\);
 \item не тратить вычислительный ресурс на H1 без нового обязательного
 сектора;
 \item возвращаться к H2 только после вывода общего parent action,
 относительной нормировки и заранее объявленной blind-карты.
\end{enumerate}
\end{protocol}

Машинный протокол сохранён в
\texttt{s2t\_revision\_r1\_route\_triage\_results.json}.

\chapter{Пакетный аудит K1--K14 и reverse-проверка высших порядков}

Расширенный список кандидатов содержит как новые конструкции, так и
маршруты, уже закрытые прежними operator-gates. Для предотвращения
повторного перебора все четырнадцать вариантов проверены в одной
замороженной ведомости.

\section{Приоритетные кандидаты K1, K2 и K6}

Для K1 eta-инвариант аддитивен. Поэтому прямая сумма двух одинаковых
блоков с \(\eta=-1/4\) формально даёт
\begin{equation}
 \eta_{\rm dbl}=-\frac12,
 \qquad
 -\frac{\eta_{\rm dbl}}2+\frac1{12}
 =\frac14+\frac1{12}=\frac13.
\end{equation}
Это точное арифметическое прохождение, но не mass gate. Два одинаковых
хиральных блока удваивают determinant-фазу, тогда как физическое
vectorlike-сопряжение сокращает eta-фазы. Вещественный mass response
по-прежнему не получает аддитивный член \(1/3\). K1 закрыт в минимальной
vectorlike-модели и может быть переоткрыт только после вывода chiral
CP-нечётного механизма.

В K2 предложенный нормировочный множитель равен
\begin{equation}
 (\sqrt2)^2\cdot\frac12=1.
\end{equation}
Следовательно, он не меняет каноническое значение
\(\mathcal J=1\) и не порождает
\(\mathcal J_{\rm req}=5.4027533071\ldots\). Дополнительное rank-деление
не задано оператором. K2 остаётся спецификацией будущей двухвершинной
алгебры, а не численно закрытым кандидатом.

K6 сохраняет наиболее сильный условно-положительный статус. Mod-2 defect
даёт один real Majorana kernel channel, а дополнение имеет ранг \(23\).
Открыт не индексный расчёт, а parent-action gate: необходимо вывести
глобальный BdG-оператор, конденсацию и включение quotient-rank в collective
stiffness из одного действия.

\section{Тензорные кандидаты K12 и K13}

Для K12 одна полностью обращённая по статистике two-form BV-система даёт
в пространственно постоянной оболочке \(+\pi^4/6\), и шесть копий
арифметически дают \(+\pi^4\). Полный четырёхмерный trace при этом
ультрафиолетово расходится, а multiplicity six и обращение статистики не
выведены из gauge symmetry. Поэтому K12 является seed новой graded
parent-theory, но не спасением текущего действия.

K13 проходит точную идентичность
\begin{equation}
 \operatorname{Var}(\bar x)
 =\frac1{6\zeta(4,1/2)}
 =\frac1{\pi^4}.
\end{equation}
Два равных источника дают арифметический контроль
\(J=\sqrt2\), \(J^2/2=1\). Однако coherent source map, ordinary trace,
bosonic half-shift и равные веса шести каналов не выведены. Статус K13:
identity pass, source/action fail.

\section{Сводный статус остальных кандидатов}

\begin{center}
\small
\begin{tabular}{>{\raggedright\arraybackslash}p{0.09\textwidth}>{\raggedright\arraybackslash}p{0.72\textwidth}>{\raggedright\arraybackslash}p{0.12\textwidth}}
\toprule
Кандидат & Результат & Статус \\
\midrule
K3 & требуемое \(X=\pi^2/(3\alpha)=450.83\ldots\) получено обращением целевой формулы & circular pruner \\
K4 & \(27/5\) и \(17/\pi\) близки только численно; оператор отсутствует & screen only \\
K5 & уровень \(k=22\) и WZW-shift не выведены из действия & open count \\
K7 & \(24\) есть порядок \(S_4\) и размерность regular representation, а не размерность группы; физическая regular representation не построена & unsupported \\
K8 & \(e^\pi\) отличается от \(23\) на \(0.612\%\) & pruner \\
K9 & finite-threshold cone алгебраически непуст, но физически пуст в доступном mass interval & closed negative \\
K10 & kinetic mixing \(B-L\) не создаёт требуемые \(SU(2)\) и \(SU(3)\) компоненты & insufficient \\
K11 & split two-loop route уже закрыт неверным знаком и anchor failure & pruner \\
K14 & множитель \(\alpha/(4\pi)\) уменьшает существующую \(m_D\)-вставку в \(1.72\cdot10^3\) раза & fails target \\
\bottomrule
\end{tabular}
\end{center}

\section{Reverse-проверка \(1/S^3\) и \(1/S^4\)}

Для остатка
\begin{equation}
 \Delta=\alpha^{-1}-S_{\rm vac}
 =3.4776519442\cdot10^{-9}
\end{equation}
обратный подбор даёт
\begin{equation}
 c_3=\Delta S_{\rm geo}^3=0.00894938\ldots,
 \qquad
 c_4=\Delta S_{\rm geo}^4=1.22639016\ldots.
\end{equation}
Эти числа корректны для двух выбранных inverse-power ansatz. Но степень
\(1/S_{\rm geo}\) не является числом фейнмановских петель. Loop expansion
должен следовать из coupling expansion и диаграмм или determinant
effective action. Аналогично, индексы оболочки \(n\) и winding \(q\)
перечисляют моды внутри одного determinant и не задают loop-order.

В coexact-башне
\begin{equation}
 \frac{T_{n=3}}{T_{\rm coex}}
 =2.43135\cdot10^{-5},
\end{equation}
то есть следующая оболочка действительно слишком мала, чтобы объяснить
данный остаток в той же absorption-нормировке. Это исключает конкретную
гипотезу о higher-shell correction, но не доказывает отсутствие
трёх- или четырёхпетлевых вкладов.

\begin{nogocriterion}[Запрет ложной loop-интерпретации]
Близость остатка к \(c/S^n\) не считается loop-выводом без:
\begin{enumerate}
 \item объявленного coupling-counting;
 \item набора диаграмм или порядка determinant expansion;
 \item схемы перенормировки;
 \item переноса того же коэффициента хотя бы на второй observable.
\end{enumerate}
Отсутствие совпадения с коротким меню констант также не доказывает полноту
формулы. Неопределённость frozen anchor в этом тесте отдельно не
пропагирована, поэтому называть остаток anchor noise преждевременно.
\end{nogocriterion}

\section{Полный BV follow-up тензорной ветви}

Последующий ответ выписал полный редуцируемый BV-комплекс twisted
двухформенного поля:
\begin{equation}
 Z_2\propto
 \det{}'(\Delta_2)^{-1/2}
 \det{}'(\Delta_1)^{+1}
 \det{}'(\Delta_0)^{-3/2}.
\end{equation}
Контроль степеней свободы даёт одну физическую bosonic моду, как и должно
быть для четырёхмерной двухформы. На spatial zero shell полный комплекс
даёт \(-\pi^4/6\), а полное обращение статистики --- \(+\pi^4/6\).
Следовательно, требуемое \(+\pi^4\) возникает только после введения шести
независимых обращённых комплексов; один BV-функционал такой multiplicity не
выводит.

Происхождение \(\zeta(4,1/2)\) корректно только для второй вариации
аддитивно деформированного оператора:
\begin{equation}
 \Gamma''(0)\propto-\Tr K^{-2},
 \qquad
 K_n=(n+\tfrac12)^2.
\end{equation}
При этом стационарность или вычитание первой вариации
\(\Tr K^{-1}\) должны быть частью объявленного generating functional, а не
последующей интерпретацией.

Полный follow-up выявил четыре окончательных разрыва:
\begin{enumerate}
 \item reality condition связывает Fourier-ветви, а BV quotient сокращает
 component count до одной физической моды; шесть independent complexes не
 следуют из \(3\times2\) кинематического разложения;
 \item Grassmann determinant не задаёт положительную вероятность и
 дисперсию commuting readout без отдельной бозонизации;
 \item нормировка источника \(J^2/2=1\) не выводится;
 \item вторая независимая normalization-sensitive величина отсутствует.
\end{enumerate}

\begin{proposition}[Закрытие тензорного parent-action кандидата]
Точная inverse-susceptibility identity сохраняется как математический
результат, но полный graded/BV generating functional не выводит физический
\(S_{\rm vac}\)-вклад. Итоговый статус:
\begin{equation}
 \text{algebraic identity pass}
 \quad/\quad
 \text{parent-action fail}.
\end{equation}
Тензорную ветвь следует архивировать до появления независимо выведенной
graded symmetry, бозонизации source readout и второго сектора.
\end{proposition}

\section{Рекомендуемый порядок}

Главным живым маршрутом остаётся K6 как derivation одного defect/BdG
parent action. K13 допустим как параллельный новый parent-action кандидат
для электромагнитного сектора. K12 разрешён только при независимом выводе
graded symmetry и полного BV-комплекса. K1, K3, K8 и K11 повторно не
открываются без новой динамики; K2 сначала требует явной vertex algebra.

После R1 внутренний статус остаётся
\begin{equation}
 R_{\rm sci}=5/10,
 \qquad
 N_{\rm closed\ physical}=0.
\end{equation}

Машинный протокол сохранён в
\texttt{s2t\_k1\_k14\_loop\_reverse\_results.json}.

\chapter{Global-consistency gate для предложенного
\texorpdfstring{\(B-L\)}{B-L} defect action}

Предложенное действие объединяет root-связность, torsion-twisted
charge-two pairing field, Majorana Yukawa vertex и half-shifted
неоднородное седло. До динамического kill-gate необходимо проверить,
являются ли эти четыре элемента глобально совместимыми.

\section{Bundle gate Majorana-вершины}

На generator \(y\) sterile line имеет holonomy
\begin{equation}
 r=i,
 \qquad
 r^2=-1.
\end{equation}
Обычное charge-minus-two pairing field имеет holonomy
\(\phi=-1\). Поэтому
\begin{equation}
 \phi r^2=+1:
\end{equation}
Majorana vertex \(\Phi N^cN^c\) является scalar, но nowhere-zero parallel
condensate запрещён.

Для предложенного twisted field
\begin{equation}
 \Phi_{\rm tw}\in
 \Gamma(L_{\rm tor}\otimes L_{B-L}^{-2})
\end{equation}
полный holonomy равен
\begin{equation}
 \phi_{\rm tw}=(-1)(-1)=+1.
\end{equation}
Parallel condensate теперь допустим, но
\begin{equation}
 \phi_{\rm tw}r^2=-1.
\end{equation}
Следовательно, записанный член
\(y_N\Phi_{\rm tw}N^cN^c\) является torsion-odd и не задаёт глобальный
scalar functional. Это ровно ранее доказанная
root--mass--condensate trilemma.

\section{Двойной учёт holonomy в спектре}

Если torsion и charge-two holonomies сокращаются, полный holonomy pairing
bundle равен \(+1\). Тогда spectrum ковариантного оператора имеет
integer-shifted вид
\begin{equation}
 k_n=\frac{2\pi n}{L}
\end{equation}
и содержит нулевую моду. Half-shifted spectrum
\begin{equation}
 k_n=\frac{2\pi n+\pi}{L}
\end{equation}
принадлежит ordinary untwisted charge-two field с total holonomy \(-1\).
Использование total holonomy \(+1\) и half-shift одновременно дважды
учитывает один и тот же знак.

Поэтому предложенное twisted действие не выводит заявленные
\(n=0,-1\) branches и порог \(1-\lambda v^2\). Эти результаты относятся к
ordinary nonuniform defect branch, в которой Yukawa vertex глобально
инвариантен, но condensate не parallel.

\section{Оставшиеся derivation gates}

Даже после выбора правильной ordinary defect branch остаются:
\begin{enumerate}
 \item Maxwell term при \(f=0\) не выбирает Wilson holonomy
 \(\oint a=\pi/2\) из непрерывного пространства flat connections;
 \item отрицательный mass-square введён определением
 \(m_\Phi^2=-\lambda v^2\), а не получен gap equation;
 \item профиль \(\tanh x\) и \(S_{\rm core}\) не выведены решением
 предъявленного действия;
 \item внешний rank-24 Majorana module не построен из перечисленных полей;
 \item равные kernel/quotient weights остаются выбором canonical metric.
\end{enumerate}

Коэффициент kinetic mixing
\(\sum Y(B-L)=8/3\) является корректным trace identity, но не второй
предсказанной величиной. Физическое mixing требует boundary value, gauge
couplings, thresholds и running, отсутствующих в предъявленном
core-complement action.

\begin{proposition}[Вердикт global-consistency gate]
Предложенное действие проваливается раньше собственного динамического
kill-gate:
\begin{equation}
 \text{twisted parallel condensate}
 \quad\Longrightarrow\quad
 \text{torsion-odd Yukawa vertex},
\end{equation}
а заявленный half-shift несовместим с total holonomy \(+1\).
Следовательно, action в записанном виде не является глобально определённым
parent action.

Выживает только ordinary nonuniform charge-two defect texture с holonomy
\(-1\). Для неё mod-2 kernel остаётся условно согласованным, но
конденсация, core profile, rank-24 embedding и relative metric должны быть
выведены заново.
\end{proposition}

\section{Follow-up: строгий выбор ordinary ветви A}

Последующий ответ корректно выбрал ordinary charge-minus-two field без
torsion twist. В этой ветви scalar Yukawa vertex глобально определена, а
holonomy \(-1\) даёт half-shifted spectrum и conjugate lowest modes.
При заранее заданных root Wilson sector и broken-symmetry potential
получается
\begin{equation}
 \rho^2=v^2-\frac{\pi^2}{\lambda L^2},
 \qquad
 \lambda v^2>\frac{\pi^2}{L^2}.
\end{equation}
Это закрывает bundle-consistency и circle-saddle gates.

Однако статус «одно действие замкнуто при данном SSB» требует ещё четырёх
уточнений:
\begin{enumerate}
 \item Maxwell term на flat connection не выбирает
 \(\oint a=\pi/2\); root Wilson sector остаётся boundary/topological input;
 \item circle plane-wave с постоянной амплитудой не выводит transverse
 radial kink \(\tanh x\) или \(\operatorname{sech}x\) zero mode; необходимо
 решить tubular vortex profile и core boundary problem;
 \item mod-2 index фиксирует нечётность числа real zero modes, но ровно один
 kernel требует minimality/gap theorem, исключающей дополнительные пары;
 \item rank \(24\) и относительный canonical вес quotient/kernel не следуют
 из предъявленного \(B-L\) action, поэтому \(23+\pi^{-1}\) остаётся
 conditional metric readout.
\end{enumerate}

\begin{remark}[Исправленный статус ветви A]
Ветвь A является единственной глобально согласованной минимальной ветвью и
условно выводит nonuniform circle texture. Она не является полностью
замкнутой даже после задания SSB: дополнительно требуются вывод root Wilson
sector, radial core solution, exact-one kernel theorem и parent-derived
rank-24 metric.
\end{remark}

Машинный протокол сохранён в
\texttt{s2t\_bl\_defect\_action\_global\_consistency\_results.json}.

\chapter{Вильсоновский селектор трёх Majorana-core мод}

Полярный аудит ordinary \(B-L\)-ветви исправил важный подсчёт: если vortex
поколенчески-слеп, то до дополнительных family-взаимодействий он создаёт
три вещественные core-моды
\begin{equation}
 \chi=(\chi_1,\chi_2,\chi_3)^T\in\mathbb R^3_{gen},
\end{equation}
а не одну. Это закрывает прежнюю попытку получить \(24-1\) одним лишь
проектированием pairing-матрицы, но открывает механизм, уже присутствующий
в части II.B.

\section{Антисимметричная связь на core}

Непрерывная вильсоновская ветвь построила правильное вращение
\(R_n(\theta_*)\in SO(3)\) на family triplet, где
\begin{equation}
 \cos\theta_*=\frac{26-9\sqrt{15}}{11}.
\end{equation}
Его главный генератор имеет вид
\begin{equation}
 K_n=\operatorname{Log}R_n(\theta_*)
     =\theta_*[n]_\times,
 \qquad
 ([n]_\times)_{ab}=\epsilon_{abc}n_c.
\end{equation}
Он вещественен, антисимметричен и при \(\theta_*\ne0\) имеет
\begin{equation}
 \operatorname{rank}K_n=2,
 \qquad
 \ker K_n=\mathbb R n.
\end{equation}

Поэтому Majorana boundary term
\begin{equation}
 S_{\rm core,fam}
 =\frac{i}{2}\int_\gamma \chi^T K_n\chi
\end{equation}
спаривает две из трёх core-мод в одну ненулевую fermionic пару и оставляет
ровно одну вещественную нулевую моду вдоль оси \(n\). Это не требует
rank-one Yukawa matrix \(uu^T\): selector возникает как kernel уже
существующего ориентированного family generator.

\section{Независимость ранга от выбора вакуума}

Factor-axis functional сокращает восемь совместных осей до двух
симметрийно-связанных минимумов. Машинная проверка для обеих осей даёт
\begin{equation}
 \Spec(iK_n)=\{-\theta_*,0,+\theta_*\},
 \qquad
 \dim\ker K_n=1.
\end{equation}
Замена ориентации \(n\mapsto-n\) меняет \(K_n\mapsto-K_n\), но сохраняет
kernel projector \(P_n=nn^T\). Следовательно, знак ориентированного
вакуумного выбора не влияет на число surviving Majorana modes.

\begin{proposition}[Условное снятие generation multiplicity]
Если family Wilson connection ограничивается на defect core именно через
\(K_n\), то поколенчески-слепой vortex сначала создаёт три локальные
Majorana-моды, после чего антисимметричная family-связь оставляет одну:
\begin{equation}
 3\longrightarrow1.
\end{equation}
При наличии ambient-to-core restriction map это восстанавливает кандидат на
rank-one kernel и complement rank \(24-1=23\) без ручной вставки \(uu^T\).
\end{proposition}

\section{Оставшийся parent-action gate}

Положительный алгебраический результат ещё не является полным замыканием.
Текущее действие не выводит в одном функционале:
\begin{enumerate}
 \item ограничение той же \(SO(3)\) family connection на \(B-L\) defect core;
 \item коэффициент и знак \(S_{\rm core,fam}\) вместе с vortex action;
 \item ambient-to-core map из \(\mathbb R^{24}\) в три локальные моды;
 \item collective stiffness, превращающую complement rank в
 \(23+\pi^{-1}\), а не только в кратность determinant.
\end{enumerate}

\begin{nogocriterion}[Wilson--defect bridge gate]
Если \(S_{\rm core,fam}\) добавляется как независимый boundary term, не
следующий из того же graded action, что и root vortex и Wilson gap-functional,
то rank-one selector имеет статус конструктивного механизма, но не вывода
S2T. Положительное закрытие требует одного covariant derivative или одного
boundary superconnection, чья редукция одновременно даёт vortex BdG operator
и \(K_n\)-связь на его трёхмерном kernel bundle.
\end{nogocriterion}

Таким образом, проблема больше не состоит в отсутствии возможного generation
selector. Конкретный selector найден внутри уже построенной Wilson-ветви.
Открытым остаётся мост между family Wilson connection и \(B-L\) defect
parent action.

Машинный аудит сохранён в
\texttt{s2t\_family\_wilson\_majorana\_core\_selector\_results.json}.

\chapter{Динамический family-order parameter и rank-one defect selector}

Постоянная антисимметричная форма на family triplet запрещена остаточной
группой, но динамическое поле
\begin{equation}
 \Sigma\in V=\chi_1\oplus\chi_2\oplus\chi_3
\end{equation}
обходит этот no-go. Поскольку \(\chi_1\chi_2\chi_3=1\), взаимодействие
\begin{equation}
 S_{\Sigma\chi}
 =\frac{ig}{2}\int_\gamma
 \epsilon_{ijk}\Sigma_k\chi_i\chi_j
\end{equation}
является \(G_{res}\)-инвариантным. Для \(\Sigma\ne0\) оператор
\begin{equation}
 A_{ij}(\Sigma)=g\epsilon_{ijk}\Sigma_k
\end{equation}
имеет rank (2) и kernel \(\mathbb R\Sigma\). Поэтому три локальные
Majorana-моды редуцируются к одной без ручной вставки rank-one projector.

\section{Исправление анализа потенциала}

Для изотропного тестового потенциала
\begin{equation}
 V(\Sigma)=m^2|\Sigma|^2+\frac u4|\Sigma|^4
 +v\Sigma_1\Sigma_2\Sigma_3,
 \qquad u>0,
\end{equation}
диагональные stationary points имеют
\begin{equation}
 \Sigma=s(\varepsilon_1,\varepsilon_2,\varepsilon_3),
 \qquad
 q=\varepsilon_1\varepsilon_2\varepsilon_3,
\end{equation}
где
\begin{equation}
 3us^2+qvs+2m^2=0.
\end{equation}
Следовательно, радиус не равен просто \(2|m^2|/u\) при \(v\ne0\).
Угловой Hessian положителен для \(qv<0\), поэтому знак \(v\) оставляет
четыре, а не восемь, энергетически предпочтительных ориентаций.

Более того, отрицательное bare \(m^2\) не является необходимым условием.
Для предпочтительной ветви первый порядок возникает при
\begin{equation}
 m^2=\frac{v^2}{27u},
\end{equation}
а ненулевой global minimum существует уже при
\begin{equation}
 m^2<\frac{v^2}{27u}.
\end{equation}
Таким образом, symmetry-allowed cubic term способен вызвать конденсацию даже
при положительной квадратичной массе. Главный динамический gate состоит не
только в доказательстве \(m^2<0\), а в выводе полного отношения
\((m^2,u,v)\) из parent action.

\section{Факторный quadratic selector}

Остаточная \(\mathbb Z_2\times\mathbb Z_2\) симметрия сама по себе разрешает
три независимых инварианта \(\Sigma_i^2\). Поэтому полностью изотропная
квадратичная масса является дополнительным предположением. Уже существующий
factor Laplacian даёт канонический кандидат
\begin{equation}
 M_\Sigma^2=m_0^2I+\kappa L(w_3,w_1),
\end{equation}
со спектром
\begin{equation}
 m_0^2+2\kappa w_3,
 \quad
 m_0^2+2\kappa w_1,
 \quad
 m_0^2+2\kappa(w_3+w_1).
\end{equation}
При \(w_1<w_3\) и \(\kappa>0\) первой теряет устойчивость уникальная
\(S^1\)-характерная компонента. Это геометрический mass selector, а не
ручной выбор оси. Открытым остаётся вывод коэффициента \(\kappa\) и его знака
из того же действия, которое создаёт defect.

\section{Вклад core Majorana pair}

Две моды, спаренные оператором \(A(\Sigma)\), имеют gap
\(g|\Sigma|\). Если разрешено минимизировать по обеим parity-ветвям
двухмодового core Hamiltonian, его zero-temperature ground-state energy
содержит
\begin{equation}
 \Delta E_f=-\frac12|g\Sigma|,
\end{equation}
а при конечном обратном температурном размере \(\beta\)
\begin{equation}
 \Delta F_f
 =-\frac1\beta\log\left[2\cosh
 \left(\frac{\beta g|\Sigma|}{2}\right)\right]
 =\mathrm{const}-\frac{\beta g^2}{8}|\Sigma|^2+O(|\Sigma|^4).
\end{equation}
Знак квадратичного вклада благоприятствует конденсации. Однако cusp
\(-|g\Sigma|/2\) не является автоматическим свойством полного defect.
При фиксированной глобальной fermion parity система может быть обязана
оставаться на одной ветви \(\pm g|\Sigma|/2\), а нечётное число локальных
Majorana operators требует bulk или удалённого parity completion. Кроме того,
при конечном \(\beta\) конденсация следует только после проверки знака полной
renormalized quadratic mass. Поэтому фермионный вклад создаёт реальный
механизм неустойчивости, но ещё не доказывает \(\Sigma\ne0\) при произвольном
bare \(m^2\).

Этот детерминант зависит только от \(|\Sigma|\) и не генерирует
\(\Sigma_1\Sigma_2\Sigma_3\); ориентационный cubic coefficient должен прийти
из Wilson/factor или другого bosonic boundary sector.

\section{Bulk--core compatibility gate}

Инвариант \(\epsilon_{ijk}\Sigma_k\chi_i\chi_j\) корректен как взаимодействие
уже спроецированных вещественных core Majorana operators. Его нельзя без
дополнительного вывода заменить bulk Yukawa-членом того же family-типа.
Для четырёхмерных Weyl/Majorana fields лоренц-скалярный bilinear
\(N_iN_j\) симметричен по family indices, поэтому contraction с
\(\epsilon_{ijk}\Sigma_k\) тождественно исчезает. Следовательно parent action
должен породить антисимметричный оператор либо после контролируемой
zero-mode projection, либо через family connection, derivative coupling или
иной bulk operator, чья core restriction равна \(A(\Sigma)\).

Это устраняет потенциальную circularity: локальные \(\chi_i\) нельзя считать
независимым фундаментальным входом до вывода самого defect spectrum.

\section{Точная граница ориентационной устойчивости}

Пусть эффективный радиальный cusp допустим и
\begin{equation}
 V_{\mathrm{eff}}=
 \sum_{i=1}^{3}\mu_i\Sigma_i^2+
 \frac u4|\Sigma|^4+v\Sigma_1\Sigma_2\Sigma_3-h|\Sigma|,
 \qquad h=\frac{|g|}{2}.
\end{equation}
Для осевого stationary point \(\Sigma=(0,s,0)\), \(s>0\), имеем
\begin{equation}
 2\mu_2s+us^3-h=0.
\end{equation}
Поперечный Hessian равен
\begin{equation}
 H_\perp=
 \begin{pmatrix}
 2(\mu_1-\mu_2) & vs\\
 vs & 2(\mu_3-\mu_2)
 \end{pmatrix}.
\end{equation}
Поэтому геометрически выбранная \(S^1\)-ось является локальным минимумом
точно при
\begin{equation}
 \mu_1>\mu_2,
 \qquad
 \mu_3>\mu_2,
 \qquad
 4(\mu_1-\mu_2)(\mu_3-\mu_2)>v^2s^2.
\end{equation}
Если последнее неравенство нарушено, cubic term уводит минимум в
многокомпонентную ветвь. Тем самым выражение ``анизотропия доминирует''
заменяется проверяемым условием.

\section{Статус rank-one утверждения}

В трёхмерном projected core subspace при любом \(\Sigma\ne0\) утверждение
точно: \(\operatorname{rank}A=2\) и \(\dim\ker A=1\). Для полного BdG operator
этого недостаточно. Нужен gap/Feshbach--Schur аргумент, исключающий
дополнительные bulk и core zero pairs и доказывающий, что coupling к высшим
модам не закрывает spectral gap. До такого доказательства rank-one является
exact projected result, но conditional full-spectrum result.

\begin{protocol}[Следующий parent-action gate]
Нужно вывести из одного действия \(m_0^2\), \(\kappa\), \(u\), разрешённые
анизотропные quartic coefficients, \(v\) и \(g\); получить антисимметричную
core coupling из допустимого bulk operator; зафиксировать global parity
completion; проверить поперечный Hessian; затем доказать full BdG gap.
Положительный исход допускает два механизма конденсации: отрицательный
fermionic mass shift или cubic-driven first-order transition. После
доказательства \(\Sigma\ne0\) rank-one kernel автоматически следует только
в projected core subspace.
\end{protocol}

\begin{nogocriterion}[Запрет ручного order parameter]
Если \(\Sigma\), знак \(\kappa\), cubic coefficient \(v\) или coupling
\(g\) назначаются после просмотра flavour/neutrino data, конструкция остаётся
новой EFT-гипотезой. Алгебраический selector при этом корректен, но физическая
ветвь не замкнута.
\end{nogocriterion}

Машинный аудит сохранён в
\texttt{s2t\_dynamic\_family\_order\_parameter\_results.json}.

\chapter{Family-connection bridge и exact-one BdG kernel}

Предыдущий gate показал, что антисимметричный family selector нельзя
реализовать четырёхмерным scalar Majorana Yukawa term. Однако такой оператор
естественно возникает из ковариантной производной вещественного
\(SO(3)\)-triplet. Это даёт локальный bulk--core bridge без ручного
boundary mass.

\section{Допустимый tubular operator}

Пусть family-blind vortex operator на поперечном диске имеет вид
\(D_\perp\) и обладает единственной нормированной вещественной zero mode
\(\psi_0\), отделённой от остального спектра gap \(\Delta_\perp>0\).
Три поколения до включения family connection образуют
\begin{equation}
 \ker(D_\perp\otimes I_3)
 =\operatorname{span}\{\psi_0\}\otimes\mathbb R^3.
\end{equation}

На трубке \(N(\gamma)\simeq D^2\times S^1_L\) введём вещественную
\(SO(3)\)-связность
\begin{equation}
 \mathcal A_s=\frac1L K_n,
 \qquad
 K_n=\theta_*[n]_\times,
 \qquad
 \operatorname{Hol}_\gamma(\mathcal A)=e^{K_n}=R_n(\theta_*).
\end{equation}
Она входит не как scalar bilinear, а через разрешённую covariant derivative
triplet Majorana field:
\begin{equation}
 S_F=\frac12\int_{N(\gamma)}
 \Psi^TC\left[
 D_{a,\Phi}\otimes I_3+
 \Gamma^s\otimes(\partial_s+\mathcal A_s)
 \right]\Psi.
\end{equation}
Поскольку \(K_n^T=-K_n\), это стандартная реальная gauge coupling. После
проекции на \(\psi_0\) она даёт
\begin{equation}
 P_0S_FP_0
 =\frac12\int_\gamma
 \chi^T(\partial_s+L^{-1}K_n)\chi,
\end{equation}
то есть прежний антисимметричный core operator теперь следует из bulk
covariant derivative.

\section{Точный separable spectrum}

Для product metric, постоянной longitudinal connection, family-blind
vortex profile и ранее выведенной полной periodic longitudinal branch
квадрат BdG operator разделяется:
\begin{equation}
 \mathcal D^2
 =D_\perp^2\otimes I_3+
 I\otimes\left[-i\partial_s+\frac{iK_n}{L}\right]^2.
\end{equation}
Так как
\begin{equation}
 \operatorname{Spec}(iK_n)=\{-\theta_*,0,+\theta_*\},
\end{equation}
полный спектр квадратов имеет вид
\begin{equation}
 \lambda_{j,k,\sigma}^2
 =\varepsilon_j^2+
 \left(\frac{2\pi k+\sigma\theta_*}{L}\right)^2,
 \qquad
 k\in\mathbb Z,
 \quad
 \sigma\in\{-1,0,+1\}.
\end{equation}
Здесь \(\varepsilon_0=0\), а \(|\varepsilon_j|\geq\Delta_\perp\) для
\(j\ne0\). Поскольку
\begin{equation}
 0<\theta_*=\arccos\!\left(\frac{26-9\sqrt{15}}{11}\right)<\pi,
\end{equation}
нулевое значение возможно только при
\begin{equation}
 j=0,
 \qquad k=0,
 \qquad \sigma=0.
\end{equation}

\begin{theorem}[Exact-one separable kernel]
В separable tubular model с unit transverse Majorana index и family Wilson
holonomy \(R_n(\theta_*)\) полный BdG operator \(\mathcal D\) имеет ровно одну вещественную
zero mode:
\begin{equation}
 \dim\ker\mathcal D=1.
\end{equation}
Она равна \(\psi_0\otimes n\). Gap над ней ограничен снизу величиной
\begin{equation}
 \delta_{\rm sep}
 =\min\left\{\Delta_\perp,\frac{\theta_*}{L}\right\}.
\end{equation}
Для систолического core \(L=\pi\) family contribution равен
\begin{equation}
 \frac{\theta_*}{\pi}=0.7979250094\ldots.
\end{equation}
\end{theorem}

Это локальная теорема на tubular neighborhood. Она не утверждает, что
постоянная связность с данным углом продолжается на весь
\(\mathbb{RP}^3\) как плоская связность.

Таким образом, в этом классе моделей не требуется ни scalar order parameter
\(\Sigma\), ни отдельный core Yukawa coefficient \(g\): selector является
longitudinal Wilson splitting того же real triplet bundle.

\section{Устойчивость вне точного product ansatz}

Пусть \(\mathcal D_E=\mathcal D+E\), где \(E\) сохраняет class-D reality,
тот же mod-two defect sector и имеет operator norm
\begin{equation}
 \|E\|<\frac12\delta_{\rm sep}.
\end{equation}
Тогда spectral projection на интервал
\((-\delta_{\rm sep}/2,\delta_{\rm sep}/2)\) сохраняет rank one. Particle--hole
symmetry не позволяет единственной непарной eigenvalue уйти от нуля, поэтому
exact-one kernel устойчив, а новый gap не меньше
\begin{equation}
 \delta_E\geq\delta_{\rm sep}-\|E\|>\frac12\delta_{\rm sep}.
\end{equation}

Для perturbation, контролируемой блоками, достаточно более локального
Feshbach-критерия. Если \(P\) проектирует на три исходные core modes,
\(Q=I-P\),
\begin{equation}
 q=\|QEQ\|<\Delta_\perp,
 \qquad
 b=\|PEQ\|,
\end{equation}
то две поперечные family modes остаются обратимыми при
\begin{equation}
 \frac{\theta_*}{L}>
 \frac{b^2}{\Delta_\perp-q}.
\end{equation}
Осевая mode остаётся точной, если perturbation ковариантна относительно той
же connection и annihilates \(n\).

\section{Что закрыто и что осталось}

Закрыты два локальных разрыва:
\begin{enumerate}
 \item антисимметричный core operator получен из допустимой bulk gauge
 coupling, а не из исчезающего scalar bilinear;
 \item для separable tube доказан full-spectrum exact-one kernel, а для
 малых class-D perturbations дана явная область устойчивости.
\end{enumerate}

\section{Торсионное препятствие для плоского ambient extension}

Систолическая геодезическая представляет элемент порядка два в
\begin{equation}
 \pi_1(\mathbb{RP}^3\times S^1)=\mathbb Z_2\times\mathbb Z.
\end{equation}
Поэтому holonomy любой плоской ambient connection обязана удовлетворять
\begin{equation}
 W_\gamma^2=I.
\end{equation}
В \(SO(3)\) это допускает только углы \(0\) и \(\pi\), тогда как
\(0<\theta_*<\pi\) и \(\theta_*\ne\pi\). Следовательно, tubular connection
\(K_n/L\) не имеет плоского ambient extension вдоль систолического core.

\begin{nogocriterion}[Flat torsion Wilson obstruction]
Нельзя одновременно считать \(V_{\rm gap}\) потенциалом плоской ambient
Wilson-моды и помещать её holonomy \(R_n(\theta_*)\) на торсионный
систолический цикл. Такое чтение несовместимо с fundamental-group relation.
\end{nogocriterion}

\section{Два согласованных продолжения}

Первый маршрут использует неторсионный внешний фактор \(S^1\). Тогда defect
в полном carrier читается как \(\gamma\times S^1\): систолическая
\(\gamma\) несёт vortex topology и periodic longitudinal zero mode, а
family Wilson connection живёт на дополнительном круге длины
\(L_1=2\pi R_1\). Separable spectrum становится
\begin{equation}
 \lambda_{j,\ell,k,\sigma}^2
 =\varepsilon_j^2+
 \left(\frac{2\pi\ell}{L_\gamma}\right)^2+
 \left(\frac{2\pi k+\sigma\theta_*}{L_1}\right)^2.
\end{equation}
При periodic total spin branch на дополнительном круге единственный ноль
снова имеет индексы \((j,\ell,k,\sigma)=(0,0,0,0)\), а
\begin{equation}
 \delta_{S^1}
 =\min\left\{
 \Delta_\perp,
 \frac{2\pi}{L_\gamma},
 \frac{\theta_*}{L_1}
 \right\}.
\end{equation}
Для \(L_\gamma=\pi\), \(R_1=1\) family gap равен
\begin{equation}
 \frac{\theta_*}{2\pi}=a_*=0.3989625047\ldots.
\end{equation}
Однако anti-periodic spin branch на внешнем \(S^1\) уничтожает осевую
нулевую моду. Поэтому этот маршрут требует periodic defect sector либо
дополнительной компенсирующей holonomy, выведенной до данных.

\section{Компенсированная spin-ветвь из charge-two Higgs sector}

Компенсация не требует нового непрерывного поля. Пусть геометрическая
spin-структура на внешнем \(S^1\) anti-periodic и \(\eta=1/2\). Для
fermion нечётного \(B-L\)-заряда \(q\) при flat gauge holonomy
\(\alpha=\oint_{S^1}a\) полный momentum равен
\begin{equation}
 p_{k,\sigma}
 =\frac{2\pi(k+\eta)+q\alpha+\sigma\theta_*}{L_1}.
\end{equation}
После конденсации charge-two field gauge group содержит residual
\(\mathbb Z_2\)-holonomy
\begin{equation}
 \alpha=0\quad\text{или}\quad\alpha=\pi.
\end{equation}
Для нечётного \(q\) ветвь \(\alpha=\pi\) сокращает anti-periodic spin sign:
\begin{equation}
 e^{2\pi i\eta}e^{iq\alpha}=(-1)(-1)=+1.
\end{equation}
При этом charge-two condensate однозначен, поскольку
\begin{equation}
 e^{2i\alpha}=+1.
\end{equation}
После целочисленной перенумерации \(k\) спектр возвращается к periodic
формуле, и осевая \(\sigma=0\) zero mode восстанавливается. Следовательно,
compensated branch совместима с уже введённым \(B-L\) field content и не
вводит непрерывного fit-параметра.

Локальное charge-two Higgs action, однако, не различает две residual flat
ветви \(\alpha=0\) и \(\alpha=\pi\). Поэтому закрыта доступность требуемой
spin\(^{G}\)-структуры, но не её динамический выбор. Для полного вывода нужен
topological term, boundary condition или determinant sign, выбирающий
\(\alpha=\pi\) до обращения к neutrino data.

Действительно, для любого минимального локального функционала, зависящего от
\(F_a\), \(|D_a\Phi|^2\) и \(V(|\Phi|)\), две конфигурации
\begin{equation}
 (a,\Phi)=(0,v)
 \quad\text{и}\quad
 (a,\Phi)=\left(\frac{\pi}{L_1}ds,
 ve^{2\pi is/L_1}\right)
\end{equation}
имеют \(F_a=0\), \(D_a\Phi=0\) и одинаковый potential energy. Они не
связаны допустимым gauge transformation при наличии unit-charge fermions,
поскольку требуемое преобразование меняет знак после обхода круга.

\begin{nogocriterion}[Residual-holonomy degeneracy]
Минимальный local charge-two Higgs action не способен динамически выбрать
compensating branch \(\alpha=\pi\) вместо \(\alpha=0\). Выбор требует
nonlocal Wilson functional, topological coupling или квантового вклада
odd-charge sector. Простое объявление condensate не закрывает spin gate.
\end{nogocriterion}

\section{Однопетлевой Hosotani sign gate}

Для massive mode с \(\rho>0\) введём конечную holonomy-dependent часть
\begin{equation}
 I_\rho(\beta)=
 \log\frac{\cosh(2\pi\rho)-\cos(2\pi\beta)}
 {\cosh(2\pi\rho)-1}.
\end{equation}
На фундаментальном интервале \(0\leq\beta\leq1/2\) она монотонно возрастает,
поскольку
\begin{equation}
 \frac{dI_\rho}{d\beta}
 =\frac{2\pi\sin(2\pi\beta)}
 {\cosh(2\pi\rho)-\cos(2\pi\beta)}\geq0.
\end{equation}
Bosonic mode даёт \(+dI_\rho/2\), fermionic mode --- \(-dI_\rho/2\).

Для anti-periodic odd-charge fermion переход \(\alpha:0\to\pi\) меняет
\(\beta:1/2\to0\), поэтому
\begin{equation}
 \Delta V_f
 =V_f(\pi)-V_f(0)
 =\frac d2 I_\rho(1/2)>0.
\end{equation}
Для anti-periodic fermion с \(q=1/3\) конечная фаза равна
\(\beta=2/3\sim1/3\), и
\begin{equation}
 \Delta V_f
 =\frac d2\left[I_\rho(1/2)-I_\rho(1/3)\right]>0.
\end{equation}
Для periodic boson с \(q=1/3\) правильный shift равен \(q/2=1/6\), а не
\(1/3\), но знак остаётся тем же:
\begin{equation}
 \Delta V_b=\frac d2I_\rho(1/6)>0.
\end{equation}
Charge-two boson не различает ветви. Следовательно, минимальный набор
anti-periodic fermions и periodic bosons действительно выбирает
\(\alpha=0\), а compensated branch \(\alpha=\pi\) является более высокой
stationary branch.

При \(\rho\gg1\) правильная асимптотика равна
\begin{equation}
 I_\rho(\beta)
 =2\left[1-\cos(2\pi\beta)\right]e^{-2\pi\rho}
 +O(e^{-4\pi\rho}).
\end{equation}
В частности,
\begin{equation}
 I_\rho(1/2)\sim4e^{-2\pi\rho},
 \qquad
 I_\rho(1/3)\sim3e^{-2\pi\rho},
 \qquad
 I_\rho(1/6)\sim e^{-2\pi\rho}.
\end{equation}
Ранее предложенные коэффициенты \(2\) и \(1/2\) были занижены вдвое;
отношение \(I(1/2):I(1/6)=4:1\) при этом сохраняется.

\section{Geometric versus total periodicity}

Sterile \(N^c\) в compensated construction имеет geometric spin shift
\(\eta=1/2\). При \(\alpha=\pi\) её \emph{total} phase становится
periodic, \(\beta=0\), но Hosotani determinant сравнивает две ветви при
фиксированной geometric \(\eta\):
\begin{equation}
 \beta(0)=\frac12,
 \qquad
 \beta(\pi)=0,
 \qquad
 \Delta V_{N^c}=+\frac d2I_\rho(1/2)>0.
\end{equation}
Следовательно, такая \(N^c\) усиливает выбор \(\alpha=0\), а не создаёт
требуемый отрицательный вклад. Периодичность defect mode вдоль
систолической \(\gamma\), полученная сокращением spin- и torsion-знаков, также
не меняет geometric spin structure на независимом внешнем \(S^1\).

Отрицательный вклад
\(-dI_\rho(1/2)/2\) требует поля с geometric \(\eta=0\), для которого
\(\beta:0\to1/2\). Такого поля существующая compensated \(N^c\)-ветвь не
предоставляет.

Для минимального \(B-L\)-контента одного поколения имеются двенадцать
anti-periodic Weyl components с \(|q|=1/3\) и четыре с нечётным целым
зарядом, включая \(N^c\). При общей \(rho\) их вклад равен
\begin{equation}
 \Delta V_{gen}
 =2I_\rho(1/2)+6\left[I_\rho(1/2)-I_\rho(1/3)\right]
 =8I_\rho(1/2)-6I_\rho(1/3)>0.
\end{equation}
Даже если искусственно перевести одну \(N^c\) на geometric periodic branch,
получается
\begin{equation}
 \Delta V_{gen}^{\rm retwist}
 =7I_\rho(1/2)-6I_\rho(1/3)>0.
\end{equation}
В large-mass limit эти коэффициенты пропорциональны соответственно
\(14e^{-2\pi\rho}\) и \(10e^{-2\pi\rho}\). Поэтому три retwisted \(N^c\)
при одинаковых masses не проходят magnitude gate.

Для существующего минимального контента common-\(\rho\) assumption не
нужна. При произвольных положительных \(\rho_i\) каждый отдельный вклад имеет
положительный знак:
\begin{align}
 \Delta V_{odd,i}
 &=\frac{d_i}{2}I_{\rho_i}(1/2)>0,\\
 \Delta V_{1/3,i}
 &=\frac{d_i}{2}\left[
 I_{\rho_i}(1/2)-I_{\rho_i}(1/3)
 \right]>0,\\
 \Delta V_{b,1/3,i}
 &=\frac{d_i}{2}I_{\rho_i}(1/6)>0.
\end{align}
Charge-two bosons дают ноль. Следовательно, сумма положительна при любом
mass spectrum существующих anti-periodic fermions и periodic bosons.

\begin{theorem}[Minimal-content multi-mass Hosotani no-go]
При фиксированных geometric spin structures минимальный \(B-L\)-контент
выбирает \(\alpha=0\) для любых \(\rho_i>0\). Полный mass spectrum не может
изменить знак, поскольку отрицательных summands в этом field content нет.
Multi-mass gate остаётся открытым только после добавления нового
geometric-periodic odd-charge sector.
\end{theorem}

\begin{nogocriterion}[Total-periodicity sign no-go]
Нельзя использовать zero mode, которая становится total-periodic только в
ветви \(\alpha=\pi\), как periodic-spin determinant, стабилизирующий эту же
ветвь. Это меняет fixed boundary problem между концами сравнения и обращает
знак Hosotani contribution. Для отрицательного вклада нужен независимо
выведенный geometric-periodic odd-charge sector.
\end{nogocriterion}

\section{Почему обычный CS/theta term не стабилизирует ветвь}

Квантованный topological term различает сектора на уровне фазы. В евклидовой
сигнатуре
\begin{equation}
 S_{CS,E}=\frac{ik}{4\pi}\int a\wedge da,
 \qquad
 S_{\theta,E}=\frac{i\theta}{8\pi^2}\int F\wedge F.
\end{equation}
На полностью flat configuration оба локальных density равны нулю. При
добавлении фиксированного magnetic flux \(m\) Chern--Simons term может дать
\begin{equation}
 \Delta S_{CS,E}=ikm\alpha,
\end{equation}
то есть относительную фазу unit modulus, но не отрицательный вклад в
\(\operatorname{Re}V_{eff}\). Его real Hessian по \(\alpha\) равен нулю.
Discrete theta term аналогично меняет знак или фазу веса topological sector,
но сам по себе не превращает Hosotani maximum в энергетический minimum.

\begin{nogocriterion}[Topological-phase stabilization no-go]
Обычный квантованный CS/theta term без дополнительного constraint sector не
может стабилизировать \(\alpha=\pi\) против положительного real Hosotani
splitting. Он способен изменить interference секторов или, в расширенной
TQFT, спроецировать часть Hilbert space, но это не эквивалентно минимуму
вещественного effective potential.
\end{nogocriterion}

Для реального выбора \(\alpha=\pi\) требуется дополнительный вклад
\begin{equation}
 \Delta V_{min}:=V_{min}(\pi)-V_{min}(0)>0,
 \qquad
 \Delta V_{new}(\pi)-\Delta V_{new}(0)<-\Delta V_{min}.
\end{equation}
Минимальный знаково подходящий пример --- periodic odd-charge fermion: он
даёт
\begin{equation}
 \Delta V_{per\,f}=-\frac d2I_\rho(1/2)<0.
\end{equation}
Его multiplicity и mass должны быть выведены из parent action и превысить
положительный вклад исходного spectrum. Альтернатива --- явная TQFT
projection, удаляющая \(\alpha=0\) sector, но тогда это новый дискретный
сектор модели, а не Hosotani stabilization.

\section{Минимальный geometric-periodic selector sector}

Минимальный anomaly-safe кандидат состоит из двух left-handed singlet Weyl
fields
\begin{equation}
 X_+:(B-L)=+1,
 \qquad
 X_-:(B-L)=-1.
\end{equation}
Они образуют vectorlike Dirac pair. Чтобы при общей anti-periodic spin
structure получить effective geometric periodicity только для этой пары,
введём flat \(\mathbb Z_4^X\)-line с charges
\begin{equation}
 z(X_+)=+1,
 \qquad
 z(X_-)=-1\pmod4,
\end{equation}
и holonomy \(-1\) на внешнем \(S^1\). Тогда spin sign и \(\mathbb Z_4^X\)
sign сокращаются, а B--L holonomy переводит
\begin{equation}
 \beta:0\longrightarrow\frac12
 \qquad (\alpha:0\longrightarrow\pi).
\end{equation}
Пара даёт требуемый real contribution
\begin{equation}
 \Delta V_X=-I_{\rho_X}(1/2)<0.
\end{equation}

Непрерывные anomaly sums сокращаются попарно:
\begin{equation}
 \sum q=\sum q^3=0.
\end{equation}
То же происходит для elementary mixed discrete sums при представлении
\(z=\pm1\):
\begin{equation}
 \sum z=\sum z^3=\sum zq^2=\sum z^2q=0.
\end{equation}
Разрешённая симметрийная Dirac mass непрерывно переводит изолированную пару
в тривиально gapped phase без нарушения \(B-L\) или \(\mathbb Z_4^X\).
Следовательно, pair не несёт остаточной discrete 't Hooft anomaly: anomaly
gate для самого нового fermionic content закрыт.

Более того, Dirac mass \(M_XX_+X_-\) разрешена, тогда как
\begin{equation}
 \Phi X_+X_+,
 \qquad
 \Phi^\dagger X_-X_-,
 \qquad
 N^cX_-
\end{equation}
имеют ненулевой \(\mathbb Z_4^X\)-charge и запрещены. Поэтому новая пара не
участвует в vortex pairing operator и не добавляет defect zero modes;
исходный exact-one kernel сохраняется.

Пусть для benchmark все три минимальных поколения имеют \(\rho=1\). Их
положительный вклад равен
\begin{equation}
 \Delta V_{min}^{(3)}
 =3\left[8I_1(1/2)-6I_1(1/3)\right]
 =0.0783386447\ldots.
\end{equation}
Одна vectorlike pair выбирает \(\alpha=\pi\), если
\begin{equation}
 I_{\rho_X}(1/2)>\Delta V_{min}^{(3)},
\end{equation}
что даёт
\begin{equation}
 \rho_X<0.6259781084\ldots.
\end{equation}
При одинаковой массе \(\rho_X=1\) требуется не менее одиннадцати таких
pairs. В выбранной \(\alpha=\pi\)-ветви даже bare-light pair получает
Wilson gap; при одной паре на threshold
\begin{equation}
 m_{phys}R_1\geq
 \sqrt{\rho_X^2+\frac14}=0.801\ldots,
\end{equation}
поэтому determinant leverage не обязательно создаёт massless spectator в
физическом вакууме.

Порог существенно зависит от объявленной reference normalization. В общем
виде \(\rho_X^{crit}(\rho_{ref})\) определяется уравнением
\begin{equation}
 I_{\rho_X^{crit}}(1/2)
 =3\left[
 8I_{\rho_{ref}}(1/2)-6I_{\rho_{ref}}(1/3)
 \right].
\end{equation}
Для двух используемых benchmark получаем
\begin{align}
 \rho_{ref}=1:
 &\quad \rho_X^{crit}=0.6259781084\ldots,
 &&m_{phys}R_1=0.8011545371\ldots,\\
 \rho_{ref}=\frac12:
 &\quad \rho_X^{crit}=0.1399290937\ldots,
 &&m_{phys}R_1=0.5192110855\ldots.
\end{align}
При уменьшении \(\rho_{ref}\) положительный minimal-content determinant
растёт, поэтому допустимый \(\rho_X\) должен уменьшаться. Значение около
\(0.74\) не удовлетворяет threshold equation для \(\rho_{ref}=1/2\).

\section{Stiffness embedding и точное \(\delta_X=0\)}

Расширенная Hilbert space имеет block form
\begin{equation}
 \mathcal H_{ext}=\mathcal H_{24}\oplus\mathcal H_X.
\end{equation}
Защитная \(\mathbb Z_4^X\) запрещает mixing между blocks. Поэтому neutrino
collective tangent продолжается нулём:
\begin{equation}
 \Xi_{ext}=\Xi\oplus0_X.
\end{equation}
Для canonical product metric
\begin{equation}
 \|\Xi_{ext}\|^2
 =\|\Xi\|^2+\|0_X\|^2
 =23+\pi^{-1},
\end{equation}
то есть
\begin{equation}
 \boxed{\delta_X=0.}
\end{equation}
Тот же вывод следует для determinant Hessian по neutrino collective
coordinate \(t\): при фиксированном выбранном \(\alpha=\pi\)
\begin{equation}
 \frac{\partial^2\Gamma_X}{\partial t^2}=0,
\end{equation}
поскольку \(D_X\) не зависит от \(t\). Если же отождествить \(t\) с Wilson
angle \(\alpha\), появляется \(\rho_X\)-dependent Hessian и denominator
теряет каноническую форму. Поэтому X-sector допустим только как orthogonal
vacuum selector, а \(\pi^{-1}\) по-прежнему приходит из исходной cycle line,
не из X determinant.

\section{Kinetic-mixing и EW/QCD second-sector gate}

Пара \(X_\pm\) является Standard-Model singlet:
\begin{equation}
 (SU(3),SU(2))_Y=(\mathbf1,\mathbf1)_0.
\end{equation}
Поэтому
\begin{equation}
 \Delta\sum Y(B-L)=0,
 \qquad
 \Delta b_1=\Delta b_2=\Delta b_3=0.
\end{equation}
При отсутствии portal couplings её one-loop logarithmic и finite
nonlogarithmic matching vector для Standard-Model gauge couplings равен
\begin{equation}
 \Delta_X^{EW/QCD}=(0,0,0),
\end{equation}
что не может воспроизвести corrected target direction
\begin{equation}
 (1,-1,-6.68221\ldots).
\end{equation}
Двухпетлевая передача через \(B-L\) kinetic mixing возможна только после
задания \(g_{B-L}\), mixing boundary value и portal structure; она не
является parameter-free вторым сектором.

\begin{nogocriterion}[SM-singlet second-sector no-go]
Минимальная vectorlike \(X_\pm\)-пара проходит vacuum-selection и сохраняет
neutrino stiffness, но не даёт независимого EW/QCD normalization-sensitive
следа. Поэтому она остаётся самосогласованным vacuum-selector extension и не
проходит two-sector gate версии III.
\end{nogocriterion}

Здесь необходимы две статусные оговорки. Во-первых, determinant пары выбирает
Wilson branch \(\alpha=\pi\), но не family axis \(n\). Ось по-прежнему должна
быть выведена из residual-orbit dynamics или boundary equations. Во-вторых,
условие на \(\rho_X=M_XR_1\) использует непрерывную массу. Пока \(M_XR_1\) не
получена из parent action, неверно говорить, что селектор не вводит
непрерывных данных: он не требует нового fitting coefficient в stiffness, но
его determinant strength остаётся mass-dependent input.

Следовательно, консолидированный статус имеет вид
\begin{equation}
 \boxed{
 \begin{gathered}
 X_\pm:\ \text{условный динамический selector ветви }\alpha=\pi,\\
 \delta_X=0,\qquad
 \text{family-axis gate открыт},\\
 \text{EW/QCD two-sector gate закрыт отрицательно},\\
 N_{\rm closed\ physical}=0.
 \end{gathered}}
\end{equation}
Положительный результат состоит в существовании anomaly-safe знакового
механизма. Он не является физическим предсказанием и не повышает научный
ledger до вывода \(M_X\), происхождения \(\mathbb Z_4^X\)-line и общего
parent action.

\begin{protocol}[Periodic selector compatibility gate]
Кандидат проходит continuous/discrete anomaly, determinant-sign и
projected-kernel tests. Для физического статуса ещё необходимо:
\begin{enumerate}
 \item вывести \(\mathbb Z_4^X\)-line и её holonomy \(-1\) из finite algebra
 или обязательного topological sector;
 \item вычислить \(M_XR_1\), а не выбирать его ниже threshold;
 \item расширить механизм так, чтобы второй normalization-sensitive sector
 возникал без назначения Standard-Model charges или portals после просмотра
 target.
\end{enumerate}
До этого это минимальная согласованная новая модель, а не восстановление
версии II.A.
\end{protocol}

\section{Collective stiffness на маршруте A}

Пусть \(U_A\) отображает каждый внутренний eigenchannel в соответствующую
нормированную plane wave на внешнем \(S^1\). Для всех spin-, gauge- и
family-shifts
\begin{equation}
 U_A^\dagger U_A=I_{24},
 \qquad
 \|\varphi_{k,\sigma}\|_{L^2(S^1)}=1.
\end{equation}
Поэтому Hilbert--Schmidt traces проекторов не меняются:
\begin{equation}
 \operatorname{Tr}\!\left[(U_AP_{ker}U_A^\dagger)^2\right]=1,
 \qquad
 \operatorname{Tr}\!\left[(U_AP_HU_A^\dagger)^2\right]=23.
\end{equation}
Дополнительный круг даёт единичный множитель нормы и не меняет intrinsic
cycle Gram metric на \(\gamma\):
\begin{equation}
 \|\widehat e_0\|^2=1,
 \qquad
 \|e_1\|^2=L_\gamma^{-1}=\pi^{-1}.
\end{equation}
Следовательно, в canonical product trace--Hodge metric
\begin{equation}
 \boxed{\|\Xi_A\|^2=23+\pi^{-1}.}
\end{equation}
Результат не зависит от \(\theta_*\), spin shift или компенсирующей
\(\mathbb Z_2\)-holonomy: они меняют phases нормированных eigenfunctions, но
не их нормы и не ranks проекторов.

Тем самым collective-stiffness gate маршрута A закрыт как точное
кинематическое тождество canonical metric. Он ещё не возвращает
\(23+\pi^{-1}\) статус физического neutrino denominator: для этого та же
parent action должна выбрать compensated branch, вывести относительную
trace--Hodge metric и пройти второй normalization-sensitive sector.

Второй маршрут сохраняет holonomy на \(\gamma\), но ambient connection
должна иметь кривизну или singular defect source. Если \(2\gamma=\partial
\Sigma_2\) и curvature редуцирована к abelian subgroup с осью \(n\), то
\begin{equation}
 \exp(2\theta_*[n]_\times)
 =\exp\left([n]_\times\int_{\Sigma_2}F_n\right),
\end{equation}
то есть требуется
\begin{equation}
 \int_{\Sigma_2}F_n=2\theta_*\pmod{2\pi}.
\end{equation}
В principal branch это flux
\begin{equation}
 2\theta_*-2\pi=-1.2696746116\ldots.
\end{equation}
Его величина и совпадение оси с factor-selected \(n\) должны следовать из
уравнений движения; иначе это два новых boundary inputs.

Итак, плоский систолический bridge закрыт отрицательно. Положительный
family-selector route теперь имеет две точные версии: плоская связь на
неторсионном \(S^1\), где compensated branch существует и сохраняет
\(23+\pi^{-1}\), а новый \(X_\pm\)-sector может условно выбрать нужную ветвь,
но не выводит axis, mass или второй сектор; либо curved/singular
connection на \(\gamma\) с flux-and-axis dynamical gate. Маршрут A является
алгебраически минимальным, но его минимальный SM-singlet selector закрыт как
two-sector extension. Дальнейшее замыкание требует единого parent action, а
не ещё одного независимого selector-а. Маршрут B остаётся резервным.

Машинный аудит сохранён в
\texttt{s2t\_family\_connection\_defect\_gap\_results.json}.

\chapter{Бинарно-тетраэдральный carrier gate}

Внешний exploratory-аудит предложил заменить
\(\mathbb{RP}^3=S^3/\mathbb Z_2\) на сферическую форму
\begin{equation}
 M_{2T}=S^3/2T,
 \qquad
 2T\simeq\operatorname{SL}(2,3),
 \qquad |2T|=24,
\end{equation}
и интерпретировать число \(23\) как размерность augmentation ideal
\(\mathbb C[2T]\ominus\mathbf1\). Проверим отдельно carrier-нормировку,
спектр и представленческий мост.

\section{Carrier-тест}

При единичном радиусе
\begin{equation}
 \Vol(M_{2T})=\frac{2\pi^2}{24}=\frac{\pi^2}{12}.
\end{equation}
Тем самым исходное условие \(Z_A=\pi^2\) нарушено в двенадцать раз.
Абелианизация \(2T\) равна \(\mathbb Z_3\), поэтому плоские абелевы
характеры имеют фазовый шаг \(2\pi/3\), а не \(\pi\). Кроме того, минимальный
геометрический сдвиг задаётся элементом с вещественной частью \(1/2\), так
что систола круглого quotient-а равна
\begin{equation}
 \ell_{sys}(M_{2T})=\arccos(1/2)=\frac{\pi}{3}.
\end{equation}
Прямая замена в прежнем геометрическом скаляре даёт
\begin{align}
 S_{geo}^{2T}
 &=4\pi^3+\frac{\pi^2}{12}+\frac{\pi}{3}
 =125.8947713058\ldots,\\
 S_{vac}^{2T}
 &=S_{geo}^{2T}-\frac{1}{24S_{geo}^{2T}}
 -\frac{1}{\pi^4(S_{geo}^{2T})^2}
 =125.8944396939\ldots.
\end{align}
Расхождение с train-якорем \(\alpha^{-1}\) равно
\begin{equation}
 S_{vac}^{2T}-\alpha^{-1}=-11.1415594831\ldots.
\end{equation}
Следовательно, candidate не является продолжением исходной numerical
architecture без нового множителя или нового определения \(S_{geo}\).

\section{Явное усреднение характеров}

Для \(g\in2T\subset SU(2)\) характер
\(\operatorname{Sym}^{\ell}(\mathbb C^2)\) равен
\(U_\ell(\Re g)\), где \(U_\ell\) --- полином Чебышёва второго рода.
Поэтому
\begin{equation}
 m_\ell^{2T}
 =\dim\operatorname{Sym}^{\ell}(\mathbb C^2)^{2T}
 =\frac1{24}\sum_{g\in2T}U_\ell(\Re g).
\end{equation}
Прямой подсчёт по двадцати четырём quaternion units даёт первые ненулевые
степени
\begin{equation}
 \ell=0,6,8,12,14,16,18,\ldots.
\end{equation}
Первая положительная scalar shell имеет
\begin{equation}
 \lambda_6=6(6+2)=48,
 \qquad d_6=(6+1)m_6^{2T}=7.
\end{equation}

Для coexact one-forms homogeneous left quotient multiplicity равна
\begin{equation}
 d_n^{2T}
 =n\,m_{n+1}^{2T}+(n+2)m_{n-1}^{2T}.
\end{equation}
Первый уровень сохраняется:
\begin{equation}
 n=1,
 \qquad \lambda_1=4,
 \qquad d_1^{2T}=3.
\end{equation}
Поэтому положительный Bessel-tail не исчезает. При \(R_3=R_1=1\)
получаем
\begin{equation}
 T_{coex}^{2T}=7.6133956135\cdot10^{-6}>0,
\end{equation}
причём первый уровень даёт более
\(0.9999999999\) полной суммы. Замена quotient-а лишь примерно делит
доминирующий residual пополам и не создаёт cancellation theorem.

\section{Почему augmentation ideal не является quotient spectrum}

Регулярное представление действительно имеет каноническое разложение
\begin{equation}
 \mathbb C[2T]=\mathbf1\oplus I_{aug},
 \qquad \dim I_{aug}=23.
\end{equation}
Однако функции и формы на \(S^3/2T\) соответствуют \(2T\)-инвариантным
состояниям на \(S^3\). Quotient-проектор сохраняет \(\mathbf1\)-isotypic
часть и удаляет nontrivial regular components. Следовательно,
\(I_{aug}\) не возникает как конечный physical spectral sector самого
quotient-а.

Утверждение, что \(Q_8\) не действует на этом пространстве, также неверно:
\(Q_8\subset2T\) действует левыми сдвигами. Ограничение регулярного
представления содержит три копии регулярного \(Q_8\)-модуля, поэтому
\begin{equation}
 \dim\mathbb C[2T]^{Q_8}=3,
 \qquad
 \dim I_{aug}^{Q_8}=2.
\end{equation}
Старый spinor-line no-go здесь непосредственно не применяется, но это не
создаёт требуемого spectral bridge.

\begin{nogocriterion}[Binary-tetrahedral direct-carrier no-go]
Прямая замена \(\mathbb{RP}^3\) на \(S^3/2T\) нарушает locked volume и
flat-phase tests, сохраняет положительный coexact tail и не реализует
augmentation ideal как quotient eigenspace. Поэтому маршрут закрыт как
геометрический источник физического ранга \(23\). Он может быть переоткрыт
только как новая модель с отдельно выведенным внутренним
\(\mathbb C[2T]\)-fiber и единым действием, связывающим его с Dirac,
калибровочным и normalization-sensitive секторами.
\end{nogocriterion}

Машинный аудит сохранён в
\texttt{s2t\_binary\_tetrahedral\_carrier\_results.json}.

\chapter{Протокол закрытия исследовательского цикла S2T--II}

Настоящая глава фиксирует завершение внутреннего конструктивного цикла
версии II по состоянию на 9 августа 2026 года. Закрытие относится не ко всем
мыслимым геометрическим или квантово-полевым моделям, а к явно объявленному
классу S2T--II и проверенным минимальным расширениям.

\section{Область действия финального вердикта}

Исчерпан следующий рабочий класс:
\begin{enumerate}
 \item круглые сферические пространственные формы с product carrier
 \(M^3\times S^1\), locked volume- и flat-phase-нормировками;
 \item проверенные линзовые, бинарно-тетраэдральные и нетривиальные
 circle-bundle варианты;
 \item canonical trace--Hodge и рассмотренные affine/positive spectral
 weight-механизмы;
 \item минимальные Maxwell--ghost, defect, \(B-L\), Wilson-selector,
 threshold и split-running реализации;
 \item перечисленные в томе finite-operator, rotor, BV/BRST и
 family-texture маршруты.
\end{enumerate}
Из отрицательных результатов в этом классе не следует теорема о невозможности
всех классических компактных многообразий, всех расслоений или всех
некоммутативных геометрий. Такое более сильное утверждение потребовало бы
отдельной классификационной теоремы, которой в версии II нет.

\begin{nogocriterion}[Запрет глобального переобобщения]
Фраза ``геометрический путь закрыт'' в версии II означает закрытие
объявленного минимального carrier-класса. Она не означает доказанной
невозможности любой геометрической унификации.
\end{nogocriterion}

\section{Закрытые и перенесённые exploratory-маршруты}

Финальный menu имеет вид:
\begin{center}
\small
\begin{tabular}{>{\raggedright\arraybackslash}p{0.24\textwidth}
                >{\raggedright\arraybackslash}p{0.20\textwidth}
                >{\raggedright\arraybackslash}p{0.45\textwidth}}
\toprule
Маршрут & Статус & Основание \\
\midrule
\(S^3/2T\) & закрыт как direct carrier & volume/phase mismatch,
positive coexact tail, augmentation ideal не является quotient eigenspace \\
Нетривиальное \(S^1\)-расслоение & закрыто для C6 и discrete selector &
тотальное пространство типа \(U(2)\) имеет \(\pi_1=\mathbb Z\), \(b_1=1\)
и непрерывную Wilson-моду \\
Higher-form extension & архивировано & требует новых gerbe/2-form данных и
не фиксирует canonical kinetic normalization в текущем menu \\
NCG spectral triple & версия III & смена категории модели, finite
KO-dimension six и новый parent action \\
Trace--Hodge incompatibility & принято в ограниченном классе & canonical
metric переводит raw lepton seed в \(8/3\), сохраняя отдельное neutrino
identity; общий no-go для всех мер не утверждается \\
\bottomrule
\end{tabular}
\end{center}

\section{Научный статус}

После всех внутренних и collaborative red-team проверок фиксируется
\begin{equation}
 \boxed{
 R_{\rm sci}=4/10,
 \qquad
 N_{\rm closed\ physical}=0.}
\end{equation}
Оценка \(4/10\) означает, что математическое ядро, вычислительные аудиты и
карта no-go предъявлены к воспроизведению, но единое действие и независимые
физические предсказания отсутствуют. Внешняя рецензия сама по себе не
повышает шкалу: она может подтвердить или понизить текущий статус. Переход к
\(5/10\) требует заранее объявленного положительного критерия, например
полного независимого computational reproduction математического ядра вместе
с устранением одного из структурных normalization gates.

Полученный red-team материал имеет внешний источник, но после итеративных
исправлений, совместного выбора тестов и включения результатов в программу
он классифицируется как \emph{collaborative external-origin audit}, а не как
окончательная независимая replication. Независимая проверка требует
исследователя, не участвовавшего в построении и корректировке маршрутов.

\section{Условия версии III}

Новая версия допустима только если до обращения к данным предъявлены:
\begin{enumerate}
 \item новый parent action и полный набор его полей, мер и симметрий;
 \item происхождение rank-one structure, family axis и относительных
 sector weights;
 \item механизм обхода canonical trace--Hodge incompatibility без
 sector-dependent fitting;
 \item не менее двух независимых normalization-sensitive тестов;
 \item собственные kill-критерии, отделяющие новую модель от продолжения
 версии II.
\end{enumerate}

\begin{protocol}[Режим архивации S2T--II]
Версия II замораживается как воспроизводимая математическая программа и
карта отрицательных результатов. Новые численные совпадения, независимые
селекторы и локальные коэффициенты в неё не добавляются. Дальнейшая работа
разрешена только в двух режимах: внешнее воспроизведение зафиксированного
ядра либо отдельная версия III с новым parent action.
\end{protocol}

\chapter{Definition of Done и текущий вердикт II.B}

\begin{protocol}[Definition of Done для inverse-susceptibility ветви]
Ветвь считается закрытой положительно только если:
\begin{enumerate}
    \item выписан local или spectral parent action нового tensor-сектора;
    \item half-integer boundary condition выведено из заранее объявленной
    \(\mathbb Z_2\)-структуры, а не выбрано ради \(\zeta(4,1/2)\);
    \item полный BV/BRST complex вычислен без изоляции одного determinant;
    \item ordinary или normalized trace выбран из action principle;
    \item после всех cancellations остаётся positive Hessian \(H=\pi^4\);
    \item source normalization даёт \(J^2/2=1\);
    \item заново воспроизводится periodic ветвь \(1/24\) либо выводится её
    согласованная замена;
    \item тот же parent action проходит независимый normalization-sensitive
    сектор до сравнения с его наблюдаемым значением.
\end{enumerate}
\end{protocol}

\begin{nogocriterion}[Провал ветви II.B-Hessian]
Если complete ghost tower меняет знак или величину \(H\), half-shift требует
невыведенной boundary condition, normalized trace уничтожает factor six,
source coefficient остаётся свободным либо новый sector не проходит второй
observable, то inverse-susceptibility ветвь закрывается как математически
интересная reconstruction без физического вывода \(S_{vac}\).
\end{nogocriterion}

Текущий статус II.B разделён:
\begin{itemize}
    \item сектор поколений: цепь \(A_3\) и её градуировка выведены, а общая
    древесная связность имеет строгий CP-no-go; уникальный двудольный
    \(C_4\)-лифт создаёт координатный нечётный цикл и Schur-хорду, но его
    градуировка не спускается на канонический affine-триплет, а ненулевой
    поток смешивает триплет с равномерным синглетом; хиральное удвоение
    сохраняет триплет и исправляет конечный оператор, но допускает
    произвольный Yukawa-блок и не создаёт selector;
    \item электромагнитный сектор: изолированный грассманов блок даёт
    \(+\pi^4\), но полный БВ-комплекс оставляет в нулевой пространственной
    оболочке только \(-\pi^4/6\) или \(+\pi^4/6\), тогда как полный
    четырёхмерный след логарифмически расходится;
    \item общая дискретная голономия: все \(48\) вариантов проверены, и ни
    один не закрывает тензорный сектор и сектор поколений одновременно;
    \item непрерывная вильсоновская ветвь: найдены восемь совместных решений,
    устойчивая стационарная точка и сокращение осей \(8\to2\to1\) после
    факторизации по остаточной симметрии;
    \item граничный ротор: найден минимальный импульсный сектор \(1^3 3^5\),
    точно воспроизводящий обратный коэффициент, но каноническая проекция
    доказала запрет однослойного получения одновременно обратного и
    логарифмического членов;
    \item APS/orbifold/inflow ветвь: точная \(Z_2/Z_4\)-проекция уже выводит
    split-direction \((17/6,1/6,2)\), а self-dual inverse susceptibility
    даёт \(\pi^{-4}\); однако предложенные eta-сумма, rank \(24-1\) и bulk
    action не проходят operator, dimension и normalization gates;
    \item phase-to-mass gate: для минимального vectorlike charged-lepton
    determinant пары \(m\pm iE\) сокращают eta-фазу blockwise, mass response
    вещественен, а circle log-determinant не равен \(|\zeta_R(-1)|\); поэтому
    критерий Ф.2 закрыт отрицательно до появления выведенного chiral
    CP-механизма;
    \item conformal-rank gate: gravitational gauge quotient действует в
    metric block и оставляет Majorana rank равным \(24\); число \(23\) не
    следует из \(\widehat A\)-знаменателя или wrong-sign conformal mode, но
    условный class-D defect сохраняет явный rank-one kernel;
    \item defect parent-action gate: square-root \(\pi\)-flux и ненулевой
    charge-two condensate строго вынуждают winding \(1\), однако текущая
    action не выводит ни singular root sector, ни отрицательный pairing
    Hessian; defect остаётся условным saddle расширенной модели;
    \item root-source menu: SU(5)-quarter element действует на sterile
    \(Y=0\) как \(+1\), spin даёт только знаки, а Nambu-quarter transition
    является следствием уже заданного vortex; обязательный root-source в
    существующем field menu отсутствует;
    \item минимальная \(B-L\)-ветвь: одно \(N^c\) на поколение сокращает
    все continuous local anomalies, holonomy \(\oint a_{BL}=\pi/2\) даёт
    sterile root-фазу \(i\), однако ordinary scalar charge \(-2\) видит
    holonomy \(-1\) на generator \(y\) и не может иметь global nonzero
    condensate; twisted section
    \(L_{\mathrm{tor}}\otimes L_{B-L}^{-2}\) делает condensate parallel,
    но оставляет Yukawa vertex torsion-odd; поэтому доказана
    трилемма «корневая голономия --- масса --- однородный конденсат», и
    выживает только неоднородная дефектная конфигурация; минимальный
    функционал Гинзбурга--Ландау показывает, что она конденсируется лишь
    при \(\lambda v^2>\pi^2/L^2\) и имеет две точно вырожденные ветви
    \(n=0,-1\), так что топология не фиксирует ни существование
    стационарной конфигурации, ни её ориентацию; проверка уже допущенных
    спектральных функций даёт три
    разных режима --- нормальную фазу, критическую точку и конденсат, а
    вещественное чётное спектральное действие сохраняет вырождение
    \(k=\pm1\); следовательно, не выведены ни неустойчивость нормальной
    фазы, ни ориентационно-нечётный член; расширение также вводит новый
    калибровочный множитель, масштаб нарушения симметрии,
    матрицу Юкавы и радиационно порождаемое кинетическое смешивание
    \(\sum Y(B-L)=8/3\) на поколение, поэтому является новой моделью, а не
    замыканием II.A;
    \item единое действие пока не построено;
    \item новые физические предсказания отсутствуют.
\end{itemize}

Следовательно, II.B впервые содержит не только список желательных свойств, но
и конкретные операторные конструкции: двенадцать направлений смежности с
полной \(M_3\)-алгеброй, изолированный полусдвинутый тензорный гессиан
\(\pi^4\), восемь единичных голономных плоскостей, квантованный импульсный
сектор \(1^3 3^5\), модульная цепь \(A_3\), точный треугольный
CP-инвариант, уникальный минимальный двудольный \(C_4\)-лифт и корректный
хирально удвоенный конечный оператор. Одновременно
закрыты наивная однослойная роторная склейка
и стандартное БВ-завершение одного двухформенного поля; координатный
\(C_4\)-лифт также закрыт как продолжение прежнего affine-механизма
поколений. Это реальное сужение пространства моделей. Но новое
градуированное действие пока не выведено.

{\raggedright
Машинные результаты части II.B сохранены в
\texttt{s2t\_family\_affine\_incidence\_exhaustive\_results.json},
\texttt{s2t\_hurwitz\_hessian\_pi4\_results.json},
\texttt{s2t\_six\_channel\_inverse\_susceptibility\_results.json},
\texttt{s2t\_selfdual\_bilaplacian\_susceptibility\_results.json} и
\texttt{s2t\_halfshift\_tensor\_ghost\_hessian\_results.json}, а прямой
двухсекторный тест --- в
\texttt{s2t\_shared\_holonomy\_two\_sector\_results.json}, а continuous
gap-action --- в
\texttt{s2t\_continuous\_wilson\_gap\_action\_results.json}, а граничный
ротор --- в
\texttt{s2t\_wilson\_boundary\_rotor\_action\_results.json}. Точная весовая
проверка сохранена в
\texttt{s2t\_wilson\_rotor\_exact\_determinant\_results.json}.
Импульсный перебор сохранён в
\texttt{s2t\_wilson\_rotor\_momentum\_sector\_results.json}.
Каноническая проекция проверена в
\texttt{s2t\_rotor\_fixed\_momentum\_projection\_results.json}.
Полный скрученный БВ-комплекс проверен в
\texttt{s2t\_twisted\_twoform\_bv\_complete\_results.json}.
Остаточная осевая орбита и минимальный критерий CKM проверены в
\texttt{s2t\_family\_residual\_axis\_orbit\_results.json} и
\texttt{s2t\_family\_relative\_orbit\_ckm\_results.json}.
Полный перебор минимального операторного меню сохранён в
\texttt{s2t\_family\_minimal\_operator\_menu\_exhaustive\_results.json}.
Путевой перебор сохранён в
\texttt{s2t\_family\_oriented\_wilson\_word\_results.json}.
Модульная градуировка и рёберная когомология проверены в
\texttt{s2t\_modular\_grading\_cocycle\_results.json}.
Физический red-team двухслойной текстуры сохранён в
\texttt{s2t\_two\_layer\_physical\_ckm\_redteam\_results.json}.
Минимальный двудольный lift проверен в
\texttt{s2t\_family\_bipartite\_c4\_lift\_results.json}.
Хиральное удвоение и физический Yukawa-перебор проверены в
\texttt{s2t\_chiral\_doubled\_triplet\_yukawa\_results.json}.
APS/orbifold/inflow red-team сохранён в
\texttt{s2t\_aps\_orbifold\_inflow\_redteam\_results.json}.
Прямой тест eta-фазы и вещественной массы сохранён в
\texttt{s2t\_eta\_phase\_mass\_gate\_results.json}.
Тест conformal mode и Majorana-rank сохранён в
\texttt{s2t\_conformal\_majorana\_rank\_gate\_results.json}.
Тест динамического происхождения Majorana-дефекта сохранён в
\texttt{s2t\_majorana\_defect\_parent\_action\_gate\_results.json}.
Исчерпание существующего root-source menu сохранено в
\texttt{s2t\_majorana\_root\_source\_menu\_results.json}.
Минимальное \(U(1)_{B-L}\)-расширение проверено в
\texttt{s2t\_bl\_root\_extension\_gate\_results.json}.
Global compatibility root holonomy и pairing condensate проверена в
\texttt{s2t\_bl\_root\_condensate\_consistency\_results.json}.
Root--mass--condensate trilemma проверена в
\texttt{s2t\_root\_mass\_condensate\_trilemma\_results.json}.
Неоднородная конфигурация наименьшего действия проверена в
\texttt{s2t\_nonuniform\_pairing\_saddle\_results.json}.
Спектральные коэффициенты жёсткости и запрет выбора ориентации проверены в
\texttt{s2t\_spectral\_pairing\_stiffness\_gate\_results.json}.
\par}

\end{document}